% 本文件由 LaTeX 排版 Agent 根据有效 translation.json 自动生成。
% author=Leo the Great, Pope (c. 395-461); work_id=3604; title=书信集

\chapter{书信集}

% structure: p0001 source=h1 target=omitted reason=page-title-already-rendered-by-parent

书信1\newline{}书信4\newline{}书信6\newline{}书信7\newline{}书信9\newline{}书信10\newline{}书信12\newline{}书信14\newline{}书信15\newline{}书信16\newline{}书信17\newline{}书信18\newline{}书信19\newline{}书信20\newline{}书信21\newline{}书信22\newline{}书信23\newline{}书信24\newline{}书信26\newline{}书信27\newline{}书信28——更广为人知的是\WorkTitle{大卷}\newline{}书信29\newline{}书信31\newline{}书信32\newline{}书信33\newline{}书信34\newline{}书信35\newline{}书信37\newline{}书信38\newline{}书信39\newline{}书信40\newline{}书信42\newline{}书信43\newline{}书信44\newline{}书信45\newline{}书信52\newline{}书信56\newline{}书信59\newline{}书信66\newline{}书信67\newline{}书信68\newline{}书信69\newline{}书信79\newline{}书信80\newline{}书信82\newline{}书信85\newline{}书信88\newline{}书信93\newline{}书信95\newline{}书信98\newline{}书信101\newline{}书信104\newline{}书信105\newline{}书信106\newline{}书信108\newline{}书信109\newline{}书信113\newline{}书信117\newline{}书信119\newline{}书信120\newline{}书信123\newline{}书信124\newline{}书信129\newline{}书信139\newline{}书信156\newline{}书信158\newline{}书信159\newline{}书信162\newline{}书信164\newline{}书信166\newline{}书信167\newline{}书信169\newline{}书信171\par

\SourceRecord{Translated by Charles Lett Feltoe.；Nicene and Post-Nicene Fathers, Second Series；Vol. 12.；Edited by Philip Schaff and Henry Wace.；Buffalo, NY: Christian Literature Publishing Co.,；1895.；<http://www.newadvent.org/fathers/3604.htm>.}

% structure: p0001 source=h1 target=section reason=page-title

\section{第一封信}

\Emph{致阿奎莱亚主教。}\par

% structure: p0003 source=h2 target=subsection reason=relative-heading-rank; base=h2

\subsection{一、因当局疏忽，佩拉奇异端在其省内蔓延。}

从我们的圣兄弟及同为主教的塞普蒂默在附函中的陈述，我们得知在贵省，某些司铎、执事及不同等级的圣职人员，因被佩拉奇异端或塞莱斯蒂异端所吸引，未经对其特有错误作任何否认，便获准进入公教共融；并且，当放哨的牧人沉睡时，披着羊皮却未弃其兽心的豺狼已进入主的羊栈。他们惯于行那些根据我们法令和规章的禁令就连未犯罪者也不允许的事：即离开他们领受或恢复职务的教会，到处散布他们的不确定，喜欢不断漂泊，从未停留在宗徒的基础上。因为未经任何考验筛分，也未受任何先前信仰宣认的约束，他们大肆强调其权利，以与教会共融为名，到一家又一家，因其虚假名声，许多人不知情而腐蚀了众多人心。但我确信，若是各教会的主事者在此类接纳之事上尽其应有的勤勉，不允许他们从一处漂泊到另一处，他们便无法行此事。\par

% structure: p0005 source=h2 target=subsection reason=relative-heading-rank; base=h2

\subsection{二、他命令召开省区教务会议，明确要求异端者作否认。}

因此，为避免此事再进一步，也避免这种因某些人的疏忽而引入的有害习惯导致许多灵魂的覆亡，我们以权威命令：兄弟，你当负责召集贵省圣职人员的教务会议，所有那些从与佩拉奇派和塞莱斯蒂派的联合中如此仓促地重新接纳到公教共融的人，无论司铎、执事或任何等级的圣职人员，之前并未被强制否认其错误，如今至少应迫使他们作真正的改正，这能有益于他们自身，也不伤害任何人，因为他们的诡诈已部分暴露。让他们通过公开宣认，谴责这僭越错误的创始人，并谴责普世教会在其教义中所弃绝的一切；让他们以亲笔签名的充分公开声明，宣告他们接受并完全赞同宗座权威为根除此异端所批准的所有教务会议法令。让他们的话语中没有任何模糊不清、模棱两可之处。因为我们知道，他们的狡猾在于，他们认为若将其可恶教义的特定条款与其余可咒观点分开，便可辩护其含义。\par

% structure: p0007 source=h2 target=subsection reason=relative-heading-rank; base=h2

\subsection{三、佩拉奇派关于天主恩宠的观点不合圣经。}

当他们假装不赞同并放弃所有定义，以便通过其完全的欺骗技艺规避，除非其意图被识破，他们便例外地保留那一教条：天主的恩宠是按领受者的功过赐予的。然而，除非白白赐予，它便不是恩赐，而是功过的酬报与补偿。因为蒙福的宗徒说：\BibleQuote{你们借着恩宠得救，是因着信德，而且这并不由你们自己，而是天主的恩赐；不是出于功行，免得有人自夸。因为我们是他所造化，在基督耶稣内受造，为行善事，就是天主所预备叫我们行在其中的。}\BibleRef{弗 2:8} 因此，一切善功的赐予都是天主的预备：因为人是因恩宠成义，而非因自身优异；因为恩宠为每一个人是正义之源、美善之源、功德之源。但这些异端者说，恩宠被人的自然良善所预期，原因在于，那种在他们看来先于恩宠便以自身善愿显著的自然，似乎不受原罪任何玷污，从而真理之言被歪曲：\BibleQuote{因为人子来，是为寻找并拯救那失落的。}\par

% structure: p0009 source=h2 target=subsection reason=relative-heading-rank; base=h2

\subsection{四、必须迅速采取措施。}

因此，亲爱的，你当注意，并以极大的勤勉作出安排，免得长久以来已被消除的冒犯借着这样的人重新兴起，也免得那已根除的教义在贵省生发出同样邪恶的种子；因为不仅会生根生长，还会以其毒害气息玷污教会未来的世代。那些希望表现出改正的人必须清除自己的一切嫌疑；并借服从我们，证明自己是属于我们的。若有人拒绝满足我们这有益的命令，无论圣职或平信徒，必须将其逐出教会团体，否则他既背叛他人的安全，也丧失自己的灵魂。\par

% structure: p0011 source=h2 target=subsection reason=relative-heading-rank; base=h2

\subsection{五、对从一教会漂泊到另一教会的圣职人员，必须执行法令。}

我们也劝告你，将教会纪律中那一部分恢复全面运作，即神圣教父们和我们以往多次下令规定：无论在司铎等级、执事序列或较低圣职中，任何人均不可随意从一教会迁移到另一教会；如此，每人当坚守他被祝圣之处，不受野心引诱、不被贪欲引入歧途、不为人的邪说所腐蚀；因此，若有人寻求自己的利益，而非耶稣基督的利益，忽视回到自己的子民和教会，那么在晋升和共融联系方面，他必被算作局外人。至于你，亲爱的，不要怀疑，如果我们认为，正如我们所不愿的，我们为维护法令和信仰完整所颁布的命令被忽视，我们必会相当严厉地动怒；因为低级圣职的过失，最应归咎于那些怠惰和散漫的统治者，他们往往因畏惧施用严厉的补救而助长了许多瘟疫。\par

\SourceRecord{Translated by Charles Lett Feltoe.；Nicene and Post-Nicene Fathers, Second Series；Vol. 12.；Edited by Philip Schaff and Henry Wace.；Buffalo, NY: Christian Literature Publishing Co.,；1895.；<http://www.newadvent.org/fathers/3604001.htm>.}

% structure: p0001 source=h1 target=section reason=page-title

\section{第四封信}

\Person{良}，罗马城主教，致候在\Place{坎帕尼亚}、\Place{皮努姆}、\Place{埃特鲁里亚}及各省所有被任命的主教，在主内问安。\par

% structure: p0003 source=h2 target=subsection reason=relative-heading-rank; base=h2

\subsection{一、引言}

由于教会和平安定使我们满意，但每当我们获悉任何违背教规和教会纪律之事被想当然或做出时，我们便深感悲伤；如果我们不尽应有的警觉加以制止，我们便无法向那指派我们为守望者的\BibleRef{则 3:17}交账——因纵容邪恶的阴谋者玷污本应保持纯洁的教会肢体，而肢体的和谐因这样的缄默而破坏。\par

% structure: p0005 source=h2 target=subsection reason=relative-heading-rank; base=h2

\subsection{二、不得祝圣奴隶和\OriginalTerm{隶农}{coloni}}

常有人被接纳进入圣秩，他们并无任何出身或品德的尊贵；甚至有些未能从主人手中获得自由的人也被提升至司铎品位，仿佛可怜的奴隶适合那样的荣誉；人们竟相信一个尚未能向自己主人证明自己的人能被天主认可。因此，在这件事上的抱怨是双重的：因为神圣的职务被如此低劣的伴侣所玷污，同时主人的权利受到侵犯，因为非法占有被轻率地施行。因此，亲爱的弟兄们，你们省的所有司铎应当远离这些人；而且不仅这些人，我们也希望你们远离那些受出身或其他服役条件束缚的人；除非那些对他们有管辖权的人提出了请求或同意。因为凡要被登记在天主服务中的人，应当不受他人约束，免得他被任何其他义务的束缚从他名字所登记的主的营中拉走。\par

% structure: p0007 source=h2 target=subsection reason=relative-heading-rank; base=h2

\subsection{三、曾两次结婚或娶寡妇者不得任司铎}

再者，当每人的出身和品行在得到认可后，什么样的人应被选入神圣祭台的服务，我们已从宗徒的教导、天主的诫命和教规的规定中得知。从这些规定我们发现许多弟兄已经偏离并完全走入了歧途。因为众所周知，娶了寡妇的人获得了司铎职；某些有过多位妻子并过着放纵生活的人，被给予一切便利并获准进入圣秩，这违背了蒙福宗徒的宣称，他对这样的人宣告说：\Quote{一妇之夫}，也违背了古代法律的诫命，它告诫说：\Quote{司铎应娶处女为妻，不可娶寡妇，不可娶被休的妇人。}因此，所有已被接纳的这类人，我们命令藉宗座权威将他们从教会职务和司铎名分中革除：因为他们将不会有任何资格去要求他们因障碍而不合格的事；我们特别为自己保留处理此事的职责，即如果有任何此类违规行为发生，应予以纠正，并不允许再次发生，且不得以无知为借口；尽管司铎从来不允许对教规所规定之事无知。因此，我们已通过我们的兄弟和同为主教的\Person{依诺森}、\Person{勒吉提默}和\Person{塞格提乌}的手，将这些信件送达你们各省：为了使那些已知萌发的恶苗被连根拔除，没有莠子破坏主的庄稼。因为如此所有真正的好种子将结出许多果实，如果那惯常杀害生长作物的东西被仔细清除。\par

% structure: p0009 source=h2 target=subsection reason=relative-heading-rank; base=h2

\subsection{四、禁止神职人员和教友的高利贷行为}

这一点我们也认为不可忽略：某些被卑劣获利之爱所占据的人，以利息放贷，并希望像高利贷者一样致富。因为令我们悲痛的是，这不仅发生在那些属于神职人员中的人身上，也发生在那些希望被称为基督徒的教友身上。我们规定，被定罪者应受严厉惩罚，以消除一切犯罪的机会。\par

% structure: p0011 source=h2 target=subsection reason=relative-heading-rank; base=h2

\subsection{五、神职人员不可借他人名义牟利，如同不可借己名义牟利一样}

以下警告我们也认为适宜发出：任何神职人员都不可借他人名义牟利，如同不可借己名义牟利一样。因为以他人的收益掩盖自己的罪行是不相宜的。不，我们应只着眼于并追求那种利息：藉着它，我们在此施予怜悯所给予的，可以从主那里收回，祂将偿还千倍、永恒之物。\par

% structure: p0013 source=h2 target=subsection reason=relative-heading-rank; base=h2

\subsection{六、任何拒绝赞同这些规条的主教必须被革职}

因此，我们的这一劝诫宣称：如果有任何我们的弟兄试图违犯这些规条，并胆敢做禁止之事，他应当知道他将被革除职位，并且他将不会与我们的共融有份，因为他拒绝参与我们的纪律。但惟恐有任何我们可能被认为遗漏之事，我们吩咐你们，亲爱的，要遵守蒙福的纪念者\Person{依诺森}的所有教令规则，以及我们所有前任所颁布的关于教会秩序和教规纪律的规则，并要如此遵守：若有任何人违犯它们，他立即知道一切宽容均被拒绝。\par

日期：十月十日，在显赫的\Person{马克西穆斯}（第二次）和\Person{帕特里乌斯}执政期间（公元443年）。\par

\SourceRecord{Translated by Charles Lett Feltoe.；Nicene and Post-Nicene Fathers, Second Series；Vol. 12.；Edited by Philip Schaff and Henry Wace.；Buffalo, NY: Christian Literature Publishing Co.,；1895.；<http://www.newadvent.org/fathers/3604004.htm>.}

% structure: p0001 source=h1 target=section reason=page-title

\section{第六封信}

\Person{良}致亲爱的弟兄\Person{阿纳斯塔西}。\par

% structure: p0003 source=h2 target=subsection reason=relative-heading-rank; base=h2

\subsection{一、他高兴\Place{伊利里库姆}的主教们就重要问题咨询了他。}

同僚们的手足之爱使我们以感恩之心阅读所有司铎的信函；因为在信中，我们仿佛面对面地在精神上彼此拥抱，通过这类书信的往来，我们相互交谈。但在目前这封信中，所显示的情感似乎比往常更为深厚：因为它告知我们各教会的情况，并敦促我们因着对自己职务的反思而保持警醒的关怀，以致我们如同按主的旨意被安置在瞭望塔上，既应赞同那些符合我们意愿的事，也应以强制性的补救措施纠正我们观察到因任何攻击而偏离正轨的事：希望如果我们不容许那些已开始滋生并破坏收成的东西增多，我们的播种将结出丰硕的果实。\par

% structure: p0005 source=h2 target=subsection reason=relative-heading-rank; base=h2

\subsection{二、效法前任先例，他任命\Person{阿纳斯塔西}为\Place{伊利里库姆}的都主教。}

因此如今，亲爱的弟兄，你的请求已通过我们的儿子、司铎\Person{尼各老}告知我们，即你愿如同你的前任们那样，从我这里依序获得对\Place{伊利里库姆}的管理权以遵守法规，我们表示同意，并恳切劝诫，在管理分布于整个\Place{伊利里库姆}的教会时，不得有任何隐瞒或疏忽，这些教会我们委托给你，代我们管理，效法荣福纪念的\Person{西里西乌}的先例，他当时首次按照固定方法，将它们托付给你的前前任，圣洁纪念的\Person{阿尼西乌}，他当时对\Place{宗座}颇有功绩，并经过后来事件的考验：以便他能援助该省内的教会，他愿这些教会保持纪律。崇高的先例必须热切效法，以使我们能在一切事上显示自己如同那些我们愿享有其特权的人一样。我们愿你效法你的前前任以及你的直接前任，后者与前一位一样，都曾享有并善用了这一特权：这样，我们便能因我们委托给你、代我们管理的教会的进展而喜乐。因为当事务托付给一个行事正确、善于履行司铎职位职责的人时，事情会进展得令人称许；反之，当权力托付给一个人，而他在使用时缺乏应有的节制时，这权力便成了他的负担。\par

% structure: p0007 source=h2 target=subsection reason=relative-heading-rank; base=h2

\subsection{三、领受圣秩者必须根据教会法规谨慎遴选。}

因此，亲爱的弟兄，你要警醒地执掌托付给你的船舵，将你心灵的目光环视你看到托付给你的一切，守护那些有助于你获得赏报的事，并抵抗那些企图扰乱教会法规纪律的人。天主的法律之制裁必须受到尊重，而教会法规的谕令更应被遵守。在你所负责的各省内，只祝圣那些因堪当的生活和在圣职人员中的地位而受称赞的司铎。不得允许人情偏袒、拉票或购买选票。对即将被祝圣者，要仔细审查他们的案件，让他们在相当长的生命时期内接受教会的纪律训练。但如果他们身上具备所有圣教父们的要求，并且他们遵守了我们读到蒙福的宗徒\Person{保禄}所命令的一切，即他只作一个妻子的丈夫，且他娶她时她是贞女，正如天主的法律所要求的，那么就可祝圣他们。我们极其渴望这一点能得到遵守，以消除一切推诿的余地，免得有人以为自己能达到司铎品位，而他在获得基督的恩宠之前娶了妻，并在她去世后，于领洗后又与另一人结合。因为前妻不能忽略，先前的婚姻不能不算数，他因前妻在领洗前所生的孩子，与他因第二任妻子在领洗后所生的孩子一样，都是孩子的父亲。因为正如罪过及明知是不合法的事情在圣洗的水中被洗净，同样，许可或合法的事情并不被消除。\par

% structure: p0009 source=h2 target=subsection reason=relative-heading-rank; base=h2

\subsection{四、都主教不得未经咨询首席主教而仓促祝圣。}

在这些教会中，不可轻率地祝圣任何人为司铎；因为倘若你的审查，弟兄，令人畏惧，这样便会对被选者形成成熟的判断。但若有任何主教违反我们的命令，由他的都主教未经你知晓而祝圣，要知道他在我们这里没有稳固的地位，而那些擅自如此行的人，必须为自己的擅专行为负责。然而，既然每位都主教被授予如此的权力，以致他有权利在他的省区中祝圣，那么我们也希望这些都主教被祝圣，但不可没有成熟而经过深思熟虑的判断。因为虽然所有被祝圣为司铎的人都应被认可并中悦天主，但我们更希望那些我们已知将要管理他们所属的同僚司铎的人具有特别的卓越。我们告诫你，亲爱的，要更勤勉谨慎地注意此事，使你能证明遵守了宗徒的那条训诫，即\BibleQuote{不可轻易给人覆手。\BibleRef{弟前 5:22}}。\par

% structure: p0011 source=h2 target=subsection reason=relative-heading-rank; base=h2

\subsection{五、在省级教务会议上无法解决的问题应提交\Place{罗马}。}

凡被召出席主教会议的弟兄，都当参加，不可拒绝临于圣会；因为在那里他尤其应当知道，凡有助于教会良好纪律的事，都必须予以解决。因为主的司铎之间若有更频繁的会议，一切过失就更能避免，而亲密的交往对进步和兄弟之爱都大有助益。在那里，若有任何问题出现，在主引领下，它们将能得到确定，以至不留任何恶感，而只存留弟兄们之间更坚固的爱。但若有更重要的问题出现，是在你主持下不能在那里解决的，弟兄，就请发送你的报告并咨询我们，以便我们能在主的启示下回信——我们能做何事全赖祂的仁慈，因为祂已向我们施恩——使我们能借由我们的裁决，按照古传传统和对宗座应有的尊重，维护我们审理的权利；因为我们既希望你在我们之位行使你的权柄，同样，我们将那些不能当场决定的事项以及向我们上诉的人，保留给我们自己。\par

% structure: p0013 source=h2 target=subsection reason=relative-heading-rank; base=h2

\subsection{六、司铎与执事不可像主教一样在平日被祝圣。}

你当安排此信送达所有弟兄知晓，以便今后无人能以无知为借口，不遵守我们所命令的这些事。我们也将劝诫信直接寄给了各省的都主教们，使他们知道必须服从宗座的训令，并因开始服从你——弟兄，我们的代表——而服从我们，正如我们所写的。我们确实听说，且不能默然不语：某些弟兄只在主日祝圣主教；但对于司铎和执事（他们的祝圣同样应当隆重），却随意在任何一日授予司祭职的尊位，这是违反教规和教父传统应受谴责的做法，因为持有这传统的那些人理当在一切圣秩中遵守此惯例：以便经过适当的时间，即将被祝圣为司铎或执事的人，能通过圣职的所有级别逐步晋升，这样一个人就有时间学习他自己将来也要教导之事。\Signature{于一月十二日，狄奥多西第十八次执政官与阿尔比努斯执政官之年（四四四年）。}\par

\SourceRecord{Translated by Charles Lett Feltoe.；Nicene and Post-Nicene Fathers, Second Series；Vol. 12.；Edited by Philip Schaff and Henry Wace.；Buffalo, NY: Christian Literature Publishing Co.,；1895.；<http://www.newadvent.org/fathers/3604006.htm>.}

% structure: p0001 source=h1 target=section reason=page-title

\section{第七封书信}

良致意大利各省主教，问候。\par

% structure: p0003 source=h2 target=subsection reason=relative-heading-rank; base=h2

\subsection{一、在罗马发现了许多摩尼教徒。}

我们邀请你们分担我们的忧虑，好使你们以牧者的勤勉，更谨慎地照看托付给你们的羊群，不让魔鬼的诡计得逞；免得那因主的慈悲眷顾，藉我们的谨慎从我们羊群中驱除的瘟疫，在你们尚未得到警告并仍对发生之事一无所知时，蔓延到你们的教会，并找到机会暗中潜入你们中间，从而我们在这罗马城中所制止的，竟在你们中暗地生根并滋长。我们的搜查已在罗马城发现了许多摩尼教不敬神的追随者和教师，我们的警醒揭露了他们，我们的权威和谴责制止了他们：凡我们能改造的，我们已加以纠正，并促使他们在教会中公开宣认，用自己的手签字谴责摩尼及其宣讲和教义，从而我们使那些承认过失的人，藉给予他们悔改的机会，从他们不义的深坑中提升出来。然而，有许多人陷得太深，任何补救都无法帮助他们，我们便依照我们基督徒君王的法律将他们交付法律制裁，免得他们以传染污染圣洁的羊群，并由公共法官将他们判为永久流放。我们还将他们著作和秘密传统中发现的一切亵渎可耻之事，公开而明确地向基督徒平信徒展示，使民众知道应当回避或远离什么；这样，连那自称是他们主教的人，也被我们审讯，并暴露了他在神秘宗教中所持的罪恶观点，正如我们审理记录可以向你们展示的那样。为此，我们也将这些记录发送给你们作为训导：阅读之后，你们将能了解我们所发现的一切。\par

% structure: p0005 source=h2 target=subsection reason=relative-heading-rank; base=h2

\subsection{二、意大利的主教们不可容许那些已离开本城的摩尼教徒逃脱或隐藏。}

因为我们知道，许多在此地受到过于紧密指控以致无法洗清自己的人已经逃走，我们便通过我们的辅祭将此信发送给你们，蒙爱者：好使你们的圣洁，亲爱的兄弟们，知晓此事，并愿以勤勉和谨慎行事，免得摩尼教错误之人有机会伤害你们的民众并宣讲他们不敬神的教义。因为除非我们以信德的热忱在主内追捕那些毁灭者和被毁灭者，并用我们所能施加的严厉，将他们与健全心智的交往隔绝，以免此瘟疫传播得更广，否则我们无法治理那托付给我们的人。因此，我劝勉你们，蒙爱者，我恳求和警告你们，要运用你们所应当和能够运用的警觉的勤勉来追查他们，免得他们在任何地方找到隐藏的机会。因为那执行有益于托付给他的人民健康之事的人，将从天主得到相应的赏报；而在主的审判座前，那不愿保护自己人民免受不敬神谬误传播者之害的人，将无法以疏忽为由为自己辩解。\newline{}444年1月30日，在著名的狄奥多西·奥古斯都（第十八次执政）和阿尔比努斯执政时期。\par

\SourceRecord{Translated by Charles Lett Feltoe.；Nicene and Post-Nicene Fathers, Second Series；Vol. 12.；Edited by Philip Schaff and Henry Wace.；Buffalo, NY: Christian Literature Publishing Co.,；1895.；<http://www.newadvent.org/fathers/3604007.htm>.}

% structure: p0001 source=h1 target=section reason=page-title

\section{第九封信}

\Person{良}主教，致亚历山大主教\Person{狄奥斯库若}，安好。\par

% structure: p0003 source=h2 target=subsection reason=relative-heading-rank; base=h2

\subsection{一、罗马教会与亚历山大教会应在一切事上合一。}

我们对您所怀的天主之爱有多深，亲爱的，您将从这一点估量出来：我们急于使您的开端建立在更稳妥的基础上，以免您的爱德似乎欠缺什么，因为您属灵恩宠的功行，如我们所验证，已为您赢得青睐。因此，圣善的弟兄，父执与弟兄般的交谈理应最令您感激，并以您所知我们向您所怀的同一精神接纳它。因为您与我们应在思想和行动上合一，正如我们所读到的那样，愿在我们身上也证明有同一颗心和同一个灵魂。既然最蒙福的伯多禄从主领受了宗徒之首的职位，而罗马教会至今仍遵守主的制度，那么认为他的圣洁门徒马尔谷——他首先治理亚历山大教会——所制定的法令遵循了不同的传统，这是邪恶的：因为毋庸置疑，门徒和师傅从同一恩宠泉源中只汲取了同一圣神，受祝圣者除了从他祝圣者那里领受的以外，不能传授任何别的东西。因此，我们不允许自己在任何事上有分歧，因为我们自认属于同一个身体和同一信德，也不允许师傅的制度显得与被教导者的不同。\par

% structure: p0005 source=h2 target=subsection reason=relative-heading-rank; base=h2

\subsection{二、应遵守固定日期祝圣司铎和执事。}

因此，我们知道先祖们非常谨慎遵守的，也希望您同样遵守，即司铎或执事的祝圣不应随意在任何一天举行：而应选择星期六之后、那一夜的开端，即一周第一日黎明之前，在那夜赐予圣洁祝福给那些将受祝圣者，祝圣者和受祝圣者都要禁食。如果确实在主日的早晨举行而不打破星期六的禁食，这一惯例就不会被违反：因为前一夜的开端构成该时期的一部分，并且无疑属于复活之日，正如关于逾越节庆典所明确规定的。因为除了我们知道基于宗徒教导的惯例之分量外，圣经也清楚表明这一点，因为当宗徒们在圣神的命令下派遣保禄和巴尔纳伯向外邦人传福音时，他们禁食祈祷为他们覆手：好叫我们知道，给予者和接受者都必须以何等虔诚谨慎，免得如此蒙福的圣事似乎被轻率地举行。因此，如果您自己也遵守这种祝圣司铎的形式，在主要求您治理的所有教会中，您将虔诚而可赞地遵循宗徒的先例：即那些将要受祝圣的人，除非在主复活之日（通常被认为从星期六傍晚开始）绝不接受祝福，这一日在天主奥妙的安排中已多次被祝圣，以至于主所有更显赫的制度都在那个崇高的日子完成。在那一天，世界开始了。在那一天，借着基督的复活，死亡被摧毁，生命开始。在那一天，宗徒们从主手中接过要向万民宣讲的福音号角，并领受他们要将之带给全世界的再生圣事。在那一天，正如蒙福的福音作者若望所见证的，当所有门徒聚集在一个地方，门关着，主进入他们中间，向他们嘘气说：\Quote{你们领受圣神：你们赦免谁的罪，谁就得赦免；你们保留谁的罪，谁就被保留}\BibleRef{若 20:22}。最后，在那一天，主所应许给宗徒们的圣神降临了：因此我们知道，这似乎是由一条天上的规则所提示和传递的，即我们应当在那些全部恩赐被赐予的日子举行司铎祝福的奥迹。\par

% structure: p0007 source=h2 target=subsection reason=relative-heading-rank; base=h2

\subsection{三、在大庆节重复举行圣体圣事并非不妥。}

再者，为使我们的习惯在各方面一致，我们也希望遵守这一点：即每逢较大的庆节吸引比平时更多的会众，且信众的聚集过于庞大，一间教堂无法同时容纳所有人时，不应犹豫重复举行祭献的奉献：否则，如果只有先到者得以参与此礼仪，后来涌入的人就会显得被拒绝；因为完全合乎虔敬与理性：每当新的会众充满正在举行礼仪的教堂时，理应照常奉献祭献。然而，如果只举行一台弥撒的习惯被保持，只有早到的人才能奉献祭献，则一部分人必被剥夺他们的敬礼。因此，我们热切而友爱地劝诫您，亲爱的，希望您的谨慎也不要忽略那些基于先祖传统而成为我们自己习惯之一部分的事，以便我们在信仰和行为上完全一致。因此，我们把这封信交给了我们的儿子波西多尼乌斯司铎，在他返回时，让他带给您，弟兄；他曾多次参与我们的礼仪和祝圣，并多次被派到我们这里，因此他非常清楚我们在一切事上拥有何种宗徒权威。日期：6月21日（？445年）。\par

\SourceRecord{Translated by Charles Lett Feltoe.；Nicene and Post-Nicene Fathers, Second Series；Vol. 12.；Edited by Philip Schaff and Henry Wace.；Buffalo, NY: Christian Literature Publishing Co.,；1895.；<http://www.newadvent.org/fathers/3604009.htm>.}

% structure: p0001 source=h1 target=section reason=page-title

\section{第十封信}

\Emph{致维埃纳省众主教。关于阿尔勒主教希拉利之事。}\par

% structure: p0003 source=h2 target=subsection reason=relative-heading-rank; base=h2

\subsection{一、建立在圣伯多禄磐石上的教会的团结必须处处维护。}

致亲爱的弟兄们，维埃纳省全体主教，罗马主教良。\par

我们的主耶稣基督，人类的救主，建立了神圣宗教的遵守，祂愿藉天主恩宠将此宗教的光辉照耀万民万邦，如此，那先前仅局限在法律和先知宣告中的真理，可藉宗徒的号角声为所有人的救恩传出，正如经上所载：\BibleQuote{他们的声音传遍了天下，他们的言语传到了地极。} 但是主愿这奥秘的职分确实为所有宗徒所关切，但祂将首要的职责放在了蒙福的伯多禄、所有宗徒之首身上：且愿从祂这位元首，祂的恩赐流向整个身体：如此，任何敢于脱离伯多禄坚固磐石的人，可以明白他在神圣奥秘中无分无份。因为祂愿那被接纳进入祂不可分割的合一中作伙伴的人，被称为祂自己所是的，当祂说：\BibleQuote{你是伯多禄，在这磐石上我要建立我的教会\BibleRef{玛 16:18}}：为使那永恒圣殿的建造，藉天主恩宠的奇妙恩赐，安息在伯多禄的坚固磐石上：如此坚固祂的教会，以致人的鲁莽不能攻击它，地狱的门也不能战胜它。但是，这块磐石的最神圣的坚固，如我们所说，是由天主建造之手所建立，一个人若以过分邪恶试图摧毁它，当他试图打破它的力量，偏向自己的欲望，不跟随他从古人所接受的，他必须愿意摧毁它：因为他相信自己不受任何法律约束，不受任何天主诫命的规则管制，并在渴望新奇的冲动中，脱离你们和我们的惯例，采取非法做法，并使他应当遵守的陷于废弃。\par

% structure: p0006 source=h2 target=subsection reason=relative-heading-rank; base=h2

\subsection{二、希拉利以其不服从而扰乱教会的平安。}

但是，我们相信，在天主的核准下，并保持对你们充满的爱，正如宗座始终如你们所记得的那样对你们所付出的，圣洁的弟兄们，我们正努力通过成熟的商议来纠正这些事，并与你们分担整顿你们教会的任务，不是靠创新，而是靠恢复旧制；好使我们能持守我们的父辈传给我们的惯常状态，并通过善工的职事除去纷争的绊脚石，以悦乐我们的天主。因此，我们愿你们记得，弟兄们，如同我们一样，宗座（因其所受的敬畏）曾无数次被你们省和其他省的司铎们咨询和请教，在各种上诉事项中，依循旧例，它曾推翻或确认判决；如此，\BibleQuote{在和平的联系中保持精神的合一\BibleRef{弗 4:3}}，并且通过书信的交流，我们可敬的作为促进了持久的友爱：因为\BibleQuote{不求自己的益处，只求基督的\BibleRef{斐 2:21}}，我们一直注意不蔑视天主赋予教会和司铎们的尊严。但是，这条在我们父辈那里一直保持得如此良好并智慧维护的道路，希拉利已经背离，并可能以其新奇的傲慢扰乱司铎们的地位和共识：他渴望将你们置于他的权力之下，却不容自己受制于蒙福的宗徒伯多禄，为自己夺取高卢各省所有教会的祝圣权，并将本属于教省主教们的尊严转移给自己；他甚至以其骄傲的言语贬低对蒙福伯多禄本身的敬意：因为不仅释放和束缚的权力比其他宗徒先给了伯多禄，而且喂养羊群的职责也特别托付给了伯多禄。然而，任何认为必须否认伯多禄为首地位的人，实际上并不能减少他的尊严：但他被自己的骄傲气息所膨胀，并将自己投入最深的深渊。\par

% structure: p0008 source=h2 target=subsection reason=relative-heading-rank; base=h2

\subsection{三、塞利多尼已被恢复主教职务，因控告他的罪名被证明为虚假。}

因此，我们的会议记录记载了我们在\Person{塞利多尼乌斯}主教一事上所采取的行动，以及\Person{希拉里}在上述主教面前和在场时所说过的话。因为当\Person{希拉里}在神圣司铎的会议中无法给出合理的答复时，\Quote{他心中的隐秘}所发出的言论，是任何平信徒都不能讲、任何司铎都不能听的。弟兄们，我承认，我们感到痛心，并努力以耐心来平息他内心的激动。因为我们不愿激化他因傲慢反驳而给自己灵魂造成的创伤，我们努力安抚这位我们视为弟兄的人——尽管是他自己的回答纠缠了自己——而不是用我们的话使他痛苦。\Person{塞利多尼乌斯}主教因此被宣告无罪，因为他在自己面前通过证人明确的回答证明了自己被不公正地剥夺了司铎职务，以致与我们同在一起的\Person{希拉里}无话可反对。因此，那份被提出并宣读的判决被撤销了，该判决声称，作为寡妇的丈夫，他不能担任司铎职务。我们维持法律的规定，希望严格遵守这一规则，不仅对司铎，也对较低等级的圣职人员：那些缔结了此类婚姻的人，或那些被证明违反宗徒训导、不止娶一个妻子的人，不应被允许进入神圣职务。但我们规定，那些因其自身行为而受谴责的人，要么根本不被接纳，要么如果已被接纳则必须被免职；同样，那些被诬告的人，经过审查后，我们必须为他们洗清罪名，不允许他们失去职务。因为如果指控属实，原判就会对他生效。因此，我们的同僚主教\Person{塞利多尼乌斯}被恢复了他的教会和他本不该失去的尊严，正如我们会议的进程和我们调查后宣布的判决所证明的。\par

% structure: p0010 source=h2 target=subsection reason=relative-heading-rank; base=h2

\subsection{四、\Person{希拉里}对待\Person{普洛耶克图斯}的方式不能为他增光。}

当此事如此结束时，接下来处理的是我们弟兄和同僚主教\Person{普洛耶克图斯}的申诉。他通过一封含泪可悲的信件向我们陈情，关于在他之上另立主教之事。我们还收到他的同城公民的一封信，附有大量个人签名，其中充满了对\Person{希拉里}最令人不快的指控：即他们的主教\Person{普洛耶克图斯}不允许生病，但在他不知情的情况下，他的主教职务已被转给他人，而继承人被\Person{希拉里}引入占据，仿佛填补空缺，而原主仍然在世。我们想听听弟兄们对此有何看法。虽然我们不应怀疑你们的感受，但当你们想象一位弟兄躺在病床上，不仅受身体虚弱折磨，更受另一种痛苦煎熬时，还有什么希望可留给他？当他人的替代被设立，而他在绝望中看到自己的主教职务被剥夺时，还有什么生机？\Person{希拉里}有一颗温柔的心，这从他相信兄弟死亡的迟缓只是对他野心计划的妨碍一事中得到了清晰证明。因为，就他而言，他熄灭了兄弟的光明；他通过另立他人来剥夺兄弟的生命，从而引起痛苦，阻碍他的康复。假若他的兄弟离世是短暂的，但按常人之道，\Person{希拉里}为何在别人的教省中为自己寻求什么？又为何声称拥有\Person{帕特罗克鲁斯}之前的任何前任都不曾拥有的东西？而那个由宗座暂时授予\Person{帕特罗克鲁斯}的地位，后来也被更明智的决定撤销了。至少应该等待公民的意愿和人民的见证；应该征询尊贵之人的意见和圣职人员的选举——那些熟悉教父法规的人在任命司铎时常遵守这些惯例：以便宗徒权威的规则在一切事上得以保持，该规则要求：将要担任教会司铎的人，不仅要有信友的见证，还要有\Quote{那些外人\BibleRef{弟前 3:7}}的见证，并且在和平与天主喜悦的和谐中，在众人完全赞同下，一位和平导师被祝圣时，不留下任何冒犯的借口。\par

% structure: p0012 source=h2 target=subsection reason=relative-heading-rank; base=h2

\subsection{五、\Person{希拉里}的行为全程应受谴责，我们已恢复\Person{普洛耶克图斯}的职务。}

但\Person{希拉里}突然来到他们中间，又同样突然离去，据我们查明，他快速完成了许多旅程，穿越遥远省份，其匆忙程度表明他渴望得到信使的迅捷名声，而非司铎的稳重。因为公民们在给我们的信中写道：\Quote{他离去之前，我们还不知道他已来到。}这不是返回，而是逃跑；不是施行牧人的有益照顾，而是盗贼和强盗的暴力，如主所说：\Quote{那不从门进入羊栈，而从别处爬进去的，便是贼，是强盗。}因此，\Person{希拉里}所急迫的与其说是祝圣一位主教，不如说是杀害那位病者，并通过不当的祝圣误导他所立在病人之上的人。然而，我们已做了——在天主审判下——我们相信你们会赞同的事：在与所有弟兄商议后，我们裁定，那位被不当祝圣的应被撤职，\Person{普洛耶克图斯}主教应继续他的主教职务；此外规定，当任何弟兄在任何教省去世时，已被承认为该教省的教省总主教应为自己要求继承人的祝圣权。\par

我们看到，这两件事已解决，尽管其中还有许多其他点似乎违反了教会的原则，应受到公正的谴责和审判。但我们不能再深究下去，因为我们被别的事务叫走，必须与你们——圣洁的弟兄们——认真商议。\par

% structure: p0015 source=h2 target=subsection reason=relative-heading-rank; base=h2

\subsection{六、必须制止\Person{希拉利}使用武装暴力的做法。}

据我们所知，一队兵士跟随司铎穿越各省，协助他依靠武力支持，暴烈地入侵那些失去自己司铎的教堂。一些对所管辖的城市完全陌生的人被拖到法庭受祝圣。因为在平安时，人们寻求的是知名且被认可的人；同样，当推出一个陌生人时，必须用暴力来确立他。弟兄们，我以天主的圣名恳求、哀求并恳请你们阻止此类事情，从你们各省中消除一切不和的缘由。无论如何，我们在天主面前为自己辩白，恳求你们不要允许此事再继续下去。应在和平与宁静中寻求那些将成为司铎的人。需要神职人员的同意、受尊敬者的见证、各等级及平信徒的认可。将要治理众人的人，应由众人选出。如前所述，每位教省主教应亲自掌控本省内发生的祝圣事宜，并与在司铎年资中居首的人协同行动——这特权是经由我们恢复给他的。任何人不得篡夺他人的权利。每个人都应守在自己的界限和范围内，并应明白他不能将自己拥有的特权转给他人。但如果有人不顾宗徒的禁止，过分注重个人恩情，愿意放弃自己的优先权，以为可以将自己的权利转给他人，那么祝圣的权力不应归于他所让与的人，而应归于在本省主教年资中居首的人。祝圣不应草率进行，而应在适当的日子举行；并应周知，凡不是在星期六晚间（即星期日黎明前夕）或主日当日受祝圣的，其身份不能确定。因为我们的祖先认为，唯独主复活的日子才配享有此荣誉，作为那些将成司铎者被献给天主的时机。\par

% structure: p0017 source=h2 target=subsection reason=relative-heading-rank; base=h2

\subsection{七、\Person{希拉利}不仅被剥夺其篡夺的管辖权，也被剥夺其合法权利，并被限制于他自己单一的主教区。}

让每个省满足于自己的教务会议，不得允许\Person{希拉利}另外召集教务会议，并以其干预扰乱主之司铎的判决。让他知道，他不仅被剥夺了别人的权利，而且被剥夺了他曾经非法占有的对\Place{维恩}省的权力。因为弟兄们，恢复古制是公平的，因为那个曾为自己索取他不负责的省的祝圣权的人，在目前的事例中也暴露出类似的行径——正如他屡次因鲁莽无礼的言语而招致定罪判决一样，现在他应按我们的命令，遵循宗座的仁慈，仅被保留于他自己城市的司铎职位。因此，他不得出席任何祝圣仪式：他不得祝圣，因为他自知罪有应得，当被要求为自己的行为答辩时，他竟以可耻的逃跑来逃脱，并使自己脱离了宗座的共融——他不配分享这共融。我们相信这是出于天主的眷顾，天主使他出乎意料地来到我们的法庭，并在调查进行中让他潜逃，以免他成为我们共融中的一员。\par

% structure: p0019 source=h2 target=subsection reason=relative-heading-rank; base=h2

\subsection{八、绝罚只应施加于犯有大罪之人，且即使如此也不应草率施行。}

任何基督徒不应轻易被剥夺共融，也不应出于愤怒的司铎的意愿而施行，而应由法官的内心在某种程度上不情愿且遗憾地执行，以惩罚大罪。因为我们发现，有人因琐碎的行为和言语而被割断了共融的恩宠，基督的宝血为之倾流的灵魂暴露于魔鬼的攻击之下，被如此残忍的惩罚所伤害、解除武装，可以说，被剥夺了一切防御，成为魔鬼的易得猎物。当然，如果出现了这样一种案件，根据所犯罪行的性质，致使一个人丧失共融是合宜的，那么只有涉及指控的人才应受惩罚；而未显示参与犯罪的人不应分受惩罚。然而，一个惯于因司铎被定罪而欢欣的人，对平信徒也表现出同样态度，又有什么奇怪呢？\par

% structure: p0021 source=h2 target=subsection reason=relative-heading-rank; base=h2

\subsection{九、指派\Person{良修}取代\Person{希拉利}的职位。}

因此，因为我们的愿望与此似乎大相径庭（因为我们切望所有教会的安定和司铎们的和谐得以维持），我们劝勉你们在爱德的联系中合一，既恳求，也本着我们的情感劝诫你们，为了你们的平安和尊严，遵守我们因天主和最蒙福的宗徒伯多禄的默感，在甄别并检验了所有争议事项之后所颁定的一切，确信我们以这种方式所决定的事，与其说是为我们自己的利益，不如说是为你们的利益。因为我们并未将你们各教省的任命权掌握在我们自己手中，正如\Person{依拉利}或许以其惯常的虚伪所言，以误导你们的心思，圣洁的弟兄们：但我们在焦虑中为你们要求，不得再有革新，并且今后也不应给篡权者任何侵犯你们特权机会。因为我们承认，只有在你们中间宗座的勤勉保持无缺，且我们在维护宗徒纪律时，不因肆无忌惮的侵犯而任由属于你们职位的东西丧失，这才只会有益于我们的声誉。并且既然年资常应受到尊重，我们愿意我们的弟兄和同为主教的\Person{莱昂提乌斯}，一位在你们中间备受称许的司铎，被提升到这一尊位，如果你们同意，没有他的同意，你们不再召集进一步的公会议，圣洁的弟兄们，并且他应按其年龄和令誉受到你们所有人的尊敬，而总主教们则在其自己的尊严和权利中得到保障。因为这是公平的，并且似乎不会对任何弟兄造成损害，如果那些在司铎年资上居首的，如其年龄所当得，在各自教省中受到其余司铎的尊重。愿天主保佑你们平安，亲爱的弟兄们。\par

\SourceRecord{Translated by Charles Lett Feltoe.；Nicene and Post-Nicene Fathers, Second Series；Vol. 12.；Edited by Philip Schaff and Henry Wace.；Buffalo, NY: Christian Literature Publishing Co.,；1895.；<http://www.newadvent.org/fathers/3604010.htm>.}

% structure: p0001 source=h1 target=section reason=page-title

\section{第十二封书信}

\Person{良}，罗马城主教，致候在\Place{非洲}的\Place{毛里塔尼亚凯撒利亚}的所有主教，在主内问安。\par

% structure: p0003 source=h2 target=subsection reason=relative-heading-rank; base=h2

\subsection{I. 省内对主教的无序任命应受谴责。}

由于那些来访者的频繁报告，曾提及你们中间在司铎祝圣方面的一些不法行为；宗教的责任要求我们，按照我们因天主命令而赐予整个教会的关怀，努力查明事实真相：因此，我们将此调查委托给我们那位动身离开的兄弟和同职司铎\Person{波腾提乌斯}；他根据我们通过他写给你们的信，将调查那些据说选举有误的主教的事实，并将一切忠实地报告给我们。因此，由于同一\Person{波腾提乌斯}已最充分地向我披露一切，并以其真实的叙述使我们清楚了解到，在（毛里塔尼亚）凯撒利亚省的某些地区，基督的一些会众是置于何种何种治理者之下，我们发现有必要通过将这封信也寄给你们各位爱者，来抒发我们心中为主羊群之危险而有的忧伤；因为我们对下述情况感到惊讶：在动乱之时，阴谋家的跋扈或民众的骚乱竟对你们有如此大的影响，以致将教会的首席牧职与治理交给了最不配的人，即那些远离司铎标准的人。这并非为人民利益筹谋，而是加害于他们；也非施行纪律，而是增加分歧。因为治理者的正直是属下之保障；那里有完全的服从，那里就有纯正的教义形式。但是，一个或由暴乱产生、或由阴谋攫取的任命，即便在道德或实践中并无过错，仅凭其起始的榜样就已造成危害；开端不良之事，难以导向善果。\par

% structure: p0005 source=h2 target=subsection reason=relative-heading-rank; base=h2

\subsection{II. 绝不可仓促祝圣主教。}

然而，如果在教会的每一等级中都需运用极大的远见与知识，以免在主房之中有任何无秩序或不当之处；那么，在选举那位于所有等级之上的人时，我们岂不更应谨慎避免错误？因为主整个家庭的平安与秩序将被动摇，如果在身体中所要求的不能在头中找到。那么，蒙福的宗徒保禄借天主之神所发出的训诫在哪里？他在\Person{弟茂德}身上教导了整个基督司铎团体，并对我们每个人说：\Quote{\BibleQuote{不可轻易给人覆手，也不要分担别人的罪过？}\BibleRef{弟前5:22}}？什么叫做\Quote{轻易给人覆手}？岂非在适当年龄之前、在未及考验之时、在未经服从而堪当之前、在尚未经纪律试炼之前，就授予司铎尊位吗？什么叫做\Quote{分担别人的罪过}？岂非祝圣者变成那本不应由他祝圣的人那样吗？因为正如一个人若在选司铎时持守正当判断，便为自己积蓄善功之果实；同样，若将不配的司铎纳入同职之列，便给自己招致严重损失。因此，我们不可在任何人的事上忽视那在一般法规中所规定的内容；违反天主法律训令的晋升，也不得视为合法。\par

% structure: p0007 source=h2 target=subsection reason=relative-heading-rank; base=h2

\subsection{III. 宗徒关于圣职人员婚姻的训诫，以基督与教会的婚姻为基础，并作为其预像。}

因为宗徒说，在选举的其他规则中，应被祝圣为主教的人须是众所周知曾为或现为\Quote{一个妻子的丈夫}；这一命令向来被如此神圣地遵守，以致在当选司铎的妻子身上也理解为必须遵守同一条件：即使他本人未有过其他妻子，他的妻子在与他缔结婚姻之前，也不应曾嫁给另一个男人。那么，谁敢容许这一伤害如此伟大圣事的行为发生？因为这一伟大而可敬的奥秘并非没有天主法律之规定的支持，其中清楚规定：司铎应娶童贞女，并成为司铎之妻者\BibleRef{肋21:13}不应认识另一个丈夫。因为即使在那些司铎身上，也预表了基督与祂教会的属神婚姻：如此，既然\Quote{男人是女人的头\BibleRef{弗5:23}，}那么圣言的配偶就应学会只认识基督，别无他人；祂只正当地拣选了她，只爱她，只接纳她进入祂的盟约。如果即使在旧约中，司铎中的这种婚姻也被坚持，那么，我们这些处于福音恩宠之下的人，岂不更应遵从宗徒的训令？这样，即使一个人被发现有良好品德，并具备神圣善工，但如果他本人并非一个妻子的丈夫，或者他的妻子并非一个丈夫的妻子，那么他就绝不可升至执事级、司铎尊位，或主教最高等级。\par

% structure: p0009 source=h2 target=subsection reason=relative-heading-rank; base=h2

\subsection{IV. 应避免过早晋升。}

但当宗徒警告说：\BibleQuote{也当先受考验，然后才可任职}\BibleRef{弟前 3:10}时，我们除了认为在这些擢升中不仅应考虑他们婚姻的贞洁，还须考虑其劳苦的功绩，免得牧职托付给那些要么刚受洗、要么突然从世俗事务中转向的人之外，还能作何理解？因为在基督教军队的所有等级中，在晋升时应当考虑一个人是否能管理更大的职责。真福教父们在论及司铎选举时，其可敬的意见正确地认为，那些经过各等级职务逐步晋升，并在其中充分证明自己的人，才适于管理神圣之事；因为在他们每个人身上，其行为的品性见证了他们的生活。如果未经时间帮助、未经历劳苦功绩便获得世俗尊荣是不当的，并且若无正直的证明而寻求职位会被指责，那么神圣职责和天上尊荣的分施应当多么勤谨、多么小心地执行，以免在何事上违反宗徒的与教规的法令，并将主的教会治理托付给那些不谙合法制度、毫无谦卑、不愿从最低等级上升而想从最高开始的人：因为让无经验者优先于有经验者、年轻者优先于年长者、新兵优先于久经沙场者，是极不公平且颠倒的。确实，如宗徒所解释，在大户人家中必须有各种器皿，有金的和银的，也有木的和瓦的；但它们的用途随材质的优劣而异，珍贵与廉价之物的使用并不相同。因为如果瓦器被优先于金器，或木器被优先于银器，一切就都会混乱。正如木器和瓦器是那些至今没有显著德行的人的象征；同样，在金器和银器中所代表的无疑是那些经过了漫长经验之火和持久劳苦的熔炉，配得上成为试炼过的黄金和纯银的人。如果这样的人得不到他们虔敬的赏报，教会的全部纪律就会松弛，一切秩序都会混乱，因为那些未经任何服务的人，因选举团的不当选择而获得不应得的擢升。\par

% structure: p0011 source=h2 target=subsection reason=relative-heading-rank; base=h2

\subsection{V. 他区分了被擢升为主教的平信徒与再婚的圣职人员，宽恕前者而不宽恕后者。}

既然在你们中间，要么是民众的热切愿望，要么是野心者的阴谋产生了如此大的影响，以致我们得知不仅平信徒，甚至再婚者或娶寡妇者也被擢升到牧职，那么难道没有最清楚的理由要求那些发生此类事情的教会以比通常更严厉的审判来洁净，不仅统治者本身，而且那些祝圣他们的人也应受到相应的惩罚？但一方面有仁慈的温良，另一方面有正义的严格。因为\Quote{上主的一切道路都是仁慈和真理}，我们基于对宗座之忠信，不得不如此调和我们的判断，以在权衡中衡量人们的过犯（当然它们并非同一尺度），并视某些在某种程度上可宽恕，而另一些则完全应受压制。因为那些要么进入第二次婚姻，要么与寡妇结合的人，无论按宗徒权威还是法律权威，都不允许持有司祭职；更何况据报告给我们，那人同时有两个妻子，或那人被妻子离弃后又娶了另一个——假设这些指控在你们判断中被证实。至于其余的人，其晋升只在此意义上受到责备：他们是从平信徒中被选为主教职务，并且在妻子的问题上没有过犯，我们允许他们保留已领受的司祭职，不损害宗座的法令，也不违背真福教父们的规则——其有益的规条是：任何平信徒，无论得到多少支持，不得升至教会的第一、第二或第三等级，除非他通过合法的步骤达到那个位置。因为我们现在允许在某种程度上可宽恕的事，今后若有人做出我们完全禁止的事，则不能不受惩罚：因为对罪的宽恕并不允许犯错，也不可因部分宽恕而重复犯罪却不受罚。\par

% structure: p0013 source=h2 target=subsection reason=relative-heading-rank; base=h2

\subsection{VI. 多纳图斯，一位归正的诺瓦提安派信徒，以及马克西穆斯，一位前多纳图派信徒，被保留在其主教职务上。}

我们希望萨拉西亚的\Person{多纳图斯}——如我们所知，他与他的民众一起从诺瓦提安派归正——管理主的羊群，条件是他记得须将一份他的信仰证书寄给我们，其中不仅谴责诺瓦提安派教义的错误，而且也毫无保留地宣认公教真理。\Person{马克西穆斯}虽然当他还是平信徒时被有罪地祝圣，但如果他现在不再是多纳图派，且已弃绝分裂的邪恶精神，我们不将其从其不正当地获得的主教尊荣上革职，条件是他通过为我们起草一份证书来宣告自己是公教徒。\par

% structure: p0015 source=h2 target=subsection reason=relative-heading-rank; base=h2

\subsection{VII. 阿加鲁斯与提贝里亚努斯（在骚乱中受祝圣）的案件被提交给主教们。}

但关于\Person{阿加鲁斯}与\Person{提贝里亚努斯}，他们的案件与其他从平信徒受祝圣的人不同，在于他们的祝圣据报告伴随着激烈的骚乱和野蛮的扰乱，我们已将整个事项托付给你们的判断，以便依靠你们对案件的调查，我们能知道对它们应作何决定。\par

% structure: p0017 source=h2 target=subsection reason=relative-heading-rank; base=h2

\subsection{VIII. 遭受暴力的少女不应与他人比较。}

那些因蛮族暴力失去贞洁的天主婢女，若不敢与无玷童贞女相比，在其谦卑和羞惭中将更值得称赞。因为虽然每一罪过都源于欲望，而意志可能保持未征服、未被肉体的堕落污染，但若她们哀叹即使在身体上失去了在精神上没有失去的，这对她们损害较小。\par

% structure: p0019 source=h2 target=subsection reason=relative-heading-rank; base=h2

\subsection{IX. 这些训令要无争论地执行。}

所以，如今你们，亲爱的，看到自己通过我们的兄弟和同主教\Person{达味}，他以个人品德和司铎价值为我们所认可，关于我们的兄弟\Person{颇登提乌斯}所叙述的几乎所有要点都得到了充分指示，兄弟们，剩下的就是你们要和谐地接受我们有益的劝勉，不作任何竞争，而是怀着全然的虔诚和热忱齐心行动，服从天主及其宗徒的规诫，决不容忍教会的明智法令受到违反。因为出于某些原因我们现在所放宽的，今后必须由古老规则来守护，以免我们现在以仁慈的宽宥所让步的，日后可能必须以适当的惩罚来追究，尤其对那些在祝圣主教时忽视了圣教父们的法规，并祝圣了本应拒绝的人，我们将以特别直接的力量行事。因此，若任何主教祝圣了不应当作司铎的人，即使他们某种程度上没有损失个人尊严，他们也不再有祝圣权，也绝不参与那圣事——他们忽视天主旨意而不当授予的圣事。\par

% structure: p0021 source=h2 target=subsection reason=relative-heading-rank; base=h2

\subsection{X. 在过小的地方任命主教是不适宜的，必须停止。}

当然，关于司铎尊严的，我们希望与所有教规保持一致遵守，即：主教不得在任何地方、任何村庄或以前未祝圣过的地方被祝圣；因为羊群小、会众少的地方，长老的关怀就足够了，而主教权威应只管辖较大的羊群和更拥挤的城市，以免违反圣教父们受天主启示的法令，将司铎职务指派给村庄、乡村庄园或偏僻人口稀少的乡镇，而使本应赋予更重要职责的尊位，因任职者众多而被轻视。\Person{雷斯提图都斯}主教报告说，在他自己的教区发生了这样的事，他有充分的理由请求，当那些本不应被祝圣的地方的主教自然去世后，这些地方本身应回归原来所属并附属于同一主教管辖。确实，由于祝圣者不谨慎的迁就，通过多余的职务倍增而削弱司铎尊严是无益的。\par

% structure: p0023 source=h2 target=subsection reason=relative-heading-rank; base=h2

\subsection{XI. 被迫失贞的童贞女要与其他童贞女稍有区别对待，但不可拒绝她们领圣体。}

现在关于那些曾发神圣贞洁誓愿（如上文第八章所说）而遭受蛮族暴力，在身体上而非精神上失去无玷纯洁的人，我们认为应当采取如此适度处理：她们既不应被贬低到寡妇的行列，也不应计入圣洁无玷童贞女之数；然而，若她们坚持童贞生活，并在心思和意念中守护贞洁的真实，则不可拒绝她们领圣事，因为她们因意愿未曾放弃、而是被敌人暴力夺走的，被指控或玷污是不公平的。\par

% structure: p0025 source=h2 target=subsection reason=relative-heading-rank; base=h2

\subsection{XII. \Person{路丕奇努斯}的事一部分已处理，一部分转交他们处理。}

关于\Person{路丕奇努斯}主教的事件，我们下令在那里审理，但由于他迫切而频繁的恳求，我们已恢复他与教会共融，理由如下：因他已向我们提出上诉，我们看到在案件悬而未决时他不应被无端暂停共融。此外还有一点：在他之上祝圣一位本不应被祝圣的人显然是鲁莽的，直到\Person{路丕奇努斯}被带到你们面前或被定罪，或至少认罪，有机会接受公正判决，以便按照教会纪律的要求，被祝圣者能接替其空缺位置。\par

% structure: p0027 source=h2 target=subsection reason=relative-heading-rank; base=h2

\subsection{XIII. 所有争端先就地处理，然后提交宗座。}

但是每当出现其他涉及教会状况和司铎和谐的案件时，我们希望你们首先以对主的敬畏自行筛选，并将所有已解决或待解决事宜的完整报告送交我们，以便那些根据教会惯例适当合理地决定的事，也能得到我们的确认批准。日期：8月10日。\par

\SourceRecord{Translated by Charles Lett Feltoe.；Nicene and Post-Nicene Fathers, Second Series；Vol. 12.；Edited by Philip Schaff and Henry Wace.；Buffalo, NY: Christian Literature Publishing Co.,；1895.；<http://www.newadvent.org/fathers/3604012.htm>.}

% structure: p0001 source=h1 target=section reason=page-title

\section{第十四封信}

\Signature{罗马城主教良，致德撒洛尼主教阿纳斯塔西。}\par

% structure: p0003 source=h2 target=subsection reason=relative-heading-rank; base=h2

\subsection{一、前言。}

如果兄弟你以真正的理性，领悟了因真福宗徒伯多禄的权威而托付给你的一切，以及因我们的恩宠而交托给你的一切，并且能公平地衡量它，我们就会因你热诚履行所托付的责任而大为欢喜。\par

% structure: p0005 source=h2 target=subsection reason=relative-heading-rank; base=h2

\subsection{二、阿纳斯塔西被指责超越其\OriginalTerm{代牧权}{vicariate}的界限，尤其是在他对阿提库的粗暴与不当处置上。}

鉴于正如我的前任对待你们的前任一样，我也效法他们的榜样，将我的权威委托给你，亲爱的：好使你效仿我们的温和，能在我们因神圣制度而主要对众教会所负的照料中协助我们，并在某种程度上代替我们亲身巡视那些远离我们的省份：因为你定期并适时地视察，就易于判断哪些事、在何种情况下，你可以凭自己的影响力解决，或留待我们决断。因为既然你可以暂缓更重大的事和更棘手的案子，等待我们的意见，那就没有理由也没有必要让你多此一举地决定超越你能力范围的事。因为我们有许多书面的警告，屡次训诫你在一切行动中要节制：用爱德的劝勉，促使托付给你的基督的众教会健康地服从。因为，虽然通常在疏忽或懈怠的弟兄中间，有些事需要强硬手段来纠正；但矫正应当如此应用，以始终保全爱德。因此，真福宗徒保禄在教导弟茂德治理教会时说：\BibleQuote{不可严责老年人，却要劝他如父；劝少年人如弟兄；劝老年妇人如母亲；劝青年妇人如姊妹，总要在纯洁中。}\BibleRef{弟前 5:1} 如果按宗徒的诫命，这种温和应当给予任何下级成员，那么对我们的弟兄和同为主教的，更应无冒犯地施予？为的是，虽然有时在司铎身上发生需要斥责的事，但仁慈对受纠正者的效力可能大于严厉；劝勉大于困扰；爱德大于权力。但那些\BibleQuote{寻求自己的事，不寻求耶稣基督的事}\BibleRef{斐 2:21}的人，很容易偏离这条律法，他们喜欢统治属下胜过关心他们的利益，因地位的骄傲而膨胀，因此原为维护和谐而设的，却成了祸害。我们不得不这样说，这使我们感到不小的悲伤。因为我发现自己某种程度上被牵扯到责备中，当我发现你如此无节制地背离了传授给你的规则。如果你不在乎自己的声誉，你至少应当顾及我的名声：免得只有你自己心中所想的，被当作得到我们的赞同。兄弟，但愿你仔细阅读我们的信件，并查阅\OriginalTerm{宗座}{Apostolic See}历任持有者寄给你前任的所有信函，你就会发现，在我们得知你所擅自行事的那件事上，我或我的前任曾有训令。\par

因为我们的兄弟，旧埃皮鲁斯的\OriginalTerm{都主教}{metropolitan}阿提库，与他省内的主教们一起来到我们这里，含泪申诉他所遭受的不应有的侮辱，在你自己的执事们面前——他们未对这些悲惨的申诉提出任何异议，表明我们所受到的印象并非缺乏真实。我们也从你带来的那些执事所持的信中读到，阿提库兄弟曾来到德撒洛尼，并且他还在一份书面声明中封印了他的同意，以致我们不能不理解为：他是出于自己的意愿和自由的热忱而来，并起草了服从的承诺声明，尽管在提及这份声明时已暴露出侮辱的迹象。因为对于已经通过自愿前来的忠诚而证明他服从的人，本不需要用书面约束他。因此，你信中的这些话证实了上述者的哀诉，通过他的直言不讳，沉默的事被揭露，即\OriginalTerm{伊利里亚}{Illyricum}的\OriginalTerm{总督区}{Præfecture}被接触，世上权贵中最尊贵的官员被调动来暴露一位无辜的主教：以致派遣了一队人来执行这可怕的行为，他们要征用所有公务员来执行命令，并从教会神圣的圣所中，将一位未受控告、至多受虚假控告的司铎拖出，在考虑他严重的病弱或残酷的冬季天气时，未给予他任何休战：但他被迫穿越无路的积雪，踏上充满艰辛和危险的旅程。这是一项如此辛劳和危险的任务，据说一些陪伴主教的人已经死去。\par

我亲爱的弟兄，我深感震惊，且极其悲痛，你竟对一位你除了得知他应召时以生病为由推脱并请求谅解之外，并无更多了解的人，如此野蛮而猛烈地动怒；尤其是，即使他本当受此对待，你也应等我答复你咨询的信函。但依我所见，你过高估计了我的习惯，并且最真切地预见到我将为维护司铎间的和睦而作出何等公正的答复：因此你急忙毫不掩饰地执行你的行动，以免在收到我们指示另一种做法的容忍信函后，你再无自由去做那已做的事。或许某种罪行传到了你耳中，身为都主教，某种新指控的重担使你感到紧迫？但实际情况并非如此，你自己因没有对他提出任何指控而确证了这一点。然而，即使他犯有某种严重且不可容忍的过失，你也应等待我们的意见：以便在知道我们的意愿之前，你不自行作出任何决定。因为我们任命你为我们的代理，亲爱的，是建立在你要分担我们的责任，而非将全权揽于自身的基础上。因此，你所施予的虔敬关怀极大地愉悦我们，而你的不当行为也深深刺痛我们。在许多事例之后，我们必须更加深谋远虑，更加谨慎防范：以使通过爱德与平安的精神，一切冒犯之事得以从主的众教会中被移除，这些教会我们已托付给你：你的主教职务在各省份中保持卓越，但所有你僭越的过分行为均被剪除。\par

% structure: p0009 source=h2 target=subsection reason=relative-heading-rank; base=h2

\subsection{三、在\Person{阿纳斯塔西乌斯}代牧权下，都主教的权利应予以遵守。}

因此，按照圣教父们所制定、并由圣神颁定且为普世尊敬所尊崇的教规，我们下令：在你的关怀（弟兄）通过我们的委托所延伸的每个省份中，都主教应保持其从古老时代传承下来的地位权利不受侵犯：但前提是他们不得因任何疏忽或傲慢而偏离现行规定。\par

% structure: p0011 source=h2 target=subsection reason=relative-heading-rank; base=h2

\subsection{四、论主教被选的不合格条件。}

在首领已逝的城市中，应依此形式填补其空缺：被祝圣者，即使其善行未经见证，也不应是平信徒、新奉教者，亦非再婚者，或虽仅有一妻但曾娶寡妇者。因为司铎的拣选如此至关重要，以至于在教会其他成员中无可指责之事，在他们身上却被视为非法。\par

% structure: p0013 source=h2 target=subsection reason=relative-heading-rank; base=h2

\subsection{五、连副执事也应守节。}

因为，尽管那些不在圣职行列中的人可以自由享受婚姻伴侣和生育子女的乐趣，但为彰显完全贞洁的纯净，连副执事也不被允许有肉体的婚姻：好叫\Quote{那些有妻子的，要像没有一样}，而那些没有的，则保持独身。但如果在这等次（从元首算起是第四级）中此事值得遵守，那么在首级、第二级或第三级中岂不更应遵守，以免任何被发现尚未克制其对妻子情欲的人，被认为适合执事的职务、或司铎的尊贵地位、或主教的卓越职位。\par

% structure: p0015 source=h2 target=subsection reason=relative-heading-rank; base=h2

\subsection{六、主教选举必须根据圣职人员和民众的意愿。}

因此，当着手选举首席司铎时，应优先选择圣职人员和民众一致同意所要求的人；但如果票数恰好在两人之间平分，则都主教的裁决应倾向于获更多票数和功绩支持的那一位：只应无人被祝圣以违背该地方的明确意愿：以免城市因未获准享有其愿望之人而轻视或憎恨其主教，并可悲地背离宗教信仰。\par

% structure: p0017 source=h2 target=subsection reason=relative-heading-rank; base=h2

\subsection{七、都主教应咨询他们的代牧：都主教的选举方式已规定。}

然而，都主教应将关于将被祝圣为主教的人选，以及圣职人员和民众的同意情况，报告给你（弟兄）；并应告知你该省份的意愿：以使祝圣典礼的适当举行也借你的权威得以加强。至于正当的拣选，你的职责是不得造成任何延迟或阻碍，以免主的羊群过长地没有牧人的照顾。\par

再者，当一个都主教去世，必须另选一位接替其位置时，该省份的主教们应集合在都主教城市中：以便在筛选了全体圣职人员和全体市民的意愿后，从同一教会的司铎或执事中选出最佳的人选，并且他的名字应由该省份的司铎们通知你，这些司铎将在确认你同意他们的选择后，执行其支持者的意愿。因为我们虽不希望正当的选举受任何延迟阻碍，但也不允许任何事情在未知会你的情况下被轻率地执行。\par

% structure: p0020 source=h2 target=subsection reason=relative-heading-rank; base=h2

\subsection{八、主教每年应召开两次教省会议。}

关于主教会议，我们给的指示不外乎圣教父们为教会的健康所制定的：即每年举行两次会议，在会议中应对所有通常在教会不同等级之间产生的投诉作出判决。但倘若在统治者们中间（天主禁止）出现关于某一大罪的争端，无法由教省审判解决，则都主教应小心地将整个事情的性质报告给你（弟兄），并且如果双方来到你面前后，事情甚至借你的判决仍不能平息，那么无论它是什么，都应转交我们的司法权。\par

% structure: p0022 source=h2 target=subsection reason=relative-heading-rank; base=h2

\subsection{九、禁止从一教区调任至另一教区。}

如果任何主教轻视自己城市的微不足道，而图谋管辖一个人口更多的地方，并以任何方式将自己转移到一个更大的羊群，他首先应被赶出他篡夺的席位，并被剥夺他自己的（席位）：这样，他既不得管理他贪婪渴望的那些人，也不得管理他傲慢轻蔑的那些人。因此，让每人满足于自己的界限，不要寻求超越其现有职位的限制。\par

% structure: p0024 source=h2 target=subsection reason=relative-heading-rank; base=h2

\subsection{X．主教不应诱拐或接收另一个教区的神职人员。}

任何（主教）不得违背其本主教之意愿而接受或邀请来自另一教区的神职人员：只有当给予者和接受者通过友好协议彼此同意时才可。因为一个人若胆敢诱拐或扣留另一兄弟教会中大有用途或极具价值的东西，便是犯下严重伤害。因此，如果此类事情发生在教省内，教省总主教应强迫离弃的神职人员返回其教会；但如果他逃得更远，则应由你那有权威的命令将他召回：以免留下任何贪利或阴谋的机会。\par

% structure: p0026 source=h2 target=subsection reason=relative-heading-rank; base=h2

\subsection{XI．当宗座代牧要求召集主教会议时，每教省派两位即足。}

在召集主教到你面前时，我们希望你要显示极大的宽容：免得在表现出很勤勉的伪装下，你似乎因兄弟的损伤而得意。因此，如果发生任何更大的案件，有理由且必要召开弟兄会议，那么，兄弟，每教省派两位主教参加即可，这些主教由教省总主教认为合适而派遣，条件是应召者不得自规定时间起被扣留超过十五天。\par

% structure: p0028 source=h2 target=subsection reason=relative-heading-rank; base=h2

\subsection{XII．若宗座代牧与主教们意见不同，应咨询罗马主教。论教会中权威的从属关系。}

但如果在你认为需要与弟兄们讨论并解决的事上，他们的意见与你自己的意愿不同，应将一切呈报给我们，附上经证实的会议记录，以便消除所有歧义，并决定天主所喜悦的事。因为我们所有的愿望和劳苦都指向此目的，即凡有助于我们和谐合一与纪律维护之事，不致因任何纷争而受损，也不因任何懈怠而被忽视。因此，亲爱的弟兄，你和那些被你过分行为所冒犯的我们的弟兄（虽然抱怨的原因不尽相同），我们劝戒并警告你们，不要以任何争论扰乱那些已经合法颁定且明智安排的事。让没有人\Quote{寻求自己的事，而是他人的事，}正如宗徒所说：\Quote{你们各人应求他人的喜悦，为建树而使他人的益处。}因为我们合一的粘合无法坚固，除非我们被爱的纽带捆绑成不可分割的整体：因为\Quote{正如在一个身体上我们有许多肢体，但所有肢体没有相同的职务；同样，我们众人在基督内是一个身体，而各自互为肢体。}整个身体的连接使所有肢体同样健康，同样美丽；而这一连接确实要求整个身体的同心一致，但它尤其要求司铎之间的和谐。虽然他们拥有共同的尊严，但并非相同的等级；因为即使在蒙福的宗徒中，尽管他们尊贵地位的相似，仍存在某种权力的区别，并且虽然他们全都被平等地拣选，但其中一位被赋予领导其余人的职责。从这个模式也产生了主教之间的区别，并且通过一项重要的规定，确保了并非每人都应为自己要求一切：而是每个教省应有一位其意见在兄弟中居首位的人；此外，那些被任命在较大城市的人应承担更全面的责任，通过他们，普世教会的关怀应汇聚于伯多禄的唯一之座，且任何地方都不应与它的头分离。因此，那知道自己被安置在某些人之上的人，不要因有人被安置在他之上而生气，而要亲自施行他向别人要求的服从：正如他不愿承担一沉重行囊，他也不敢将难堪的重担放在别人的肩上。因为我们是谦逊而良善的师傅的门徒，他说：\BibleQuote{向我学习，因为我是良善心谦的，你们将找到灵魂的安息。因为我的轭是易于负的，我的担子是轻的。} \BibleRef{玛 11:29}而我们又如何体验到这一点，除非我们也想起同一主所说的话：\Quote{你们中更大的，要成为你们的仆役。但自高自大的，必被贬抑；而自谦自卑的，必被高举。}\par

\SourceRecord{Translated by Charles Lett Feltoe.；Nicene and Post-Nicene Fathers, Second Series；Vol. 12.；Edited by Philip Schaff and Henry Wace.；Buffalo, NY: Christian Literature Publishing Co.,；1895.；<http://www.newadvent.org/fathers/3604014.htm>.}

% structure: p0001 source=h1 target=section reason=page-title

\section{第十五书}

\Emph{致\Person{图里比乌斯}，\Place{阿斯图里亚}主教，论\OriginalTerm{普里西利安派}{Priscillianists}之谬误。}\par

% structure: p0003 source=h2 target=subsection reason=relative-heading-rank; base=h2

\subsection{一、引言}

\Person{良}主教致\Person{图里比乌斯}主教，问安。\par

弟兄，你那值得称赞的对公教信仰真理的热忱，以及你在牧职中为主羊群所付出的勤勉奉献，已由你托付我们、由你的执事转交的信函所证明；你在信中留心向我们陈明那从古老瘟疫的残余中在你地区爆发的疾病的性质。因为你的信函措辞、详细陈述以及你那小册子的文本清楚说明，普里西利安派的污秽水坑仍在你们中间活跃。没有一种污秽不是从各种异端者的观念流入这教义的，因为他们从世俗意见的泥淖中搜罗了杂色残渣，为自己调制成一种唯独他们能全盘吞下、而他人仅尝点滴的混合物。\par

事实上，若将普里西利安之前兴起的一切异端仔细研究，几乎找不到一种谬误不曾沾染这不虔敬；因它不满足于接受那些假基督之名离开福音者的谎言，还投身于外教的黑暗，藉着魔术技艺的亵渎秘密和占星家的虚谎，将宗教信仰和道德行为置于恶魔的能力和星辰的影响之上。若此可被相信和教导，则德行之功、过犯之罚将不复存在，不仅人律、连神律的一切诫命都将瓦解；因若致命必然性将心灵冲动驱向任何一方，而人所行的一切非藉人而是藉星辰，则善恶行为便不可能有标准。这疯狂属于那将整个人体按十二星座划分的怪诞分类，使每部分受不同势力支配；那按自己肖像造人的受造者，被星辰支配如同其肢体彼此相连一样。因此，我们父辈在这可憎异端兴起之时，迅速在全界追讨它，使这亵渎之误处处被逐出教会；连世间统治者都如此厌恶这渎神的愚蠢，以国法之剑诛灭了其创始人及多数门徒。因他们看见，若允许这类人持此信条在任何地方生存，一切对正直行为的渴求将消失，一切婚姻纽带将解除，神律与人律同时被摧毁。这种严苛处理长久以来有助于教会温和的法律——此法律虽依靠司铎的判断，避免血腥的报复，但有时亦藉基督徒君主的严厉法令得到辅助，当那些畏惧肉体惩罚的人求助于纯粹灵性纠正之时。但由于许多省份被敌人的入侵所占据，法律的执行也因这些风暴般的战争而中断。而当天主司铎间的交往变得困难、聚会稀少时，秘密的背信便在普遍混乱中自由行动，并被那些本应遏止它的灾祸本身激动，去搅乱许多人的心灵。但在这地区，如亲爱的弟兄你所指出的，甚至某些司铎的心灵也沾染了这致命疾病，有多少民族、多少民族未受此瘟疫污染？那些本应被视为真理之必要抑制者和捍卫者的人，却成了使天主的福音受制于普里西利安教义的人；以致神圣典籍的忠实被扭曲为亵渎之意，在先知和宗徒的名字下所宣讲的，不是圣神所教导的，而是魔鬼的仆人所插入的。因此，如你所爱者，以你力所能及的一切信德勤勉，已将这些已被谴责的意见细分为16条，我们也再次对其加以严格审查；唯恐这些亵渎之论中有任何被视作可容忍或可疑。\par

% structure: p0007 source=h2 target=subsection reason=relative-heading-rank; base=h2

\subsection{二、(一) 驳斥普里西利安派对天主圣三的否认}

因此，在第一项下，显示了他们对天主圣三所持的不圣洁观点：他们声称圣父、圣子及圣神的位格是同一且相同的，仿佛同一天主有时称父、有时称子、有时称圣神；仿佛生者、所生者及自两者发出者并非各异，而应理解为不可分的单一，虽以三名言之，却不由三位格组成。此类亵渎借自\Person{撒伯里乌}，其追随者亦被正当地称为\OriginalTerm{父受苦论者}{Patripassians}，因为若子与父同一，则子的十字架便是父的受苦（\OriginalTerm{父受苦}{patris-passio}），而且父自取了子以奴仆形象所取的一切，并服从于父。这无疑违反公教信仰，公教信仰承认天主性圣三为\OriginalTerm{同一本质}{\GreekText{ὁμοούσιον}}，并相信父、子、圣神不可分而无混淆，无时间而永恒，无差别而平等；因为不是同一位格，而是同一本质充满圣三中的合一。\par

% structure: p0009 source=h2 target=subsection reason=relative-heading-rank; base=h2

\subsection{三、(二) 驳斥他们关于从天主发出的某种德能的幻想}

在第二项下，显示了他们关于从天主发出的某些德能——祂开始拥有它们，且它们在本性上后于天主自身——的愚蠢空想。于此他们又支持亚流派的错误，他们说父先于子，因为有一段时间祂没有子，并在生子时成为父。但正如天主教会憎恶亚流派，也憎恶这些认为同一本质者曾对天主有所欠缺的人。因为说祂进步与说祂可变化的同样邪恶。因为增加意味着变化，正如减少也意味着变化一样。\par

% structure: p0011 source=h2 target=subsection reason=relative-heading-rank; base=h2

\subsection{四 (3) 他们关于\Quote{独生子}称号的论述被驳斥}

第三点同样涉及这些人的亵渎断言：他们称天主之子之所以被称为\Quote{独生子}，是因为唯有祂是由童贞女所生。的确，他们若非饮下了\Person{撒摩撒特人保禄}和\Person{佛提努斯}的毒液，绝不敢这样说；此二人声称我们的主耶稣基督直到由童贞玛利亚诞生之前并不存在。但如果他们希望以自己的主张表达别的含义，而不将基督的起源定于祂母亲的子宫，那么他们必然要断言：天主子不止一位，其他也有从至高之父受生的，而这位是由女子所生，因此被称为独生子，因为天主的其他儿子中没有一个经历过这样出生的境况。因此，无论他们往何处逃遁，都会陷入大不敬的深渊，只要他们要么坚持主基督起源于祂的母亲，要么不相信祂是天主圣父的独生子：因为那本是天主的却由母亲所生，而除了圣言之外，没有谁由圣父所生。\par

% structure: p0013 source=h2 target=subsection reason=relative-heading-rank; base=h2

\subsection{五 (4) 他们于主诞及主日守斋被否定}

第四点涉及以下事实：主诞，即祂取了真正人性的场合，天主教会对此十分重视，因为\Quote{圣言成了血肉，寄居在我们中间\BibleRef{若 1:14}，}但这些人并未真正尊敬它，尽管他们装出尊敬的样子，因为他们那一天守斋，如同也在主日——基督复活之日——守斋一样。毫无疑问，他们这样做，是因为他们不相信主基督是以真正的人性诞生的，反而坚持认为通过某种幻觉，呈现了并非真实的假象，追随\Person{策尔多}和\Person{玛尔基翁}的观点，并与他们的亲属摩尼教徒完全一致。因为正如我们的审查已揭示并使他们明白的，他们在悲哀的斋戒中拖长了主日，而这一日因着我们救主的复活而被祝圣：他们将这种斋戒，按照解释，奉献给对太阳的崇拜：以至于他们处处与我们的信仰合一不相和谐，那一日我们以喜乐度过，他们却以自苦度过。因此，这些基督十字架与复活的敌人接受一种与他们所选教义一致的意见是合适的。\par

% structure: p0015 source=h2 target=subsection reason=relative-heading-rank; base=h2

\subsection{六 (5) 他们关于灵魂是天主性一部分的观点被驳斥}

第五点涉及他们断言人的灵魂是天主性的一部分，并且我们人类境况的本性与造物主的本性并无不同。这一不敬的观点源于某些哲学家的见解，以及摩尼教徒，而公教信仰谴责它：因为公教信仰知道，没有受造物如此崇高卓越，以致其本性本身即是天主。因为属于祂自身的一部分就是祂自身，除了圣子和圣神之外，别无其他。除了这至高的天主圣三的同体、永恒、不变的天主性之外，在一切受造物中，没有一样在其起源上不是从无中创造出来的。此外，任何超越同类受造物的事物并不因此就是天主，而如果某物伟大奇妙，它也不等同于\Quote{惟独祂行伟大奇事}。没有\Emph{人}是真理、智慧、正义；但许多人分享真理、智慧、正义。然而，唯有天主免除任何分享：任何在某种程度上可恰当地用来述说祂的，都不是属性，而是祂的本质本身。因为在不变者中，无所添加，无所减损：因为\Quote{自有}是祂常有的特性，这就是永恒。因此，祂恒常不变，更新万物\BibleRef{默 21:5}，不接受任何祂未曾赐予的东西。因此，那些说灵魂是天主性一部分的人狂妄自大、眼瞎如石，不明白他们实际上是在断言天主是可变的，且祂自身会遭受任何可能加于其本性的事物。\par

% structure: p0017 source=h2 target=subsection reason=relative-heading-rank; base=h2

\subsection{七 (6) 他们关于魔鬼从不善，因而非天主所造的观点被驳斥}

第六项注意到他们说魔鬼从不善，他的本性不是天主的造化，而是出自混沌和黑暗：因为我想他没有唆使者，而是他自己是一切恶的根源和本体：然而真实的信仰，即公教信仰，承认一切精神或物质受造物的本体是善的，而恶没有正面存在；因为天主，宇宙的创造者，没有造出任何不是善的事物。因此，魔鬼若保持他受造时的状态，也会是善的。但由于他滥用了他本性的优越，并\BibleQuote{没有存留在真理中\BibleRef{若 8:24}，}他并没有转入相反的本体，而是从他所应忠于的至善中叛离了：正如那些作出如此断言的人自己从真理一头栽入谎言，并为其自愿的乖张受到谴责一样：当然，这种恶在他们之内，但它本身不是一种本体，而是加于本体的一种惩罚。\par

% structure: p0019 source=h2 target=subsection reason=relative-heading-rank; base=h2

\subsection{八 (7) 他们拒绝婚姻被谴责}

第七点紧接着的是他们对婚姻的谴责以及对生儿育女的憎恶：在这一点上，他们几乎在所有方面都与摩尼教徒的不敬一致。但正如他们的行为所证明的，他们之所以憎恶婚姻纽带，正因为一旦保存了婚姻与后代的贞洁，就没有放荡的自由了。\par

% structure: p0021 source=h2 target=subsection reason=relative-heading-rank; base=h2

\subsection{九 (8) 他们不信肉体复活已被教会谴责}

他们的第八点是：人的身体形成是魔鬼的伎俩，而\OriginalTerm{受孕的种子}{seed of conception}是在妇女子宫中藉魔鬼的帮助而成形的；因此，不可相信\OriginalTerm{肉身的复活}{resurrection of the flesh}，因为构成身体的质料与灵魂的尊严不相符。这谎言无疑出自魔鬼的工作，而如此怪异的见解是恶魔的伎俩，恶魔并不在妇女腹中塑造人，而是在异端者的心中炮制这些谬误。这从不洁的毒源——尤其从\OriginalTerm{摩尼教}{Manichæan}邪恶的源泉——流出的毒液，早已被公教信仰提出控诉并定罪。\par

% structure: p0023 source=h2 target=subsection reason=relative-heading-rank; base=h2

\subsection{X. (9) \Quote{\OriginalTerm{恩许的子女}{children of promise}被圣神怀孕}的观念全然不合圣经且不合公教。}

第九项通知宣称，他们说\OriginalTerm{恩许的子女}{sons of promise}确实由妇女所生，但由\OriginalTerm{圣神}{Holy Spirit}怀孕；以免由\OriginalTerm{血肉种子}{carnal seed}所生的后裔似乎分享天主的产业。这与公教信仰相悖，公教信仰承认每一个人在灵魂和身体的实质上都是由宇宙的创造者所形成，并在母腹中接受生命的气息：尽管罪的玷污和必死的责任从第一父母传到后代；直到\OriginalTerm{重生的圣事}{sacrament of Regeneration}来救助他，藉此我们因圣神而重生为恩许的子女，不是在肉体的子宫中，而是在\OriginalTerm{圣洗}{baptism}的能力中。因此，达味——他确实是恩许的子女——对天主说：\BibleQuote{你的手造了我，塑造了我。} 对耶肋米亚，主说：\BibleQuote{我尚未在母腹中形成你，我已认识你，在你母亲胎中我已祝圣了你。} \BibleRef{耶 1:5}\par

% structure: p0025 source=h2 target=subsection reason=relative-heading-rank; base=h2

\subsection{XI. (10) 他们关于\OriginalTerm{灵魂}{souls}在进入人身之前具有\OriginalTerm{前在存在}{previous existence}的理论被驳斥。}

在第十项标题下，他们被报道主张，被安置在人身中的灵魂先前没有身体，并在天上的居所中犯了罪，为此原因，从高位堕落到低位，降落在各种品质的统治灵体上，并经过一系列空气和星辰的力量——有些更猛烈，有些更温和——后被封闭在不同种类和条件的身体中，以至于今生分配给我们的任何种类和不平等，都似乎是先前原因的结果。这个\OriginalTerm{亵渎的寓言}{blasphemous fable}是他们从许多人的错误中为自己编织出来的；但所有这些，\OriginalTerm{公教信仰}{Catholic Faith}从与它身体的合一中切断，持续而真实地宣告：人的灵魂在注入身体之前并不存在，并且它们的植入不是出于任何别的，而是出于天主，他是灵魂和身体的创造者。因为通过第一个人的过犯，整个人类种族被玷污，除非通过基督的\OriginalTerm{圣洗圣事}{sacrament of baptism}，没有人能从旧亚当的状态中被释放；在圣洗中，重生者之间没有区别，正如宗徒所说：\BibleQuote{你们凡受洗归于基督的，就是披戴基督了：不再有犹太人或希腊人；不再有为奴的或自由的；不再有男的或女的；因为你们众人在基督耶稣内已合而为一} \BibleRef{迦 3:27}。那么星辰的轨迹或命运的计谋有何相干？世俗事物的变化状态及其不安的差异呢？请看\OriginalTerm{天主的恩宠}{grace of God}如何使所有这些不平等者变得平等；这些人无论今生中有什么劳苦，只要保持忠信，就不会不幸，因为他们能在每一次试炼中与宗徒一同说：\BibleQuote{谁能使我们与基督的爱分离？是困苦、是窘迫、是迫害、是饥饿、是赤身、是危险、是刀剑吗？正如经上所载：\InnerQuote{为了你，我们终日被处死，被看作待宰的羊} \BibleRef{咏 43:22}。然而，靠着爱我们的那一位，我们在这一切事上大获全胜} \BibleRef{罗 8:35}。因此，教会——\OriginalTerm{基督的身体}{body of Christ}——不惧怕世界的不平等，因为她不渴望暂时的财物；也不害怕被命运空洞的威胁所压倒，因为她知道自己在患难中藉忍耐而坚强。\par

% structure: p0027 source=h2 target=subsection reason=relative-heading-rank; base=h2

\subsection{XII. (11) 他们的\OriginalTerm{占星观念}{astrological notions}被定罪。}

第十一项亵渎是，他们认为人的灵魂和身体都受制于\OriginalTerm{命中注定的星辰}{fatal stars}的影响：这种愚昧迫使他们纠缠于外教人的一切错误，并努力吸引他们认为有利的星辰，而缓和那些不利于他们的星辰。但对于追随这些追求的人，在\OriginalTerm{天主教会}{Catholic Church}中没有位置；一个投身于这类信念的人完全将自己从\OriginalTerm{基督的身体}{body of Christ}中分离。\par

% structure: p0029 source=h2 target=subsection reason=relative-heading-rank; base=h2

\subsection{十三、(12) 他们认为\OriginalTerm{某些权能}{certain powers}统治\OriginalTerm{灵魂}{soul}，而\OriginalTerm{星辰}{stars}统治\OriginalTerm{身体}{body}，这是不合圣经且荒谬的。}

这些十二点中的第十二项是：他们将灵魂的各部分划分到某些权能之下，将身体的肢体划分到其他权能之下；并且他们通过用圣祖的名字来命名统治灵魂的内在权能，相反地，他们将星宿的标记归属于那些他们置于身体之下的东西。他们用所有这些把自己纠缠在无法解开的迷宫中，不听宗徒的话，他说：\Quote{你们要小心，恐怕有人用哲学和虚妄的欺骗，照着人的传统，照着世界的元素，而不照着基督，把你们掳去；因为天主性圆满的全部有形地居住在他内，你们在他内也得到了圆满，他是一切执政掌权者的元首 \BibleRef{哥 2:8}}。又说：\Quote{不要让任何人用自愿自卑和敬拜天使骗取你们的赏报，他却窥探他所没有看见的，妄自尊大，凭着肉体的感受，不持守元首；其实全身靠着关节和韧带，获得供应和联系，以天主的增长而增长。}那么，把法律没有教导、先知没有歌唱、福音真理没有宣告、宗徒道理没有传授的东西接受在心里，有什么用处呢？然而，这些事适合那些宗徒所说的人的心意：\Quote{因为时候要到，人必不忍受健全的道理，反而耳朵发痒，就随从自己的情欲，为自己堆积许多教师；并且掩耳不听真理，偏向虚构的故事 \BibleRef{弟后 4:3}}。所以，我们与那些竟敢教导或相信这样的事，并且竭尽全力说服人相信肉身的实体与复活的希望无关，从而摧毁基督降生成人的全部奥迹的人，毫无共同之处；因为基督要是不该解救完整的人性，那么他就不该取了完整的人性。\par

% structure: p0031 source=h2 target=subsection reason=relative-heading-rank; base=h2

\subsection{十四（十三）他们臆想的圣经划分被驳斥。}

第十三位，他们声称全部正典圣经要以圣祖的名字被接受：因为那十二种革新内在之人的美德在他们的名字中被指明，并且没有这种知识，灵魂就不能实现其革新，并回归到它原先出来的实质。但基督信仰的智慧鄙视这种邪恶的错觉，因为它知道真实天主性的本质是不可侵犯且不可改变的；而灵魂，无论活在身体内还是与身体分离，都受许多情欲的支配；而当然，如果它是神圣本质的一部分，就不会遭遇任何苦难。因此，两者之间无法比较：一位是造物主，另一位是受造物。因为他始终如一，不受任何变化；而灵魂是可变的，即使它没有改变，因为其不变的能力是恩赐，而非本性。\par

% structure: p0033 source=h2 target=subsection reason=relative-heading-rank; base=h2

\subsection{十五（十四）他们以为圣经认可他们将身体屈服于星体影响下的主张被驳斥。}

在第十四项下，陈述了他们关于身体状态的观点，即：由于身体的属地特性，它被置于星辰和星座的权下，并且在圣书中发现许多提及外在之人的事情，其目的是在圣经本身中可以看到神圣本性与属地本性之间某种对立在起作用；并且治理灵魂的权能所要求的东西，可以区别于塑造身体者所要求的东西。编造这些故事是为了主张灵魂是神圣实质的一部分，并相信肉身属于恶劣的本性：因为他们认为世界本身及其元素不是慈善天主的工程，而是邪恶作者的产物；并且为了用漂亮的外衣掩盖这些亵渎的谎言，他们几乎用自己邪恶观念的着色污染了所有神圣的言论。\par

% structure: p0035 source=h2 target=subsection reason=relative-heading-rank; base=h2

\subsection{十六（十五）他们篡改的圣经抄本以及他们的伪经被禁止。}

关于这一点，你在第十五项下的言论提出了抱怨，并表达了对其魔鬼般僭越的应有憎恶，因为我们也从可信见证人的叙述中确认了这一点，并且发现他们的许多抄本——尽管被冠以正典之名——极度败坏。因为他们若不将毒杯用少许蜂蜜调甜，如何能欺骗头脑简单的人，以免那致人死命的东西因过度难闻而被察觉？因此必须谨慎，并尽司铎之最大勤勉，以防止使用那些与纯净真理不符的伪造抄本进行诵读。而那些以宗徒之名行世的伪经，是诸多谎言的温床，不仅应被禁止，而且应完全拿走，在火中烧成灰烬。因为尽管其中有某些看似敬虔的东西，但它们绝无毒害，并且透过故事的诱惑，先以奇迹叙述迷惑人，然后将他们套入某种错误的圈套。因此，若有任何主教没有禁止人家里拥有伪经著作，或容许那些被\Person{普里西利安}伪造的篡改所败坏的抄本在教会中诵读，就让他知道他将被视为异端者，因为那不解救别人脱离错误的人，表明他自身已经误入歧途了。\par

% structure: p0037 source=h2 target=subsection reason=relative-heading-rank; base=h2

\subsection{十七（十六）关于\Person{狄克提尼乌斯}的著作。}

在最后一点上，有人正当抱怨，\Person{狄克提纽斯}与\Person{普里西利安}教义一致所写的论著，被许多人恭敬阅读：因为若他们认为对\Person{狄克提纽斯}的纪念应予尊重，他们应当钦佩他的悔改，而非他的堕落。因此，他们所读的不是\Person{狄克提纽斯}，而是\Person{普里西利安}；他们赞同他误写的，而非他悔改后所偏爱的。但任何人不得冒此风险而不受惩罚，也不得有人被视为天主教徒，若他利用那些不仅被天主教会，且被作者本人所谴责的著作。不要让那些迷失的人得以伪装，在基督徒名号的掩饰下规避帝国法令的规定。因为他们心中怀着所有这些分歧而依附于天主教会，目的是既尽可能使人归附，又冒充我们的人以逃避法律的严厉。\OriginalTerm{普里西利安派}{Priscillianists}和\OriginalTerm{摩尼派}{Manichæans}都这样做；因为二者之间有如此紧密的联系，以至于他们只在名称上不同，在亵渎上却是一致的：因为虽然摩尼派拒绝旧约，而另一派假装接受旧约，但两者的目的都趋向同一结局，即一方在接受时败坏，而另一方攻击并拒绝它。\par

但在他们那可憎的奥秘中，越是污秽，就越被小心隐藏，他们的罪行是同一的，他们的污秽思想是同一的，他们的丑恶行径是相似的。虽然我们羞于如此直言，但我们通过最痛苦的追查，并借着被捕\OriginalTerm{摩尼派}{Manichæans}的供认，揭露了出来，并公之于众：免得有任何怀疑的余地，尽管这些事已经在我们的法庭上由那些犯下罪恶的人亲口披露；当时不仅有大批司铎聚集，还有名望和尊严的人，以及若干元老和平民在场，正如我们写给主内亲爱的你们的信函所显示的那样。关于\OriginalTerm{普里西利安派}{Priscillianists}的不道德行为，所发现和广泛传播的，与关于摩尼派的污秽邪恶所发现的完全相同。因为那些在思想上全然处于同一堕落水平的人，不可能在宗教事务上有所不同。\par

因此，我们已详尽驳斥了所有内容，而他们观点的备忘录与之并无矛盾；我想，我们已经充分表明了对于兄弟你向我们询问的那些事的意见，以及如果这样亵渎性的错误甚至在一些司铎心中得到接纳，或者说得温和些，没有被他们积极反对，那是多么难以容忍。他们怎能凭良心维持赋予他们的尊贵地位，却不为托付给他们的灵魂劳苦？野兽闯入，他们却不关羊栈。强盗埋伏，他们却不设警戒。疾病增多，他们却不寻求药方。但此外，当他们对行事更谨慎的人表示不同意，并回避通过书面宣认来绝罚那些已被全世界谴责的亵渎时，他们想让人们理解为，他们不属于弟兄的行列，而是站在敌人那边？\par

% structure: p0041 source=h2 target=subsection reason=relative-heading-rank; base=h2

\subsection{十八、基督的身体确实安息在坟墓中，也确实复活了。}

此外，在你机密信函中最后提到的事情上，我感到惊讶，任何有智识的基督徒竟然会困惑：当基督下降到阴府时，他的肉身是否安息在坟墓中？因为正如他真实地死了并被埋葬，他也确实在第三日复活了。因为主自己也曾预告，对犹太人说了：\BibleQuote{你们拆毁这座圣殿，三天内我要把它重建起来。} \BibleRef{若 2:19} 而福音作者补充评论说：\BibleQuote{但祂所说的是指祂身体的圣殿。} 达味先知也以主和救主的口吻预言了这真理，说：\BibleQuote{而且我的肉身也要安息于希望中，因为祢必不将我的灵魂遗弃在阴府，也不让祢的圣者见到腐朽。} 从这些话中显然可见，主的肉身被埋葬，确实安息了，并未腐朽：因为藉着灵魂的返回，它迅速复苏并复活了。不信这一点是极为亵渎的，无疑与\Person{摩尼}和\Person{普里西利安}的教义如出一辙，他们以亵渎的观念假装承认基督，但只是为了破坏祂降生成人、死亡和复活的真实性。\par

因此，应在你们中间召开一次主教会议，并让邻近省份的司铎们在一个对所有人都适合的地方聚会；如此，按照我们对你咨询请求的答复，可以全面调查是否有任何主教感染了这种异端的瘟疫：因为他们如果拒绝谴责这个最不义、充满错误理念的派别，无疑必须被断绝共融。因为绝不允许一个承担宣讲信德职责的人胆敢坚持与基督的福音以及普世教会的信经相反的意见。在这样的师傅教导的地方，会有怎样的门徒呢？当反对人类社会利益时，贞洁的神圣被根除，婚姻的纽带被推翻，生育被禁止，肉身的本质被谴责，并且与对真天主的真实敬拜相反，天主圣三被否认，位格的独特性被混淆，人的灵魂被宣称是神圣本质，并被魔鬼的意志禁锢在肉身中，天主子被宣告为因生于童贞女而是独生者，而非由圣父所生，同时又被主张既不是天主的真后裔，也不是童贞女的真儿子：以致在虚假的苦难和不真实的死亡之后，甚至从坟墓中重新取回的肉身的复活也应被视为虚构。但他们自称天主教会之名是徒劳的，因为他们不反对这些亵渎之言；如果他们能如此耐心地倾听这些言论，他们必定相信它们。因此，我们已写信给我们在\Place{塔拉戈}、\Place{迦太基}、\Place{卢西塔尼亚}和\Place{加里西亚}各省的弟兄和同僚主教们，敦促召开全体主教会议。亲爱的，你们将负责安排，将我们的权威指示传达给上述各省的主教们。但如果有什么（天主禁止的事）阻碍了加里西亚全体会议的召开，至少让司铎们聚会，我们的弟兄\Person{伊达西乌}和\Person{塞波尼乌斯}将负责召集他们，并借助你们自己的辛勤努力，也通过一个省级会议来加快对这些严重创伤的治疗。日期为七月二十一日，在显赫的\Person{卡利皮乌斯}和\Person{阿尔达布里}担任执政官期间（447年）。\par

\SourceRecord{Translated by Charles Lett Feltoe.；Nicene and Post-Nicene Fathers, Second Series；Vol. 12.；Edited by Philip Schaff and Henry Wace.；Buffalo, NY: Christian Literature Publishing Co.,；1895.；<http://www.newadvent.org/fathers/3604015.htm>.}

% structure: p0001 source=h1 target=section reason=page-title

\section{第十六封信}

主教利奥问候全西西里的诸位主教，在主内致意。\par

% structure: p0003 source=h2 target=subsection reason=relative-heading-rank; base=h2

\subsection{一、引言。}

因天主的诫命和宗徒的劝诫，我们受激励要谨慎看守所有教会的情况；若发现任何需要责备之处，就当以迅速的关怀使人们从愚昧无知或鲁莽狂妄中醒悟。因为我们受主亲自的命令所告诫，圣宗徒伯多禄曾被再三以奥秘的嘱托压在身上，即爱基督的人应牧养基督的羊群，我们因敬畏那以天主恩宠的丰盈所占有的宗座，不得不尽可能逃避懈怠的危险：免得首生宗徒的宣认——他证明自己爱天主——不被发现在我们身上；因为若他（藉着我们）漫不经心地牧养那屡次托付给他的羊群，便证明他并不爱首生牧者。\par

% structure: p0005 source=h2 target=subsection reason=relative-heading-rank; base=h2

\subsection{二、圣洗应在逾越节期间举行，而非在主显节。}

因此，当我以可靠的证据听闻（我早已对你们的作为怀着兄弟般的关爱焦虑，亲爱的弟兄）在教会的主要圣事之一，你们偏离了宗徒规制的做法，在主显节的庆典中施行的圣洗人数比逾越节更多时，我感到惊讶，你们或你们的前任竟会引入如此不合理的革新，以致混淆两个节期的奥迹，以为基督受贤士朝拜的日子与他从死者中复活的日子没有区别。你们本不该陷入这种错误，如果你们从你们受祝圣为主教职的源头领受你们的一切礼规；如果圣宗徒伯多禄的宗座——那是你们司铎尊严之母——是教会方法的公认导师，就不会如此。我们本可以更不宽容地忍受你们背离其规则，若你们先前曾受过我们任何形式的警戒责问。但现在我们并未绝望于纠正你们，所以必须显示温和。虽然以无知为托词的借口在司铎身上几乎不可容忍，但我们宁愿缓和必要的责备，并明确教导你们教会的正确方法。\par

% structure: p0007 source=h2 target=subsection reason=relative-heading-rank; base=h2

\subsection{三、必须区分不同节期在尊严和缘由上的差异。}

人类的恢复确实始终不变地预先定于天主永恒的计划中；但藉着我们的主耶稣基督在时间中要完成的事件序列，始于圣言降生成人。因此，有一个时刻：当天使传报时，蒙福的童贞玛利亚相信自己将因圣神怀孕并确实受孕；另一个时刻：当她无损童贞纯洁时，男婴诞生，因天朝随从的欢跃而显给牧人们；又一个时刻：婴孩受割损；又一个时刻：按法律要求的祭品为他奉献；又一个时刻：三位贤士受新星光辉吸引从东方来到白冷，以神秘的礼品奉献朝拜这婴孩。\par

同样，日子并不相同：他藉天主的安排逃往埃及脱离邪恶的黑落德，以及追捕者死后他从埃及被召回加里肋亚。在这些不同境遇中，也包括他身体的成长：正如福音作者所见证，主在年龄和恩宠上增进；逾越节时他与父母上耶路撒冷圣殿，当他在回程的人群中不见时，被发现坐在长老中，在惊讶的经师中辩论，并解释他留下的原因：\BibleQuote{为什么你们找我？难道你们不知道我必须在我父亲的地方吗？}表明他是那圣殿所属之父的儿子。再一次，后来当他更公开地显明并寻求前驱若翰的圣洗时，关于他是天主的身份还能存疑吗？主耶稣受圣洗后，圣神如同鸽子降下并停留在他身上，圣父的声音从天上传来：\BibleQuote{你是我的爱子，我喜悦你。}\BibleRef{玛 3:17}这一切我们尽可能简洁地提及，为叫你们知道，亲爱的弟兄，虽然基督一生中所有日子都因他的许多大能事迹而圣化，并且在他所有的行动中闪耀着奥妙的圣事，但有些时候事件由预兆暗示，有些时候实现应验；而且记载的所有救主事迹并不都适合圣洗的时间。因为如果我们不加区分地同样纪念那些我们知道主在受圣若翰圣洗后所做的事，他的一生就得连续不断地庆祝节期，因为他所有的作为都充满奇迹。但因为智慧与知识之神教导整个教会的宗徒和教师们，不允许在我们基督徒的礼规中存在任何无序或混乱，我们必须辨别各种庆典的相对重要性，并在我们祖先和统治者的所有制度中遵守合理的区分：否则我们无法成为\BibleQuote{一个羊群，一个牧人}\BibleRef{若 10:17}，除非如宗徒教导我们：\BibleQuote{我们大家说同样的话，并在同一心意和同一判断上圆满合一}\BibleRef{格前 1:10}。\par

% structure: p0010 source=h2 target=subsection reason=relative-heading-rank; base=h2

\subsection{四、解释为什么逾越节和五旬节是圣洗的适宜时期。}

因此，尽管与基督的贬抑有关的事和与祂的受举扬有关的事都集中于同一且同一的位格，并且祂内一切天主性的能力和人性的软弱都有助于完成我们的恢复；然而，适当的做法是，在受难者的死亡之日和死者的复活之日，圣洗的力量应将旧人变为新人：使基督的死亡和祂的复活在重生者身上运作，正如蒙福的宗徒所说：\Quote{难道你们不知道，我们这受过洗归于基督耶稣的人，就是受洗归于祂的死亡吗？我们藉着洗礼已与祂同葬于死亡，好使基督藉着父的光荣从死者中复活，我们也能同样在新生活中行走。如果我们已与祂死亡的形状结合，我们也将与祂复活的形状结合，}以及外邦人的导师在推荐圣洗圣事时进一步讨论的其余内容：好使人从这教义的精神看出，那一日就是所选定的日子，时候到了，重生人子并收养他们为天主的儿子，藉着奥妙的象征和仪式，在肢体中所行的与在元首本身所行的相符，因为在圣洗礼仪中，死亡经由杀死罪恶而产生，三次浸入模仿了在坟墓中三日，而从水中升起则相似于从坟墓中复活的那一位。因此，行为本身的性质教导我们，那一日就是普遍领受恩宠的公认日子，恩赐的能力和行动的特质都源自于此。这一点由以下考虑有力地证实了：主耶稣基督自己从死者中复活后，将圣洗的形式和能力传授给了祂的门徒，而所有教会的首领都从他们身上领受了训示，用这些话：\BibleQuote{你们去使万民成为门徒，因圣父、圣子及圣神之名给他们授洗。\BibleRef{玛 28:19}}。当然，祂本可以在受难之前就教导他们，但祂特别希望人明白，重生的恩宠始于祂的复活。确实还必须补充，由圣神降临所圣化的庄严五旬节时期也被允许，如同逾越节庆的延续和完成。尽管其他庆节在一周的其他日子举行，但这个庆节（五旬节）总是在主复活所标记的那一天发生：可以说是伸出援助恩宠的手，邀请那些因疾病、长途旅行或航行困难而无法参加复活节庆的人，藉着圣神的恩赐获得他们渴望的目的。因为天主的独生子自己愿意在信徒的信德和祂工程的效能中，自己和圣神之间没有差异被感受到：因为他们的本性没有差异，正如祂所说：\BibleQuote{我要请求父，祂要赐给你们另一位护慰者，使祂永远与你们同在，就是真理之神。\BibleRef{若 14:16}；}又说：\Quote{但那护慰者，就是圣神，父因我的名要派遣来的，祂要教导你们一切，并使你们想起我对你们所说的一切；}又说：\Quote{当祂，真理之神，来到时，祂要引导你们进入一切真理。}因此，既然基督是真理，圣神是\OriginalTerm{真理之神}{Spirit of Truth}，而\OriginalTerm{护慰者}{Comforter}之名适用于两者，那么这两个庆节并非不相同，因为圣事是相同的。\par

% structure: p0012 source=h2 target=subsection reason=relative-heading-rank; base=h2

\subsection{五、\Person{圣伯多禄}的例子作为\OriginalTerm{五旬节}{Whitsuntide}洗礼的权威。}

我们并不是凭自己的信念主张这一点，而是持守宗徒的权威，我们用一个十分恰当的例子来证明，追随蒙福的宗徒\Person{伯多禄}，他在圣神降临应许之日，使信者的人数得以满足，将他讲道所归化的人中三千人藉着圣洗圣事奉献给天主。\WorkTitle{宗徒大事录}在其忠实叙述中教导此事，说：\Quote{他们听见这话，心中刺痛，遂向\Person{伯多禄}和其余宗徒说：\InnerQuote{诸位弟兄，我们该作什么？}\Person{伯多禄}对他们说：\InnerQuote{你们悔改，并各人因耶稣基督的名受洗，好使你们的罪得赦，并领受圣神的恩惠。因为这恩许是给你们和你们的子女，以及一切远方的人，凡是主我们的天主所召叫的。}他还用许多别的话向他们作证并劝勉说：\InnerQuote{你们应救自己脱离这邪恶的世代。}于是接受他话的人都受了洗，那一天约增添了三千人。\BibleRef{宗 2:37}}\par

% structure: p0014 source=h2 target=subsection reason=relative-heading-rank; base=h2

\subsection{六、在紧急情况下，其他时间也可以许可行洗礼。}

因此，既然我们所说的这两个时节显然是教会中为选民施洗的正确时节，我们劝告你们，亲爱的，不要在此规定之外增添其他日子。因为，尽管还有其他庆节也应在天主的尊荣中受到极大的尊敬，但我们仍应理性地守护这首要且最伟大的圣事，将其视为深奥的奥迹，而非日常常规的一部分：但并不禁止在任何时候为处于危险中的人施行洗礼以提供救助。因为，尽管我们将那些没有患病压力且生活在平安中的人的誓愿推迟到这两个紧密相连且相关的庆节，但我们绝不拒绝在任何时间将真正救恩的唯一保障给予那些面临死亡危险、围城危机、迫害痛苦或船难恐惧的人。\par

% structure: p0016 source=h2 target=subsection reason=relative-heading-rank; base=h2

\subsection{七、主由\Person{若翰}所受的洗礼与信徒的洗礼大不相同。}

但若有人以为主显节（虽然理当按适当程度予以应有的尊敬）因为有些人推定主在那一天来到若翰的圣洗前，而声称此节日享有圣洗的特权，便应知那圣洗的恩宠及其理由大不相同，并且与那些由圣神重生的人所领受的能力（经上说：\BibleQuote{他们不是由血、也不是由肉欲、也不是由人意，而是由天主生的}\BibleRef{若 1:13}）并不平等。因为主无需罪过的赦免，也不寻求重生的救治，祂愿意受洗，正如祂愿意受割损，并愿意为祂的洁净献上祭牲：祂曾\BibleQuote{生于女人}\BibleRef{迦 4:4}，如宗徒所说，也为要成为\Quote{属于法律}的那位，祂来\BibleQuote{不是为废除，而是为成全}\BibleRef{玛 5:17}，并以成全来终结，正如蒙福的宗徒所宣示的：\BibleQuote{因为基督是法律的终向，使凡信祂的人都获得正义}\BibleRef{罗 10:4}。但圣洗的圣事是祂亲自建立的，因为\BibleQuote{祂在一切事上居首位}\BibleRef{哥 1:18}，祂教导说祂自己就是元始。祂在那时批准了重生的能力，当时从祂的肋旁流出了赎价的血和圣洗的水。因此，正如旧约为新约作证，\Quote{法律是藉梅瑟传授的，恩宠和真理却是由耶稣基督而来的}；正如各种祭献预表了那唯一的牺牲，许多羔羊的宰杀因那一位的奉献而结束，经上论及祂说：\BibleQuote{请看，天主的羔羊，除免世罪者}\BibleRef{若 1:29}；同样，若翰——不是基督，而是基督的先行者；不是新郎，而是新郎的朋友——他如此忠信地寻求\BibleQuote{不是自己的利益，而是耶稣基督的利益}\BibleRef{斐 2:21}，以至于自认不配解祂的鞋带：因为他自己确实是\Quote{以水施洗，叫人悔改}，但那一位将以双重能力既恢复生命又毁灭罪过，祂要\Quote{以圣神和火施洗}。因此，亲爱的弟兄们，既然所有这些明确的证据摆在你们面前，使你们消除一切疑虑，认识到在给那些按宗徒规条须藉驱魔得以净化、藉斋戒得以圣化、并藉多次讲道得以训导的被选者施行圣洗时，只应遵守两个季节，即复活节和五旬节；我们嘱咐你，弟兄，不要再进一步偏离宗徒制度。因为今后，任何认为可以藐视宗徒规条的人，都不会不受惩罚。\par

% structure: p0018 source=h2 target=subsection reason=relative-heading-rank; base=h2

\subsection{八、西西里主教们应每半年派遣三人出席在罗马的主教会议。}

因此，为保持完美的和谐，我们最首要要求：按照圣教父们健全的规条，每年应举行两次主教会议，你们中的三人务必在每次9月29日出席弟兄们的会议；因为这样，藉天主的恩宠，我们将更容易防范基督教会中罪行和错误的兴起；而且此会议必须在蒙福的宗徒伯多禄面前常会集并商议，使他的所有宪章和教会法令，在主的所有司铎中，得以保持不可侵犯。\par

我们藉蒙主默感认为有必要训导你们的这些事项，愿通过我们的弟兄和同为主教的兄弟巴基鲁斯及帕斯卡西努斯传达给你们。愿我们藉他们的报告得知宗座的制度被你们虔敬地遵守。\Signature{日期为10月21日，在显赫的阿里比乌斯和阿尔达布里乌斯执政之年（447年）。}\par

\SourceRecord{Translated by Charles Lett Feltoe.；Nicene and Post-Nicene Fathers, Second Series；Vol. 12.；Edited by Philip Schaff and Henry Wace.；Buffalo, NY: Christian Literature Publishing Co.,；1895.；<http://www.newadvent.org/fathers/3604016.htm>.}

% structure: p0001 source=h1 target=section reason=page-title

\section{第十七封信}

% structure: p0002 source=h2 target=subsection reason=relative-heading-rank; base=h2

\subsection{禁止出售教会财产，除非有利于教会}

教宗\Person{良}，致\Place{西西里}众主教。\par

特定的投诉事件引起了我们的关注，因为我们负有\Quote{关心所有教会}的责任，因此我们应当颁布一项永久性法令，杜绝所有主教将那些在你们省内的两个教会中曾毫无顾忌地提出并错误执行的做法奉为惯例。\Place{陶罗梅尼乌姆}教会的圣职人员哀叹他们因主教出售、赠与或处理而挥霍了所有产业而陷入贫困；\Place{帕诺尔穆斯}教会的圣职人员最近任命了一位新主教，在由我们主持的神圣主教会议上，也提出了类似的对前任主教治理不善的投诉。因此，尽管我们已经就两个教会的利益发出了指示，但为了防止这种可恶掠夺的邪恶榜样将来被援引为先例，我们愿意将此项正式命令永久地约束于你们，亲爱的弟兄们。因此，我们规定：任何主教都不得擅自赠予、交换或出售其教会的任何财产；除非他预见这样做可能带来利益，并且在征询全体圣职人员的意见并取得他们同意后，他才可以决定做那些无疑对该教会有利的事。因为司铎、执事或任何级别的圣职人员如果纵容教会的损失，必须知道他们将丧失级别和共融：因为完全公平的是，亲爱的弟兄们，不仅主教，而且整个圣职人员都应当促进他们教会的利益，并完好地保存那些为自己灵魂的救恩而将自身财物奉献给教会的人所捐赠的礼物。发于10月20日，在显赫的\Person{卡勒皮乌斯}执政之年（447年）。\par

\SourceRecord{Translated by Charles Lett Feltoe.；Nicene and Post-Nicene Fathers, Second Series；Vol. 12.；Edited by Philip Schaff and Henry Wace.；Buffalo, NY: Christian Literature Publishing Co.,；1895.；<http://www.newadvent.org/fathers/3604017.htm>.}

% structure: p0001 source=h1 target=section reason=page-title

\section{第十八封书信}

\Emph{罗马城主教良，致阿奎莱亚主教雅奴雅里。}\par

% structure: p0003 source=h2 target=subsection reason=relative-heading-rank; base=h2

\subsection{凡弃绝异端与分裂而回归教会者，必须极其明确地宣明其反悔；属于神职人员者，可保留其品级，但不得晋升。}

弟兄，披阅你的来信，我们认出你信德的活力——这原是我们早已知道的——并且祝贺你以牧者般的警醒，守护基督的羊群：免得那些披着羊皮潜入的豺狼，以野兽般的凶残撕碎纯朴的人，不仅自己肆无忌惮地横行，也毁坏那些健全的。又为防止蛇蝎般的诡计得逞，我们以为宜于提醒你，挚爱的，提醒你：凡是从我们中间堕落到异端或分裂派别之中，并以异端共融的玷污无论程度如何地污染了自己的人，若在悔悟时，未经合法且明确的补赎，就被接纳进入天主教会共融，那是对他灵魂有危险的。因为，最为有益且充满一切属灵医治之益的，是那些司铎、执事、副执事或任何品级的神职人员——他们希望显得已经改过，并恳求重新回到他们久已失去的公教信仰——应当首先毫不含糊地承认，他们弃绝自己的错误以及这些错误的始作俑者，好使他们的卑劣见解被彻底摧毁，不留任何重蹈覆辙的希望，并且没有任何肢体因与他们接触而受损，每一项都已以其应有的反悔予以应对。关于他们，我们也命令遵守此项教规：他们若被允许留在现品级中不得再行晋升，应视为极大的宽免；但仅以他们未经第二次圣洗玷污为条件。主不会轻易放过那判断任何此类人适于晋升圣秩的人。倘若连无可指摘者，也须在全面审查后才予晋升，那么对于受怀疑者，岂不更应拒绝晋升？因此，挚爱的弟兄——我们因你的虔诚而喜乐——请关心我们的指示，并安排谨慎而迅速地执行这些有益的建议与健康的训令，它们关系到整个教会的福祉。但请不要怀疑，挚爱的，倘若我们所规定的为遵守教规及信仰纯正之事被忽略（我们并不预料此事），我们将强烈动怒：因为低级人员的过失，最应归咎于怠惰而不尽责的治理者，他们往往因拒绝施用必要的补救而助长许多病症。颁于杰出的卡勒丕与阿尔达布里任执政之年（447年）十二月三十日。\par

\SourceRecord{Translated by Charles Lett Feltoe.；Nicene and Post-Nicene Fathers, Second Series；Vol. 12.；Edited by Philip Schaff and Henry Wace.；Buffalo, NY: Christian Literature Publishing Co.,；1895.；<http://www.newadvent.org/fathers/3604018.htm>.}

% structure: p0001 source=h1 target=section reason=page-title

\section{第十九封书信}

\Person{利奥}主教，致其爱弟\Person{多鲁斯}\EditorialNote{本味芬托主教}。\par

% structure: p0003 source=h2 target=subsection reason=relative-heading-rank; base=h2

\subsection{一、他责备多鲁斯容许一位资浅的司铎被擢升到资深者之上，且第一和第二资深者竟默许此事。}

我们痛心，因为我们获悉您所行之事以其应受责难的新奇侵犯了整个教会纪律体系，这使我们本来对您所怀的期待落空了：尽管您完全清楚，我们何等关切地希望教规的规定在主的诸教会中得以遵守，且万民的所有司铎都应将以一切合理的权威防止因任何越轨行为而违反神圣典章的规定视为自己的特殊职责。因此我们感到惊讶，您这位本应严格持守宗座命令的人，竟如此疏忽大意，甚至可说是悖逆，以致于成为传给您的法律的破坏者而非守护者。因为我们从您的司铎保禄的附随报告中得知，司铎团在您那里的秩序因奇怪的阴谋和卑劣的勾结而被扰乱；一个未经考虑、过早地被提升，而另一些其年龄本应推荐他们晋升且未被指控有任何过失的人却被忽略。然而，如果阴谋家的急切或支持者的无知热心要求那种习惯从未允许的事，即新人优先于经验丰富者，稚嫩少年优先于年长者，那么您有责任以一切合理的权威，通过勤勉和教导来抑制请求者的不当愿望，以免您仓促擢升到司铎职位的人进入他的职务时损害与他共事的人，并因他内心增长的不是谦卑之德而是自大之恶而腐化。因为您并非不知道主曾说过：\Quote{自谦者必被高举: 自高者必被贬抑,}又曾说：\Quote{但你们寻求从微小中增长，从伟大中变得微小。}因为这两种行为都是不合顺序、不合时宜的；人劳作的一切果实都丧失，其功德的一切尺度都归于无效，如果获得尊严与奉承的程度成正比：这样，渴望出人头地不仅贬低了渴望者本身，也贬低了纵容他的人。但如果，如所声称的，第一和第二司铎如此同意埃皮卡皮乌斯居于他们之上，以致为了自己的羞辱而要求他得荣耀，那么当他们自愿贬低自己时，本不应答应他们所愿：因为反对而不是屈从于这种可悲的愿望对您来说更值得。但他们卑贱怯懦的顺从不能损害其他良心清白且没有藐视天主恩宠的人；因此，无论他们以何种操作将优先权让给他人，都不能降低他们之后的人的地位，也不能因为他们将最后者置于自己之上，那人就领先于其余的人。\par

% structure: p0005 source=h2 target=subsection reason=relative-heading-rank; base=h2

\subsection{二、让步的司铎应与篡位者一同降到末位，其余人保持原位。}

因此，上述司铎们，尽管他们甚至应被剥夺司铎职，但他们既已自认不配其应有等级，然而为显示宗座对他们的宽容，他们应被放在教会所有司铎之末；为使他们承担自己的判决，他们也将低于那个他们凭自己判断所优先于己的人；所有其他司铎则保持其按晋铎时间所指定的顺序。除上述两人外，任何人不得遭受尊严上的损失，但此耻辱只归于那些选择将自己置于一位刚被祝圣的资浅者之下的人：好让他们感到福音的那句话适用于他们，经上说：\BibleQuote{你们用什么判断来判断，你们也要受什么判断；你们用什么量器来量，也用什么量器量给你们。}但司铎保禄应保留其位置，他未曾以值得称赞的坚定离开此位；不得再有任何侵害他人权益的行为：这样，您，挚爱的，您并非无端地在此事中承担全部恶名，应迅速采取措施，至少通过执行我们的这些命令来治愈此事；免得再次有正当的申诉向我们提出，我们被迫采取更强烈的愤怒：因为我们宁愿通过纠正错误来恢复纪律，而不愿加重刑罚。须知我们将命令的执行托付于我们的弟兄及同僚主教\Person{尤利乌斯}，以便一切按我们所规定立即确立。发于三月八日，在显赫的\Person{波斯图米亚努斯}执政之年（448年）。\par

\SourceRecord{Translated by Charles Lett Feltoe.；Nicene and Post-Nicene Fathers, Second Series；Vol. 12.；Edited by Philip Schaff and Henry Wace.；Buffalo, NY: Christian Literature Publishing Co.,；1895.；<http://www.newadvent.org/fathers/3604019.htm>.}

% structure: p0001 source=h1 target=section reason=page-title

\section{第二十封书信}

% structure: p0002 source=h2 target=subsection reason=relative-heading-rank; base=h2

\subsection{他感谢对方提供关于聂斯多略主义复兴的信息，并称赞其热忱}

\Person{良}主教致其爱子\Person{欧提克斯}司铎。\EditorialNote{欧提克斯是\Place{君士坦丁堡}的一位隐修院院长。}\par

亲爱的，你通过来信让我们知悉，有些人活动起来，使\Person{聂斯多略}的异端再次复兴。我们回复说，你对此事的关切令我们欣慰，因为从我们所收到的言语中可以看出你的心意。因此，你不要怀疑，主、大公信仰的创立者，必将在一切事上眷顾你。而当我们能够更充分地查明此事是因谁的恶行而发生时，我们必须在天主的助佑下设法彻底根除这种早已受到谴责的毒草。愿天主保佑你平安，我亲爱的儿子。签署于显赫的\Person{波斯图米亚努斯}与\Person{芝诺}执政之年（448年）六月一日。\par

\SourceRecord{Translated by Charles Lett Feltoe.；Nicene and Post-Nicene Fathers, Second Series；Vol. 12.；Edited by Philip Schaff and Henry Wace.；Buffalo, NY: Christian Literature Publishing Co.,；1895.；<http://www.newadvent.org/fathers/3604020.htm>.}

% structure: p0001 source=h1 target=section reason=page-title

\section{第二十一封信}

\Emph{从\Person{欧提基}致\Person{良}}\par

% structure: p0003 source=h2 target=subsection reason=relative-heading-rank; base=h2

\subsection{一、他陈述了会议经过。}

天主圣言在万有之先就是我的见证，祂确知我在主基督、万有的天主内的望德和信德，并看出我在此事上持守真理的证据；但我亦呼求你的圣德为我的心、我的见解和言语的合理性作证。然而邪恶的魔鬼已对我那本应摧毁其权势的热忱和决心施加了邪恶的影响。于是，他竭尽其所有邪恶之力，挑起了多利莱城的主教欧瑟伯来反对我；欧瑟伯向君士坦丁堡教会的主教圣弗拉维安，以及他在同一城中因各自事务而聚集的若干其他人提出了一项指控。在指控中，他称我为异端者，并非提出任何真实的控告，而是为我图谋毁灭，为天主的教会制造纷扰。\par

他们的圣德们传唤我回应他的控告。但我因重病及年迈而耽搁，仍前来为自己辩白，深知有人已结党危及我的安全。事实上，我连同附有签名的上诉状，向他们递交了一份表明我对神圣信仰之宣认的声明。然而圣弗拉维安并未接受该文件，也未命人宣读，却听到我逐字答覆了神圣尼西亚大公会议所颁布、并在厄弗所大公会议所确认的信仰。他们要求我承认两个性体，并诅咒否认此说的人。但我畏惧大公会议的定断，不愿对神圣尼西亚大公会议所颁布的信仰增减一言，又知我们神圣而蒙福的教父和主教们——尤利乌斯、斐理克斯、阿塔纳修、额我略——均拒绝\Quote{两个性体}的说法，且我不敢讨论天主圣言的本性——祂在末日无改变地、按祂所愿所知的进入童贞玛利亚的胎中，真实地成为人，而非幻觉——也绝不敢诅咒上述教父们。因此我请求他们让您的圣德知晓此事，以便您判断何者合宜；我定当全然追随您的裁决。\par

% structure: p0006 source=h2 target=subsection reason=relative-heading-rank; base=h2

\subsection{二、他的解释未获理睬。}

但他们不听我说的任何话，就解散了会议，并公布了对我的革职判决——这在审理之前就已准备好。他们如此结党对我造谣中伤，以致若非天主因您圣德的转求而迅速帮助我脱离武力的袭击，我的性命也难保全。然后他们开始强迫其他隐修院的院长签字同意革除我的职务（这种事从未加诸于自称异端者，甚至对奈斯多略本人也未曾做过）。以致当我试图为安抚人心而对我的信仰加以陈明时，那些策划上述反我党争的人不仅不让我陈词被听，还立刻夺走它，好使我在众人面前成为异端者。\par

% structure: p0008 source=h2 target=subsection reason=relative-heading-rank; base=h2

\subsection{三、他向良求助。}

因此我投奔于您——信仰的捍卫者、此类党争的厌恶者，至今绝未带来任何违背最初传授给我们的信仰异说；相反，我诅咒亚波里纳里、瓦伦提诺、摩尼、奈斯多略，以及那些说我们的主耶稣基督救主的肉身是从天降下，而非由圣神及童贞女所生的人；连同从西满．玛古以来的一切异端。然而我仍如异端者一般，生命危在旦夕。我恳求您，不要因他们对我的阴谋诡计而对我存有偏见，而要按您认为合宜的，对信仰作出判决；今后更不要让这伙党争对我说任何毁谤的话；也不要让一位在节制与一切洁德中度过了七十年隐修生活的人被逐出正统者的行列，以致在生命的末刻遭受破产。我在信后附上了两份文件：一是我的控告者在会议中提出的，二是我递交但未被接受的；此外还有我的信仰声明，以及我们神圣教父们关于两个性体的决定。\par

\Emph{\Person{欧提基}的信仰宣认}\par

我在赐生命给万有的天主，及在本丢．比拉多面前作了那美好宣认的基督耶稣面前，恳求你勿徇情面。因为我与我的祖先持守同一信仰，自幼受那由全世界聚于尼西亚的三百一十八位最蒙福主教的神圣大公会议所订立的同一信仰光照，此信仰又由在厄弗所举行的神圣大公会议重新确认并批准为唯一当接受的信仰；我从未有异于正确而唯一真实的正统信仰所命令的。我同意同一神圣大公会议关于同一信仰所规定的一切；该大公会议的领袖与主持人是蒙福的记忆中的阿历山大的主教济利禄——他与那些圣者及天主所拣选者——大额我略、另一位额我略、巴西略、阿塔纳修、阿提古、普若克鲁——是同道的伙伴，在宣讲与信仰上的共享者。我视他和他们所有人皆为正统而忠信的，尊敬他们为圣人，尊崇他们为我之师。我诅咒奈斯多略、亚波里纳里及直至西满的一切异端者，以及那些说我们的主耶稣基督的肉身从天降下的人。因为祂——天主的圣言——无肉身从天降下，并在童贞女的胎中无变化、无改变地成了肉身，正如祂自己所知所愿。祂在世世代代之前常是完美的天主，在末日为我们并为了我们的救恩，也成了完美的人。求您圣德审阅我的这份完整宣认。\par

我，\Person{欧提基}，司铎兼总隐修院院长，亲手签署了此声明。\par

\SourceRecord{Translated by Charles Lett Feltoe.；Nicene and Post-Nicene Fathers, Second Series；Vol. 12.；Edited by Philip Schaff and Henry Wace.；Buffalo, NY: Christian Literature Publishing Co.,；1895.；<http://www.newadvent.org/fathers/3604021.htm>.}

% structure: p0001 source=h1 target=section reason=page-title

\section{第22封信}

\Emph{第一封：\Place{君士坦丁堡}主教\Person{弗拉维安}致教宗\Person{莱奥}。}\par

% structure: p0003 source=h2 target=subsection reason=relative-heading-rank; base=h2

\subsection{一、魔鬼的计谋使\Person{欧迪奇}误入歧途。}

致最神圣、热爱天主的父和同僚主教\Person{莱奥}，\Person{弗拉维安}在主内致候。\par

没有什么能阻止魔鬼的邪恶，那\BibleQuote{不安的邪恶，充满致死的毒素}\BibleRef{雅 3:8}。它上下\Quote{巡行}，寻找\Quote{可}击、惊吓并\Quote{吞吃}\BibleRef{伯前 5:8}。因此，警醒、祈祷、亲近天主、避免愚昧的争论、追随教父、不越过永恒界限，这些都是我们从圣经中学到的。所以我放下过度悲痛和大量的眼泪，因我手下一位圣职人员的被掠，我无法救他脱离豺狼，尽管我已准备好为他舍命。他是如何被捕获，如何跳开，憎恨呼召者的声音，背离对教父的记忆，彻底厌恶他们的道路。因此我继续我的陈述。\par

% structure: p0006 source=h2 target=subsection reason=relative-heading-rank; base=h2

\subsection{二、异端的诱惑捕获不警惕的人。}

有些人在\BibleQuote{外披羊皮，内里却是贪婪的狼}\BibleRef{玛 7:15}；凭他们的果实我们能认出他们。这些人起初似乎属于我们，但却不属于我们：\BibleQuote{因为他们若属于我们，必会留在我们中}\BibleRef{若一 2:19}。但当他们吐出他们的不敬，抛出内心的诡计，抓住较弱的人，以及那些对神圣话语尚未操练感官的人，他们就把这些人带向毁灭，曲解并侮辱教父的教义，正如他们曲解圣经以自取灭亡一样。我们必须预先警告并留心，免得有人被他们的邪恶误导，动摇其坚定。因为他们\Quote{磨利舌头如同蛇，虺蛇的毒汁在它们嘴唇之下}，如同先知曾论他们呼喊的。\par

% structure: p0008 source=h2 target=subsection reason=relative-heading-rank; base=h2

\subsection{三、\Person{欧迪奇}的异端陈述。}

因此，现在在我们中间出现了一个这样的人，\Person{欧迪奇}，多年担任司铎和隐修院院长，假装持有与我们相同的信仰，并具有正确的信德：实际上他抗拒\Person{奈斯多利}的亵渎，并假装与他辩论，但由318位圣教父所订立的信经，以及神圣的\Person{济利禄}写给\Person{奈斯多利}的信，还有同一作者就同一主题写给东方人的另一封信，这些所有同意过的著作，他却试图推翻，并复活亵渎的\Person{瓦伦提诺}和\Person{阿波利纳里}的古老邪说。他不畏惧真王的警告：\BibleQuote{谁使这些小子中的一个跌倒，倒不如把磨石拴在他的颈上，沉在海的深处。}但他抛去一切羞耻，抖掉遮盖他错误的斗篷，公开在我们神圣的主教会议上坚持说，我们的主耶稣基督在降生成人后，不应被我们理解为具有两个本性，在一个实体和一个位格内；也不应认为主的肉体与我们同性同体，如同从我们中取来并与天主圣言\OriginalTerm{位格性地}{hypostatically}结合；而他说，童贞女（生他的）确实按肉体与我们同性同体，但主自己并没有从她那里取用与我们同性同体的肉体；但主的身体不是人的身体，尽管从童贞女所出的成了人的身体，拒绝所有圣教父的阐述。\par

% structure: p0010 source=h2 target=subsection reason=relative-heading-rank; base=h2

\subsection{四、他将会议记录送给\Person{莱奥}，让他看到所有细节。}

但为避免过于冗长而详述一切，我们已将早些时候就此事所采取的会议记录送于您的圣座：在会上我们因他被证实有这些罪名而剥夺了他的司铎职、他隐修院的管理权以及我们的共融：为使您的圣座也知道他案件的事实，能使他的邪恶显明于您座下所有热爱天主的各主教；免得若他们不知道他所持且被公开定罪的观点，他们或许会通过书信或其他方式与他通信，视他为同道。我和与我同在的人向您及基督内所有弟兄致以诸多问候。愿主保佑您平安，并为我们祈祷，最热爱天主的父。\par

\SourceRecord{Translated by Charles Lett Feltoe.；Nicene and Post-Nicene Fathers, Second Series；Vol. 12.；Edited by Philip Schaff and Henry Wace.；Buffalo, NY: Christian Literature Publishing Co.,；1895.；<http://www.newadvent.org/fathers/3604022.htm>.}

% structure: p0001 source=h1 target=section reason=page-title

\section{第二十三书}

致其至爱之弟，主教弗拉维安\EditorialNote{君士坦丁堡主教}，主教良。\par

% structure: p0003 source=h2 target=subsection reason=relative-heading-rank; base=h2

\subsection{一、 抱怨弗拉维安未将欧迪克的案件详尽告知}

我至虔敬与仁爱的皇帝，因其神圣可赞的信德及对天主教会平安的关切，已将你处引起扰攘之事函告我等。兄弟，我们惊异于你竟能对此引起的丑闻缄默不语，未曾即刻设法以你之报告告知我们，使我们不致对事实真相存疑。因为我们收到了司铎欧迪克的一份文件，他控诉因尤西比乌主教的指控而被不公地剥夺了共融，尽管他说他曾应你的召叫，并未拒绝到场；且声称他在法庭上提交了上诉书，但未被接受；因此他被迫在君士坦丁堡城发出了辩护书。此事悬而未决，我们尚不知他被剥夺教会共融是否公正。但鉴于此事之重要，我们愿知你行动的理由，并让整件事情为我们所知；因为我们这些渴望主之司铎的审判慎重者，不能在未获知全部程序之前做出决定，直至我们准确掌握所有情况。\par

% structure: p0005 source=h2 target=subsection reason=relative-heading-rank; base=h2

\subsection{二、 现在要求知悉}

因此，兄弟，请通过最合适且胜任之人，以详尽报告告知我们，那需以如此严厉判决纠正的、针对古老信仰的新奇之事是什么。因为教会的温和与最虔诚君主的信德，都要求我们极为关切基督信仰之平安：以使纷争得以澄清，天主教信仰保持完整，使经过考验的信德者因我们的权威而坚固，同时使坚持错误者从谬误中召回。此事不会产生困难，因为该司铎已在其陈述中自称，若发现他有何可责备之处，愿意接受纠正。因为在此类事情上，我们宜采取一切预防措施，保持爱德，并免于争吵而捍卫真理。因此，因你看到，亲爱的，我们对如此重大之事焦虑，请尽快以尽可能详尽明了的方式告知我们一切（这本应早已做），以免因双方对立的陈述而陷入某种不确定，并使本应在萌芽时即扑灭的纷争滋长；因为我们的心被天主的默感所触动，需要保护可敬教父们那些已获神圣批准且属于信仰基础的宪章，免受任何人的曲解。愿天主保佑你平安，亲爱的兄弟。日期：449年2月18日，在显赫的阿斯图里乌斯与普罗托革内斯执政之年。\par

\SourceRecord{Translated by Charles Lett Feltoe.；Nicene and Post-Nicene Fathers, Second Series；Vol. 12.；Edited by Philip Schaff and Henry Wace.；Buffalo, NY: Christian Literature Publishing Co.,；1895.；<http://www.newadvent.org/fathers/3604023.htm>.}

% structure: p0001 source=h1 target=section reason=page-title

\section{第二十四封信}

\Person{良}主教，致\Person{德奥多西}·奥古斯都。\par

% structure: p0003 source=h2 target=subsection reason=relative-heading-rank; base=h2

\subsection{一、他赞颂皇帝的虔诚并提及\Person{欧提克}的上诉。}

主通过你的仁慈和信德赐予祂的教会多少保护，你又通过你寄给我的这封信显示出来了：以致我们喜乐于你内不仅有王者的心，也有司铎的心。看到你除了帝国的公务和公众事务外，对基督宗教怀有最虔诚的焦虑，唯恐在天主子民中产生分裂、异端或其他冒犯。因为当人们藉着宣认一个天主性事奉永恒而不变的天主圣三时，你的国度便处于最佳状态。在君士坦丁堡教会中发生了什么扰动，以致能如此激动我的弟兄兼同为主教\Person{弗拉维安}，使他剥夺了司铎\Person{欧提克}的共融，我至今未能清楚理解。因为虽然上述司铎曾用书面将其困扰寄给宗座，但他只是简要触及了一些要点，声称他遵守了尼西亚会议的法令，并因信德差异而被无端指责。\par

% structure: p0005 source=h2 target=subsection reason=relative-heading-rank; base=h2

\subsection{二、他指责\Person{弗拉维安}的沉默。}

但是他的控告者\Person{欧瑟比}主教的陈述——上述司铎曾将其副本寄给我们——关于他的反对意见没有任何清楚的内容，尽管他指控一位司铎犯有异端，却没有明确说出他反对他什么意见：虽然该主教本人也声称他坚持尼西亚会议的法令：因此我们无法了解更多情况。并且因为我们信德的方法和你的虔敬所表现出来的值得称赞的焦虑要求了解案件的实质，现在绝不允许有任何欺骗的余地，但我们必须被告知他认定他不正统的那些点，以便在充分了解情况后作出正确的判决。我已经给上述主教寄了一封信，他可以从信中看出，我对他在如此严重的事情上依然对所发生的事情保持沉默感到不满，他本应一开始就主动向我们披露一切：我们相信，即使在此提醒之后，他也会把全部情况告诉我们，以便现在看似模糊的事情被带到光明中时，可以按照福音和宗徒的教导作出判决。日期：二月十八日，在显赫的\Person{阿斯图里乌斯}和\Person{普罗托革尼斯}执政之年（449年）。\par

\SourceRecord{Translated by Charles Lett Feltoe.；Nicene and Post-Nicene Fathers, Second Series；Vol. 12.；Edited by Philip Schaff and Henry Wace.；Buffalo, NY: Christian Literature Publishing Co.,；1895.；<http://www.newadvent.org/fathers/3604024.htm>.}

% structure: p0001 source=h1 target=section reason=page-title

\section{第二十六书}

\Emph{第二封，弗拉维安致良一世。}\par

% structure: p0003 source=h2 target=subsection reason=relative-heading-rank; base=h2

\subsection{一、欧缇克斯异端重述。}

致最神圣、可敬的父及同工\Person{良}，\Person{弗拉维安}在主内问安。\par

最蒙天主所爱者，如你所知，对司铎而言，无有比虔诚与正确分施真理之言更为宝贵。因为我们的一切希望、安全及所许美物的酬报皆系于此。为此，我们必须竭尽全力维护真信仰，以及圣教父们所陈明并制定的各项规定，使之无论在何时何地均得以完整无损地保存与持守。因此，在当前情况下，我们见正统信仰遭受损害，而阿波利纳里与瓦伦提努的异端被邪恶隐修士欧缇克斯重新煽起，实属必要不可忽视，并当公开揭露以保全民信德。因为此人，欧缇克斯，将他那病态且有害的见解隐藏于心中，竟敢冲撞我们的温和，并恬不知耻地将其亵渎之言灌输给许多人，声称：我们的救主耶稣基督在降生成人之前确实有两个性体，即天主性与人性；但结合之后，二者成了单一的性体。他不知自己所说的是什么，也不知自己为何如此武断。因为，如你的虔诚所知，即使在基督内结合的两个性体，其结合过程并未使它们的属性混淆；相反，两个性体的属性在结合中仍完整保留。他还添加了另一项亵渎之言，说：由玛利亚所出之主的身躯并非出自我们的人性，亦非出于人的质料；尽管他称之为人的身躯，却拒绝承认这身躯与我们及按肉身生育他的那位同质。\par

% structure: p0006 source=h2 target=subsection reason=relative-heading-rank; base=h2

\subsection{二、欧缇克斯规避主教会议的手段。}

尽管如此，在厄弗所会议的文件中，神圣普世主教会议致邪恶而被罢免的聂斯多略的书信明确写道：\Quote{为形成真正合而为一的性体确实彼此不同；然而从二者中仅有一位基督、一个圣子。不是因结合而消除了两个性体间的差异，而是这些相同的性体，即他的天主性与人性，藉着那不可言说、不可理解、归于合一的相遇，为我们成就了唯一的主耶稣基督。}这并非你之圣德所不知，你无疑已阅过厄弗所会议记录。然而这位欧缇克斯不重视这些话，认为自身不受该神圣普世主教会议所定罚则的约束。因此，我们发现许多淳朴人士因他的争辩而信德受损，当虔诚主教欧瑟伯控告他，他出席神圣会议，并亲口向与会主教们陈述其见解后，我们因他背离真信仰而将他罢免，一如你的圣德将从关于他的决议中所知；这些决议我们随此信呈上。此外，我认为亦应告知你：这位欧缇克斯在遭受公正且合乎教规的罢免后，并未以随后的行为弥补先前之过，也未以审慎的补赎和多泪的忏悔平息天主，并以真心悔改安慰我们因他跌倒而深感忧伤的心；他不仅没有这样做，反而竭力使本地的至圣教会陷入混乱：公开张贴充满侮辱与咒骂的告示，此外还向最虔诚、爱基督的皇帝呈递祈求，这些祈求充满傲慢与无礼，企图以此完全凌驾于神圣教规之上。\par

% structure: p0008 source=h2 target=subsection reason=relative-heading-rank; base=h2

\subsection{三、他确认收到良的书信。}

但在这一切发生之后，你的圣德的书信由最尊贵的伯爵\Person{潘索斐}转交给我们；从信中我们得知同一位欧缇克斯曾给你送去一封充满谎言与诡计的书信，声称在审判时他曾向我们及当时在场的圣主教会议提交了上诉书，并向你的圣德提出上诉：他对此事从未做过，而且在这一事上，他也犯了欺骗之罪，如同谎言之父一般，妄想博得你的垂听。因此，至圣之父，我们被他所行的一切以及仍在对我们和至圣教会所做的一切所激怒，请用你身为司铎所惯有的敏捷，在保卫神圣教会的公益与和平中，惠允以你的书信批准已合乎教规对他所下的决议，并坚固我们最虔诚、爱基督之皇帝的信仰。因为此事只需你的权威与支持，藉你的智慧，将立刻带来普遍的和平与安宁。这样，藉天主的助佑，通过你的一封书信，已兴起的异端及其所激起的混乱将轻易平息；那传闻中的主教会议也将得以避免，从而普世至圣教会无需受扰。我和所有同我在一起的人向你处所有弟兄致敬。愿你在主内安然无恙，仍为我们祈祷，至敬爱天主、至圣之父。\par

\SourceRecord{Translated by Charles Lett Feltoe.；Nicene and Post-Nicene Fathers, Second Series；Vol. 12.；Edited by Philip Schaff and Henry Wace.；Buffalo, NY: Christian Literature Publishing Co.,；1895.；<http://www.newadvent.org/fathers/3604026.htm>.}

% structure: p0001 source=h1 target=section reason=page-title

\section{第二十七封信}

% structure: p0002 source=h2 target=subsection reason=relative-heading-rank; base=h2

\subsection{对\Person{弗拉维安}第一封信的确认，并承诺更详尽的答复}

\Person{良}致\Place{君士坦丁堡}主教\Person{弗拉维安}。\par

亲爱的，我们一有机会，即我们的可敬之子\Person{罗达努斯}前来，就确认收到了你的信函（这信函是要向我们报告那由错误异见在你们中间不幸挑起的事件）。此人长久以来似乎热心宗教，却对信仰作了不正当的表达，虽然他本不该背离公教传统，而应坚守众人所共持的同一信仰。然而，关于此事，我们通过那位将你的信带给我们的使者，亲爱的，作了更充分的答复：好让我们就整个事件给你一切必要的指示。因为我们不允许他坚持其错误的固执；也不允许你，亲爱的，以如此忠信的热忱抵抗他错误和愚蠢的异见，长久被敌对者的反对所困扰。我们希望你以应有的爱心接待我们的前述之子（我们托他送去这封信），并在其返回我们这里时给予答复。公元449年5月21日，\Person{阿斯图里乌斯}与\Person{普罗托格内斯}执政之时。\par

\SourceRecord{Translated by Charles Lett Feltoe.；Nicene and Post-Nicene Fathers, Second Series；Vol. 12.；Edited by Philip Schaff and Henry Wace.；Buffalo, NY: Christian Literature Publishing Co.,；1895.；<http://www.newadvent.org/fathers/3604027.htm>.}

% structure: p0001 source=h1 target=section reason=page-title

\section{第二十八封信——\WorkTitle{论集}}

\Emph{致\Person{弗拉维安}，通常称为《\WorkTitle{论集}》。}\par

% structure: p0003 source=h2 target=subsection reason=relative-heading-rank; base=h2

\subsection{一、\Person{欧提克}因狂妄无知而陷入谬误。}

亲爱的，读了你的来信（我们对这封迟到的信感到惊讶），并仔细审阅了主教们行动的详细报告后，我们终于发现了那在你们中间兴起、危及信德的纯洁的丑闻是什么：先前看似隐藏的事，如今已被揭示并展现在我们眼前。由此可见，那位曾因其司铎职务而显得值得尊敬的欧提克，实是极其鲁莽、极端无知，甚至先知论及他时所说的话也可用在他身上：\Quote{他拒绝明智，为行善；他躺在床上图谋不义。}然而，还有什么比坚持亵渎的意见，不肯向比自己更明智、更有学识的人让步更为不义的呢？现在那些因某种晦暗而无法认识真理的人，不求助先知的言论、宗徒的书信和福音的训令，只依赖自己，从而陷入了这种愚昧：他们成为谬误的导师，是因为他们从未做过真理的门徒。因为这个人，连信经的基本原理都未曾掌握，又能从新旧约的篇章中学到什么呢？而那种为普世每个即将重生的人所口口宣认的信经，这位老人心中却尚未领会。\par

% structure: p0005 source=h2 target=subsection reason=relative-heading-rank; base=h2

\subsection{二、论基督的双重诞生与本性。}

因此，他既不知道应当如何思考天主圣言的降生成人，又不愿通过研读圣经的博大精深来获得知识之光，他至少应当仔细聆听那种普遍的、一致的宣认，即全体信徒宣认他们\Quote{相信全能的天主圣父，和耶稣基督、祂的独生子、我们的主，祂由圣神和童贞玛利亚诞生。}这三句话推翻了几乎所有异端的诡计。因为不仅相信天主既是全能者又是圣父，而且圣子与祂同永恒，在各方面与圣父没有任何区别，因为祂是\Quote{出自天主的天主}，出自全能者的全能者，由永恒者所生而与祂同永恒：在时间上不晚，在能力上不低，在光荣上不异，在本质上不分；然而，永恒圣父的独生子在时间内由圣神和童贞玛利亚永恒地诞生。这个在时间内的诞生，对那神圣永恒的诞生既无减损也无添加，而是完全致力于恢复那被欺骗的人，以便他能战胜死亡，并以自己的力量推翻那掌握死亡权势的魔鬼。因为除非祂取了我们的本性并使之成为自己的，我们决不可能战胜罪恶与死亡的创始者——那位既不能被罪恶污染也不能被死亡拘禁的。无疑，祂是在童贞母亲的子宫中由圣神受孕，她生祂时无损其童贞，正如她受孕时无损其童贞一样。\par

但是，如果他不能从基督徒信仰的这个纯洁源泉中获得正确的理解（因为他用自己特有的盲目面纱遮蔽了清晰真理的光辉），他本应服从福音的教导。而玛窦提及\BibleQuote{耶稣基督、达味之子、亚巴郎之子的族谱\BibleRef{玛 1:1}}时，他也可以寻求宗徒言论所提供的教导。读到《罗马人书》中\BibleQuote{保禄，耶稣基督的仆人，蒙召为宗徒，被选拔为传天主的福音——这福音是天主先前借自己的先知在圣经中所预许的，论及祂的儿子，按肉身是生于达味的后裔\BibleRef{罗 1:1}}，他本应对先知的书卷付出忠诚的细心。发现天主对亚巴郎所作的许诺：\BibleQuote{万民都要因你的后裔蒙受祝福\BibleRef{创 12:3}}，为了消除对这后裔所指的一切疑虑，他本应跟随宗徒所说的：\BibleQuote{所许诺的话原是向亚巴郎和他的后裔说的，并不说后裔们，如同为许多，而是为那一个，就是你的后裔，那后裔就是基督\BibleRef{迦 3:16}。}依撒意亚的预言他也本应通过更仔细地留意他所说的\BibleQuote{看，一位贞女将怀孕生子，人要称祂的名字为厄玛奴耳，意思是：\InnerQuote{天主与我们同在。}}而领会。同一先知的话他也本应忠信地诵读：\Quote{有一个婴孩为我们诞生了，一个儿子赐给了我们；权势负担在祂肩上，祂的名字要称为伟大谋略的天使、奇妙、策士、全能的天主、和平之王、永世之父。}那么他就不会如此错误地说，圣言成了血肉，以致于从童贞女子宫出生的基督具有人的形状，却没有祂母亲身体的真实。或者，难道他可能认为我们的主耶稣基督不属于我们的本性，因为那被派遣到永远童贞的荣福玛利亚那里的天使说：\BibleQuote{圣神要临于你，至高者的能力要庇荫你，因此那要诞生的圣者，将称为天主的儿子\BibleRef{路 1:35}}，以为是因着童贞女的受孕是一种神圣行动，被受孕的血肉就不分享受孕者的本性？但是，那种独特奇妙而奇妙独特的诞生，不应当理解为因受造的革新而消除了祂种类的属性。因为虽然圣神赐予了童贞女生育能力，但真实的肉体是从她的肉体中接受了：\BibleQuote{智慧建造房屋\BibleRef{箴 9:1}}，\BibleQuote{圣言成了血肉，居住在我们中间}，就是说，居住在那从人取来并用更高生命的呼吸使之活起来的血肉之中。\par

% structure: p0008 source=h2 target=subsection reason=relative-heading-rank; base=h2

\subsection{三、论天主的计划和旨意，关乎圣言的降生成人。}

因此，两性及其本质在一位格内结合，各不损伤其属性：尊威取了卑微，能力取了软弱，永恒取了有死；为偿付我们境况所负的债，不可损伤的本性结合于可受苦的本性，以便——如我们的情形所需——天主与人之间的唯一中保、降生成人的基督耶稣，能依其一性而死，依其另一性而不死。如此，在真实人性的完整与圆满中，真天主诞生了：在自己的方面完全，在我们的方面也完全。所谓\Quote{我们的}，是指造物主从起初在我们内所塑造的，以及祂前来修复的。因为那欺骗者所引入、受骗的人所犯下的，在救主身上毫无痕迹。祂虽分担了人的软弱，却并未因此分享我们的过犯。祂取了奴仆的形体，而无罪恶的瑕疵，提升了人性，并未减少天主性：因为祂的\Quote{空虚自己}，使不可见者成为可见，且作为万物的造物主和主，愿成为有死的，这是怜悯的俯就，而非能力的亏损。因此，祂在保持天主的形体时，成为人；在奴仆的形体中，也成为人。因为两性各自保留其固有的特性而不丧失：正如天主的形体并未废除奴仆的形体，奴仆的形体也未损害天主的形体。因为魔鬼曾夸耀：人受其诡计欺骗，丧失了天主的恩赐，被剥夺了不朽的恩惠，接受了死亡的判决；且他在困苦中因有犯罪的同伴而得到些许安慰；并且，天主因公义原则的要求，也改变了对那曾受如此尊荣而造的人的计划。为此，需要一项隐秘的决议：使那不可改变的天主，祂的意志不会被剥夺其固有的仁慈，能以更隐秘的奥秘实现祂对我们最初的慈父关怀；并使那因魔鬼诡诈的狡猾而陷入过犯的人，不致违背天主的旨意而丧亡。\par

% structure: p0010 source=h2 target=subsection reason=relative-heading-rank; base=h2

\subsection{四、论基督双重诞生与本性的属性相互权衡。}

于是，天主子进入这世界的低处，从天上的居所降下，却未离弃圣父的光荣，以新的秩序、藉新的诞生而受生。所谓新的秩序：因为祂在自己的本性中不可见，却在我们内成为可见；那无可包容者，甘愿受限制；祂在万世之前就已存在，却开始存在于时间中；万有之主，祂遮掩了其无限的尊威，取了奴仆的形体；那不能受苦的天主，不嫌弃成为能受苦的人，且作为不朽者，甘愿服从死亡的规律。主取了母亲的性体，而未取其过犯：在主耶稣基督内，由童贞女胎中出生，其诞生的奇妙并未使其本性不同于我们。因为祂是真天主，也是真人：在这结合中，没有虚假，因为人性的卑微与天主性的崇高在此相遇。因为天主并未因显示怜悯而改变，人性也未被尊威所吞没。因为每一形性都互相合作，履行其固有之事：即圣言行属于圣言的事，血肉实行属于血肉的事。其一以奇迹闪耀，另一则承受伤害。并且，正如圣言并未停止与圣父的光荣平等，血肉也未放弃我们族类的本性。因为必须一再重申：同一位真是天主子，也真是人子。祂是天主，因为\Quote{在起初已有圣言，圣言与天主同在，圣言就是天主\EditorialNote{若 1:1}}；祂是人，因为\Quote{圣言成了血肉，住在我们中间}。祂是天主，因为\Quote{万物是藉着祂造的；凡受造的，没有一样不是藉着祂造的}；祂是人，因为\Quote{由女人所生，生于法律之下\BibleRef{迦 4:4}}。肉身的诞生是人性之显示；童贞女的生育是天主德能之明证。婴孩的幼年显露于摇篮的卑微中；至高者的伟大由天使的歌声宣告。黑落德奸诈地试图消灭的祂，与我们在最初阶段相似；但贤士们俯伏崇拜的祂，是万有之主。同样，当祂来到前驱若翰那里受洗时，为免因遮盖祂天主性的血肉帷幕而未被认出，圣父的声音从天上发出，说：\Quote{这是我的爱子，我所喜悦的\BibleRef{玛 3:17}。}因此，那被魔鬼的诡计作为人攻击的，有天使的服役作为天主服事。饥饿、干渴、疲倦和睡眠，明显是属于人的；但以五饼使五千人饱足，赐给撒玛黎雅妇人活水，喝后永不再渴，以不沉的脚行走海面，以斥责风平息浪涛，毫无疑问是天主性的。因此，略过许多其他事例：为死去的朋友动怜悯之泪，与从四日坟墓的石后以命令之声使同一朋友复活，不是同一本性的事；被钉在十字架上，且使白昼变为黑夜，震动万物；被钉子穿透，却为盗贼的信德开启天堂之门，也不是同一本性的事；说\Quote{我与父原是一体}，又说\Quote{父比我大}，也不是同一本性的事。因为虽然在主耶稣基督内，天主和人是一个位格，但二者共有的卑微之源是一个，二者共有的光荣之源是另一个。因为祂的人性，小于父，来自我们这一边；祂的天主性，与父相等，来自父。\par

% structure: p0012 source=h2 target=subsection reason=relative-heading-rank; base=h2

\subsection{五、论基督的血肉由圣经证实为真实的。}

因此，由于这同一主体（位格）在两种性体内被理解的结果，我们也读到人子从天降下，当天主子从生祂的童贞女取了肉身时。同样，天主子据说被钉十字架并被埋葬，尽管实际上不是在祂的天主性内（独生子藉此与圣父同永恒、同实体），而是在祂软弱的人性内受了这些苦。因此，我们在信经中也全体宣认天主的独生子被钉十字架并被埋葬，正如宗徒的那句话：\BibleQuote{因为若知道，决不会把光荣的主钉在十字架上。}\BibleRef{格前2:8}但是当我们的主和救主亲自藉祂的提问教导门徒的信德时，祂说：\BibleQuote{人们说人子是谁？}当他们把别人的各种意见记录下来后，祂说：\BibleQuote{但\Emph{你们}说我是谁？}我，即我是人子，你们在奴仆的形体内、在真实肉身中所看见的，你们说我是谁？于是蒙福的\Person{伯多禄}，其由天主启示的宣认注定有益于万民，说：\BibleQuote{你是基督，生活天主之子。}\BibleRef{玛16:13}他并非不配地被主称为有福，从那角石汲取了权能的稳固（他的名也表达了这点），他藉圣父的启示宣认祂既是基督又是天主子：因为接受两者之一而不要另一个，对救恩毫无裨益，相信主耶稣基督要么只是神而没有人性，要么只是人而没有神性，同样危险。但在主的复活之后（当然，那是祂真实身体的复活，因为祂复活时与祂死时和被埋葬时一样），四十天的延迟除了除去我们信德纯净的一切黑暗之外，还有什么作用呢？为此，祂与门徒交谈，与他们同住、同食，允许那些受疑惑困扰的人以勤勉而好奇的触摸来触摸祂，在门徒们关着门时进入，并藉向他们嘘气赐予他们圣神\BibleRef{若20:22}，赐予他们领悟之光，开启了圣经的奥秘。\BibleRef{路24:27}同样，祂又显示了祂肋旁的创伤、钉痕，以及祂最近受苦的一切迹象，说：\BibleQuote{看我的手和脚，就是我。摸摸我，看看，神灵没有肉和骨头，如同你们看到我有一样。}好让祂天主性和人性的属性得以认作仍然不可分离：并使我们知道圣言与肉身并无不同，以致我们也宣认同一的天主子既是圣言又是肉身。论到这信德的奥秘，您的对手\Person{欧提克斯}想必知之甚少，如果他在天主的独生子内，既没有通过祂必须受死的屈辱，也没有通过祂复活的荣耀，认出我们的人性。他也不怕蒙福的宗徒兼福音作者\Person{若望}的宣告——他说：\BibleQuote{凡承认耶稣基督在肉身内而来的，就是出于天主；凡消灭耶稣的，就不是出于天主，这就是假基督。}但所谓\Quote{消灭耶稣}，除了从祂身上剥夺人性，并用最无耻的虚构使那唯一使我们得救的奥秘失效，还能是什么呢？事实是，既然他对基督身体的性体处于黑暗之中，他也必在祂受苦的事上被同样的盲目愚弄。因为如果他不认为主的十字架是虚构的，也并不怀疑祂为拯救世界所受的惩罚是真实的，那么就让承认祂的死亡的人，承认祂的肉身吧：不要不信祂是与我们一样有身体的人，既然他承认祂能够受苦：因为否认祂真实的肉身，也是否认祂身体的苦难。因此，如果他接受基督信仰，并且不侧耳不听福音的宣讲：就让他看看被钉在木十字架上、被钉子刺穿的到底是什么性体，当被钉者的肋旁被兵士的长矛刺开时，让他明白血和水是从哪里流出的，好使天主的教会从洗礼池和圣爵得到浇灌。也让他听蒙福的宗徒\Person{伯多禄}宣告：圣神的圣化藉基督宝血的洒布而实现。并且让他不要草率地阅读同一位宗徒的话，他说：\BibleQuote{你们要知道，不是用可朽坏的东西，如金银，从你们由祖先传下来的虚妄生活中被赎出来的，而是用耶稣基督——无玷无瑕的羔羊——的宝血。}\BibleRef{伯前1:18}也不要让他抗拒蒙福的宗徒\Person{若望}的见证，他说：\BibleQuote{天主子耶稣的血，洗净我们的一切罪过。}\BibleRef{若一1:7}又说：\BibleQuote{这就是战胜世界的胜利——我们的信德。}以及\BibleQuote{这经由水和血而来的，就是耶稣基督：不仅凭着水，而是凭着水和血。并且作证的是圣神，因为圣神就是真理，原来那作证的有三个：圣神、水和血，而这三个是一致的。}这圣神即圣化之神，这血即救赎之血，这水即洗礼之水：因为三者是一，并且不可分，没有哪一个与这联系分离；因为天主教会凭此信德生存并且前进，以致在基督耶稣内，既不信仰不带真实神性的人性，也不信仰不带真实人性的神性。\par

% structure: p0014 source=h2 target=subsection reason=relative-heading-rank; base=h2

\subsection{VI. \Person{欧提克斯}错误且有害的让步；他得以恢复共融的条件；派遣代表前往东方。}

但在你们的讯问中，当欧提克回答并说：\Quote{我承认我们的主在结合前有两个性体，但结合后我只承认一个，}我感到惊讶，他如此荒谬且错误的陈述竟然没有受到审判者的批评和斥责，这样一句达到愚蠢和亵渎顶峰的话竟被允许通过，仿佛没有听到任何冒犯之言。因为说天主子在降生成人前有两个性体，其不虔敬程度等同于断言在\BibleQuote{圣言成了血肉}之后祂只有一个性体。为使欧提克不因你们未曾驳斥而认为这是正确或可辩护的意见，我们劝告你们，亲爱的弟兄，要殷切注意，倘若因天主的仁慈感动，此案得以圆满解决，他的无知心灵也能从这有害观点及其他错误中被洁净。事实上，根据议程记录，他刚开始从错误的信念中退却，当被你们的劝诫围困时，他同意说出之前未说的话，并同意接受他先前反对的信仰。然而，当他拒绝同意谴责他那亵渎性的教义时，弟兄，你们明白他仍坚持其背信，理应受到判决的宣告。但是，如果他真诚且有益地为此痛悔，并在迟来的时刻承认主教们的权威被正当地运用；或者他亲自在你们面前口手并用撤销其错误观点，那么当他悔改时向他显示的仁慈无可指摘。因为我们的主，那位真实且\BibleQuote{善牧}，祂为羊群舍弃性命，并来为拯救而非丧失人的灵魂\BibleRef{路 9:50}，愿意我们效法祂的仁慈；如此，当我们犯罪时正义约束我们，当我们回头时仁慈使我们不致被弃绝。因为那时，真正的信德才最有益地得以维护，即当错误的信仰甚至被其支持者所谴责之时。\par

为了忠诚且忠信地执行整件事，我们委派我们的弟兄朱利乌斯主教、勒纳特司铎\Place{圣克肋孟领衔堂}以及我的儿子希拉里执事代表我们。我们还联合了我们的公证人杜尔奇丢弃斯，他的信德经得起考验：确信天主的助佑将赐予我们，以致那犯错者能在其错误观点被谴责后得救。愿天主保佑你平安，亲爱的弟兄。\par

449年6月13日，最显赫的阿斯图留斯与普罗托格内斯执政期间。\par

\SourceRecord{Translated by Charles Lett Feltoe.；Nicene and Post-Nicene Fathers, Second Series；Vol. 12.；Edited by Philip Schaff and Henry Wace.；Buffalo, NY: Christian Literature Publishing Co.,；1895.；<http://www.newadvent.org/fathers/3604028.htm>.}

% structure: p0001 source=h1 target=section reason=page-title

\section{第二十九封书信}

% structure: p0002 source=h2 target=subsection reason=relative-heading-rank; base=h2

\subsection{他通知在厄弗所大公会议任命代表一事}

\Emph{致凯撒·德奥多西，最虔诚、最敬主的奥古斯都，罗马城天主教会教宗良。}\par

天主眷顾如何垂顾人间的利益，这由你出于仁慈的关怀所显示：你受天主圣神的激励，不愿有任何纷争或分歧；因为绝对合一的信德，在任何事上都不能与自身不同。因此，尽管欧提克斯，正如主教会议的记录所揭示的，已被揭露陷于愚昧无知的错误，并且本应放弃他那被正判定罪的固执己见，然而，由于你的虔诚出于对天主荣誉的热忱而热爱公教真理，决定在厄弗所举行一次主教会议的判决，以便使那在他眼中盲目的真理能传达给这位无知的老人，我派遣了我的弟兄\Person{尤利乌斯}主教、\Person{雷纳图斯}司铎和我的儿子\Person{希拉里}执事，作为我的代表，按事情的需要行事；他们将带着如此公正与仁慈的精神，以致于当整个误入歧途的错误被谴责（因为对基督徒信德的完整性毫无怀疑）时，倘若那误入歧途者悔改并恳求宽恕，他就能获得司铎仁慈的援助。因为在他寄给我们的上诉中，他保留了自己通过承诺纠正我们意见在他意见中所不认可的一切来赢得我们宽恕的权利。但是，天主教会关于主降生成人的奥迹所普遍相信和教导的，在我已寄给我的弟兄、同为主教的\Person{弗拉维安}的书信中有更详尽的阐述。日期：于显赫的\Person{阿斯塔里乌斯}与\Person{普洛托革内斯}执政之年六月十三日（四四九年）。\par

\SourceRecord{Translated by Charles Lett Feltoe.；Nicene and Post-Nicene Fathers, Second Series；Vol. 12.；Edited by Philip Schaff and Henry Wace.；Buffalo, NY: Christian Literature Publishing Co.,；1895.；<http://www.newadvent.org/fathers/3604029.htm>.}

% structure: p0001 source=h1 target=section reason=page-title

\section{第三十一封信}

\Person{良}致\Person{普尔喀丽亚·奥古斯塔}。\par

% structure: p0003 source=h2 target=subsection reason=relative-heading-rank; base=h2

\subsection{一、他提醒普尔喀丽亚她从前对教会的服务，并建议她介入欧提克斯争议。}

主曾如何藉你的仁厚保护祂的教会，我们屡次以许多标记得以验证。我们在当代对抗公教真理攻击者时，司铎付出的所有努力，也主要归功于你的荣耀；因为你从圣神的教导中学会，将你的权威在一切事上服从于祂——你藉祂的恩宠并在祂的保护下治国。因此，由于我从我的兄弟和同为主教\Person{弗拉维安}的报告得知，在君士坦丁堡教会中，\Person{欧提克斯}引发了一场针对基督信仰完整性的争议（而公会议记录全文已向我展示了整件事的确切性质），这值得你的伟大令名去消除那——在我看来更多源于无知而非智谋——的错误，免得它因愚昧者的支持而固执成见、获得力量。因为无知有时也会陷入严重错误，头脑简单者往往因不慎坠入魔鬼的陷阱：我相信，谎言之灵就是这样潜入\Person{欧提克斯}的；他幻想自己更虔诚地敬仰天主圣子的威严，竟否认祂内我们本性的真实临在，最终断定那\BibleQuote{成了血肉}的圣言整个是同一本质。正如\Person{聂斯多略}因主张基督是由祂母亲所生的纯粹之人而大大偏离真理，此人同样从公教道路偏离不少，因为他不相信我们的本体是从同一童贞女所取；他自然希望将其理解为只属于祂的天主性；以致那取了奴仆的形体、与我们相似且同形的那一位，不过是某种影像，而非我们本性的真实。\par

% structure: p0005 source=h2 target=subsection reason=relative-heading-rank; base=h2

\subsection{二、人的救恩要求基督内两性的结合。}

然而，若不信我们的主、童贞玛利亚之子是从福音所归给祂的那一血统而为人，则说祂是真实而完美的人便是无用的。因为\Person{玛窦}说：\BibleQuote{耶稣基督的诞生书：亚巴郎之子，达味之子} \BibleRef{玛 1:1}：并追溯祂人类起源的谱系，将先祖各代一直下达到\Person{若瑟}——主母曾许配给他。而\Person{路加}则一步步回溯，将祂的传承上溯至人类始祖本身，以表明第一个\Person{亚当}和最后一个\Person{亚当}具有同一本性。毫无疑问，全能的天主之子为了教导和使人成义，本可完全以祂向圣祖和先知显现的血肉形象的方式出现；例如，当祂与\Person{雅各伯}角力、交谈时，或当祂不拒绝款待、甚至食用摆上的食物时。但这些显现只是那一位的预表，奥秘预言已宣告祂将从历代先祖的血统中取人性。我们的救赎奥迹从永远就被预定，其完成并未借助任何预象，因为圣神尚未降临于童贞女，至高者的能力也尚未荫庇她；以至于\BibleQuote{智慧为自己建造房屋}，在她无玷的身体内，\BibleQuote{圣言成了血肉}；天主的形体和奴仆的形体结合于一位，时间的创造者在时间中诞生；万物藉祂而受造的那一位，在万物中间被生出。因为若新的人没有成为有罪的肉身的样式，没有取我们旧的本性，且与圣父同体而屈尊与祂的母亲同体，独自无罪而将我们的本性结合于祂自己，整个人类就仍被束缚于魔鬼的轭下，我们也无法利用得胜者的胜利——若这胜利是在我们本性之外取得的。\par

% structure: p0007 source=h2 target=subsection reason=relative-heading-rank; base=h2

\subsection{三、从两性的结合涌流圣洗的恩宠。他直接向普尔喀丽亚求助。}

但藉着基督两性奇妙的相通，重生的奥迹光照了我们：使我们这些从肉情而生的人，也能藉着同一圣神——基督就是由祂而受孕、诞生的——从灵性根源重生。因此，福音作者论及信徒时说：\BibleQuote{不是由血，也不是由肉情，也不是由人意，而是由天主而生} \BibleRef{若 1:13}。这不可言喻的恩宠，凡是从信仰中排除我们救恩主要途径的人，都无从分享，也不能被算作天主的义子。因此，我甚为烦恼和悲伤，这个人——先前在谦卑方面看似值得称许——竟敢于对我们及我们祖先的唯一信仰发动这些空洞而愚蠢的攻击。当他看到自己无知的见解冒犯了公教徒的耳朵时，他本应撤回自己的意见，而不应如此扰乱教会的统治者，以致当得定罪判决；当然，若他执意坚持己见，无人能撤销此判决。因为宗座的温和是这样使用其宽容的：对顽固者严厉处置，同时愿意向接受矫正者提供宽恕。既然我对你崇高的信德和热忱怀有极大信任，我恳请你的卓越仁慈，正如公教信仰的宣讲一直由你的圣善热忱所辅助，如今也请你维护其自由行动。或许主允许它如此受攻击，正是为此：好让我们发现那隐藏在教会内的是些什么人。显然，我们切不可忽视照料这样的人，免得我们因他们实际的丧失而痛苦。\par

% structure: p0009 source=h2 target=subsection reason=relative-heading-rank; base=h2

\subsection{四、他出席大公会议须获谅解。争议的问题非常严重。}

但最尊贵且虔诚的皇帝，急于尽快平息骚动，为主教会议指定了一个过于短促的日期，他希望会议在\Place{厄弗所}举行，并定于八月一日开会：因为从五月五日我们接到陛下的诏书之日起，剩余的大部分时间都必须用来为我所委派的能够代表我的司铎们的行程做充分准备。至于陛下所提议的我也必须亲自出席会议之必要，即使曾有先例，现在也无论如何无法实现：因为目前极其不稳定的局势不允许我离开这座大城的子民；而心怀不轨者可能会铤而走险，如果他们以为我借教会事务之机舍弃我的国家与宗座。既然陛下明白，因着您的仁厚宽容，我不应拒绝我子民的热切祈祷，这关乎公共利益，那么请视这些我派遣代我出席的弟兄为我本人与众位一起临在；我已向他们清楚详尽地说明了应当坚持的立场，鉴于详细报告以及被告亲自向我陈述已对案情作了令人满意的说明。因为问题并非关乎我们信仰的某个细小部分，尚未有明确声明；而是所提出的愚蠢异议竟敢攻击我主不愿教会中任何男女对此一无所知的事情。因为那简短而完整的公教信经，包含十二位宗徒的十二句条文，装备有天上的全副武装，以致一切异端者的见解皆可因这一武器而遭致命打击。如果\Person{欧提克斯}满足于以纯朴之心完整地接受那信经，他绝不会丝毫偏离最神圣的\Place{尼西亚}大公会议的法令，他会明白圣教父们如此规定，为的是任何心智或修辞上的机巧都不得凌驾于绝对唯一的宗徒信仰之上。那么，请陛下以您惯有的虔敬，尽力而为，将这一亵渎且愚昧的、针对人类救恩唯一圣事的攻击从众人心中驱除。如果那人本人，即陷入这诱惑的人，恢复理智，以书面悔改谴责自己的错误，就不应拒绝他与自己品级的共融。陛下当知，我也以同样口吻写信给圣\Person{弗拉维亚诺}主教：好使仁爱不被忽略，如果错误被消除。于光荣的\Person{阿斯图里乌斯}与\Person{普洛托格内斯}执政之年（449年）六月十三日。\par

\SourceRecord{Translated by Charles Lett Feltoe.；Nicene and Post-Nicene Fathers, Second Series；Vol. 12.；Edited by Philip Schaff and Henry Wace.；Buffalo, NY: Christian Literature Publishing Co.,；1895.；<http://www.newadvent.org/fathers/3604031.htm>.}

% structure: p0001 source=h1 target=section reason=page-title

\section{第三十二封书函}

% structure: p0002 source=h2 target=subsection reason=relative-heading-rank; base=h2

\subsection{他称赞他们的热忱并将他们引向\WorkTitle{大卷}}

致其爱子\Person{福斯托}、\Person{玛尔蒂诺}及众隐修院院长，主教良。\par

由于我正为那被欧提基斯试图搅扰的信仰起见，派遣\OriginalTerm{特使}{legates de latere}以协助维护真理，我认为也应该写一封信给你们，亲爱的诸位：我确实知道你们在宗教事业上如此热心，以致绝不能平静地听取这类亵渎而亵圣的言论；因为宗徒的命令铭记在你们心中，其中说道：\BibleQuote{若有人给你们宣讲的福音，与你们所接受的不同，他就当受诅咒。\BibleRef{迦 1:9}}我们也决定，上述欧提基斯的意见应被摒弃，因为我们从审阅案件记录中得知它已被正义地谴责；这样，如果那愚蠢的坚持者执迷不悟，他就可能与那些他追随其错误的人为伍。因为谁说基督没有人的——即我们的——本性，他就应被逐出基督的教会。但是，如果他藉天主圣神的怜悯而悔改，承认自己的邪恶错误，毫无保留地谴责公教徒所拒绝的一切，我们愿意不拒绝施予仁慈，以免主的教会遭受损失：因为悔改者始终可被重新接纳，惟有错误必须被排除。关于伟大虔敬的奥迹，即藉天主的圣言降生成人而成就我们的成义与救赎，按我的判断，在我已寄给我们的兄弟弗拉维安主教的信中，已根据教父的传统充分解释；这样，通过你们首领的宣言，你们可以知道，按照我们的主耶稣基督的福音，我们愿所有信友铭记在心的是什么。\Signature{公元449年6月13日，在杰出的\Person{阿斯图里乌斯}与\Person{普罗托格内斯}执政之时。}\par

\SourceRecord{Translated by Charles Lett Feltoe.；Nicene and Post-Nicene Fathers, Second Series；Vol. 12.；Edited by Philip Schaff and Henry Wace.；Buffalo, NY: Christian Literature Publishing Co.,；1895.；<http://www.newadvent.org/fathers/3604032.htm>.}

% structure: p0001 source=h1 target=section reason=page-title

\section{第三十三封书信}

% structure: p0002 source=h2 target=subsection reason=relative-heading-rank; base=h2

\subsection{一、他称赞皇帝对\Person{伯多禄}之座的求援}

\Person{良}主教，致在\Place{厄弗所}召集的神圣主教会议\par

我们最仁慈的君王的虔诚信仰，深知防止任何错谬的种子在天主教会内萌芽尤关乎他的荣耀，因而如此尊重神圣的制度，以致向宗座的权威提出请求以作适当的解决：仿佛他希望由最蒙福的\Person{伯多禄}本人宣示，在他\OriginalTerm{宣认}{confession}中受赞扬的是什么，因为主说：\Quote{人们说\OriginalTerm{人子}{Son of man}是谁？} 门徒们提到了各种人的意见；但当主问他们自己相信什么时，宗徒之长在简短的一句话中包含了信德的全貌，说：\Quote{你是基督，\OriginalTerm{永生天主之子}{Son of the living God}。} 这就是说：你真是人子，也真是永生天主之子；我说，你在\OriginalTerm{天主性}{Godhead}上真实，在\OriginalTerm{人性}{flesh}上也真实，且完全合而为一，两\OriginalTerm{性体}{natures}属性保持完整。如果\Person{欧提克斯}曾明智而彻底地相信这一点，他就绝不会从这信德的道路上退却。因为\Person{伯多禄}为自己的宣认从主那里得到了这个答复：\BibleQuote{你是有福的，\Person{西满·巴尔约纳}，因为不是血肉启示了你，而是我在天之父。我也告诉你，你是伯多禄，在这磐石上，我要建立我的教会，阴间的门不能战胜她。}\BibleRef{玛 16:17} 但是，那既拒绝蒙福的\Person{伯多禄}的宣认，又否认基督福音的人，离这建筑的合一甚远；因为他表明自己从未有任何理解真理的热忱，只拥有崇高声望的空虚外表，他没有以任何成熟的内心判断来装饰苍苍的白发。\par

% structure: p0005 source=h2 target=subsection reason=relative-heading-rank; base=h2

\subsection{二、\Person{欧提克斯}的异端应受谴责，尽管他若完全悔改，可恢复其地位。}

但是，因为即使对这样的人的医治也不可忽视，而且最基督化的皇帝虔诚而热切地希望召开一次主教会议，以便通过更完整的判断消灭一切错误，我派遣了我们的兄弟\Person{尤利乌斯}主教、\Person{雷纳图斯}司铎和我的儿子\Person{希拉里}执事，以及\Person{杜尔西提乌斯}公证人——我们已验证了他们的信德——代表我出席你们的神圣会议，弟兄们，并与你们共同决定什么是符合主旨意的。也就是说，首先应谴责那致命的错误，然后讨论他的恢复——那个如此不明智地犯错的人——但条件是：他接受真道，并用自己的声音和签名完全而公开地谴责那些使他的无知被诱入的异端意见；因为他在寄给我们的申诉中承诺了这一点，保证在一切事上服从我们的判断。在收到我们的兄弟和同为主教的\Person{弗拉维安}的信后，我们就他似乎在向我们咨询的问题给他作了相当长的回复：即，当这个似乎已兴起的错误被消除后，全世界可有同一个信德和同一个宣认，以赞美和光荣天主，并使\BibleQuote{因耶稣之名，天上、地上和地下的一切，无不屈膝，并且一切唇舌无不明认主耶稣基督在圣父的光荣中。}\BibleRef{斐 2:10} \Signature{公元449年6月13日，在显赫的阿斯图里乌斯与普罗托革内斯执政期间。}\par

\SourceRecord{Translated by Charles Lett Feltoe.；Nicene and Post-Nicene Fathers, Second Series；Vol. 12.；Edited by Philip Schaff and Henry Wace.；Buffalo, NY: Christian Literature Publishing Co.,；1895.；<http://www.newadvent.org/fathers/3604033.htm>.}

% structure: p0001 source=h1 target=section reason=page-title

\section{书信三十四}

\Person{良}主教致其爱弟\Person{犹利安}主教\EditorialNote{科斯主教}。\par

% structure: p0003 source=h2 target=subsection reason=relative-heading-rank; base=h2

\subsection{一、\Person{欧提克思}显然已背离信仰。}

你最近寄来的信函，亲爱的，显示出你对公教信仰怀有多么热切的神爱；这令我非常欣喜，因为虔诚的心都在同一意见上一致，以致按照圣神的教导，在我们身上实现了宗徒所说的：\BibleQuote{如今我藉着我们的主耶稣基督的名劝告你们，各位都说同样的话，在你们中间不要有分裂；但要同心合意，坚守同一判断。\BibleRef{格前 1:10}}但\Person{欧提克思}若坚持他的乖谬，仍不明白魔鬼捆绑他的绳索，并认为凡赞同他的无知与狂妄的人都可算在主的司铎之列，那么他就完全置身于这一合一之外。一段时间以来，我们不清楚他在何事上使公教信徒不悦；当我们没有收到我们的兄弟\Person{弗拉维安}的来信，而\Person{欧提克思}本人在其信中抱怨聂斯多略异端正在复活时，我们无法完全了解如此狡猾指控的来源或动机。但主教们的会议记录一送到我们手中，所有隐藏在他欺骗性抱怨之幕下的可憎之事便显露无遗。\par

% structure: p0005 source=h2 target=subsection reason=relative-heading-rank; base=h2

\subsection{二、他宣布从身边派遣代表。}

因为我们最仁慈的皇帝心存仁爱与虔诚，希望对一位至今似乎享有崇高声誉者的情况进行更审慎的判决，并为此认为应召开主教会议，因此经由我们的兄弟\Person{犹利乌斯}主教、\Person{雷纳特}司铎以及我以本人名义\OriginalTerm{从身边}{ex latere}派遣的儿子\Person{希拉里}执事之手，我已向我们的兄弟\Person{弗拉维安}致信一封，以适应案情需要；由此你，亲爱的，以及整个教会，可以了解那古老而唯一的信仰（这位无知的反对者曾攻击它），即我们如何持守那由天主所传授的，并如何毫不更改地宣讲它。然而，因为我们不可忘记仁慈的本分，我们认为符合我们作为司铎的温和之道：若被定罪的司铎毫无保留地改正自己，则应撤销束缚他的判决；但若他选择躺在他愚昧的泥沼中，则让判决继续，并让他与他所追随其错误的人同命运。日期：六月十三日，在显贵\Person{阿斯图里乌斯}和\Person{普罗托杰内斯}执政之年（449年）。\par

\SourceRecord{Translated by Charles Lett Feltoe.；Nicene and Post-Nicene Fathers, Second Series；Vol. 12.；Edited by Philip Schaff and Henry Wace.；Buffalo, NY: Christian Literature Publishing Co.,；1895.；<http://www.newadvent.org/fathers/3604034.htm>.}

% structure: p0001 source=h1 target=section reason=page-title

\section{第三十五封信}

\Emph{致科斯主教尤利安}\par

% structure: p0003 source=h2 target=subsection reason=relative-heading-rank; base=h2

\subsection{一、欧提克斯的异端涉及许多其他异端。}

罗马城主教良，致其爱弟尤利安主教。\par

虽然我们已通过为信仰从罗马城派遣的弟兄之手，向我们的弟兄弗拉维安送交了一份对欧提克斯极端异端的最全面驳斥，然而，由于我们通过我们的儿子巴西略收到了爱卿您的来信，这封信因其公教精神的热忱使我们非常喜悦，所以我们又添加了这页与另一文件一致的内容，使您能团结一致、奋力抵抗那些企图败坏基督福音的人，因为圣神的智慧和训导在您内与在我们内是同一的；凡不接受它的，就不是基督身体的肢体，也不能以那首级为荣——在他内，他否认自己本性的临在。那位极不明智的老者，以奈斯多利异端之名撕裂那些人的信仰，而那些人最虔诚的信德是他无法粉碎的，这对他有何益处呢？当他宣称天主的独生儿子从荣福童贞的胎中诞生时，只是穿着人身的形影，而没有真正的人性血肉与圣言结合，他偏离正道之远，正如奈斯多利将圣言的天主性与祂所取的人性实体分离一样。从这巨大的虚妄中，谁看不到会生出何等怪异的见解呢？因为谁否认耶稣基督的真实人性，必然充满许多亵渎之词，被阿波利纳里认作己有，被瓦伦提努抓住，或为摩尼所紧握——他们中没有一人相信基督内有真实的人性血肉。然而，当然，若不接受这一点，不仅否认祂——那具有天主性体却取了奴仆形象者——按血肉和理性灵魂成为人而诞生，而且也否认祂被钉、死、埋葬，第三日复活，坐在圣父右边，将在那曾受审判的身体中审判活人死人；因为如果我们不信基督具有真实而完整的人性，这些我们救赎的保证就归于无效。\par

% structure: p0006 source=h2 target=subsection reason=relative-heading-rank; base=h2

\subsection{二、在基督内可找到两种本性。}

或者，因为祂天主性的迹象是毋庸置疑的，就应假定祂有人体的证明是假的，从而接受两种本性的迹象以证明祂是创造者，却不接受它们为受造物的得救吗？不，因为血肉并未减少属于祂天主性的事物，天主性也未摧毁属于祂血肉的事物。因为祂同时是：由圣父而来永恒的，由母亲而来时间上的；力量上不可侵犯的，软弱中可受苦的；在三位一体的天主性中，与圣父和圣神同一实体，但在取了人性时，不是同一实体，而是同一[位格，以致祂在贫穷中富足，在顺从中全能，在苦难中不受苦，在死亡中不死]。因为圣言没有在任何部分变为肉或灵魂，因为天主性绝对不变的本性在其本质中永远是完整的，不受损失也不增加，如此祝福了它所取的本性，使那本性在荣耀者的位格中永远受光荣。[但是，为何圣言、血肉和灵魂成为一位耶稣基督，天主子与人之子成为一位，会显得不合适或不可能呢？既然即使没有圣言的降生成人，不同本性的血肉和灵魂也能形成一个人；那么，天主性的能力造成祂自己与人的合一，比人性自身的软弱在它自己的实体中造成合一，岂不是容易得多吗？]因此，圣言并未变为血肉，血肉也未变为圣言；但两者同存于一，一存在于两者之中，不被差异所分，不被混合所乱：祂并非因圣父而为一位，因母亲而为另一位，而是同一的，一面在万有之先由圣父所生，一面在末期由母亲所生；这样，祂是\BibleQuote{天主与人之间的中保，乃人基督耶稣，}\BibleRef{弟前 2:5} 在祂内住着\BibleQuote{天主性的一切丰盈，有形有体地住在祂内，}\BibleRef{哥 2:9} 因为被提升的是那被取的（本性），而非取者（本性），因为天主\BibleQuote{高举了祂，赐予祂那超越万名的名字，使耶稣之名，天上的、地上的和地下的万膝都屈拜，万口都宣认耶稣基督是主，归于天主圣父的光荣。}\BibleRef{斐 2:9}\par

% structure: p0008 source=h2 target=subsection reason=relative-heading-rank; base=h2

\subsection{三、基督的灵魂和身体在完全的人性意义上都是真实的，尽管祂诞生的环境是独特的。}

[但是关于\Person{欧提克斯}在主教的法庭上竟敢说的\Quote{在降生成人之前，基督内有二性，但在降生成人之后只有一性}，审判官本应用频繁而焦虑的问题迫使他对自己的承认作出解释，以免它被当作小事忽略过去，尽管它被视为出自与他其他有毒意见相同的源头。因为我认为，他说这话是因为他确信：救主所取的灵魂，在祂由童贞玛利亚诞生之前，已经居住在诸天之上，并且圣言在母胎中就将它结合于自己。但这对公教的思想和耳朵是无法容忍的：因为那从天降下的主，并没有带来任何属于我们状况的事物：因为祂既没有接受一个先前已存在的灵魂，也没有接受一个不属于祂母亲身体的血肉。毫无疑问，我们的本性并非以先被创造然后被取的方式被取，而是藉着这取本身而被创造。因此，那在\Person{奥力振}身上应受谴责的，也必须在\Person{欧提克斯}身上受惩罚，除非他愿意放弃自己的意见，即主张灵魂在被放入人体之前，不仅已有生命，而且还有不同的活动。因为虽然主按肉体的诞生具有某些特性，使其超越人类存在的通常开端，既因为祂独自由一位贞洁的童贞女受孕和诞生而无私欲偏情，又因为祂从母胎中出生时，母亲的生育能力使祂诞生而不损及童贞：然而祂的血肉并非与我们的本性不同；灵魂吹入祂内，也并非来自与所有其他人不同的源头，并且祂超越他人不是在种类的差异上，而是在能力的优越上。因为祂在祂的血肉中没有对立，[欲望的斗争也没有引起意愿的冲突]。祂的身体感官活动而无罪的律法，祂情绪的现实受祂的天主性和祂的理智控制，既不受诱惑的攻击，也不屈服于有害的影响。但真实的人与天主结合，并且就预先存在的灵魂而论并非从天降下，就血肉而论也非从无中创造：它在圣言的天主性中具有同一的位格，并在身体和灵魂上拥有与我们相同的本性。因为祂不会成为\BibleQuote{天主与人之间的中保}，除非天主和人以两种本性共存，形成一个真实的位格。主题的重大促使我们进行冗长的讨论：但对于像你这样有学识的人，无需如此详尽的论述，尤其是我们已经通过我们的代表向我们的兄弟\Person{弗拉维安}寄出了一封充分的信函，以坚固众心，不仅是司铎的心，也包括平信徒的心。我们相信，天主的仁慈将提供，使健全者能免受魔鬼的诡计，受伤者得医治，而不失去任何一个灵魂。日期：六月十三日，在显赫的\Person{阿斯蒂里乌斯}和\Person{普罗托吉尼斯}执政年（449年）。]\par

\SourceRecord{Translated by Charles Lett Feltoe.；Nicene and Post-Nicene Fathers, Second Series；Vol. 12.；Edited by Philip Schaff and Henry Wace.；Buffalo, NY: Christian Literature Publishing Co.,；1895.；<http://www.newadvent.org/fathers/3604035.htm>.}

% structure: p0001 source=h1 target=section reason=page-title

\section{书信第三十七封}

% structure: p0002 source=h2 target=subsection reason=relative-heading-rank; base=h2

\subsection{信德一致是必要的，但争议之点如此明确，几乎无需召开大公会议}

良致德奥多西皇帝。\par

收到陛下的信件后，我意识到普世教会大有理由欢喜，因您将使基督信仰——天主圣三借此受尊崇与钦崇——在任何事上既不致相异，亦不致自相矛盾。因为在呼求天主仁慈时，还有什么比全体向祂的威严献上同一声感恩与同一宣认的祭献更能有力地支持人间事务呢？其中，司铎与全体信徒的热诚终将达到圆满，只要对于天主圣言、天主的独生子为我们所作的救赎，所信的与所宣讲的，不外乎祂自己所命令的。因此，尽管诸多考虑使我无法按陛下为各地主教会议所定的日期亲自出席——因并无先例，且时代的需求不容我离开此城，尤其所争议的信仰问题如此明确，原本更有理由避免召集会议——然而，凡主惠赐助佑之处，我已尽心遵命，指派了能胜任处理此类事务、消除障碍的弟兄们代我出席。因为所产生的问题并无任何可以或应当怀疑之处。公元449年6月21日，在显赫的阿斯图里乌斯与普罗托革内斯执政时期。\par

\SourceRecord{Translated by Charles Lett Feltoe.；Nicene and Post-Nicene Fathers, Second Series；Vol. 12.；Edited by Philip Schaff and Henry Wace.；Buffalo, NY: Christian Literature Publishing Co.,；1895.；<http://www.newadvent.org/fathers/3604037.htm>.}

% structure: p0001 source=h1 target=section reason=page-title

\section{第三十八封书信}

% structure: p0002 source=h2 target=subsection reason=relative-heading-rank; base=h2

\subsection{他确认收信，并建议如果\Person{欧提克斯}愿意悔改，则应施以仁慈。}

\Person{良}致\Place{君士坦丁堡}主教\Person{弗拉维安}。\par

当我们的弟兄们已经出发，我们为信德的缘故派遣他们到你那里，我们收到了你的来信，亲爱的，由我们的儿子\Person{巴希尔}执事带来，你在信中对我们共同焦虑之事说得非常少，既因为已经到达的报告已使我们充分了解了一切，也因为私下询问可以与前述巴希尔交谈，他现在因着天主的恩宠，我们所信赖的，我们在此回复中恳请你，亲爱的，使用宗徒的话语说：\BibleQuote{在任何事上不要被敌人惊骇；这对他们是灭亡的原因，对你们却是得救的。}\BibleRef{斐 1:28} 因为有什么比否认基督降生成人的真实性而想要摧毁人类救恩的一切希望，并反驳宗徒明确说的\Quote{大哉！虔诚的奥秘，这奥秘曾显现于肉身}更悲惨的呢？有什么比为了福音的信德，对抗基督诞生和十字架的仇敌更光荣的呢？关于祂最纯净的光芒和不可战胜的大能，我们在已经送给你、亲爱的的信中，已经揭示了心中的想法：免得在我们之间，对于我们所学的、并按照天主教教义所教导的那些事，有任何怀疑。但是，既然真理的见证如此清晰有力，一个人若不立即在理性的光明中摆脱谎言的迷雾，就必被视为完全盲目和顽固；我们希望你使用恒久忍耐的药方来治愈无知的疯狂，好使那些年岁虽长，心智却如婴儿的人，通过你慈父般的劝诫，学习服从长者。如果他们放弃无知的虚妄自负，恢复理智，并且谴责他们所有的错误，接受唯一的真信仰，请不要拒绝他们主教仁慈心胸的慈悲：尽管如果他们应受谴责的不敬持续在邪恶之中，你的判决必须维持。发于7月23日，在显赫的\Person{阿斯图里乌斯}和\Person{普罗托格内斯}执政之年（449年）。\par

\SourceRecord{Translated by Charles Lett Feltoe.；Nicene and Post-Nicene Fathers, Second Series；Vol. 12.；Edited by Philip Schaff and Henry Wace.；Buffalo, NY: Christian Literature Publishing Co.,；1895.；<http://www.newadvent.org/fathers/3604038.htm>.}

% structure: p0001 source=h1 target=section reason=page-title

\section{\WorkTitle{第三十九封信}}

% structure: p0002 source=h2 target=subsection reason=relative-heading-rank; base=h2

\subsection{他斥责弗拉维安因多次回信未获答复。}

主教良致弗拉维安主教。\EditorialNote{弗拉维安是君士坦丁堡宗主教。}\par

你的缄默加深了我们的焦虑，因为自我们收到你的来信已过去许久，亲爱的弟兄。我们这些因担忧信仰辩护而分担你忧虑的人，曾多次按时机写信给你：为能用我们劝慰的帮助来支援你，使你在捍卫信仰时不向敌人的攻势屈服，并感到我们是你在辛劳中的同伴。我们相信，一段时间以前，我们的信使已到达你处，弟兄，你从我们的信件和指示中得到了充分的指导，我们也已按你的请求将巴西略送回你那里。如今，为免你以为我们遗漏了任何与你沟通的机会，我们便通过我们的儿子欧普西基乌斯（一位我们极为尊敬和爱戴的人）送交此函，请你尽快回复我们的信，并立即告知我们关于你本人和我们的代表的行动，以及整个事件的完成情况：好让我们能用更振奋的消息缓解我们此刻为捍卫信仰而感受到的焦虑。日期：光荣的阿斯图里乌斯与普罗托格内斯执政之年（449年）八月十一日。\newline{}\Signature{良}\par

\SourceRecord{Translated by Charles Lett Feltoe.；Nicene and Post-Nicene Fathers, Second Series；Vol. 12.；Edited by Philip Schaff and Henry Wace.；Buffalo, NY: Christian Literature Publishing Co.,；1895.；<http://www.newadvent.org/fathers/3604039.htm>.}

% structure: p0001 source=h1 target=section reason=page-title

\section{第四十书}

\Emph{致高卢地区阿尔勒省的各位主教。}\par

% structure: p0003 source=h2 target=subsection reason=relative-heading-rank; base=h2

\subsection{他赞同他们一致推选\Person{拉文尼乌斯}为\Place{阿尔勒}主教。}

致他钟爱的弟兄们：\Person{君士坦丁}、\Person{奥登提乌}、\Person{鲁斯提库}、\Person{奥斯皮奇}、\Person{尼西塔}、\Person{奈克塔里}、\Person{弗洛鲁}、\Person{阿斯克勒皮}、\Person{尤斯图}、\Person{奥古斯塔利}、\Person{伊南提}、\Person{赫里萨菲}，教宗良。\par

当我们得知主的司铎所做的既符合教父法规的准则，又符合宗徒的制度时，我们有正当合理的理由欢欣。因为整个教会身体必然随着健康的成长而增长，如果治理的成员在其权威的力量与和平的管理上出类拔萃。因此，弟兄们，我们以我们的认可批准你们的善行：在\Person{圣希拉里}去世后，根据圣职人员的愿望、首要市民和平信徒的一致意见，你们在\Place{阿尔勒}城祝圣了我们所认可的\Person{拉文尼乌斯}弟兄为主教。因为一个促进和平与和谐的选举，既不缺少个人功德也不缺少会众的善意，我们相信这不仅表达了人的选择，也表达了天主的默感。所以，亲爱的弟兄们，让上述司铎运用天主的恩赐，明白他所应尽的献身精神，好使他以勤恳和审慎的方式履行委托给他的职务，证明自己无愧于你们的见证，完全配得上我们的眷顾。愿天主保佑你们平安，亲爱的弟兄们。公元449年（\Person{阿斯图里乌斯}与\Person{普罗托基内斯}执政）八月廿二日。\par

\SourceRecord{Translated by Charles Lett Feltoe.；Nicene and Post-Nicene Fathers, Second Series；Vol. 12.；Edited by Philip Schaff and Henry Wace.；Buffalo, NY: Christian Literature Publishing Co.,；1895.；<http://www.newadvent.org/fathers/3604040.htm>.}

% structure: p0001 source=h1 target=section reason=page-title

\section{\WorkTitle{书信四十二}}

\Emph{致亚尔勒主教拉文纽斯。}\par

% structure: p0003 source=h2 target=subsection reason=relative-heading-rank; base=h2

\subsection{求他处置某位佩特罗尼亚努斯的欺骗行为。}

教宗良致蒙爱的主内兄弟拉文纽斯。\par

愿你谨慎小心，免得任何应受责备的僭妄行径提出不当要求：因为，一旦它以狡诈的隐秘方式找到入口，就会以所窃取的尊荣之名，扩张为更大的鲁莽。我们从你的神职人员的可靠见证得知，某个流浪无定的佩特罗尼亚努斯在你们高卢各省自称是我们的执事，并以此身份出入当地的各教会。蒙爱的兄弟，我们愿你如此遏制他可憎的厚颜无耻：揭露他的欺骗行为，警告全区的各位主教，并将他逐出所有教会的共融，免得他继续妄称。愿主保佑你平安，蒙爱的主内兄弟。\Signature{颁发于八月二十六日，显赫的阿司图留斯与普罗托格内斯执政之年（449年）。}\par

\SourceRecord{Translated by Charles Lett Feltoe.；Nicene and Post-Nicene Fathers, Second Series；Vol. 12.；Edited by Philip Schaff and Henry Wace.；Buffalo, NY: Christian Literature Publishing Co.,；1895.；<http://www.newadvent.org/fathers/3604042.htm>.}

% structure: p0001 source=h1 target=section reason=page-title

\section{第四十三封信}

% structure: p0002 source=h2 target=subsection reason=relative-heading-rank; base=h2

\subsection{一、他控诉\Person{迪奥斯库禄斯}在\Place{厄弗所}大公会议上的行为}

致最光荣、最宁静的皇帝\Person{狄奥多西}，\Person{良}主教。\par

我们自始已在所举行的会议中，从至圣的宗徒之长\Person{伯多禄}获得了如此自由的发言权，既能为了平安而维护真理，又能不容许任何人动摇其坚定地位，而立即击退破坏。既然你下令为\Person{弗拉维亚诺}在\Place{厄弗所}城召开的主教会议，对信德本身造成了损害，并给所有教会带来创伤——；此事并非透过不可靠的传信者，而是由我们派出的最可敬的主教们及我们最忠实的执事\Person{希拉留}亲自告知我们，他们向我们叙述了所发生的事。这些事件应归咎于与会者的过失，因为他们没有照惯例以纯洁的良心和正确的判断就信德及偏离信德者作出明确声明。因为我们得知，并非所有应当与会的人都聚集到会议中，有些人被驱逐，另一些人被接纳；他们因上述司铎的阴谋而被诱骗签署了一项不敬神的文件。因为他所造成的这一宣言具有伤害所有教会的性质。因为当我们派去的人看到它极端不敬神且敌对信德时，他们便通知了我们。\par

% structure: p0005 source=h2 target=subsection reason=relative-heading-rank; base=h2

\subsection{二、他请求他恢复古老的天主教教义。}

因此，最爱好和平的君王，请为信德的缘故恩准将这一危险从你敬天的良心中除去，不要让人间的僭越对基督的福音施以暴力。本着我和与我同在的主教们真诚的愿望，即愿你，最基督徒、最可敬的君王，首先在一切事上中悦天主——整个教会同心合意为你的帝国向祂献上祈祷——我劝告你，恐怕如果我们对如此重大的事保持沉默，我们将在基督的审判台前受到惩罚。因此，我在同一天主神性的不可分的天主圣三（这圣三因这些恶行而受损害，且是你王国的守护者）并在基督的圣天使面前恳求你：所有事物保持审判前的原状，等待大公会议更权威的决定——全世界所有主教聚集一起；并且不要让你自己承担他人恶行的重担。我们因紧迫需要的畏惧而被迫使如此明白地说出来。但要常将真福伯多禄的光荣，所有宗徒与他共享的冠冕，以及殉道者的喜乐放在眼前——他们受苦的动机无非是对真神性的宣认和在基督内的完美持守。\par

% structure: p0007 source=h2 target=subsection reason=relative-heading-rank; base=h2

\subsection{三、并请求召集另一次大公会议。}

现在，既然这一宣认正受到一些人的不敬神攻击，我们辖区的所有教会和所有司铎都含泪恳求你的仁慈，按照\Person{弗拉维亚诺}在他的上诉中提出的请求，命令在\Place{意大利}召集一次特别大公会议，以便排除或平息一切反对，使信德没有偏离或含糊；所有东方省份的主教也应前来，以便任何人若已偏离真理之路，可通过有益的疗法被召回忠诚，而受更严重指控的人或通过劝导归顺，或从唯一的教会中被剪除。这样，我们既须遵守尼西亚教会法的规定，也须遵守全世界主教们按照天主教会习惯所规定的定义，并维护我们祖传信德的自由——你的安宁基于此。因为我们祈求，当那些伤害教会的人被驱逐，你的各省享有正义，对这些异端者执行了刑罚时，你的王权也得以由基督的右手护卫。\par

\SourceRecord{Translated by Charles Lett Feltoe.；Nicene and Post-Nicene Fathers, Second Series；Vol. 12.；Edited by Philip Schaff and Henry Wace.；Buffalo, NY: Christian Literature Publishing Co.,；1895.；<http://www.newadvent.org/fathers/3604043.htm>.}

% structure: p0001 source=h1 target=section reason=page-title

\section{第四十四封书信}

\Person{良}主教与在\Place{罗马}集会的圣会议，致\Person{狄奥多西·奥古斯都}。\par

% structure: p0003 source=h2 target=subsection reason=relative-heading-rank; base=h2

\subsection{一、他揭露了\Place{以弗所}会议的无耻性质。}

从陛下致蒙福的宗徒\Person{伯多禄}之座的信函中，我们获得了对您维护真理与和平的如此信心，以致于我们认为，在如此清晰且井然有序的事情上，不会发生任何有害的事；尤其是，当那些被派往您命令在\Place{以弗所}召开的主教会议的代表们，得到了如此充分的指示，以至于，如果\Place{亚历山大里亚}的主教允许将他们带给圣会议或主教\Person{弗拉维安}的信函在主教们耳中宣读，那么藉着最纯洁信仰的宣言——这宣言是神圣所默启、我们既领受又持守的——一切争论的喧嚣都将完全平息，以至于无知再也不能自娱，嫉妒再也找不到行恶的机会。但是，因为在宗教的掩盖下，私利得到了考虑，少数人的不忠造成了必将伤害整个教会的事。因为，我们从一位最忠实的目击者——我们的执事\Person{希拉里}那里得知了所行之事，他为了避免被迫签署他们的决议而历尽艰辛逃脱——我们获悉，许多司铎聚集在会议中，其人数无疑会帮助辩论和决定，如果那自诩为首位的人同意保持司铎的节制，以便按照惯例，在所有人都自由表达意见之后，经过安静而公正的审议，所颁布的既能符合信仰，又有助于那些在谬误中的人。但据我们被告知，所有来的人都没有出席实际的决定：因为我们得知，有些人被拒绝，而另一些人被接纳，他们应前述司铎的要求，屈服于一项不义的签署，知道若不服从他的命令就会遭受损害，并且他提出了一项这样的决议，以致于在攻击一个人的同时，他对整个教会发泄了他的狂怒。我们宗座代表看到这是如此亵渎且违背公教信仰，以至于没有任何压力能迫使他们同意；因为在同一会议中，他们像应该做的那样坚决抗议，宗座永远不会接受正在通过的东西：因为整个基督徒信仰的奥秘被绝对摧毁了（愿上智在你皇恩的统治下避免此事），除非这超越一切先前亵渎的可憎邪恶被废止。\par

% structure: p0005 source=h2 target=subsection reason=relative-heading-rank; base=h2

\subsection{二、并恳请皇帝帮助撤销他们的决定。}

但是，因为魔鬼以邪恶的诡计欺骗不谨慎的人，并以一种虔诚的外表如此愚弄一些人的轻率，以至于说服他们去做有害而非有益的事，我们恳求陛下，弃绝一切在这危及宗教与信仰中的同谋，并在处理神圣事物时，赐予您法律公正所允许的世俗事务中同样的东西，以使人的僭越不致对基督的福音施暴。看，最基督徒与受尊敬皇帝，我与我的同僚司铎，向您尊敬的仁爱履行真诚爱德的职责，并渴望您事事中悦天主——教会为您祈祷——以免我们在主基督的审判台前因沉默而被判有罪——我们恳求您在一体不可分的天主圣三面前（这一行为冒犯了祂，因为祂是您帝国的主保与创造者），并在基督的圣天使面前，命令一切恢复到决定前的状态，直到从世界各地聚集更多的司铎。不要让自己因他人的罪而担负重担，因为（我们不得不说），我们害怕那位其宗教正被摧毁的天主被激怒。请您将蒙福的\Person{伯多禄}的荣耀、所有宗徒与他共享的冠冕、以及所有殉道者的棕榈枝放在眼前，并以您全部的心神虔敬地凝视他们，他们受苦的唯一理由就是在基督内宣认真天主与真人。\par

% structure: p0007 source=h2 target=subsection reason=relative-heading-rank; base=h2

\subsection{三、他要求在\Place{意大利}召开一次大公会议。}

并且，因为这一奥秘如今正被少数无知之人不虔敬地反对，我们各地的所有教会及所有司铎，含泪哀叹，看到我们的代表忠实地抗议，且主教\Person{弗拉维安}以书面方式向他们提出了上诉，恳求陛下下令在\Place{意大利}举行一次大公会议，这会议将以这样的方式解决或平息一切争论，使得信仰中不再有任何疑问、爱德中不再有分裂。当然，东方各省的主教必须前来，如果他们中有人因威胁与伤害而屈服、偏离了真理之道，可由康复性的措施完全恢复；而那些本人情况更艰难的人，如果他们同意更明智的劝告，可能不会从教会的合一中跌落。在上诉提出之后，这一请求有多么必要，由\Place{尼西亚}全体主教通过的教规法令（附于下方）可以证明。请按照您自己及您祖先的惯例，恩待公教徒。赐予我们如此捍卫公教信仰的自由，以致没有任何暴力、任何世俗的恐惧，在您尊敬的仁爱安全之时，能够夺走它。因为我们所辩护的不仅是教会的事业，也是您王国与繁荣的事业，使您可以和平地统辖您的省份。捍卫教会以不可动摇的和平对抗异端，好使您的帝国也由基督的右手来捍卫。日期：十月十三日，在显赫的\Person{阿斯图里乌斯}与\Person{普罗托革内斯}执政之年（449年）。\par

\SourceRecord{Translated by Charles Lett Feltoe.；Nicene and Post-Nicene Fathers, Second Series；Vol. 12.；Edited by Philip Schaff and Henry Wace.；Buffalo, NY: Christian Literature Publishing Co.,；1895.；<http://www.newadvent.org/fathers/3604044.htm>.}

% structure: p0001 source=h1 target=section reason=page-title

\section{\WorkTitle{第45封信}}

主教\Person{良}、以及在\Place{罗马}城聚集的圣主教会议，致\Person{普尔喀丽亚}奥古斯塔。\par

% structure: p0003 source=h2 target=subsection reason=relative-heading-rank; base=h2

\subsection{I. 他送去一封未能到达她手中的前信副本。}

如果先前经由我们圣职人员之手呈递至陛下圣驾的、关于信仰的信函确实抵达了您处，那么主帮助之下，您无疑早已能够为那些违背信仰所行之事提供补救。因为您何时曾辜负过司铎、宗教或基督的信仰？但当那些被派遣之人受到如此完全的阻碍，无法抵达陛下尊前，以致只有我们的一位执事\Person{希拉里}艰难逃脱并返回时，我们觉得有必要重写我们的信函；并且，为使我们的恳求能获得更多重视，我们附上了那份未送达陛下的文件副本，比从前更恳切地请求您维护那您所精通的宗教——您愈是尊贵，便愈需要对付的罪行之严重性，您的皇家信仰要求您采取行动，以免基督教信仰的完整受人计谋的侵犯。因为那些认为需要藉由在\Place{厄弗所}召开的主教会议来解决和治愈的事情，非但没有平息纷争，反而导致了更大的和平扰乱，而且（更令人痛惜的是）甚至推翻了那使我们成为基督徒的信仰本身。\par

% structure: p0005 source=h2 target=subsection reason=relative-heading-rank; base=h2

\subsection{II. 他也送上一封致皇帝的信函副本并解释其内容}

那些被派遣之人——其中一位逃脱了那凡事自专的\Place{亚历山大}主教的暴行，忠实地向我们报告了主教会议上所发生的事——他们按照身份反对（我姑且称之为）一个人的狂乱而非判断，抗议说那些藉暴力与恐惧推行之事不能推翻教会的奥迹和宗徒们所撰写的信经本身，并且任何伤害都不能使他们脱离那从蒙福宗徒\Person{伯多禄}之座完整陈述并阐述后带至圣主教会议的信仰。既然在主教请求下不允许宣读这份声明，为的是（实在是）藉着拒绝那曾使宗祖、先知、宗徒及殉道者得冠冕的信仰，来推翻我们的主耶稣基督按肉身的诞生以及宣认祂真实的死亡与复活（我们颤栗地说出这一点），我们根据我们的能力就此事写信给我们最荣耀的、而且（更重要的）是我们信奉基督的君主，同时附上一封给您的信函副本，使他决不允许那信仰——他是在其中藉天主的恩宠而重生并执政的——被任何创新所腐化，既然\Person{弗拉维安}主教仍与我们全体共融，并且那无视正义、违反所有教规教导而行之事，无论如何都不能被视为有效。由于\Place{厄弗所}主教会议没有消除反而增加了分歧的丑闻（我曾请求他）指定一个在\Place{意大利}境内召开会议的地点和时间，暂停双方所有争执与偏见，以便更仔细地重新考虑所有引起冒犯之事，并在不伤害信仰、不损害宗教的情况下，使那些因犹豫而被迫签署的司铎们能返回基督的平安，并只除去他们的错误。\par

% structure: p0007 source=h2 target=subsection reason=relative-heading-rank; base=h2

\subsection{III. 他请求她协助向皇帝陈情。}

为使我们堪当获得此事，愿您那久经考验的信德与保护——您一直在劳苦中帮助教会——俯就藉着蒙福宗徒\Person{伯多禄}的特殊委托，向我们最仁慈的君主推进我们的请求；好在这内战和破坏性的战争在教会内获得力量之前，他能藉天主的助佑赐予恢复合一的机会，深知他帝国的力量将因他善意所影响的每一份公教自由的扩展而增长。\par

颁于公元449年10月13日，在显赫的\Person{阿斯图里乌斯}和\Person{普罗托格内斯}执政官年份。\par

\SourceRecord{Translated by Charles Lett Feltoe.；Nicene and Post-Nicene Fathers, Second Series；Vol. 12.；Edited by Philip Schaff and Henry Wace.；Buffalo, NY: Christian Literature Publishing Co.,；1895.；<http://www.newadvent.org/fathers/3604045.htm>.}

% structure: p0001 source=h1 target=section reason=page-title

\section{书信五十二}

\Emph{从居鲁士主教德奥多勒致良。（见本丛书第三卷，第293页。）}\par

% structure: p0003 source=h2 target=subsection reason=relative-heading-rank; base=h2

\subsection{一、如果\Person{保禄}曾求助于\Person{伯多禄}，那么普通信徒更应求助于他的继承人。}

致\Place{罗马}主教\Person{良}。\par

如果\Person{保禄}，真理的宣告者，圣神的号角，曾求助于伟大的\Person{伯多禄}，以便为那些在\Place{安提约基雅}就遵守法律问题争论的人从他获取决定，那么我们这些卑微渺小的人更应投奔宗座，从你那里获得医治教会创伤的药方。因为你在一切事上理应居首位，因为你的宗座拥有许多特殊的特权。因为其他城市因大小、美丽或人口而扬名，有些缺乏这些优势的城市因某些神恩而得到补偿，但你的城市从一切美善的施与者那里获得了最丰满的美物。因为她是所有城市中最伟大、最著名的，是世界的主母，人口稠密。此外，她还创建了一个至今仍然主导的帝国，并将她自己的名号强加于其臣民。但她主要的装饰是她的信德，\Person{保禄}宗徒对此是一个确实的见证，当他高呼说：\Quote{你们的信德传遍了天下\BibleRef{罗 1:8}；} 如果她在接受救恩福音的种子之后立即结出如此奇妙果实的重量，那么有什么话足以表达如今在她内所找到的虔诚呢？她也拥有我们共同的父亲和真理导师、\Person{伯多禄}和\Person{保禄}的坟墓，以光照信徒的灵魂。这一对蒙福而神圣的夫妇确实在东方升起，向各方放射光芒，但自愿在西方经历生命的日落，从那里如今却光照全世界。这些使你的宗座如此荣耀：这是你所有美物中的首要。而他们的宗座仍然因他们天主的临在之光而蒙福，因为祂在其中安置了你的圣座，以传播唯一真信德的光芒。\par

% structure: p0006 source=h2 target=subsection reason=relative-heading-rank; base=h2

\subsection{二、他赞扬\Person{良}对摩尼教徒的热忱，以及后来对欧提基主义的热忱，尤其是在\WorkTitle{书函}中所证明的。}

的确，关于此事，虽然可以找到许多其他证据，但你对臭名昭著的摩尼教徒的热忱足以证明，这一热忱是你的圣座近年所显示的，从而揭示了你在神圣之事上对天主的虔诚之深。同样，你现今所写的也足以证明你的宗徒品格。因为我们看到了你的圣座关于我们的天主及救主降生成人所写的，并钦佩那作品的认真细致。因为它同样很好地证明了两个要点，即：从永生之父所生的独生子的永恒天主性，以及同时祂从\Person{亚巴郎}和\Person{达味}的后裔所取得的人性，以及祂接受了一个在各方面都与我们相似的本性，除了一点：祂完全免除一切罪过，因为罪不是生于本性，而是生于自由意志。你的信中也包含了这一点，即：天主的独生子是唯一的，祂的天主性不受苦难、不可逆转、不可改变，如同生了祂的父和全体至圣圣神一样。既然天主性不能受苦，祂便取了能受苦的本性，为的是藉着祂自己肉体的苦难，使信仰祂的人免于苦难。这些要点及所有与之相关的内容，信中都包含了。我们钦佩你的属灵智慧，赞颂藉着你说话的圣神的恩宠，并祈求、恳求你的圣座来救援如今遭遇风暴的天主教会。\par

% structure: p0008 source=h2 target=subsection reason=relative-heading-rank; base=h2

\subsection{三、他控诉\Person{迪奥斯哥鲁}对自己的苛待。}

因为当我们期望通过你的圣座派往\Place{厄弗所}的人平息风浪时，我们反而陷入了更糟的风暴。因为那位最正直的\Place{亚历山大}首席主教并不满足于非法且极其不义地将主的至圣且爱慕天主的\Place{君士坦丁堡}主教\Person{弗拉维安}免职，同样其他主教的惨遭罢黜也未平息他的愤怒。而且他还用笔在我不在场的情况下杀死了我，既未传唤我受审，也未亲自对我进行审判，更没有就我对我们的天主及救主降生成人的立场加以质询。但即使是杀人犯、破坏坟墓者、夺人妻室者，审判者也不会定他们的罪，除非他们自己通过口供证实指控，或者被他人明确证实有罪。然而我们，在相距三十五天的行程之外，他尽管深受神律培养，却凭自己的意愿定了我们的罪。不仅如此，他不仅在现在这样做，而且在去年，当两个染有阿波利拿里派病症的人来到这里，对我们提出虚假指控之后，他在教堂里站起来，诅咒了我们，而当时我已写信给他，并在信中表达了我的立场。\par

% structure: p0010 source=h2 target=subsection reason=relative-heading-rank; base=h2

\subsection{四、这种苛待发生在他于\Place{居鲁士}教区二十年善工之后。}

我悲叹教会的苦难，渴望它的和平。因为藉着你们的祈祷，我治理了宇宙的天主托付给我的教会二十年，无论是在真福\Person{德奥多托}（东方总督）的时代，还是在那些继承他\Place{安提约基雅}宗座的人的时代，我从未受到丝毫责备，反而，天主的恩宠与我合作，使一千多个灵魂从\Person{玛尔基翁}的病症中解脱出来，并使许多其他灵魂从\Person{亚略}和\Person{欧诺米}的同党归向主\Person{基督}。我必须牧养八百座教堂：因为那是\Place{居鲁士}教区的堂区数目，在其中，藉着你们的祈祷，没有一根杂草存留。但我们的羊群已从一切异端谬误中得解放。那洞察万物的天主知道我如何被派来反对我的臭名昭著的异端者用石头击打，以及我在东方的许多城市中与希腊人、犹太人和一切异端谬误进行过何等斗争。而在这一切辛劳和困难之后，我不经审讯就被定罪了。\par

% structure: p0012 source=h2 target=subsection reason=relative-heading-rank; base=h2

\subsection{五、他满怀信心地向宗座上诉。}

然而，我等待着您的\OriginalTerm{宗座}{Apostolic See}的裁决，并恳求并祈求您的圣座，在我向您公正正义的法庭上诉时，救助我，并命我来到您面前，表明我的教导遵循宗徒的踪迹。因为我有一些著作，有些是大约二十年前写的，有些是十八年前，有些是十五年前，有些是十二年前，有些反对亚略派信徒和\OriginalTerm{欧诺米派信徒}{Eunomians}，有些反对犹太人和希腊人，有些反对\OriginalTerm{波斯术士}{Magi}，有些也关于普世眷顾，其他关于天主的本质和天主的降生成人。我藉着天主的恩宠，解释了宗徒著作和先知言论，从中容易看出，我是否坚定不移地持守了信德的标准，还是偏离了它的正直道路。我恳求您不要拒绝我的请求，也不要忽视堆在我可怜的白发上的侮辱。\par

% structure: p0014 source=h2 target=subsection reason=relative-heading-rank; base=h2

\subsection{六、他应当默从于对他的撤职吗？}

首先，我恳求您告诉我，我是否应当默从于这不义的撤职。因为我等待着您的裁决，如果您命我顺服于对我的判决，我就顺服，今后不再烦扰任何人，只等待我们天主和救主无误的裁决。我实在，主天主是我的见证，我毫不关心荣誉和荣耀，只关心放在人路上的绊脚石。因为许多较单纯的民众，尤其是那些曾被我们从各种异端中拯救出来的人，会相信那些谴责我们的人，或许会把我们视为异端者，因为他们不能分辨教义的确切真理。并且因为，在我漫长的主教职任之后，我既没有获得房屋，也没有田地，没有\OriginalTerm{欧波尔}{obol}，没有坟墓，只有自愿贫困，甚至立即将从我父辈死后得来的财产分施了，所有生活在东方的人都知道。\par

% structure: p0016 source=h2 target=subsection reason=relative-heading-rank; base=h2

\subsection{七、他本人受阻，便差遣代表为他辩护。}

在一切事之前，我恳求您，圣善而蒙天主爱戴的兄弟，助成我的祈祷。我已将这些事呈报于您的圣座，经由最虔诚而蒙天主宠爱的司铎，\Person{希帕提乌斯}和\Person{阿布拉米乌斯}，\OriginalTerm{乡村主教}{chorepiscopi}，以及\Person{阿利皮乌斯}，我们地区隐修院的监督：因我被皇帝的阻留信所阻碍，不能亲自前来，其他人也一样。我恳求您的圣座，以慈父之心眷顾他们，以真挚的仁慈倾听他们，并认为我那被诽谤和妄加攻击的处境值得您的保护，尤其要全力捍卫那如今正遭阴谋反对的信德，为众教会保持教父传承的完整，如此您的圣座必将从慷慨之主那里得到圆满的赏报。\EditorialNote{日期约为449年底。}\par

\SourceRecord{Translated by Charles Lett Feltoe.；Nicene and Post-Nicene Fathers, Second Series；Vol. 12.；Edited by Philip Schaff and Henry Wace.；Buffalo, NY: Christian Literature Publishing Co.,；1895.；<http://www.newadvent.org/fathers/3604052.htm>.}

% structure: p0001 source=h1 target=section reason=page-title

\section{第五十六封书信}

% structure: p0002 source=h2 target=subsection reason=relative-heading-rank; base=h2

\subsection{\Person{加拉·普拉西狄亚}·奥古斯塔致\Person{狄奥多西}}

给吾主\Person{狄奥多西}，常胜皇帝，她的永远尊敬的皇儿，\Person{加拉·普拉西狄亚}，最虔诚与最昌盛，永久的奥古斯塔及母亲。\par

当我们抵达这座古老城市之际，正于最蒙福的宗徒\Person{伯多禄}殉道者的祭台前献上我们的敬礼，最可敬的\Person{利奥}主教在礼仪后稍作停留，向我等悲叹公教信仰，并同样以我们刚刚亲近的宗徒之长作证，身边围绕着许多他凭其地位之权威与尊严从\Place{意大利}众多城市召集而来的主教，他还在话语中加上眼泪，吁请我们将我们的哀叹与他的哀叹联合。因为从那些事件中产生了不小的损害，以致自我们最神圣的父亲\Person{君士坦丁}——他是皇宫中首位以基督徒身份出现者——的时代以来长久守护的公教信仰的准则，近日被一个人的僭越所扰乱，此人在\Place{厄弗所}召开的会议中被指控反而煽动仇恨与纷争，以士兵在场恐吓\Place{君士坦丁堡}主教\Person{弗拉维亚诺}，因他曾通过最可敬的\Place{罗马}主教所委派的代表之手（这些代表一贯按\Place{尼西亚}大公会议的规定如此出席），向宗座及本地区所有主教发出上诉，最神圣的吾主、吾子及敬爱的皇帝啊。为此，我们恳求陛下的仁慈，以真理反对这样的纷扰，并命令保持公教信仰的准则无玷，以便依照宗座的准则与决定——我们同样尊崇其为至高无上的——\Person{弗拉维亚诺}在其司铎职位上完全不受伤害，并将此事呈交宗座的大公会议，那确然是首先享有首席权者，就是被视为配得接受天堂的钥匙的那一位；因为我们在一切事上都应当维持对这座普世之主的伟大城市的敬重；并且我们也必须最谨慎地确保这一点，即我家从前所守护的，在我们时代似乎不被侵犯，且通过当前的例子，不在主教们或至圣教会之间助长分裂。\par

\SourceRecord{Translated by Charles Lett Feltoe.；Nicene and Post-Nicene Fathers, Second Series；Vol. 12.；Edited by Philip Schaff and Henry Wace.；Buffalo, NY: Christian Literature Publishing Co.,；1895.；<http://www.newadvent.org/fathers/3604056.htm>.}

% structure: p0001 source=h1 target=section reason=page-title

\section{第五十九封书信}

\Emph{致\Place{君士坦丁堡}城的神职人员和信众。}\par

% structure: p0003 source=h2 target=subsection reason=relative-heading-rank; base=h2

\subsection{一、他祝贺他们勇敢地抵抗错误。}

\Person{良}主教致居住在\Place{君士坦丁堡}的神职人员、显贵和信众。\par

虽然我们对最近在\Place{厄弗所}司铎会议中所报告的事甚感悲伤，因为正如一贯传闻且也由结果所证明，那里既未遵守应有的节制，也未遵守信仰的严格，然而我们在你们的虔诚和圣洁子民的欢呼中喜悦，这些例子已引起我们的注意，我们赞同你们所有人的正确情感；因为良善子民对卓越父亲应有的孝爱仍然存留和居住，你们不容许公教教导的圆满在任何部分被败坏。因为毫无疑问，正如圣神已向你们启示，那些否认天主的独生子取了我们的真实人性，并坚持祂所有身体的行为都是虚假幻影的行为之人，是结盟于摩尼派异端的。并且唯恐你们在任何事上同意此亵渎，我们现在已由我的儿子\Person{埃皮法尼乌斯}和\Place{罗马}教会公证人\Person{狄奥尼修斯}送给你们劝勉信件，在其中我们自愿向你们提供了你们所寻求的援助，以使你们不会怀疑我们倾注所有父亲般的关怀给你们，并藉天主仁慈的帮助，竭尽全力摧毁无知和愚昧之人所树立的所有绊脚石。任何人若被证实持有如此可憎的亵渎意见，就不可炫耀其司铎尊位。因为若无知在平信徒中似乎难以容忍，那么在治理者中，它就更不可原谅或可宽恕了；尤其当他们竟敢为其虚假和乖谬的观点辩护，并用威吓或哄骗说服不坚定者同意他们时。\par

% structure: p0006 source=h2 target=subsection reason=relative-heading-rank; base=h2

\subsection{二、凡否认基督血肉之真实者，即每一位在圣体圣事中领受者所不断宣认之真理，都当被弃绝。}

让这样的人被基督奥体的圣洁肢体所弃绝，不要让公教的自由承受不忠信者的轭。因为他们被算作在天主恩宠之外和人类救恩的奥秘之外，那些否认基督内我们血肉的本性的人，他们反对福音，反对信经。他们也没有认识到他们的盲目使他们陷入如此深渊，以致在主的苦难或复活的真实中都没有稳固的立足点：因为如果我们的血肉本性不被相信在祂内，这两者在救主内都不可信。他们迄今一定在多么无知的黑暗、多么彻底的惰怠中躺卧，没有从听闻中学习，也没有从阅读中理解，那在天主教会中不断地在人们口中，甚至在婴儿的舌头也不沉默于圣体圣事礼仪中基督身体和圣血的真理？因为在那神秘分施属神滋养的过程中，所给予和所接受的具有这样的性质：领受天上食物的德能，我们进入祂的肉中，祂成为我们的肉。因此，为坚固你们——可爱的——在你们值得称赞的忠于真理抵抗仇敌中，我将充分而合时宜地使用宗徒的语言和情感，说：\BibleQuote{因此，我一听见你们对主耶稣的信德和对众圣徒的爱德，就不断为你们感谢天主，在祈祷中纪念你们，使我们主耶稣基督的天主，荣耀的父，赐给你们智慧和启示的神恩，使你们认识祂，光照你们的心眼，使你们知道祂的召唤有什么希望，在圣徒中祂荣耀的基业有何等丰盛，以及祂对我们这些相信的人所施展的、根据祂大能的力量，这力量何等超越，祂已在基督内运行，使祂从死者中复活，并让祂坐在自己的右边，远超一切执政的、掌权的、大能的、统治的及一切可称呼的名字，不仅在今世，也在来世：并将万有置于祂脚下，使祂作教会万有之元首，教会是祂的身体，是那充满万有者的圆满。}\BibleRef{弗 1:15}\par

% structure: p0008 source=h2 target=subsection reason=relative-heading-rank; base=h2

\subsection{三、真天主和真人在基督内结合。}

在本文中，让真理的敌对者说，全能圣父何时或按什么本性高举了祂的圣子于万有之上，或使万有服从于什么实体？因为圣言的天主性在一切方面都与圣父相等、与圣父同质，生者与受生者的权能始终是同一且不变的。当然，万有的创造者，既然\BibleQuote{万物是藉着祂造的；凡受造的，没有一样不是藉着祂造的}\BibleRef{若 1:3}，祂当然超越祂所创造的一切，祂所造之物也从未不服从于它们的创造者，祂永恒的属性就是：除了出自圣父，与圣父无任何差异。如果祂被赐予更大的权能、更崇高的尊严、更超越的卓越，那么这位被提升者就比那提升祂者为小，并且祂从祂的圆满中领受时，并不拥有本性的圆满。但如此思想的人便被卷入亚略的团体中，其异端因这亵渎而大为得助，这亵渎否认圣言中人性之存在，因此，拒绝在天主内谦卑与威严的结合，它要么断言基督内有一个虚假的幻影身体，要么声称祂一切肉身的行动和受难属于天主性而非属于血肉。但祂敢于维护的一切完全是愚蠢的：因为我们的宗教信仰或奥迹的范围都不允许天主性受苦，也不允许真理在任何事上自相矛盾。因此，不能受苦的天主子，祂永恒地与圣父及圣神同在，在不变之天主圣三的同一本体中始终是祂所是的，当时期一满（这时期已由永恒旨意预定，并由言语和行为的预言意义所应许），祂降生成人，不是藉转变祂的本体，而是藉取了我们的本性，并\BibleQuote{来寻找并拯救那失落的}\BibleRef{路 19:10}。但祂的来临不是藉地方性的接近或身体的移动，好像祂之前不在，现在才在场，或是离开祂已到达之处；而是藉那看得见且与人相同的方式，为旁观者所显示，就是说，在童贞母亲的胎中接受了人的肉身和灵魂，如此，祂仍居于天主的形象内，将奴仆的形象和罪恶肉身的形状与自己结合，藉此，祂并未因人性而减少神性，反而因神性而提升了人性。\par

% structure: p0010 source=h2 target=subsection reason=relative-heading-rank; base=h2

\subsection{IV. 圣洗圣事象征并实现这一结合于每一位信徒。}

因为所有必死者的情况从我们的始祖开始就是这样：在原罪传给后裔之后，若不是圣言成了血肉并居住在我们中间——即居住在那属于我们血统和种族的本性中，无人能逃脱定罪的惩罚。因此，宗徒说：\BibleQuote{就如因一人的过犯，（审判）众人都被定罪，同样，因一人的义行，（恩宠）众人都获得生命的成义。因为正如因一人的悖逆，许多人成为了罪人，同样，因一人的顺服，许多人将成为义人}\BibleRef{罗 5:18}；又说：\BibleQuote{因为死亡是藉着人（来了），死人的复活也是藉着人（来了）。就如因亚当众人都死了，同样，因基督众人都要复活}\BibleRef{格前 15:21}。也就是说，所有虽生于亚当却在基督内重生的人，都有确凿的证据证明他们藉恩宠而成义，以及基督分享了他们的本性；因为凡不相信天主的独生子在达味后裔的童贞女胎中取了我们的本性的人，便无分于基督宗教的奥迹，既不认识新郎，也不认识新娘，不能在婚宴中有份。因为基督的肉是圣言的帷幕，凡无保留地承认祂的人皆以此帷幕为衣。但以之为耻并拒绝它的人，将不会从祂得到装饰，虽然他来到王宴，不合时宜地参加神圣的筵席，但这闯入者必逃脱不了君王的辨察，反而如主亲自所说，将被捉住，手脚被捆绑，丢在外面的黑暗中，那里有哀号和切齿。因此，凡不承认基督内有人性肉身的人，必须知道他不配降生成人的奥迹，也无分于宗徒所说的神圣结合：\BibleQuote{因为我们都是祂的肢体，是祂的肉、祂的骨。为此，人要离开父亲和母亲，依附自己的妻子，两人成为一体。}接着他解释这话的意义，又说：\BibleQuote{这奥迹是伟大的，但我是指基督和教会说的。}因此，从人类初始，基督就被宣布要带着肉身来临。在这肉身中，如所说，\BibleQuote{两人成为一体}，确实有两位：天主和人，基督和教会，教会出自新郎的肉，当她从被钉者肋旁流出的血和水中接受了救赎与重生的奥迹。因为新受造物的状态——在圣洗中，人脱去的不是真实肉身的覆盖，而是旧定罪之玷污——就是：人成了基督的身体，因为基督也是人的身体。\par

% structure: p0012 source=h2 target=subsection reason=relative-heading-rank; base=h2

\subsection{V. 重申降生成人的真道并嘱托他们持守。}

所以我們稱基督不僅只是天主，如同摩尼教異端者，也不僅只是人，如同佛提諾異端者，也不是在祂內缺少任何確實屬於人性的東西，無論是靈魂、理性心靈或不是從女人而來的、而是由聖言轉變成肉體；這三種虛假而空洞的主張曾被亞波里拿留異端者的三個派別以不同方式提出。我們也不說真福童貞瑪利亞懷了一個沒有天主性的人，這人是聖神所創造，後來被聖言所收納；我們正確地並且恰當地譴責了聶斯多略宣講這樣的道理：但我們稱基督是天主子、真天主，由天主聖父所生，沒有時間上的開端，同樣也是真人，由一位人類母親所生，在預定的時期圓滿之時，並且我們說，祂的人性——父因此更大——絲毫不減損祂與父同等的那一性體。但這兩種性體形成一位基督，祂最真實地既按祂的天主性說過：\BibleQuote{我與父原是一體}，又按祂的人性說過：\BibleQuote{父比我大}。這種真實而不可摧毀的信德，親愛的諸位，唯有它使我們成為真正的基督徒；我們欣然聽聞，你們正以忠貞的熱忱和可稱許的摯愛保衛這一信德，望你們堅定持守並勇敢維護。此外，除天主的助佑外，你們還必須贏得公教王公們的善意，因此請謙卑而明智地懇求最仁慈的皇帝恩准我們的祈求：即召開一次全體主教會議；如此，藉著天主仁慈的助佑，健康者可以增加勇氣，患病者若願意接受治療，也能得到補救。（公元四四九年十月十五日，在顯赫的\Person{Asturius}與\Person{Protogenes}執政期間。）\par

\SourceRecord{Translated by Charles Lett Feltoe.；Nicene and Post-Nicene Fathers, Second Series；Vol. 12.；Edited by Philip Schaff and Henry Wace.；Buffalo, NY: Christian Literature Publishing Co.,；1895.；<http://www.newadvent.org/fathers/3604059.htm>.}

% structure: p0001 source=h1 target=section reason=page-title

\section{Letter 66}

\Emph{\Person{良}的答复致第六十五封信。}\par

教宗良，致亲爱的弟兄\Person{康斯坦提努斯}、\Person{阿尔门塔里乌斯}、\Person{奥登提乌斯}、\Person{塞维里亚努斯}、\Person{瓦莱里亚努斯}、\Person{乌尔苏斯}、\Person{斯特凡努斯}、\Person{内克塔里乌斯}、\Person{康斯坦提乌斯}、\Person{马克西穆斯}、\Person{阿斯克勒皮乌斯}、\Person{狄奥多鲁斯}、\Person{尤斯图斯·英格努斯}、\Person{奥古斯塔利斯}、\Person{苏佩尔文托尔}、\Person{伊纳提乌斯}、\Person{丰泰乌斯}和\Person{帕拉迪乌斯}。\par

% structure: p0004 source=h2 target=subsection reason=relative-heading-rank; base=h2

\subsection{一、\Place{维也纳}主教已抢先回应他们的申诉。他提议以公正态度进行仲裁。}

当我们读到你们的来信（亲爱的弟兄），这封信是由我们的儿子、司铎\Person{佩特罗尼乌斯}和执事\Person{雷古卢斯}带来的，我们认识到你们对我们的弟兄和同主教\Person{拉文尼乌斯}怀有多么深切的爱戴：因为你们的请求是，他前任因过分僭越而理应失去的，可以归还给他。但你们的请求，弟兄们，被\Place{维也纳}主教抢先了一步，他派遣了信函和使节，控诉\Place{阿尔勒}主教非法主张祝圣\Place{瓦萨}主教之权。因此，由于我们既要尊重教父们的教规，又要尊重你们对我们的善意评价，以致在教会特权问题上不应允许任何侵犯或剥夺，我们有责任通过运用如此公正的节制来维护\Place{维也纳}行省内的和平，既不忽略古老惯例，也不忽视你们的愿望。\par

% structure: p0006 source=h2 target=subsection reason=relative-heading-rank; base=h2

\subsection{二、\Place{维也纳}主教保留对邻近四城的管辖权，其余城市归属\Place{阿尔勒}。}

因为，在考虑了出席的双方圣职人员提出的论据之后，我们发现贵省内的\Place{维也纳}和\Place{阿尔勒}两城历来如此著名，以至于在教会特权的某些事务上，时而这座，时而那座，交替占据优先地位，尽管国家传统是它们从前共享权利。因此，我们不容许\Place{维也纳}城完全失去尊荣，尤其是在教会管辖权方面，特别是它已拥有我们谕令的权威以享受其特权：即当从\Person{希拉里}手中收回时，我们以为适宜授予\Place{维也纳}主教的权力。为了使他不显得突然且不适当地被降级，他应管辖邻近四城，即\Place{瓦伦提亚}、\Place{塔兰塔西亚}、\Place{日内瓦}和\Place{格拉提亚诺波利斯}，连同\Place{维也纳}本身作为第五城，所有这些教会的事务应由\Place{维也纳}主教负责。但同一行省的其他教会应置于\Place{阿尔勒}主教的权威和管理之下，从他那温和节制的处事方式我们相信，他将如此热切地追求友爱与和平，以至于绝不认为自己被剥夺了他看到已让与其弟兄的东西。日期为五月五日，\Person{瓦伦提尼安}皇帝第七次执政，以及最著名的\Person{阿维恩}（公元450年）任执政官之时。\par

\SourceRecord{Translated by Charles Lett Feltoe.；Nicene and Post-Nicene Fathers, Second Series；Vol. 12.；Edited by Philip Schaff and Henry Wace.；Buffalo, NY: Christian Literature Publishing Co.,；1895.；<http://www.newadvent.org/fathers/3604066.htm>.}

% structure: p0001 source=h1 target=section reason=page-title

\section{第六十七封信}

致亲爱的弟兄\Person{拉文努斯}\EditorialNote{阿尔勒主教}，教宗\Person{良}。\par

我们将我们的儿子——司铎\Person{佩特罗尼乌斯}与执事\Person{雷古鲁斯}——久留于罗马，既因他们在我们眼中的宠爱配得如此，亦因信仰的需要——它现今正受某些人谬误的攻击——如此要求。因为我们希望他们出席我们讨论此事，并查明一切我们愿通过你、亲爱的，使你得以传达到所有我们的弟兄和同为主教者的知识；特将此委托给你，亲爱的弟兄，好藉着你的警醒勤勉，使我们为捍卫信仰而发往东方的信函，或真福\Person{济利禄}完全符合我们观点的信函，得以被所有弟兄知晓；好使他们在获得论据后，能以属神的力量武装自己，对抗那些胆敢以谬误信仰侮辱主的降生成人的人。亲爱的弟兄，你现在有了绝佳机会，将你主教职的开始推荐给各教会及我们的天主，如果你按照我们托付和命令你的方式执行这些事。但那些不宜付诸笔墨的事，你当倚靠天主的助佑，如我们所说，在从我们上述儿子口中得知后，有效地、可赞地执行。愿天主保佑你平安，最亲爱的弟兄。日期：五月五日，在最荣耀的\Person{瓦伦提尼安}（第七次）与著名的\Person{阿维努斯}执政之年（450年）。\par

\SourceRecord{Translated by Charles Lett Feltoe.；Nicene and Post-Nicene Fathers, Second Series；Vol. 12.；Edited by Philip Schaff and Henry Wace.；Buffalo, NY: Christian Literature Publishing Co.,；1895.；<http://www.newadvent.org/fathers/3604067.htm>.}

% structure: p0001 source=h1 target=section reason=page-title

\section{第六十八封信}

\Emph{来自三位高卢主教致圣良}\par

\Person{塞雷提乌斯}、\Person{萨洛尼乌斯}和\Person{维拉努斯}致圣主、最蒙福的教父、最堪当宗座的教宗\Person{良}。\par

% structure: p0004 source=h2 target=subsection reason=relative-heading-rank; base=h2

\subsection{一、他们为\WorkTitle{大卷}向\Person{良}道贺并致谢。}

在仔细阅读了阁下的信函——您为教导信德而撰写，并寄给\Place{君士坦丁堡}主教——之后，我们认为自己有责任，既然已被如此丰富的教义所充盈，至少通过给您写信来偿还我们感恩的债务。因为我们感激您对我们慈父般的关怀，并承认我们更应归功于您那预防性的照顾，因为我们如今在经历恶事之前便已受益于补救。因为知道那些伤口造成后才施行的补救几乎为时已晚，您以充满爱心的远见之声劝诫我们要以那些宗徒性的防御措施武装自己。最蒙福的教宗，我们坦率地承认，您以何等非凡的仁慈向我们传达了您内心的最深处，藉此您确保了其他人的安全：当您从他人的心中拔出古老毒蛇注入的毒液时，您仿佛站在爱的高塔上，以宗徒的关顾和警觉高声呼喊，免得敌人出其不意地攻击我们，免得粗心的安逸使我们暴露于袭击，哦，圣主、最蒙福的教父、最堪当宗座的教宗。此外，我们这些特别属于您的人，充满了巨大而难以言表的喜悦，因为您这一特殊的教义陈述在各教会聚集之处受到如此高度的重视，以至于一致认为宗座的首要地位正当地归属于那里，因为宗徒之神的训谕仍然在那里获得诠释。\par

% structure: p0006 source=h2 target=subsection reason=relative-heading-rank; base=h2

\subsection{二、他们请求他订正或增补其\WorkTitle{大卷}抄本。}

因此，若您认为值得，我们恳请圣座通读并订正这份抄写员在这部工作中——无论现在还是将来都如此珍贵，我们出于保存它的愿望已将其抄录在羊皮纸上——的任何错误，或者，若您出于热忱已设想了任何将有益于所有读者的事，请以您慈爱的关怀下令将其增补到这份抄本中，以便不仅\Place{高卢}各省的许多神圣主教我们的弟兄，而且许多作为您平信徒儿子的那些人——他们极愿看到这封信以揭示真理——当这封信经您圣手订正后送回我们时，得以抄写、阅读和保存它。如您认为合适，我们渴望我们的信使能尽快返回，以便我们更快地得到您健康的消息而欢喜：因为您的安康就是我们的喜乐和健康。\par

愿主基督长久保持尊驾对我们的卑微心怀顾念，哦，圣主、最蒙福的教父、最堪当宗座的教宗。\par

\Signature{我，\Person{塞雷提乌斯}，您的义子，向宗座致敬，并将自己托付于您的祈祷。}\par

\Signature{我，\Person{萨洛尼乌斯}，您的崇拜者，向宗座致敬，恳求您祈祷的帮助。}\par

\Signature{我，\Person{维拉努斯}，您宗座的敬拜者，向您的圣座致敬，并恳求您为我祈祷。}\par

\SourceRecord{Translated by Charles Lett Feltoe.；Nicene and Post-Nicene Fathers, Second Series；Vol. 12.；Edited by Philip Schaff and Henry Wace.；Buffalo, NY: Christian Literature Publishing Co.,；1895.；<http://www.newadvent.org/fathers/3604068.htm>.}

% structure: p0001 source=h1 target=section reason=page-title

\section{书信六九}

\Person{良一世}主教，致永远至尊的\Person{德奥多西}皇帝。\par

% structure: p0003 source=h2 target=subsection reason=relative-heading-rank; base=h2

\subsection{一、暂缓对\Person{阿纳多略}任命发表意见，直至他公开宣认大公信仰}

在您所有虔诚的信函中，虽然我们为信仰而受苦，您却以如此忠信支持\Place{尼西亚}大公会议，不容许主的司铎们偏离它（正如您已多次写信告知我们），从而给了我们安全的希望。但为了避免我似乎对捍卫大公信仰有任何不利之处，我认为目前对那位开始在\Place{君士坦丁堡}教会担任牧职者之任命，不宜采取任何冒然行动回复——这并非缺乏关爱，而是等待着大公真理得以显明。我恳请陛下宽容对待此事：当他证明自己在关乎大公信仰方面符合我们的期望时，我们便能更完全、更安全地为他的真诚而喜乐。但为使任何恶意的怀疑不因我们对他的态度而困扰他，我除去一切困难的理由，不要求任何看似苛刻或可争议之事，而是发出任何大公信徒都不会拒绝的邀请。因为在我们之前、无论是在希腊语还是拉丁语中宣讲大公真理的人，举世闻名，他们学识与教导甚至为今日一些人所采用，从他们的著作中产生了统一而多样的信理陈述；这陈述既推翻了\Person{纳斯托略}的异端，也切除了如今重新萌发的这一错误。因此，让他再读一读圣教父们在主降生成人一事上的信仰，这信仰一向被守护并同样地宣讲；当他认识到，荣福纪念的\Place{亚历山大}主教\Person{基利尔}的书信（\EditorialNote{在此信中，他愿纠正并医治纳斯托略，驳斥其错误论述，更清晰地陈明尼西亚所界定的信仰；这封信由他发送并存放在宗座的图书馆中}）与先前诸教父的观点一致时，让他进一步重新考虑\Place{厄弗所}大公会议的记录——其中，荣福纪念的\Person{基利尔}收录并维护了关于主降生成人的大公司铎们的证言。他也请不要轻视阅读我的书信，他会发现这封信与教父们的虔诚信仰完全一致。当他认识到，所要求于他的正是服务于同一善果时，让他衷心赞同大公信徒们的判断，以便在全体圣职人员和全体人民面前，毫无保留地宣认他对共同信仰的真诚认同，并将此认同通报给宗座、所有主的司铎与教会。这样，世界因同一信仰而和睦，我们众人便能如天使在救主由童贞玛利亚诞生时所歌唱的那样说：\Quote{在天上光荣归于天主，在地上和平归于善人\BibleRef{路 2:14}。}\par

% structure: p0005 source=h2 target=subsection reason=relative-heading-rank; base=h2

\subsection{二、若\Person{阿纳多略}作此宣认，他承诺接受他；并请求在意大利召开大公会议以最终界定信仰}

但是由于我们和荣福教父们（我们所敬畏并追随其教导）在同一信仰上同心合意，正如所有各省主教所证实的，愿陛下的虔诚信仰确保这位\Place{君士坦丁堡}主教能以一位经过认可的大公司铎之身份，尽快向我们提交应有的文件——即公开、明确地宣称：凡相信或主张关于天主圣言降生成人的任何其他观点，而不符合我和所有大公信徒所陈明的道理的，他将与其断绝共融；这样我们才能合宜地在基督内赐予他兄弟之爱。为使我们有益的心愿能藉着天主助佑、通过陛下的信仰产生更迅速、更圆满的成效，我已派遣我的兄弟和同僚主教\Person{阿邦迪}与\Person{阿斯泰里}，同司铎\Person{巴西略}与\Person{塞纳托尔}一同前往陛下处；他们的虔敬已在我面前得到充分证明。当他们展示了我们送去的训令后，您便能恰当了解我们信仰的标准；如此，若这位\Place{君士坦丁堡}主教衷心赞同同一宣认，我们便能安妥地、合宜地为教会的和平而喜乐，也不会再有任何歧义隐藏在后，以可能无端的猜疑困扰我们。但是，若有人背离我们信仰的纯正和教父们的权威，则为此目的而在\Place{罗马}召集的宗教会议将与我一同恳请陛下，准许在\Place{意大利}境内召开一个普世大公会议；这样，所有因无知或恐惧而跌倒的人都聚集一处，便可以采取措施纠正并医治他们，再不许任何人引用\Place{尼西亚}大公会议的方式以证明自己反对其信仰；因为无论对整个教会还是对您的统治，都将有益：如果整个世界以同一宣认持守同一的天主、同一的信仰、同一的人类救恩奥迹。\par

日期：七月十七日，在显赫的\Person{瓦伦提尼安}第七次执政和\Person{阿维努斯}（四五〇年）执政年间。\par

\SourceRecord{Translated by Charles Lett Feltoe.；Nicene and Post-Nicene Fathers, Second Series；Vol. 12.；Edited by Philip Schaff and Henry Wace.；Buffalo, NY: Christian Literature Publishing Co.,；1895.；<http://www.newadvent.org/fathers/3604069.htm>.}

% structure: p0001 source=h1 target=section reason=page-title

\section{\WorkTitle{第七十九封書信}}

\Person{良}，羅馬城主教，致\Person{普爾喀麗亞·奧古斯塔}。\par

% structure: p0003 source=h2 target=subsection reason=relative-heading-rank; base=h2

\subsection{一、他因\Person{普爾喀麗亞}對抗\Person{聶斯多略}和\Person{歐迪奇}的熱忱而歡喜。}

關於\OriginalTerm{恩主}{Your Grace}的神聖意圖，我們始終所預期的，如今已完全證實，即：無論惡人如何以多種方式攻擊基督徒信仰，但當\OriginalTerm{恩主}{Your Grace}在場，並由主預備予以捍衛時，信仰便不致動搖。\newline{} 因為天主不會放棄祂仁慈的奧蹟，也不會放棄\OriginalTerm{恩主}{Your Grace}勞苦的功績：藉此功績，\OriginalTerm{恩主}{Your Grace}早已將我們神聖宗教的狡猾仇敵從教會的心臟驅逐出境。那時，\Person{聶斯多略}的不敬未能維持其異端，因為\OriginalTerm{恩主}{Your Grace}，真理的婢女和學生，沒有被那油嘴滑舌者的彩色謊言所瞞過——他給單純的人們注入了多少毒藥！\newline{} 那場偉大鬥爭的後續是：藉著\OriginalTerm{恩主}{Your Grace}的警醒，魔鬼通過\Person{歐迪奇}所策劃的事未能隱藏，而那些在雙重異端中選擇一方的人，被大公信仰的統一而不可分割的力量所推翻。\newline{} 因此，這是\OriginalTerm{恩主}{Your Grace}對\Person{歐迪奇}錯誤之毀滅的第二次勝利。\newline{} 如果\Person{歐迪奇}有任何清醒的頭腦，那麼這個錯誤既然在\OriginalTerm{其煽動者}{his instigators}身上早已被擊敗和混亂，他本應很容易避免企圖復燃那些冒煙的灰燼，從而只是分擔那些他所追隨範例之人的命運，最光榮的\Person{奧古斯塔}。\newline{} 因此，我們想歡躍，並為\OriginalTerm{恩主之福}{your clemency's prosperity}向天主獻上應有的許諾，祂已在全世界各處——在傳揚主福音的地方——賜予\OriginalTerm{恩主}{Your Grace}雙重的棕櫚枝和冠冕。\par

% structure: p0005 source=h2 target=subsection reason=relative-heading-rank; base=h2

\subsection{二、他感謝\Person{普爾喀麗亞}對大公事業的援助，並說明他對恢復失足主教的期望。}

\OriginalTerm{恩主}{Your clemency}必須知道，全羅馬教會對\OriginalTerm{恩主}{Your Grace}所有忠信的行為深懷感激，無論是\OriginalTerm{恩主}{Your Grace}以虔誠的熱忱協助了我們的全權代表，並將那些被不義判決逐出教會的大公司鐸領回；還是\OriginalTerm{恩主}{Your Grace}促使那位無辜而神聖的司鐸——神聖記憶的\Person{弗拉維安}——的遺體，以應有的榮譽歸還給他管理有方的教會。\newline{} 在所有這些事上，\OriginalTerm{恩主}{Your Grace}的光榮無疑倍增，因為\OriginalTerm{恩主}{Your Grace}按聖人們的功績尊敬他們，並關切地將荊棘和雜草從主的田地中除去。\newline{} 但是，我們也從我們代理人的報告以及我的兄弟和同僚主教\Person{阿納托利}（他蒙\OriginalTerm{恩主}{Your Grace}恩准推薦給我）的報告中得知，某些主教請求為那些似乎已同意異端之事的人尋求和解，並渴望他們獲得大公的共融。\newline{} 我們准許他們的請求，條件是：在我們的代理人與上述主教協同行事，他們受到糾正，並親手譴責自己的惡行之前，不應賜予他們和平的恩惠；因為我們基督教信仰要求：真正的正義應當約束頑固的人，而愛不應拒絕懺悔者。\par

% structure: p0007 source=h2 target=subsection reason=relative-heading-rank; base=h2

\subsection{三、他將其某些主教和教會託付給\Person{普爾喀麗亞}照顧。}

因為我們知道\OriginalTerm{恩主}{Your Grace}願以多大的虔誠關懷大公司鐸，我們已下令告知\OriginalTerm{恩主}{Your Grace}：我的兄弟和同僚主教\Person{歐瑟伯}正與我們同住，並分享我們的共融；我們將其教會託付給\OriginalTerm{恩主}{Your Grace}，因為據說不當被選接替他位置的人正在蹂躪該教會。\newline{} 此外，我們請求\OriginalTerm{恩主}{Your Grace}——我們不懷疑\OriginalTerm{恩主}{Your Grace}會自願這樣做——將應有的恩惠同樣賜予我的兄弟和同僚主教\Person{猶利安}，以及那些以忠誠依附神聖\Person{弗拉維安}的君士坦丁堡神職人員。\newline{} 關於所有必須執行或安排的事，我們已通過我們的代理人向\OriginalTerm{恩主}{Your Grace}說明。\newline{} 日期：\OriginalTerm{顯赫的亞德斐烏斯}{illustrious Adelfius}執政官之年（451年）四月十三日。\par

\SourceRecord{Translated by Charles Lett Feltoe.；Nicene and Post-Nicene Fathers, Second Series；Vol. 12.；Edited by Philip Schaff and Henry Wace.；Buffalo, NY: Christian Literature Publishing Co.,；1895.；<http://www.newadvent.org/fathers/3604079.htm>.}

% structure: p0001 source=h1 target=section reason=page-title

\section{第八十封书信}

\Emph{致君士坦丁堡主教\Person{安纳多略}。}\par

% structure: p0003 source=h2 target=subsection reason=relative-heading-rank; base=h2

\subsection{一、他因\Person{安纳多略}证明自己为正统而欢喜。}

主教\Person{良}，致\Person{安纳多略}主教。\par

我们因主而喜乐，并因祂的恩宠之赐而荣耀，祂已显示您为福音教导的追随者，正如我们从您信函及我们派往君士坦丁堡的弟兄们的报告中所发现的：因为如今藉着司铎的认可信德，我们正有理由相信交付给他的整个教会将没有皱纹或谬误的斑点，如宗徒所说：\Quote{因为我曾把你们许配给一个丈夫，把你们如同贞洁的童女献给基督}，见：\BibleRef{格后 11:2}。因为那位童女就是教会，是唯一丈夫基督的新娘，她不容许自己被任何谬误所败坏，以致普世我们拥有一完整而纯洁的共融，在此我们现今欢迎您为同道，亲爱的，并批准我们已收到的程序记录，并以必要的签署予以认可。因此，为使您的心灵，亲爱的，也藉我们的话得到坚固，我们在逾越节庆后，将我们的儿子，司铎\Person{卡斯特里}、执事\Person{帕特里西}和\Person{阿斯克勒庇亚}送回，他们曾将您的信函带来给我们，告知您，如我们前面所述，我们为君士坦丁堡教会的和平而喜乐，我们为此如此劳心，希望它不被任何异端欺骗所污染。\par

% structure: p0006 source=h2 target=subsection reason=relative-heading-rank; base=h2

\subsection{二、背教主教中的忏悔者，应按\Person{安纳多略}与\Person{良}的代表所商定的计划，被重新接纳至圆满共融。}

但关于那些弟兄，我们从您的信函和我们代表的报告中得知，他们渴望与我们共融，因为他们后悔未能坚忍抵抗暴力和威吓，却因恐惧而慌乱，草率同意，参与了公教会且无罪的主教（\Person{弗拉维亚诺}）的定罪，并接纳了可憎的异端（\Person{欧提基}），我们赞同在我们代表在场并合作下所决定的，即他们应暂时满足于自己教会的共融，但希望我们派往您处的代表与您商议，并作出安排，使那些以充分补赎保证谴责自己恶行，宁愿控告自己而非辩护自己的人，能因与我们和平共融而欢喜；条件是他们首先完全绝罚反对公教会信德的一切。否则，在天主之教会——基督的身体内，既无有效的司祭职，也无真实的祭献，除非在祂本性的真实中，真正的大司祭为我们赎罪，无玷羔羊的真正宝血洁净我们。因为尽管祂坐在圣父右边，但在祂从童贞玛利亚所取的同一天主性肉身上，祂继续着赎罪的奥迹，如宗徒所说：\Quote{基督耶稣已死，且已复活，祂在天主右边，也为我们转求}，见：\BibleRef{罗 8:34}。因为我们的仁爱在任何情况下都不能被责备，当我们接纳那些保证补赎的人，我们曾因他们的迷惑而忧伤。因此，与我们共融的恩赐既不应严苛地拒绝，也不应轻率地赐予，因为正如以基督般的爱德对待受压迫者完全符合我们的信仰，同样，将全部责任归于骚乱的肇事者也是公平的。\par

% structure: p0008 source=h2 target=subsection reason=relative-heading-rank; base=h2

\subsection{三、\Person{狄奥斯库鲁斯}、\Person{朱文纳尔}与\Person{尤斯塔修}之名不应在圣祭台前宣读。}

关于在圣祭台前宣读\Person{狄奥斯库鲁斯}、\Person{朱文纳尔}和\Person{尤斯塔修}的名字，亲爱的，您应当遵守我们当时在场的朋友们所说应当做的事，这也与圣\Person{弗拉维亚诺}的可敬纪念相符，并且不会使教友的心背离您。因为那些以攻击骚扰无辜公教徒的人，与圣人的名字混杂不分，是极不恰当和不相称的，因为他们不放弃自己被定罪的异端，反而以他们的乖戾定自己的罪；这样的人或者应为他们的不忠而受责罚，或者应竭力寻求宽恕。\par

% structure: p0010 source=h2 target=subsection reason=relative-heading-rank; base=h2

\subsection{四、关于个别人士的一两点指示。}

然而，我们的弟兄和同为主教的\Person{尤利安}，以及那些忠于圣善记忆的\Person{弗拉维亚诺}、忠信服事他的圣职人员，我们希望他们也依附于您，亲爱的，好使他们知道那位我们确信凭其信德之功与我们天主同活的人，在您身上临在他们中间。我们还希望您知晓此事，亲爱的：我们的弟兄和同为主教的\Person{优西比乌}，他曾为信德忍受许多危险和辛劳，现暂居我们这里，继续与我们共融；愿您的关怀保护他的教会，使他在外期间一无损失，且无人敢在任何事上伤害他，直到他持我们的信函到达您处。为使您在我们——更确切地说，全体基督徒人民——心中的爱激发更大的程度，我们希望这封信，亲爱的，为众人所知，好使事奉我们天主的人，因宗座与您的和平圆满完成而感恩。至于其他事项和人士，您将由我们代表带给您的信函中获得更充分的指示。于阿德尔菲乌斯显贵执政之年（451年）4月13日。\par

\SourceRecord{Translated by Charles Lett Feltoe.；Nicene and Post-Nicene Fathers, Second Series；Vol. 12.；Edited by Philip Schaff and Henry Wace.；Buffalo, NY: Christian Literature Publishing Co.,；1895.；<http://www.newadvent.org/fathers/3604080.htm>.}

% structure: p0001 source=h1 target=section reason=page-title

\section{第八十二封书信}

\Emph{致\Person{马尔西安}皇帝。}\par

% structure: p0003 source=h2 target=subsection reason=relative-heading-rank; base=h2

\subsection{一、在祝贺皇帝的崇高行为之后，他反对对信德道理进行随意探究。}

虽然我已藉\Place{君士坦丁堡}圣职人员之手回复了陛下，但当我藉我儿\Person{塔提安}——这位杰出的\Place{君士坦丁堡}城市长官——获得陛下的仁慈眷顾时，我更加感到值得祝贺，因为我了解到您对教会和平的强烈热忱。这一神圣愿望，正如其公正所应得的，为您的帝国带来了与您为宗教所寻求的同样幸福的状态。因为当天主圣神在基督信仰的君主之间建立和谐时，普世便产生双重的信心：爱德与信德的进步使他们的武力在两方面都不可战胜，以致天主因一致的宣认而息怒，异端者的虚妄与蛮族的敌意同时被推翻，最光荣的皇帝。因此，藉着皇帝的友谊，天助的希望增加，我更加自信地恳求陛下，为了人类救恩的奥秘，不容许任何人以虚妄和放肆的诡计探究那些必须持守的道理，仿佛它们是不确定的。尽管我们在一个字上也不可偏离福音与宗徒的教导，也不可对圣经持有与蒙福的宗徒和我们的教父所学习与教导的不同的见解，但在这末后的日子里，无知而亵渎的探究被提了出来——从前圣神藉真理的门徒粉碎了这些探究，当魔鬼一在适合其目的的心灵中激起它们时。\par

% structure: p0005 source=h2 target=subsection reason=relative-heading-rank; base=h2

\subsection{二、待决的问题只是哪些失足者应被恢复，以及条件为何。}

但最不合时宜的是，因少数人的愚蠢，我们再次陷入危险的见解和血肉的争论中，仿佛争论要重新开始，我们必须解决\Person{欧提克斯}是否持有亵渎的观点，以及\Person{迪奥斯库鲁斯}是否作出了错误的判决——他在定罪神圣记忆的\Person{弗拉维安}时，给了自己致命一击，并使较单纯的人也卷入同样的毁灭。如今，据我们了解，许多人已采取补救措施，并为其软弱的轻率恳求宽恕，我们要决定的不是信德的性质，而是接受谁的祈祷，以及以何种条件。因此，您为召集主教会议所愿怀有的那种最虔诚的关切，将在我认为与案件需要相关的一切问题上，通过将尽快（如果天主许可）到达陛下身边的代表们，得到充分而及时地陈述。日期：在杰出的\Person{阿德尔菲乌斯}执政时期，四月二十三日（451年）。\par

\SourceRecord{Translated by Charles Lett Feltoe.；Nicene and Post-Nicene Fathers, Second Series；Vol. 12.；Edited by Philip Schaff and Henry Wace.；Buffalo, NY: Christian Literature Publishing Co.,；1895.；<http://www.newadvent.org/fathers/3604082.htm>.}

% structure: p0001 source=h1 target=section reason=page-title

\section{第八十五封书信}

\Emph{致\Place{君士坦丁堡}主教\Person{亚纳多略}。}\par

% structure: p0003 source=h2 target=subsection reason=relative-heading-rank; base=h2

\subsection{一、\Person{亚纳多略}偕同\Person{良}的使者，处理那些曾随\Person{欧提克斯}暂时迷失者重新接纳的问题。}

主教\Person{良}致主教\Person{亚纳多略}。\par

虽然我希望，亲爱的，你致力于一切善工，但为使你的行动更有效，有必要并适宜派遣我的弟兄\Person{路森休}主教和\Person{巴西略}司铎，正如我们所承诺的，与你结合，亲爱的，好使关于普世教会利益的事务，无所犹豫或懈怠地办理；因为只要你亲临现场，我们已将执行我们旨意之事委托于你，一切皆能如此有节制地处理，既不忽略仁慈，也不忽略正义，且不徇情面，处处顾及天主的判断。但为使此事正确遵守并保持，必须首先维护公教信仰的完整，因为在所有情况下，\BibleQuote{通往生命的道是窄而狭的}\BibleRef{玛 7:14}，不可偏离其轨迹，或左或右。因为福音及宗徒信仰须对抗一切谬误，一面推翻\Person{奈斯多略}，一面粉碎\Person{欧提克斯}及其同伙，请记住遵守这一规则的必要性：所有那些在不能、也不配被称为会议的会议中——\Person{迪奥斯哥鲁}在那里表现其恶意，\Person{尤维纳}表现其无知——据我们从你的报告中所知，亲爱的，他们因恐惧而被征服，并在惊骇中被逼同意那不义的判决，如今渴望获得公教共融者，应在充分保证悔改后接受弟兄们的和平，条件是他们必须以毫不含糊的措辞，诅咒、谴责并弃绝\Person{欧提克斯}及其邪说和追随者。\par

% structure: p0006 source=h2 target=subsection reason=relative-heading-rank; base=h2

\subsection{二、较重罪者的案件须留待当前处理。}

至于那些在此事中犯罪更重，且在同样不幸的会议中自居高位，以邪恶的傲慢激怒卑微弟兄纯朴心灵的人，若他们恢复理智，停止辩护自己的行为，转而谴责自己的特定错误，并给出无可争议的悔改保证，则应将他们的案件保留，等待宗座更成熟的审议，以便在一切筛检权衡后，对其真实行为得出正确结论。在你所治理的、主愿意你管理的教会中，凡我们已写明的这类人，不得在祭台上宣读他们的名字，直到事件发展表明当如何决定他们。\par

% structure: p0008 source=h2 target=subsection reason=relative-heading-rank; base=h2

\subsection{三、请\Person{亚纳多略}忠实配合\Person{良}的使者。}

至于你的圣职人员呈交给我们的文书，亲爱的，无需将我的意见写入信函；只需将一切托付给我的代表，他们的话语将细致地指示你每一要点。所以，最亲爱的弟兄，请与这些我们选为如此重大事务的合适代理人的弟兄们一道，忠诚而有效地执行符合天主教会的事；尤其是因为事情本身的性质和天主助佑的许诺激励你，我们最仁慈的君主们表现出如此神圣的信德和虔诚的热忱，以致我们在他们身上不仅看到基督徒普遍的同情，甚至看到司祭职的精神。他们无疑会依照那使他们自夸为天主仆人的虔诚，以相称的精神接受你对公教信仰有益的一切建议，以便藉他们的帮助，基督的和平得以恢复，邪恶的谬误得以消除。若在某些点上需要更多建议，请迅速传信给我们，以便在调查事情性质后，我们审慎地规定正当的措施。于显赫的\Person{阿德尔斐乌斯}执政之年（451年）六月九日。\par

\SourceRecord{Translated by Charles Lett Feltoe.；Nicene and Post-Nicene Fathers, Second Series；Vol. 12.；Edited by Philip Schaff and Henry Wace.；Buffalo, NY: Christian Literature Publishing Co.,；1895.；<http://www.newadvent.org/fathers/3604085.htm>.}

% structure: p0001 source=h1 target=section reason=page-title

\section{第八十八封信}

莱奥，主教，致帕斯卡西诺，利利贝主教。\par

% structure: p0003 source=h2 target=subsection reason=relative-heading-rank; base=h2

\subsection{寄送《大卷》副本并进一步阐明欧提克斯的异端性。}

虽然我毫不怀疑，弟兄，关于吾主耶稣基督降生成人所引起的东方教会中的一切丑闻根源，你都已充分得知，但恐有偶然疏漏之处，我已将我们致真福弗拉维安的、对此事作最详尽论述的信函寄去供你仔细研读；这信函已为普世教会所接纳。如此，你便能明白，这一整个亵渎性的错误如何在天主的助佑下被彻底摧毁，而你自己也当本着对天主的爱展现同样的精神，并清楚知道那些追随欧提克斯的亵渎与疯狂、竟敢说我们的主、天主的独生子、那位为拯救人类而自甘降卑者内没有两个性体（即完整的天主性与完整的人性）的人，是绝对可憎的。他们以为说他们相信圣言的一个性体降生成人，就能麻痹我们的警惕；然而天主的圣言在圣父、自身及圣神的天主性内固然只有一个性体，但当祂取了我们血肉的真实性时，我们的人性也与祂不变的实体结合了：因为除非圣言取了血肉，否则连\Quote{降生成人}也不可说。这取血肉形成了如此完整的结合，以致不仅在真福童贞女的生产中，而且在她的怀孕中，都不可想象天主性与被赋予生命的肉体之间有丝毫分离，因为在位格的合一中，天主性与人性在童贞女的怀孕与生产中结合为一。\par

% structure: p0005 source=h2 target=subsection reason=relative-heading-rank; base=h2

\subsection{欧提克斯本可从以往异端者的命运得到警告。}

因此，在欧提克斯身上所应憎恶的亵渎，正如从前在教父们手中曾被谴责并推翻的异端者一样；这些先例本应使这个愚昧之徒获益，提醒他提防自己凭理智所无法把握的事，以免因否认吾主耶稣基督内真实的人性而毁坏我们救恩的无与伦比的奥迹。因为，如果祂内没有真实而完美的人性，那么祂就没有取了我们的人性；按照他的异端，我们全部的信仰与教导便成了虚空与谎言。然而，因为真理不说谎，天主性不受苦难，故在天主圣言内两种实体共存于一个位格中；教会如此宣认她的救主：既在祂的天主性中不受苦难，又在祂的肉体中能受苦，正如宗徒所说：\Quote{祂虽然由于（我们的）软弱被钉在十字架上，却由于天主的能力仍然活着\BibleRef{格后 13:4}。}\par

% structure: p0007 source=h2 target=subsection reason=relative-heading-rank; base=h2

\subsection{寄送教父语录，并宣布东方各教会已接纳《大卷》。}

为了使你在一切事上得到更充分的教导，亲爱的弟兄，我已给你寄送了我们圣教父们的一些语录，以便你清楚了解他们对主降生成人奥迹的感悟以及在教会中的宣讲；这些语录连同我们的信函，也在君士坦丁堡由我们的代表们呈上。你必须明白，整个君士坦丁堡教会，连同所有隐修院及许多主教，都已表示赞同，并以签字弃绝了聂斯多略和欧提克斯连同他们的教义。你还必须明白，我最近收到了君士坦丁堡主教的信，其中提到安提约基雅主教已向他所辖各省的所有主教发出指示，使他们赞同我的信函，并同样弃绝聂斯多略和欧提克斯。\par

% structure: p0009 source=h2 target=subsection reason=relative-heading-rank; base=h2

\subsection{请他通过咨询权威人士，解决亚历山大与罗马在455年复活节计算上的分歧。}

此外，我们认为有必要嘱咐你的谨慎：在你确知消息的地方，应勤奋查问关于复活节计算中的那个点——我们在特奥斐洛的表格中发现了它，且这使我们极为困扰——并与精通这类计算的人讨论，四年后主的复活日应在何时举行。因为，靠着天主的仁慈，下一个复活节定于3月23日举行，第二年在4月12日，第三年在4月4日，而真福特奥斐洛将455年的复活节定为4月24日，我们发现这完全违反教会的规则；但在我们的复活节周期中，如你所熟知，那年的复活节定于4月17日。因此，为消除我们所有疑虑，我们请你与最权威的人士仔细讨论这一点，以便今后避免此类错误。发于6月24日，显赫的阿德尔菲乌斯执政之年（451年）。\par

\SourceRecord{Translated by Charles Lett Feltoe.；Nicene and Post-Nicene Fathers, Second Series；Vol. 12.；Edited by Philip Schaff and Henry Wace.；Buffalo, NY: Christian Literature Publishing Co.,；1895.；<http://www.newadvent.org/fathers/3604088.htm>.}

% structure: p0001 source=h1 target=section reason=page-title

\section{第九十三封信}

\Emph{致加采东大公会议}\par

% structure: p0003 source=h2 target=subsection reason=relative-heading-rank; base=h2

\subsection{一、他为缺席会议致歉，并介绍其代表。}

\Person{良}，\Place{罗马}城主教，致在\Place{尼西亚}集会的圣教会议。\par

我确实曾祈祷，亲爱的诸位，为我亲爱的同僚们，愿主的一切司铎恒心持守对公教信仰的一致虔诚，并且无人因世俗权力的偏袒或恐惧而误入歧途，背离真理之路；但因许多事常引发忏悔，天主的仁慈超越犯者的过犯，且报复被推迟以使改过得以发生，我们必须重视我们最仁慈君主的虔诚目的之会议，他在会议中渴望你们神圣的兄弟团体聚集，以摧毁魔鬼的陷阱，恢复教会的和平，且如此尊重最有福的宗徒伯多禄的权利和尊严，以致也写信邀请我们惠临你们可敬的会议。这确因时代需要或任何先例都不允许。然而，在这些弟兄——即\Person{帕斯卡西努斯}和\Person{卢森提乌斯}主教，\Person{波尼法爵}和\Person{巴西略}司铎，由宗座派遣——身上，愿你们的兄弟团体认为我亲自主持会议；因为我的临在并未离开你们，如今由我的代表所代行，且长久以来，在宣告公教信仰上，我实在与你们同在；因此，你们既无法不知道我们按古代传统所信的是什么，也就不能怀疑我们所愿望的是什么。\par

% structure: p0006 source=h2 target=subsection reason=relative-heading-rank; base=h2

\subsection{二、他恳请他们重申\WorkTitle{书信}中所陈述的信仰。}

因此，最亲爱的弟兄们，务要彻底平息一切攻击天主所默启之信仰的企图，并使异端者虚妄的不信止息；也不可辩护那不可信之理；因为依据福音的权威论述、先知们的发言和宗徒们的教导，在我们寄给有福纪念的\Person{弗拉维安}主教的信中，已极其充分而清晰地规定了关于我等主耶稣基督降生成人的奥迹的忠诚而纯洁的宣认。\par

% structure: p0008 source=h2 target=subsection reason=relative-heading-rank; base=h2

\subsection{三、被驱逐的主教必须复职，且奈斯多略教规继续有效。}

但因我们深知，由于恶毒的嫉妒，许多教会的情况已受扰乱，大批主教因不接受异端而被逐出教座、流放远方，而另一些人在他们尚在世时被安插其位；对这些创伤，首先必须实施正义的医治，绝不可剥夺一人的所有，使他人得以享用；因为如果我们所愿，众人皆放弃其谬误，则无人需失去其现有品秩，而那些为信仰辛劳者，应恢复其权利，并享有所有特权。愿先前由圣善记忆的\Person{基利尔}主教主持的厄弗所大公会议上特别针对\Person{奈斯多略}的法令继续有效，以免那时被谴责的异端因\Person{欧迪奇}受到应得之咒骂而在任何事上自鸣得意。因为我们在同一精神下与我们的圣教父们所宣告的信仰与教义的纯洁，同样谴责并抨击\OriginalTerm{奈斯多略}{Nestorian}和\OriginalTerm{欧迪奇}{Eutychian}的谬误及其支持者。诸位最亲爱的弟兄，愿主内平安。日期：六月二十六日，显贵阿德尔菲乌斯执政之年（四五一年）。\par

\SourceRecord{Translated by Charles Lett Feltoe.；Nicene and Post-Nicene Fathers, Second Series；Vol. 12.；Edited by Philip Schaff and Henry Wace.；Buffalo, NY: Christian Literature Publishing Co.,；1895.；<http://www.newadvent.org/fathers/3604093.htm>.}

% structure: p0001 source=h1 target=section reason=page-title

\section{第九十五封信}

\Emph{致普尔喀丽亚·奥古斯塔，由长官特奥克提斯图斯转交。}\par

% structure: p0003 source=h2 target=subsection reason=relative-heading-rank; base=h2

\subsection{一、他告知皇后，他已忠诚地承认了她所命令召开的会议，并派遣代表携信件赴会。}

主教利奥致普尔喀丽亚·奥古斯塔。\par

陛下对公教信仰的虔诚关怀，您不懈地倾注其中，我凡事都能认出，并感谢天主，看到您对普世教会如此关心，以致我可以自信地提出我认为合乎正义与仁爱之事，从而您那通过基督的怜悯迄今无可指责的热忱，能更迅速地达成我们感谢的结果，最高贵的奥古斯塔。因此，您命令在\Place{尼西亚}召开主教会议，并温和地拒绝了我请求在\Place{意大利}召开的要求（这样我们地区的所有主教都可以被召集并集会，如果情况允许他们的话），我以远远不是轻视的态度接受，并分别提名两位主教同僚和两位司铎同僚代表我，同时也向可敬的会议发出了一封合适的信函，使聚集在那里的弟兄团体能够了解他们决定时必须遵循的标准，以免任何轻率之举损害信仰的规范、教规的条款或仁爱精神所需的补救措施。\par

% structure: p0006 source=h2 target=subsection reason=relative-heading-rank; base=h2

\subsection{二、在处理此事时必须保持节制，这种节制在以弗所完全缺失。}

因为，正如我从这件事一开始就在信函中多次说明的，我希望在分歧的意见和肉体的嫉妒中保持这种节制，以便既不允许从信仰的纯洁中夺去或增添任何东西，又应赐予那些回归合一与和平的人宽恕的补救。因为当人心被召回对天主和邻人的爱时，魔鬼的作为就被更有效地摧毁了。然而，当时他们的行为与我的警告和恳求何其相反，说来话长，也无需要在这封信中记载在\Place{以弗所}那次不是审判者而是强盗的会议上所允许发生的一切；在那次会议上，主教会议的首领们既不饶恕那些反对他们的弟兄，也不饶恕那些赞同他们的人，因为为了摧毁公教信仰和加强可憎的异端，他们剥夺了一些人的合法品级，并玷污了另一些人使其与罪有染；他们的残忍无疑对那些他们通过说服使之与纯洁分离的人，比对那些他们通过逼迫使之成为蒙福的宣信者的人，更为严重。\par

% structure: p0008 source=h2 target=subsection reason=relative-heading-rank; base=h2

\subsection{三、对于那些撤回错误的人必须宽容对待。}

然而，因为这样的人因自己的恶行而伤害了自己最大，而且伤势越重，就越须谨慎施治，我从未在任何信函中坚持认为，即使他们悔改自新，也不应给予宽恕。虽然我们坚定地憎恶他们的异端，这是基督教信仰最大的敌人，但这些人本身，如果他们毫无疑问地改正自己的行为，并通过悔改的充分保证澄清自己，我们并不认为他们被排除在天主不可言喻的仁慈之外；相反，我们与哀哭的人同哭，\BibleQuote{我们与哭泣的人一同哭泣，}\BibleRef{罗 12:15}并且服从公正的要求，在罢免时遵守正义，同时不忽视仁爱的补救措施；正如陛下所知，这不仅是口头承诺，也由我们的行动所证实，因为几乎所有那些曾被误导或被迫同意主持主教的人，通过撤销他们所颁布的、谴责他们所写的，已完全获得罪过的赦免和宗座和平的恩赐。\par

% structure: p0010 source=h2 target=subsection reason=relative-heading-rank; base=h2

\subsection{四、即使是祸首，也有可能通过悔改获得宽恕。}

因此，如果陛下屈尊反思我的动机，便会满意我自始至终的行动都是为了在不丧失一个灵魂的情况下废除异端；并且对于那些残酷骚乱的制造者，我稍作调整了我的做法，以便他们缓慢的心灵能被一些悔悟之情所激发，从而请求宽大的处理。因为尽管自从他们作出既亵渎又不公正的决定以来，他们不能再被公教弟兄团体像以前那样尊敬，但他们仍然保留着他们的主教座和主教品级，前景要么是在真正和必要的悔改迹象后接受整个教会的和平，要么是（如果天主禁止）坚持他们的异端，从而收获他们妄信的报应。日期：451年7月20日，在显赫的阿德尔菲乌斯执政期间。\par

\SourceRecord{Translated by Charles Lett Feltoe.；Nicene and Post-Nicene Fathers, Second Series；Vol. 12.；Edited by Philip Schaff and Henry Wace.；Buffalo, NY: Christian Literature Publishing Co.,；1895.；<http://www.newadvent.org/fathers/3604095.htm>.}

% structure: p0001 source=h1 target=section reason=page-title

\section{第九十八封书信}

\Emph{自加采东大公会议，致良}\par

伟大的、圣洁的普世大公会议，因天主的恩宠和我们最虔诚、爱基督的皇帝们的谕令，在\Place{彼提尼亚}省\Place{加采东}都会召开，致最圣洁、蒙福的罗马总主教良。\par

% structure: p0004 source=h2 target=subsection reason=relative-heading-rank; base=h2

\subsection{一、他们祝贺良在维护信德方面起了最重要的作用。}

\Quote{我们的口充满了喜乐，我们的舌头洋溢着欢欣。}这句预言，恩宠恰当地适用于我们，因为宗教的安全得到了保障。因为还有什么比信德更能激励喜悦呢？还有什么比神圣的知识更能激起欢欣呢？这知识是救主亲自从上赐给我们，为得救恩，说：\BibleQuote{你们去使万民成为门徒，因父、子及圣神之名给他们授洗，教导他们遵守我所吩咐你们的一切}\BibleRef{玛 28:19}。这条从命令的制定者延伸至我们的金链，您本人坚定地保存着，作为蒙福的\Person{伯多禄}对众人的代言人，并将他的信德之福分施给众人。因此，我们明智地以您为一切美善的向导，向教会的子女展示了他们承继的真理，不是各自私下传授，而是\OriginalTerm{在自负中}{in conceit}，同心合意地公开宣认我们的信德。我们都非常欢喜，如同在皇帝的宴席上，享受基督透过您的信函赐予我们的灵性食粮；我们似乎实际看见了天上的新郎与我们同在。因为如果\Quote{哪里有两三个人因祂的名聚在一起}，祂说\Quote{我就在他们中间}，那么，当有五百二十位司铎聚集时，他们宁愿传播关于祂的知识，而不顾自己的国家和安逸，祂岂不更特别临在吗？在这些司铎中，您以肢体之头的地位为首，通过那些代表您的人显示了您的善意；而我们虔诚的皇帝们主持了会议以促进秩序，邀请我们重建教会的教义结构，如同\Person{则鲁巴贝耳}邀请\Person{耶叔亚}重建耶路撒冷一样。\par

% structure: p0006 source=h2 target=subsection reason=relative-heading-rank; base=h2

\subsection{二、他们详述\Person{狄奥斯库若}的恶行。}

那仇敌本会像一头野外的野兽，独自咆哮，无力捉住任何人，若不是已故的亚历山大里亚主教投身作他的猎物。此人虽然先前做过许多可怕的事，但后来的行为更超过了前者；因为他违反所有教规的命令，撤职了君士坦丁堡那位蒙福的圣人牧者\Person{弗拉维安}——他显示了如此宗徒般的信德——以及最虔诚的主教\Person{欧瑟伯}，并通过恐吓获得的票数，开释了因异端被定罪的\Person{欧迪奇}，恢复了他被您的圣洁视为不配而剥夺的尊严。他如同最奇怪的野兽，扑向一棵状态极佳的葡萄树，将其连根拔起，并带回了那棵被丢弃为不结果子的树。他砍杀了那些作为真正牧者的人，让那些显为豺狼的人管理羊群。除此之外，他甚至向那位被救主托付看管葡萄树的人——我们指的是您的圣洁——发泄他的愤怒，并意图将教会合一的维护者逐出教会。他非但没有为此忏悔，没有含泪恳求怜悯，反而像做了善事一样欢欣，拒绝您的圣洁的信函，抗拒所有真理的教义。\par

% structure: p0008 source=h2 target=subsection reason=relative-heading-rank; base=h2

\subsection{三、我们已撤职\Person{欧迪奇}，尽可能仁慈地对待他。}

我们本应让他留在自己所处的位置；但是，既然我们宣讲救主的教导\BibleQuote{祂愿意所有的人得救，并得以认识真理}\BibleRef{弟前 2:4}，事实上，我们尽力对他实行这种仁慈的政策，以兄弟的方式传唤他受审，并非试图割断他，而是给他辩护和痊愈的空间。我们祈祷他能胜过他们对他提出的众多指控，以便我们能和平快乐地结束会议，使\Person{撒殚}无法占我们的便宜。但他因自己的良心备受定罪，通过回避审判，支持了指控，并拒绝了他收到的三次合法传唤。因此，我们以尽可能温和的方式确认了他因自己的错误对自己作出的判决，剥去了他长期以来被证明是伪装披着的牧羊皮。此后，我们的麻烦停止了，立即迎来了可喜的快乐时光。拔掉了一棵莠子，我们满心欢喜地用纯净的麦子充满了整个世界。领受了拔除和栽种的全权，我们将拔除限于一人，并精心栽种了一棵好果子。因为是天主在工作，得胜的\Person{欧弗弥亚}像为婚礼加冕一样加冕了会议，她将我们的信德定义作为自己的宣认，通过我们最虔诚的皇帝和爱基督的皇后献给新郎，平息了所有反对者的骚动，并确立我们对真理的宣认为祂所悦纳，用手和舌为我们在宣告中的全体投票盖印。这些是我们所做的一切，您以精神临在，并认同我们是兄弟，几乎通过您代表的智慧可见于我们。\par

% structure: p0010 source=h2 target=subsection reason=relative-heading-rank; base=h2

\subsection{四、他们宣布决定，\Place{君士坦丁堡}应在\Place{罗马}之后享有优先权，并请求\Person{良}的同意。}

我们还通知您，为了教会事务的良好管理和稳定，我们也决定了一些其他事项，我们相信，当您得知后，您的圣座将会接受并批准它们。长久以来的惯例，即天主在君士坦丁堡的神圣教会一直为亚细亚、本都和色雷斯各省任命都主教，我们现在已通过主教会议的投票予以批准，这与其说是赋予君士坦丁堡宗座特权，不如说是为了那些城市的良好治理，因为主教去世时经常发生混乱，神职人员和教友那时都没有领袖，扰乱了教会的秩序。这一点并未逃脱您的圣座注意，尤其是在厄弗所的情况，它常使您烦恼。我们还批准了在神圣记忆的大德奥多西时期于君士坦丁堡召开的150位圣教父的教规，该教规规定，在您最神圣的宗座之后，君士坦丁堡宗座应位列第二；因为我们相信，由于您平常对别人的关怀，您经常将属于您的宗徒权威也扩展到君士坦丁堡的教会，凭借您极大的无私，将您所有美好的事物与您的属灵亲属分享。因此，最神圣和蒙福的父亲，请垂允接受您自己所愿的，并有助于良好治理的这些决定，我们为消除一切混乱和确认教会秩序而作的决定。因为您的圣座的代表，最虔诚的主教帕斯加西诺和路森提乌斯，以及与他们一起的极敬虔的司铎波尼法爵，曾强烈反对这些决定，出于强烈的愿望，希望这一善工也从您的远见开始，以便良好秩序和信仰的确立都归功于您。因为我们恰当地考虑到我们最虔诚和爱基督的皇帝们，他们对此感到喜悦，以及显赫的元老院和几乎整个帝都，我们认为利用这次大公会议来批准您的荣誉是适时的，并确信这决定如同由您以惯常的培育热情所发起，知道子女的每一个成功都归于父母的荣耀。因此，我们恳求您，以您的同意来尊重我们的决定，正如我们已向元首屈服，在荣誉之事上达成一致，同样愿元首也为子女成就合宜之事。因为这样，我们虔诚的皇帝们将受到应有的尊重，他们已批准您的圣座判断为法律，而君士坦丁堡宗座也将因其在宗教事务上始终对您表示如此忠诚，并如此热心地与您完全合而为一而得到报偿。但为使您知道我们并非出于偏爱或仇恨行事，而是受天主圣意引导，我们已将我们全部行动的范围告知您，以巩固我们的立场，并批准和确立我们所做的一切。\par

\SourceRecord{Translated by Charles Lett Feltoe.；Nicene and Post-Nicene Fathers, Second Series；Vol. 12.；Edited by Philip Schaff and Henry Wace.；Buffalo, NY: Christian Literature Publishing Co.,；1895.；<http://www.newadvent.org/fathers/3604098.htm>.}

% structure: p0001 source=h1 target=section reason=page-title

\section{第一○一书}

\Emph{来自\Place{君士坦丁堡}主教\Person{安纳托利乌斯}，致\Person{莱奥}。}\par

\EditorialNote{本函论述与安纳托利乌斯自述的第98书大致相同的内容；此处翻译第三章，以说明第98书第三章。}\par

% structure: p0004 source=h2 target=subsection reason=relative-heading-rank; base=h2

\subsection{三、他描述大公会议制定降生成人教义的情形。}

但既然在对他作出判决后，我们必须以祈祷和眼泪就正确信德的定义达成一致；因为那是皇帝召集神圣大公会议的主要原因，你的圣德在精神上与我们同在，并通过你派来的敬畏天主的人与我们同工；我们有了至圣荣美的殉道者圣\Person{欧斐米亚}的保护，便全体极其审慎地投入这件重要的事。而由于时机要求所有聚会的圣主教们发表一致的决定，为清晰并明确宣认我们的主耶稣基督、主天主的信德，祂甚至为不寻求祂的人所寻见和启示，是的，也为不询问祂的人所显明\BibleRef{依 65:1}，尽管最初有一些抗拒的尝试，但祂依然向我们显示了祂的真理，并命令所有人都一致无异议地写下并宣告它，从而坚固了坚强者的灵魂，并邀请所有偏离真理的人走上真理之道。事实上，在一致签署这份文件后，我们这些聚集在此届大公会议上的人，以同一至圣得胜的殉道者圣\Person{欧斐米亚}的信德，以及我们最虔诚、爱基督的皇帝\Person{马尔基安}，和我们最虔诚、在一切事上最忠实的女儿皇后\Person{普尔喀丽亚}·奥古斯塔的名义，在祈祷、喜乐和幸福中，将按照您的圣书所写的定义放在圣祭台上，以确认我们父辈的信德，并将其呈献给他们虔诚的照顾；因为他们曾请求接受它，而接受之后，他们与我们一同光荣了他们的主人基督，祂驱散了异端的一切迷雾，并恩赐地显明了真理的言语。如此，藉着救主的仁慈，教会及司铎们关于纯正信德的一致意见同时得以建立。\par

\SourceRecord{Translated by Charles Lett Feltoe.；Nicene and Post-Nicene Fathers, Second Series；Vol. 12.；Edited by Philip Schaff and Henry Wace.；Buffalo, NY: Christian Literature Publishing Co.,；1895.；<http://www.newadvent.org/fathers/3604101.htm>.}

% structure: p0001 source=h1 target=section reason=page-title

\section{第一〇四封信}

致\Person{马尔基安·奥古斯都}，关于\Person{阿纳托利}的僭越，经由主教\Person{路济安}与执事\Person{巴西略}之手。\par

% structure: p0003 source=h2 target=subsection reason=relative-heading-rank; base=h2

\subsection{一、他祝贺皇帝在公教信仰胜利中的分享。}

\Emph{\Person{莱奥}，主教，致\Person{马尔基安·奥古斯都}}\par

借着天主慈恩的浩大恩赐，当通过您的圣洁光荣之热忱，那最有害的谬误在我们中被铲除时，整个天主教会的喜乐增加了；因此我们的劳苦更迅速地达到了预期的目的，因为您事奉天主的陛下如此忠信而有力地协助了它们。因为虽然福音的自由必须借着圣神的能力和经由\OriginalTerm{宗徒的宗座}{Apostolic See}的中介来捍卫，以对抗某些异议者，然而天主的恩宠比我们所期望的更显明地彰显了，祂恩赐世界：在真理的胜利中，只有违反信仰的始作俑者应当灭亡，而教会恢复其健全。因此，那由我们和平之敌所挑起的战争，主用右手为我们作战，如此幸福地结束了，以致当基督凯旋时，祂所有的司铎都分享了同一胜利，当真理之光显现时，只有谬误的阴影及其拥护者被驱散。因为正如在相信主的复活时，为坚固信仰的开端，由于某些宗徒怀疑了我主耶稣基督的身体实在，并通过看见和触摸检查了钉痕和枪伤，因怀疑而消除了众人的怀疑，信心大增；同样现在，当一些人的错误信仰被驳斥时，所有犹豫者的心得到坚固，那曾使少数人瞎眼的事，却有助于整个身体的启蒙。在这项工作中，您的仁慈理所当然地欢喜，因为您忠信而恰当地提供了防范，使得魔鬼的陷阱不会伤害东方教会，反而随处献上更中悦天主的全燔祭；因为借由天主与人之间的中保，就是那为人的基督耶稣，人民、司铎和君王都持守同一信经。啊，最光荣的儿子和最仁慈的奥古斯都。\par

% structure: p0006 source=h2 target=subsection reason=relative-heading-rank; base=h2

\subsection{二、综合考虑所有情况，\Person{阿纳托利}本应表现出更多的谦逊。}

而今，关于这些聚集了如此众多司铎的大事，已经得到良好而可欲的结局，我惊讶且痛心的是，那原已由天主恢复的普世教会的和平，竟又被一种自求私利的精神所扰乱。因为，虽然我的弟兄\Person{阿纳托利}似乎必然地为了自身利益而离弃了那些祝圣他之人的错误，并以有益的转变接受了天主教信仰，但他本应小心，不要因任何卑劣的欲望而玷污那经您手段所获得的成果。因为，我们顾念到您的信德和中介，尽管因着那些祝圣他的人，他的前科可疑，我们宁愿对他仁慈而非公正，以便通过医治措施平息所有因魔鬼的运作而激起的骚动；而这本应使他谦卑，而非相反。因为，即使他因卓越的功绩、最明智的选拔而被合法且按常规祝圣，然而若不尊重\OriginalTerm{教父们的教规}{canons of the Fathers}、圣神的法令和古代的成例，任何选票也不能对他有利。我在一位基督徒、真正虔诚、真正正统的君王面前说：\Person{阿纳托利}主教因贪求不当的显赫，极大地减损了他本应有的功绩。\par

% structure: p0008 source=h2 target=subsection reason=relative-heading-rank; base=h2

\subsection{三、\Place{君士坦丁堡}城，尽管是皇城，却绝不能升到\OriginalTerm{宗徒的等级}{Apostolic rank}。}

让\Place{君士坦丁堡}城，如我们期望的那样，保持其高位，并在天主的右手保护下，长久享受您陛下的统治吧。然而，世俗之事与神圣之事根基不同；除非建立在主所安放为根基的那磐石上，否则没有稳固的建筑。凡贪图非其应得之物者，必失其所有。对\Person{阿纳托利}来说，借由您的虔诚之助和我的恩准与赞同，他获得了如此重要城市的主教职，这应已足够。让他不要轻视这座皇城，尽管他不能使它成为\OriginalTerm{宗徒的宗座}{Apostolic See}；并且让他绝不要希望，他可以通过损害他人而上升。因为由圣教父们的教规所决定、由\OriginalTerm{尼西亚大公会议}{Nicene Synod}的法令所固定的教会特权，不能被任何无耻的行为推翻，也不能被任何革新所扰乱。并且，在忠实执行这项任务中，借由基督的助佑，我必表现出坚定不移的奉献；因为这是一项交托给我的职责，如果那些由教父们所批准、在天主圣神引导下在\Place{尼西亚}大公会议上为整个教会管理而制定的规则，因我的默许而被违反（愿天主禁止），或如果单个弟兄的愿望比主整个家庭的公益对我更有分量，这就会导致我的定罪。\par

% structure: p0010 source=h2 target=subsection reason=relative-heading-rank; base=h2

\subsection{四、他请求皇帝表示对\Person{阿纳托利}自寻私利精神的不悦。}

因此，深知您光荣的仁慈关怀教会的和平，并保护和支持那些促进和平统一的措施，我恳切地祈求并哀求您，拒绝给予这些违背基督徒合一与和平的肆无忌惮的企图任何支持与保护，并对我的兄弟 \Person{阿纳托略} 的欲望加以有益的约束——如果他坚持，只会损害他自己：愿他不再渴望那些有损于您荣耀和时代需要的事物，不愿超越他的前任，并让他自由地尽其在德行上显赫，而他只有在以爱德装饰而非野心膨胀时，才能分享德行。这种过分愿望的构想，他本不应在心中接受，但当我代表在场的兄弟和同为主教们反对他时，他至少可以在他们有益的反对中放弃他那不合法的自我追求。因为您仁慈的陛下和他的信都证实，宗座的代表们以最正当的抵抗反对了他，他的傲慢因此更不可原谅，因为即使在受责备时，他也没有克制自己。\par

% structure: p0012 source=h2 target=subsection reason=relative-heading-rank; base=h2

\subsection{五、并力求使他回心转意。}

因此，因为您光荣的信德相称于，正如异端已被推翻，天主藉您行事，如今一切自我追求也应被击败，请您做那符合您基督徒和君王良善之事，以使所述主教服从教父，促进和平事业，不再认为他有权为\Place{安提约基雅}教会祝圣主教，如同他毫无先例且违反教会法规擅自所为：我们出于重建信仰的热望和为了和平，决定不取消这一行为。因此，让他停止藐视教会规矩，避免不法过犯，以免在试图破坏和平之事时，将自己与普世教会隔绝。我宁愿因他行为无过而爱他，而不愿发现他坚持这种可能使我们所有人分离的傲慢心态。我的兄弟和同为主教 \Person{路奇安}，与我的儿子、执事 \Person{巴西略}，将陛下仁慈的信函带给了我，他以全部虔诚履行了他作为代表所承担的任务：因为不能认为他的使命失败了，而是事态发展辜负了他。\Signature{于光荣的 \Person{赫尔库拉努斯} 执政之年五月二十二日（452年）。}\par

\SourceRecord{Translated by Charles Lett Feltoe.；Nicene and Post-Nicene Fathers, Second Series；Vol. 12.；Edited by Philip Schaff and Henry Wace.；Buffalo, NY: Christian Literature Publishing Co.,；1895.；<http://www.newadvent.org/fathers/3604104.htm>.}

% structure: p0001 source=h1 target=section reason=page-title

\section{第一〇五封信}

\Emph{（致普尔喀丽亚·奥古斯塔，关于阿纳托利乌斯的自私自利。）}\par

% structure: p0003 source=h2 target=subsection reason=relative-heading-rank; base=h2

\subsection{一、祝贺皇后战胜异端，但遗憾教会中又引入新争论。}

\Person{良}主教致\Person{普尔喀丽亚·奥古斯塔}。\par

我们与陛下同享不可言喻的喜乐，因为公教信仰已战胜异端，并且通过陛下虔诚而中悦天主的热情，和平已恢复于整个教会：为此我们感谢仁慈全能的天主，祂不允许任何人——除非那些喜爱黑暗胜过光明的人——被剥夺福音真理：以致藉着谬误迷雾的消散，最纯净的光辉能在众人心中升起，而那位喜爱黑暗的仇敌不能战胜某些软弱的灵魂，他们不但被未受伤的人击败，而且也被那些他曾使之动摇的人击败；藉着谬误的废除，真信仰得以统治世界，并且\Quote{万口都要承认主耶稣基督在父天主的光荣中\BibleRef{斐 2:11}。}但当全世界已在福音的合一中稳固，所有司铎的心已导向同一信念时，除了那些为神圣大公会议而召集并经由陛下热情达成圆满协议的事宜之外，本不应引进任何对抗如此巨大利益的事物，也不应使主教会议成为不适时推进非法欲望的借口。\par

% structure: p0006 source=h2 target=subsection reason=relative-heading-rank; base=h2

\subsection{二、尼西亚典章不可更改，对普世有约束力。}

因为我的弟兄和同僚主教\Person{阿纳托利乌斯}，没有充分考虑到陛下的仁慈和我赞同的恩惠——他藉此获得了君士坦丁堡教会司祭职——反而对所获不知感恩，被超乎其等级的过度欲望所灼烧，相信他那无节制的自私自利能通过声称某些人已用强迫签署表示同意来推进；尽管我的弟兄和同僚主教们——他们代表我——对这些注定迅速失败的企图，忠实而可赞地表示反对。因为没有人敢做出与教父典章相悖的事，这些典章在许多年前于\Place{尼西亚}城，基于圣神的法令而确立，以致任何人若想通过不同决议，只会伤害自己，而非损害它们。如果所有主教都如其所当保持它们不受侵犯，那么所有教会中必有完美的和平与完全的和谐：将不会有关于等级的争执，不会有关于祝圣的争论，不会有关于特权的争议，不会有关于夺取他人之物的争斗；而是藉着爱的公平律法，在行为与职务上保持合理秩序，而真正伟大者将是那种远离一切自私自利的人，正如主说：\Quote{谁若愿意在你们中成为大的，就作你们的仆役；谁若愿意在你们中为首，就作你们的奴仆。就如人子来不是为受服事，而是为服事人\BibleRef{玛 20:26}。}然而这些训诫当时是赐给那些希望从卑微地位上升，从最低处走向最高处的人；但君士坦丁堡教会的统治者还要贪求什么超过他所获得的呢？或者什么才能使他满足，如果这么伟大城市的辉煌与声名仍不够？如此越出一切适当界限，践踏古老习惯，想夺取他人权利，实在过于傲慢无度了：以损害这么多都主教首席权为代价来提高一个人的尊严，并在早已由神圣尼西亚大公会议法令平静的省份中，引入一场新的混乱战争：破坏可敬教父们的法令，借口某些主教的同意，而这些同意甚至经过这么多年也未能生效。因为有人夸耀说这已被睁一只眼闭一只眼将近六十年了，而该主教认为这对他有利；但他妄图从那些即使有人敢期望也永远无法获得的东西中获得帮助，是徒劳的。\par

% structure: p0008 source=h2 target=subsection reason=relative-heading-rank; base=h2

\subsection{三、只有效法其前任，他才能重新赢得良的信任：主教的同意被宣告无效。}

让他认清他继承的是怎样一个人，驱除一切骄傲精神，效法\Person{弗拉维安}的信仰、\Person{弗拉维安}的谦逊、\Person{弗拉维安}的谦卑，这已将他提升至宣认者的光荣。如果他能以其德行发光，他必获得一切赞美，并在各方赢得丰厚的爱，不是通过寻求人的升迁，而是通过配得天主的恩宠。藉着这种谨慎的作风，我保证也会将我的心与他联结，而宗座对君士坦丁堡教会一贯施予的爱，绝不会因任何改变而受损。因为即使统治者有时因缺乏节制而陷入错误，但基督的教会并不失去其纯洁。至于那些反对在\Place{尼西亚}与陛下忠信的合作下所制定神圣典章规定的主教同意，我们不承认，并以荣福宗徒\Person{伯多禄}的权威全面彻底地予以废除，在所有教会事务中，我们遵守圣神藉三百一十八位主教为所有司铎的和平遵守所颁布的法律，其方式为：即使有更多数的人通过与他们不同的决议，凡是反对他们法令的，都必须视为无效。\par

% structure: p0010 source=h2 target=subsection reason=relative-heading-rank; base=h2

\subsection{四、请求皇后惠予考虑此信。}

因此，我恳请阁下以相称的精神接纳这封长信，我在此信中不得不阐述我的观点。这封信将由我的兄弟及同为主教者路济亚诺（\Person{路济亚诺}）呈递——他尽其所能，忠实地履行了作为我的代表所肩负的焦虑职责——以及我的儿子执事巴西略（\Person{巴西略}）呈递。又因您素来致力教会的和平与合一，为使他的灵魂获益，请将我的兄弟阿纳托利主教（\Person{阿纳托利}）保持在对他有益的限度内——我曾因您的建议而向他表达了我的爱——以便您阁下仁爱的荣耀因信德的恢复已然增光，也因克制私欲而传扬在外。于赫库拉努斯（\Person{赫库拉努斯}）任执政官之年（452年）五月二十二日。\par

\SourceRecord{Translated by Charles Lett Feltoe.；Nicene and Post-Nicene Fathers, Second Series；Vol. 12.；Edited by Philip Schaff and Henry Wace.；Buffalo, NY: Christian Literature Publishing Co.,；1895.；<http://www.newadvent.org/fathers/3604105.htm>.}

% structure: p0001 source=h1 target=section reason=page-title

\section{第一〇六封书信}

\Emph{致君士坦丁堡主教\Person{阿纳托利}，责备其自利之心。}\par

% structure: p0003 source=h2 target=subsection reason=relative-heading-rank; base=h2

\subsection{一、他赞赏\Person{阿纳托利}的正统信仰，但谴责其狂妄。}

\Person{良}主教，致\Person{阿纳托利}主教。\par

如今，藉天主的恩宠，福音真理的光照已如我们所愿地彰显，最恶毒的谬误之夜已从普世教会驱散，我们在主内不胜喜乐，因为托付给我们的艰巨使命已如愿达成，正如你的信函所宣告，按照宗徒的教训：\BibleQuote{我们都说一样的话，我们中间不可有分裂；但要一心一意，彼此相合。}\BibleRef{格前 1:10} 在这项工作中，我们称赞你，亲爱的，因你参与了：你藉自己的活动帮助了需要纠正的人，并使自己摆脱与犯法者的一切同谋。当你的前任、荣福弗拉维盎因捍卫天主教真理而被罢免时，人们理所当然地认为，祝圣你的人似乎祝圣了一个与他们相似的人，违反了神圣法令的规定。但天主的仁慈在此引导并坚固了你，使你善用不良的开端，表明你被提升不是出于人的判断，而是出于天主的慈爱：这可以被视为真实，只要你不再因另一种冒犯的缘由而丧失这神圣恩赐的恩宠。因为天主教信徒，尤其是主的司铎，不仅不可陷入任何错误，也不可被任何贪婪所腐蚀；正如圣经所说：\Quote{不要顺从你的欲情，要远离你的私欲。} 必须抵制许多世俗的诱惑和许多虚荣，才能达到真正自律的圆满，其首要污点是骄傲——过犯的开端和罪恶的起源。因为贪图权柄的心灵既不知道禁止的不可行，也不知道允许的不可享，只要过犯不受惩罚，便陷入无纪律和邪恶的过度行为，恶行倍增，这些只是我们为恢复信仰的热忱与和谐之爱而忍受的。\par

% structure: p0006 source=h2 target=subsection reason=relative-heading-rank; base=h2

\subsection{二、尼西亚法令不可取消或修改。}

因此，在你祝圣的并非无可指责的开端之后，在你擅自为安提约基雅主教行祝圣礼（违反法令条例）之后，我深感忧伤，亲爱的，你竟陷入这等事：试图推翻尼西亚法令最神圣的规程，仿佛这一机会特意向你提供，为使亚历山大里亚宗座失去第二位的特权，安提约基雅教会放弃其第三位尊严的权利，以便这些地方置于你的管辖之下后，所有都主教都被剥夺应有的尊荣。凭借这种前所未闻且从未尝试过的越分行为，你竟超越自身，将神圣会议拖入自利的机会，并强迫这次由我们最虔诚的基督信仰君王为消灭异端、确认天主教信仰而召集的会议默许；仿佛多数人的非法愿望不可拒绝，且圣神在尼西亚法令中真正制定的秩序可由任何人部分推翻。任何议会会议不可因大会规模而自夸，任何数量更大的司铎不可敢于将自己与那三百一十八位主教相比或自认优于他们，因为尼西亚大公会议被天主尊以如此特权，以致无论藉更多或更少的人支持教会判决，凡反对其权威的，全然失去一切权威。\par

% structure: p0008 source=h2 target=subsection reason=relative-heading-rank; base=h2

\subsection{三、加采东大公会议为单一目的召开，绝不应被用于其他目的。}

因此，那些与最神圣法令相悖的事是极其无原则和谬误的。这种傲慢自大导致整个教会的动荡，它意图如此滥用一次议会会议，藉邪恶的论据说服或藉恐吓强迫弟兄们同意它，而他们被召集仅为信仰事务，并已就应关注的主题作出决定。正是基于此，我们的兄弟——由宗座派遣、以值得称赞的坚定主持我们在此次会议者——对违法企图提出了抗议，公开反对任何违背尼西亚大公会议法规的应受责备的革新。他们的反对无疑存在，因为你在书信中抱怨他们企图违背你的尝试。你在信中如此向我推荐他们，从而大大称赞了他们，而你在拒绝服从他们关于你不法图谋的事上控告自己，徒然寻求无法获准的事，渴求有害你灵魂健康、且永不能赢得我们同意的事。愿我永不有罪助长如此错误的愿望，这愿望更应藉我及所有不图高远之事、而与卑微者同心之人的协助而被推翻。\par

% structure: p0010 source=h2 target=subsection reason=relative-heading-rank; base=h2

\subsection{四、尼西亚法令普遍适用，不得曲解为私人解释。}

这些神圣可敬的教父们在尼西亚城谴责亵渎者亚略及其不敬之后，为教会制定了直至世界终结生效的法令全书；他们的宪章不仅在我们中间，也在全人类中永存。若有人胆敢违反其法令，则\OriginalTerm{事实上}{ipso facto}无效；因此，为我们的永久利益而普遍制定的事，永不可因任何改变而修改，也永不可将原本为共同利益而设的事歪曲为私人利益。如此，只要教父们所定的界限存留，无人可侵犯他人权利，每人要在适当合法的范围内，按自己的能力，在爱的广博中操练自己；君士坦丁堡主教若能充分收获其果实，只要他依靠谦卑之德，而不因自利之心自高自大。\par

% structure: p0012 source=h2 target=subsection reason=relative-heading-rank; base=h2

\subsection{五、据称六十年前核准君士坦丁堡对亚历山大和安提约基雅至高权的文件毫无价值。}

\Quote{不可心高气傲，} 弟兄， \Quote{但要恐惧\BibleRef{罗 11:20}，} 并且停止以无理的要求扰乱基督徒君王虔敬的耳朵；我相信，你的谦逊比你的骄傲更能取悦他们。因为你的意图绝得不到某些主教书面同意的支持，如你所称，这同意是六十年前给出的，且从未由你的前任们呈报宗座；这项从一开始就注定失败且早已失效的交易，你现在却想用为时已晚且无用的手段来支撑，即从弟兄们那里榨取同意的假象，他们的谦逊因厌烦而损害自身。请记住主对那使一个小子跌倒之人的警告，并要有智慧去理解，那不怕使如此多的教会和如此多的司铎跌倒的人，将承受何等天主的审判。因为我承认，我被对全体弟兄的爱如此束缚，以致我不会同意任何人提出违背自身利益的要求；因此你可以清楚看出，我的反对，亲爱的，是出于善意，要以更明智的劝告约束你，免得你扰乱普世教会。省区首席主教的权利不可推翻，都主教也不能被剥夺基于传统的特权。\Place{亚历山大}宗座不可失去它通过\Person{圣马尔谷}——圣史和\Person{圣伯多禄}的门徒——所赢得的任何尊严，如此伟大教会的荣光也不可被他人的阴云遮蔽，因为\Person{狄奥斯哥罗}因坚持不虔敬而跌倒。\Place{安提约基雅}教会，在那里因\Person{圣宗徒伯多禄}的宣讲，基督徒的名号首次出现\BibleRef{宗 11:26}，必须继续持守教父所赋予的地位，既被置于第三位，绝不可从那里降低。因为宗座与它的拥有者处于不同层面；每个人的主要尊荣在于自身的正直。既然正直在任何地方都不会失去其应有的价值，那么当它被置于\Place{君士坦丁堡}城的壮丽之中时，该是多么光荣啊！在那里，许多司铎既能找到教父教规的辩护，也能在遵守你的过程中找到正直的榜样？\par

% structure: p0014 source=h2 target=subsection reason=relative-heading-rank; base=h2

\subsection{六、基督徒的爱德要求舍弃自我，而非追求私利。}

弟兄，我这样写信给你们，是在主内劝勉和告诫你们，要放下一切野心，反而珍爱爱德的精神，并以爱德的德行装饰自己，使自己得益，按宗徒的训诲。因为爱德\Quote{含忍而慈祥，不嫉妒，不行不义，不自大，不贪慕，不求己益\BibleRef{格前 13:4}}。因此，如果爱德不求己益，那么贪图他人之物的人犯了多大的罪！我愿你们完全保守自己，并记住那句话说：\Quote{你要持守你所有的，免得别人夺去你的冠冕\BibleRef{默 3:11}}。因为如果你寻求不允许的事，你将因自己的行动和判断，剥夺自己普世教会的和平。我们的弟兄和同席主教\Person{路济安}，以及我们的儿子执事\Person{巴西略}，以他们所拥有的一切热忱遵行了你的训令，但正义拒绝使他们的请求生效。日期：显赫的\Person{赫尔库拉诺斯}执政官任期（四五二年）五月二十二日。\par

\SourceRecord{Translated by Charles Lett Feltoe.；Nicene and Post-Nicene Fathers, Second Series；Vol. 12.；Edited by Philip Schaff and Henry Wace.；Buffalo, NY: Christian Literature Publishing Co.,；1895.；<http://www.newadvent.org/fathers/3604106.htm>.}

% structure: p0001 source=h1 target=section reason=page-title

\section{第一〇八书信}

\Person{良}，主教，致\Person{德奥多禄}，\Place{尤里乌姆广场}主教。\par

% structure: p0003 source=h2 target=subsection reason=relative-heading-rank; base=h2

\subsection{一、\Person{德奥多禄}本不应未经其都主教而接近他。}

你首先应当做的，是在焦虑时，就那似乎需要查问的事项，咨询你的都主教；如果他也不能帮助你，亲爱的，你们二人应请求（由我们）予以指教；因为凡涉及全体主的司铎的事务，未经首席主教不得查问。但为使咨询者的疑虑无论如何得以消除，我不隐瞒教会关于忏悔者处境的规则。\par

% structure: p0005 source=h2 target=subsection reason=relative-heading-rank; base=h2

\subsection{二、补赎的恩宠是为那些在领洗后跌倒的人。}

天主的丰盛仁慈如此帮助跌倒的人，以致不仅藉着圣洗的恩宠，也藉着补赎之药，永生的希望得以重振，好使那些玷污了第二次诞生的恩赐、凭自己的判断定自己罪的人，能获得对其罪过的赦免；神圣美善的安排如此规定：不藉司铎的恳求，便不能获得天主的宽恕。因为天主与人之间的中保，基督耶稣那人，将这种权能传授给了那些治理教会的人，使他们既应许给告罪者一条补赎的道路，又在他们藉有益之惩戒洁净后，通过和好之门接纳他们参与圣事的共融。在这项工作中，救主本身确实无间地参与，从不远离那些祂托付给祂的仆人去执行的事，祂说：\BibleQuote{看，我同你们天天在一起，直到今世的终结：}\BibleRef{玛 28:20}因此，凡藉我们的服务按正当秩序并获满意成果而完成的一切，我们确信是藉圣神所恩赐的。\par

% structure: p0007 source=h2 target=subsection reason=relative-heading-rank; base=h2

\subsection{三、补赎只有在今生才确定。}

但是，若我们所哀求天主为他代祷的人中，有人因某些阻碍被耽误，失去立即赦免的恩惠，并在未获指定的补救之前，按自然进程结束此生，那么，当他脱去肉体时，便不能获得那尚在肉身时未曾领受的。我们也不必权衡这样死去之人的功德与行为，因为主我们的天主，祂的判断不可揣测，已将我们司铎职务未能完成之事留给祂自己定夺：因为祂愿意祂的权能被如此敬畏，以致这敬畏造福众人，每个人都畏惧那发生于冷淡或疏忽之人身上的事。因为最有益且必要的是，罪的债应在生命最后一日之前，藉司铎的恳求得解除。\par

% structure: p0009 source=h2 target=subsection reason=relative-heading-rank; base=h2

\subsection{四、但补赎与和好不可拒绝给临终之人。}

但那些在需要之时和紧迫危险中首先恳求补赎之助，而后恳求和好之助的人，既不可禁止他们改过自新的方法，也不可禁止和好：因为我们不能为天主的仁慈设定界限，也不能为祂规定时间——对祂而言，真正的皈依不会使赦免延迟，正如天主的圣神藉先知所说：\BibleQuote{你若回头哀哭，就必得救；}又在别处说：\BibleQuote{你当预先陈明你的过犯，好使你成义；}又说：\BibleQuote{因为上主富于仁慈，祂有丰富的救赎。}因此，在分施天主的恩赐时，我们不可苛刻，也不可忽视自责者的眼泪与叹息，因为我们认为，补赎的情感本身出自天主的感动，正如宗徒所说：\BibleQuote{或许天主赐予他们悔改，使他们摆脱魔鬼的圈套，因为魔鬼已俘虏他们，使之顺从其意愿。}\BibleRef{弟后 2:25}\par

% structure: p0011 source=h2 target=subsection reason=relative-heading-rank; base=h2

\subsection{五、临终悔改虽有危险，但即使只能以记号请求时，也不得拒绝赦罪的恩宠。}

因此，每个基督徒都应听从自己良心的判断，以免他日复一日地推迟归向天主，而将改过的时间定在生命之末；因为人性的软弱与无知若将自己局限于这样的境况，以致沦为数小时的不确定性，并以更充分的改过来赢得宽恕，反倒选择那几乎没有空间容忏悔者告解或司铎赦罪的时间的狭隘界限，这是极其危险的。但是，如同我已说过的，纵使是这些人的需要也必须如此协助，即使他们在声音已失的情况下，通过仍然清晰的理智的记号来要求，也不可拒绝他们补赎的自由行动和共融的恩宠。如果他们被疾病的压力所压倒，以致不能在司铎面前表示他们此前所请求的，那么，在场信徒的证言必须为他们奏效，好使他们同时获得补赎和和好的恩惠，只要在那些因背弃信仰而得罪天主的人身上，遵守教父的教规。\par

% structure: p0013 source=h2 target=subsection reason=relative-heading-rank; base=h2

\subsection{六、他应将此信呈报都主教。}

兄弟，这些是我对你问题的答复，为使任何人不得以无知为借口而行事不同，你当将此呈报你的都主教；倘若此前有任何弟兄认为这些点上有疑问，他们可由他依据我所写给你的内容得到指教。日期：\Person{赫古拉努斯}阁下执政之年（452年）六月十一日。\par

\SourceRecord{Translated by Charles Lett Feltoe.；Nicene and Post-Nicene Fathers, Second Series；Vol. 12.；Edited by Philip Schaff and Henry Wace.；Buffalo, NY: Christian Literature Publishing Co.,；1895.；<http://www.newadvent.org/fathers/3604108.htm>.}

% structure: p0001 source=h1 target=section reason=page-title

\section{第一〇九封书信}

\Emph{致科斯主教尤利安}\par

% structure: p0003 source=h2 target=subsection reason=relative-heading-rank; base=h2

\subsection{一、他为\Place{巴勒斯坦}近期的骚乱而哀叹。}

教宗\Person{良}致\Person{尤利安}主教。\par

贤弟，你关于假隐修士骚乱行径的报告是严重的，并且令人不胜哀伤；因为这些骚乱源于邪恶的\Person{欧提克斯}藉着欺骗者的疯狂对福音和宗徒的宣讲所发动的战争，尽管这战争将以他自己及其追随者的毁灭而告终；但因天主的宽容而延迟，为要显明基督十字架的仇敌是如何深深地被魔鬼所奴役；因为异端邪说撕破其古老的伪装面纱，不能再约束于伪善的界限之内，已经倾泻出其长期隐藏的毒液，不只用笔杆，而且用暴力行为来攻击真理的门徒，以迫使无知的单纯或被恐惧所慑的信德屈服。但是光明之子不应如此惧怕黑暗之子，以致明智的人同意疯狂者的想法，或认为应对这种人表示任何尊重；因为，如果他们宁愿灭亡也不恢复理智，就必须作出安排，以免他们逃脱惩罚会造成更广泛的伤害，并且长期容忍他们会导致许多人的毁灭。\par

% structure: p0006 source=h2 target=subsection reason=relative-heading-rank; base=h2

\subsection{二、必须将首要分子迁往远处。}

我并非不知道，我们那些神圣而真正的隐修士——他们不放弃修道的节制，并实践他们誓愿所承诺的——理应受到何等爱戴和恩惠。但是这些傲慢的扰乱者，他们以对司铎的侮辱和伤害为荣，不应被视为基督的仆人，而是假基督的兵卒，并且主要应在他们的领袖身上被羞辱，这些领袖煽动无知的暴众支持他们的不顺服。因此，鉴于我们极其仁慈的君侯以虔敬心灵的整个热忱热爱天主教信仰，并且对这帮叛逆异端者的厚颜无耻深感愤怒，正如各处所报道的，我们必须恳求他的仁慈，使这些骚乱的煽动者从他们的疯狂集会中移除；不仅\Person{欧提克斯}和\Person{迪奥斯库若}，而且任何曾积极支持他们错谬疯狂之人，都应被安置在不能与他们的亵渎同伙联络的地方；因为这样或可治愈某些人的单纯，并且如果人们从有害教师的煽动中解脱出来，将更容易被召回健全的心智。\par

% structure: p0008 source=h2 target=subsection reason=relative-heading-rank; base=h2

\subsection{三、他寄送圣\Person{亚大纳削}的一封书信，以表明当前的异端不过是昔日被摒弃异端的复活。}

但是，为使必要教导——或为坚固忠实的心灵，或为驳斥异端——不致缺乏或表达不明，我已寄送了神圣记忆的主教\Person{亚大纳削}致主教\Person{埃皮克托}的书信，神圣记忆的\Person{济利禄}曾在\Place{厄弗所}大公会议上引用该书信反对\Person{奈斯多利}，因为它如此清晰而详尽地阐述了圣言的降生成人，以致预先推翻了那些时代的异端中的\Person{奈斯多利}和\Person{欧提克斯}。让\Person{欧提克斯}和\Person{迪奥斯库若}的追随者胆敢指责如此权威的无知或异端吧——他们声称我们的宣讲背离了圣教父的教导和知识。但是，这应当有助于坚固所有主的司铎的心灵，他们已经被发现并因所追随的权威而被判为异端，如今开始更公开地陈述他们亵渎的教条，免得，若其含义隐藏在沉默的外衣之下，是否\Person{阿波林}的三重谬误和\Person{摩尼}教徒的疯狂念头真的在他们身上复活仍存疑问。既然他们不再寻求隐藏，而是大胆起来攻击基督的教会，我们岂能不设法摧毁他们一切企图的力量，正如我所说的，持守这样的辨别力，以便将不可救药者与较温顺的心灵分开？因为\Quote{恶劣的交谈败坏良好的品行\BibleRef{格前15:33}，} 并且\Quote{明智者将比受惩罚的瘟疫之人更锋利；} 这样，无论邪恶者的团体以何种方式被打破，有些器皿或可从魔鬼手中被夺回？因为我们不应被粗鄙的空话所激怒，以致不顾他们的改正。\par

% structure: p0010 source=h2 target=subsection reason=relative-heading-rank; base=h2

\subsection{四、他表示希望\Person{尤文纳尔}及时认错的态度会为其余人所效仿。}

但是主教\Person{尤文纳尔}（他的伤害是可悲的）过于轻率地加入了那些亵渎的异端者，并因接纳\Person{欧提克斯}和\Person{迪奥斯库若}而使许多无知之人因他的榜样而鲁莽行事，尽管他后来以更明智的劝告纠正了自己。然而，那些更贪婪地饮下邪恶毒药的人，成了他们先前曾为弟子之人的敌人，以至于他所供应的食物反过来成了他自己的毁灭；尚且希望他们会在改正缺点方面效仿他，只要他们居住的邻近圣洁团体帮助他们恢复理智。但是，篡夺了仍在世主教职位之人的品性，从其所行的事迹中无可置疑，并且受到信仰攻击者所爱之人必定是个不信者。同时，贤弟，不要犹豫，以焦虑的关怀通过更频繁的书信让我了解事态的进展。写于\Person{赫库兰努斯}执政年十一月廿五日（主历四五二年）。\par

\SourceRecord{Translated by Charles Lett Feltoe.；Nicene and Post-Nicene Fathers, Second Series；Vol. 12.；Edited by Philip Schaff and Henry Wace.；Buffalo, NY: Christian Literature Publishing Co.,；1895.；<http://www.newadvent.org/fathers/3604109.htm>.}

% structure: p0001 source=h1 target=section reason=page-title

\section{第一一三封信}

\Emph{致科斯主教尤利安}\par

% structure: p0003 source=h2 target=subsection reason=relative-heading-rank; base=h2

\subsection{一、感谢尤利安的同情后，他控告阿埃提乌斯被从总执事之职免去。}

\Person{良}，罗马主教，致\Person{尤利安}，\Place{科斯}主教。\par

我于来信中，亲爱的，认出了兄弟之爱的情愫，因你真诚地同情我们所承受的诸多惨痛祸患。但我们祈祷，主所允许或愿意让我们承受的这些事，能有助于那些经历之人的改正，并因过犯的停止而使逆境终止。这两种结果都将因天主的仁慈而实现，只要祂除去鞭笞，并将祂子民的心转向祂。但是，兄弟，你为我们周围的敌对行动感到悲伤，同样，我也焦虑，因为正如你的来信所表明，异端者的诡诈攻击在\Place{君士坦丁堡}教会仍未平息，人们寻找机会迫害那些公教信仰的捍卫者。因为只要\Person{阿埃提乌斯}以升迁为借口被从总执事之职免去，而\Person{安德肋}——他曾因与异端者交往而被遗弃——取代了他的位置；只要神圣记忆的\Person{弗拉维亚诺}的控告者受到尊重，而那位最虔诚的宣信者的同伴或门徒被压制，这就太清楚不过地表明了该教会主教的心意。对于此人，我推迟采取行动，直到我听到案件的实情，并等待他在我儿\Person{阿埃提乌斯}告知我他将送来的信中亲自与我处理，给予自愿改正的机会，我愿借此平息我的恼怒。尽管如此，我已就这些关乎教会和平之事写信给我们最仁慈的君王和最虔诚的\OriginalTerm{奥古斯塔}{Augusta}；我不怀疑他们将因信德的虔诚而留意，以免一项已被谴责的异端成功复燃，损害他们自己的光荣事业。\par

% structure: p0006 source=h2 target=subsection reason=relative-heading-rank; base=h2

\subsection{二、他请尤利安代他行事，因阿纳托利乌斯缺乏活力。}

那么，亲爱的兄弟，请你将必要的思考倾注于宗座的忧虑，宗座以其作为你母亲的权柄将维护公教真理以对抗聂斯脱利派和欧提基派的事委托于你——你曾在其怀中受哺育——以便你受神圣助佑的支持，不停止关注\Place{君士坦丁堡}城的利益，免得错误的暴风雨随时在其中兴起。又因为我们光荣的君王的信德是如此之大，你可以自信地向他们提出必要之事，请为普世教会的益处使用他们的虔诚。但是，亲爱的，如果你就你认为是疑惑的事咨询我，我的答复必不缺少指导，以便除了应由各别教会主教的调查决定的事案外，你可以作为我的特使，并承担特殊责任，防止聂斯脱利派或欧提基派异端在任何地方复燃；因为\Place{君士坦丁堡}主教不具有公教的活力，对人类救恩的奥迹或对他自己的名声也不甚热心：然而，如果他有一点灵性活动，他本应考虑由谁祝圣了他，他继承了谁的位置，以致应追随有福的\Person{弗拉维亚诺}而非那些助他晋升的工具。因此，当我们最虔诚的君王肯按我的恳求就那些应受责备的事情训诫我们的兄弟\Person{阿纳托利乌斯}时，亲爱的，请将你的勤勉与他们的结合，以便所有过犯的起因能通过最彻底的改正而被移除，并且他停止伤害我儿\Person{阿埃提乌斯}。因为对于一位有公教情怀的主教，即使总执事身上有某种看似会激怒人的事，也应为尊重信德而放过，而不是让最无价值的异端者取代一位公教徒的位置。因此，当我得知其余的故事后，我将更清楚地知道该如何行事。因为在此期间，我认为最好克制我的恼怒并操练忍耐，以便为宽恕留有余地。\par

% structure: p0008 source=h2 target=subsection reason=relative-heading-rank; base=h2

\subsection{三、他请求关于\Place{巴勒斯坦}和\Place{埃及}骚乱的进一步信息。}

但是关于\Place{巴勒斯坦}的隐修士，据说他们长期处于叛乱状态，我不知道他们目前被什么精神所驱动。也还没有人向我解释他们似乎为自己的不满提出什么理由：例如，他们是希望用如此疯狂来服侍欧提基异端呢，还是他们不可调和地恼怒，因为他们的主教竟被误导进入那种亵渎中，尽管圣地的关联——从那里发出了给全世界的教导——他仍使他自己与主的降生成人的真理疏离，而在他们看来，在别人身上必须通过赦免洗去的罪，在他身上却是不可宽恕的。因此，我希望更充分地了解这些事，以便采取适当措施加以纠正；因为用邪恶武装自己反对信德是一回事，而为信德而过度不安是另一回事。你也必须知道，\Person{阿埃提乌斯}司铎先前告诉我说已经寄出的文件，以及你说已经寄出的信德概要，至今尚未到达。因此，如果有机会找到更快的信使，我将乐于尽快收到任何似乎合宜的信息。我很想知道\Place{埃及}的隐修士是否恢复了他们的平安和信德，以及关于\Place{亚历山大}教会，你得到了什么可靠的音讯：我希望你知道我写给其主教或其祝圣者或神职人员的信，因此我已送给你一封信的副本。你也能从我送去的副本中得知我对我们最仁慈的君王和最虔诚的皇后所说的话。\par

% structure: p0010 source=h2 target=subsection reason=relative-heading-rank; base=h2

\subsection{四、他请求将\Place{加采东}大公会议决议译为拉丁文。}

我愿意知道，我托执事\Person{巴西略}带给你的那封信，是否已送达你手，弟兄？那封信是关于主降生成人的信仰，当时圣善的\Person{弗拉维亚诺}还在世。因为我想你从未对其内容发表过任何意见。关于\Place{加采东}大公会议期间所制定的会议法案，由于语言差异，我们没有十分清楚的信息。因此，我特别嘱咐你，弟兄，请你将全部内容收集成册，当然要准确地翻译成拉丁文，以便我们对会议的任何一个部分都不再有疑问，并且在你费心使我完全理解之后，不会有任何不确定之处。日期：三月十一日，显赫的\Person{欧皮利奥}执政之年（四五三年）。\par

\SourceRecord{Translated by Charles Lett Feltoe.；Nicene and Post-Nicene Fathers, Second Series；Vol. 12.；Edited by Philip Schaff and Henry Wace.；Buffalo, NY: Christian Literature Publishing Co.,；1895.；<http://www.newadvent.org/fathers/3604113.htm>.}

% structure: p0001 source=h1 target=section reason=page-title

\section{\WorkTitle{第一一七封信}}

\Emph{致\Place{科斯}主教\Person{尤利安}。}\par

% structure: p0003 source=h2 target=subsection reason=relative-heading-rank; base=h2

\subsection{一、他愿众人知晓他对\Place{加采东}法案的赞同。}

\Person{良}致\Person{尤利安}主教。\par

弟兄，你的信函显示了你是何等警醒且热诚地守护公教信仰，而它所包含的消息使我极大的焦虑得以缓解；此外，我们虔敬皇帝的最虔敬孝爱之情也进一步证实了这一点，这显然是由主为整个教会的坚固而预备的；如此，当基督徒君王以神圣的热忱为信仰行事时，主的司铎可以满怀信心地为他们的国度祈祷。\par

因此，我最仁慈的皇帝认为必要的事，我欣然遵行，向所有出席\Place{加采东}主教会议的弟兄们发送了信函，在其中表明我赞同我们圣洁的弟兄们关于信仰准则所作的决议；这主要是为了那些以掩饰自身背信弃义为目的的人，他们假称未经我认可的主教会议法令是无效或可疑的；尽管在代表我出席会议的弟兄们返回后，我已向\Place{君士坦丁堡}主教致函；若他有意公布，那封信本可充分证明我多么欣然地赞同了主教会议关于信仰的决议。但是，由于该信函包含的答复与他的自利追求相悖，他宁愿不让众人知道我接受了弟兄们的决议，以免我的答复同时因\Place{尼西亚}法令的绝对权威而为人所知。因此，亲爱的，你要注意，通过屡次提醒来告诫我们最仁慈的君王，请他将其话语加于我们的话语之上，并命令将宗座书信发送给每个省份的司铎，以免日后真理的仇敌敢于以我的沉默为借口来自辩。\par

% structure: p0007 source=h2 target=subsection reason=relative-heading-rank; base=h2

\subsection{二、他对皇帝与皇后所表现的热忱表示感谢。}

至于最基督徒皇帝的法令，其中表明某些隐修士的无知愚行应得何种处罚，以及最仁慈的奥古斯塔的答复，她在此答复中斥责了隐修院的院长们，我希望我的巨大喜悦为人所知，因为我确信这信仰的热忱是来自神圣的灵感，为的是使所有人承认他们的卓越不仅在于其王权，也在于其司祭的圣德；我过去和现在都请求他们对你完全信任，深信他们的善意，并且他们不会拒绝倾听必要的建议。\par

% structure: p0009 source=h2 target=subsection reason=relative-heading-rank; base=h2

\subsection{三、他欲知他致\Person{尤多西亚}皇后信函的效果。}

再者，由于最仁慈的皇帝乐于通过我们的儿子\Person{保禄}秘密委托我劝诫我们的女儿、最仁慈的奥古斯塔\Person{尤多西亚}，我已照他所愿行事，以便她从我信中得知，若她拥护公教信仰的事业将对她多么有益；我还设法让她也通过其子——那位最仁慈的君主的信函受到劝诫；我毫不怀疑她本人也将以虔诚的热忱着手工作，使叛乱的首领们知晓其行为的后果，并且，若他们不明白教导他们之人的话语，至少使他们惧怕惩罚他们之人的权力。因此，亲爱的，我愿立即通过你的信函得知我们这一关切产生了何种效果，以及他们无知的顽固是否终于平息了；关于此事，若他们认为对我们的教导有任何疑问，至少他们不应拒绝接受\Person{阿塔纳修斯}、\Person{提奥菲鲁斯}和\Place{亚历山大}的\Person{济利禄}等如此圣洁的司铎的著作，我们的信仰声明与之完全一致，以致凡声称赞同他们的人，在各方面都与我们毫无分歧。\par

% structure: p0011 source=h2 target=subsection reason=relative-heading-rank; base=h2

\subsection{四、\Person{艾提乌斯}目前必须满足于皇帝的恩宠。}

对于我们的儿子、司铎\Person{艾提乌斯}的悲伤，我们深表同情；既然一位先前已被认为应受责难的人取代了他的位置，这无疑对公教徒造成损害。但与此同时，这些事情必须耐心忍受，以免我们被认为超出了惯常的节制尺度；目前，\Person{艾提乌斯}必须满足于我们最仁慈君王的鼓励恩宠，我不久前刚在信函中如此称赞他，以致我毫不怀疑他的美誉在他们最宗教的心灵中已得到提升。\par

% structure: p0013 source=h2 target=subsection reason=relative-heading-rank; base=h2

\subsection{五、\Person{阿纳托利乌斯}在其后的行为中未显示忏悔。}

我们也愿你知道，\Person{阿纳托利乌斯}主教在我们禁止之后，仍如此坚持其鲁莽的妄为，以致召集\Place{依利里库姆}的主教们签名；这消息是由\Place{德撒洛尼基}主教派来宣告其祝圣的主教带给我们的。我们拒绝就此致函\Person{阿纳托利乌斯}，尽管你可能期待我们这样做，因为我们察觉他并无悔改之意。我制作了我致主教会议信函的两个版本，一个附有我致\Person{阿纳托利乌斯}信函的副本，另一个没有；我将判断权留给你，你认为应提交给我们最仁慈君王的那一份，就交付，而保留另一份。写于3月21日，在显贵\Person{奥皮利奥}执政时期（453年）。\par

\SourceRecord{Translated by Charles Lett Feltoe.；Nicene and Post-Nicene Fathers, Second Series；Vol. 12.；Edited by Philip Schaff and Henry Wace.；Buffalo, NY: Christian Literature Publishing Co.,；1895.；<http://www.newadvent.org/fathers/3604117.htm>.}

% structure: p0001 source=h1 target=section reason=page-title

\section{\WorkTitle{第一一九书}}

\Emph{致安提约基雅主教马克西默书，由司铎玛里安及执事敖林比乌带交}\par

% structure: p0003 source=h2 target=subsection reason=relative-heading-rank; base=h2

\subsection{一、信德是欧提克斯与聂斯多略两个极端之间的中道}

良致安提约基雅的马克西默书。\par

亲爱的，你有多么渴望我们共同信德的最神圣合一及教会和平的宁静和谐，你的信函内容足以表明，该信由我们的儿子司铎玛里安及执事敖林比乌带交我们。我们格外欢迎此信，因为借此我们仿佛能进行交谈，于是天主恩宠越发显明，因公教真理的启示，整个世间洋溢着更大的喜乐。然而，我们仍为一些人深感忧痛（你的使者如此表示），他们依旧喜爱自己的黑暗；尽管白昼的光芒已普照各处，他们仍耽于盲目之暗，背弃信德，仅存虚名之为基督徒，无力辨别此错与彼错，也无法区分聂斯多略的亵渎与欧提克斯的不虔。因为他们的任何妄想都不能被视为可原谅，因他们在乖僻中自相矛盾。虽然欧提克斯的门徒憎恨聂斯多略，聂斯多略的追随者诅咒欧提克斯，但在公教信徒的判断中，双方均被定罪，两种异端一并从教会身体上切除：因为没有任何虚假能与我们一致。他们既在主的降生成人真理上以亵渎方向分歧，这一点并不重要，因为他们的错误见解既不持守福音的权威，也不符合奥迹的含意。\par

% structure: p0006 source=h2 target=subsection reason=relative-heading-rank; base=h2

\subsection{二、马克西默务要使东方各教会远离这两种相反异端}

因此，亲爱的弟兄，你当全心省思，主命你治理的是哪一座教会，并记起那套教义体系，这体系由众宗徒之长、真福伯多禄奠定根基，不仅通过他在普世一致的宣讲，尤其通过在安提约基雅及罗马两城的训诲：好使你明了，凡被置于他自己荣名之家的首领，他要求他持守那些从他接受自自己所宣认的真理本身的传授。在东方各教会，特别是在尼西亚太神圣教父之教规划归安提约基雅主教区者，你绝不可允许无赖异端攻击福音，也不可容任何人维护聂斯多略或欧提克斯的教条。因为，如我所说，公教信德的\OriginalTerm{磐石}{petra}，真福宗徒伯多禄由主手中得名之石，拒绝这两种异端的任何痕迹；因为它公开明确地诅咒聂斯多略，因他在真福贞女受孕时分离圣言与血肉之性，将独一基督分裂为二，愿意区分神性位格与人性位格：因为祂是全然同一的，那在永恒神性中从无时间生于圣父者，在真实血肉中随时间生于祂的母亲；同样，它也摒弃欧提克斯，因他忽略主耶稣基督内真实人性肉身，主张圣言本身转化为血肉，以致祂的诞生、抚养、成长、受苦、死亡、埋葬、第三日复活，全部仅属于祂的神性，祂所取的奴仆形象并非实体，而是幻影。\par

% structure: p0008 source=h2 target=subsection reason=relative-heading-rank; base=h2

\subsection{三、安提约基雅作为基督教世界第三席位应保留其特权}

因此，你当竭尽警醒，以防这些堕落异端胆敢自高自大；因为你当以司铎的全部权威抵抗他们，并常常通过你的报告告知我们教会进展之事。因为理当由你与宗座同担此责任，并明了基督教世界第三席位的特权赋予你行动的全部信心，这些特权任何阴谋都不得丝毫损害：因我对尼西亚教规的尊重如此之深，我从未允许，也永不允许圣教父的制度被任何革新所侵犯。因为个别主教的长处虽有不同，但其主教席的权利却是永久的；尽管竞争或偶尔引发它们的一些骚动，却不能损害其尊严。所以，弟兄，无论何时，你若认为应采取任何行动来维护安提约基雅教会的特权，务必以你自己的信件阐明，好使我们能充分且适当地回复你的请求。\par

% structure: p0010 source=h2 target=subsection reason=relative-heading-rank; base=h2

\subsection{四、阿纳多略颠覆尼西亚决议的企图是徒劳的}

但此时，在各方面作一般性的宣告就足够了：若在任何大公会议中，有人试图或似乎借机攻击尼西亚教规的规定，他无法玷污其不可侵犯的法令；任何阴谋的协议被破坏，也比上述教规的任何部分被废弃更容易。因为阴谋不会错过窃取利益的机会，每当事务的发展导致司铎们的全体集会时，无耻之徒的贪婪很难不试图谋取某种不正当利益：正如在推翻亵渎者聂斯多略及其教义的厄弗所大公会议中，儒文纳尔主教相信他能担任巴勒斯坦省的主席，并冒险用一封谎言的函件召集不服者。对此，亚历山大主教、真福西里禄感到适当地惊愕，他在致我的信中指出对方的贪婪导致了他何等狂妄：并怀着焦虑的恳求，竭力请求我不要同意他的非法企图。因为应让你知道，我们在我们的书柜中找到了所寻的西里禄信函的原始文件，你曾寄给我们抄本。然而，我特别强调这一点：虽然由于某些人的狡诈，多数司铎作出了与三百一十八位教父的规程相悖的决定，但按正义原则，该决定应被视为无效：因为整个教会的平安无法以其他方式保存，除非始终对教规表示应有的尊重。\par

% structure: p0012 source=h2 target=subsection reason=relative-heading-rank; base=h2

\subsection{五、如果良的使节以任何方式超越了他们的训令，其所作所为视为无效。}

当然，若那些我所差遣代替我出席神圣大公会议的弟兄们，在超出与信德相关之事上做了任何事，那都将毫无效力：因为他们受宗座差遣，只为铲除异端，维护大公信仰。因为提交给主教们审查的，超出大公会议所处理的特定主题之外的事项，若圣教父们在尼西亚对此未作任何规定，则可允许一定程度的自由讨论。但任何与他们的规则和规程不一致之事，永远无法得到宗座的认可。然而，应如何以极大的勤奋遵守这一规则，你将从我们致君士坦丁堡主教信函的抄本中得知，该信函旨在遏制他的贪欲；你应安排，使该函为所有我们的弟兄和同司铎们所知。\par

% structure: p0014 source=h2 target=subsection reason=relative-heading-rank; base=h2

\subsection{六、除司铎外，无人获准宣讲。}

亲爱的，这也需要你提防：除主的司铎外，无人敢声称拥有教导或宣讲的权利，无论是隐修士或平信徒，自诩有些学问者。因为尽管教会的所有子女应明白正确健全之事，但无人可在司铎品级之外担任宣讲的职务，因为在天主的教会中，一切应有秩序，使在基督的同一身体内，更高贵的肢体应履行其本分，而低微的不得反抗高贵的。\Signature{签发于六月十一日，在尊贵的奥皮利奥担任执政官之年（四五三年）。}\par

\SourceRecord{Translated by Charles Lett Feltoe.；Nicene and Post-Nicene Fathers, Second Series；Vol. 12.；Edited by Philip Schaff and Henry Wace.；Buffalo, NY: Christian Literature Publishing Co.,；1895.；<http://www.newadvent.org/fathers/3604119.htm>.}

% structure: p0001 source=h1 target=section reason=page-title

\section{第一二〇封信}

\Emph{致居鲁士主教德奥多雷：论在信德中的恒心。}\par

% structure: p0003 source=h2 target=subsection reason=relative-heading-rank; base=h2

\subsection{一、他祝贺德奥多雷他们共同的胜利，并表达对诚实探讨带来美好结果的赞同。}

主教良，致其爱弟德奥多雷主教。\par

蒙福的伯多禄宗座派遣至神圣大公会议的我们的兄弟和同侪司铎返回后，我们得知，亲爱的，你我与我们一起，藉着上天的助佑，战胜了聂斯多略的亵渎和欧提克斯的疯狂。因此我们以主夸耀，与先知同唱：\Quote{我们的救助是在上主的名里，祂创造了天地；}祂容许我们在我们弟兄身上毫发无伤，却藉着全体弟兄不可撤销的同意，证实了祂已藉着我们的职务所规定之事：以表明那首先由基督教世界的首席宗座制定，然后由整个基督教世界的判断所接受的事，确实出自祂自己：使在此事上，肢体也与元首合一。在此，我们喜乐的原因更大，因为我们看见，敌人越猛烈攻击基督的仆人，他反而越加折磨自己。因为，惟恐其他宗座对万有之主所指定居于首位者之同意，似乎仅是逢迎，或恐其他邪恶的猜疑潜入，有几个人在我们之决定被最终接受前，便出来争议。而当一些人，受不和之源的煽动，冲入矛盾的战争时，藉着一切美善之源的引导，因他的跌倒反而产生了更大的善。因为天主的恩宠之赐予，在藉着巨大努力获得时，对我们更为甜蜜；而不受干扰的和平，往往显得不如藉辛劳恢复的和平那样美善。而且，真理本身更加明亮地照耀，当信德早已教导的事随后被进一步的探讨所证实，也更为勇敢地持守。此外，司铎职务的美名更添光彩，当最高权柄得以保存，而低级阶层的自由不被认为受到任何侵犯时。而讨论的结果有助于天主更大的光荣，当辩论者满怀信心地努力克服反对者时：使得那本身显为错误的事，不至于在有害的沉默中被忽略。\par

% structure: p0006 source=h2 target=subsection reason=relative-heading-rank; base=h2

\subsection{二、基督的胜利赢回许多人归向信德。}

因此，亲爱的弟兄，欢跃吧，是的，在独生子天主内欢跃得胜。藉着我们，祂已为自己征服了那否认祂肉体真实性的人。藉着我们，为我们，祂已得胜；我们为祂的事业而得胜。这幸福的日子，对于世界而言，仅次于主的来临。那强盗被击倒了，神圣降生成人的奥迹被重新归还给我们的时代，这奥迹曾被人类的敌人用诡计蒙蔽，因为事实不容他真正毁灭它。不，那不朽的奥迹已从不信者心中消亡，因为如此伟大的救恩对不信者毫无益处，正如真理本身对祂的门徒所说：\BibleQuote{信而受洗的，必要得救；但不信的，必被判罪。}\BibleRef{谷 16:16} 公义之太阳的光线，曾在东方被聂斯多略和欧提克斯的乌云所遮蔽，却从西方明亮地照耀出来，在那里，它在宗徒和教会导师身上达到了顶点。然而，即使在东方，也不可相信它曾被遮蔽，因为你们中间曾出现高贵的宣信者：所以，当旧敌人再次藉着一个现代法郎的顽固之心，试图抹去信德的亚巴郎的子孙和恩许之子时，他因天主的仁慈而疲惫，除了他自己之外，不能伤害任何人。关于他，全能者还行了这奇事：祂并没有将那些与他一起屠杀以色列人民的同伙与暴政的创始人一同覆灭，而是将他们聚集到自己的子民中；并且，正如一切仁慈的源泉所知道的那样，值得祂自己且唯独祂能行，祂使他们成为与我们一同的征服者，我们曾是他们的征服者。因为谎言之灵是人类唯一的真正敌人，无疑，所有被真理争取到祂一边的人，都分享祂对那敌人的胜利。确实，现在清楚显示，我们的救赎主的话是何等神圣地授权，这些话如此适用于信德的敌人，以至于无人怀疑它们是对他们说的：\Quote{你们，}祂说，\Quote{是出于你们的父亲魔鬼，你们父亲的私欲，你们偏要行。他从起初就是凶手，不站在真理中，因为在他内没有真理。他说话时，出于自己说谎，因为他是撒谎者，也是谎言之父。}\BibleRef{若 8:44}\par

% structure: p0008 source=h2 target=subsection reason=relative-heading-rank; base=h2

\subsection{三、狄奥斯库若，在他的疯狂中甚至攻击了罗马主教，显示出自己是撒殚的工具。}

因此，那些对我们在真天主内的本性接受了谬误的人，竟在这些问题上也随从他们的父，坚持认为那在独生子内所看见、所听见、并确实藉福音的见证所触摸和感受的，并不属于那被证明所归属的实体，而属于一个与圣父同永恒、同本体的本质，好像天主性的本性能在十字架上被刺透，好像不变者能从婴孩长成成人，或永恒的智慧能在智慧上增进，或天主——祂是神体——以后能充满圣神。同样，他们的极端狂妄暴露了其根源，因为它尽其所能地试图伤害每一个人。因为那以迫害折磨你们的人，也误导了别人，驱使他们同意他的邪恶。甚至对我们也是如此，尽管他在每一位弟兄身上伤害了我们（因为他们是我们的肢体），但他并没有免除我们特别的烦恼，竟以奇怪、前所未闻且难以置信的厚颜无耻，试图伤害我们的元首。但愿他在所有这些暴行之后能恢复理智，而不因他的死和永罚使我们悲伤。他所行的恶没有限度：他既不饶恕生者也不饶恕死者，弃绝真理而与谎言结盟，将早已玷污的手浸入一个无罪的天主教司铎的血中，这对他来说还不够。因为经上记载：\BibleQuote{恨自己弟兄的，就是杀人者：}\BibleRef{若一 3:15}他确实实现了他在仇恨中已声称做过的事，好像他从未听过这话，也从未听过主所说的：\BibleQuote{你们跟我学习吧，因为我是良善心谦的，这样你们必会找到你们灵魂的安息；因为我的轭是柔和的，我的担子是轻省的。}\BibleRef{玛 11:29}这个埃及的掠夺者被证明是魔鬼谬误的合格宣讲者，他像教会所遭遇过的最残忍的暴君一样，通过暴民骚乱的暴力和士兵沾满鲜血的手，将他的邪恶亵渎强加给可敬的弟兄们。当我们的救赎主的声音向我们保证，杀人的和撒谎的出于同一个源头时，他两者都同样实现了：好像这些事记载下来不是为了避免，而是为了实行。因此，他将天主子的有益警告用于完成他的毁灭，并对同一位主所说的充耳不闻：\BibleQuote{我所说的，是我从圣父那里看见的；而你们所作的，却是从你们的父那里听见的。}\BibleRef{若 8:38}\par

% structure: p0010 source=h2 target=subsection reason=relative-heading-rank; base=h2

\subsection{四、那些受命以权威谈论教义的人，必须在极端之间保持平衡。}

因此，当他竭力切断可纪念的\Person{弗拉维安}在现世的生命时，他使自己失去了真生命的光明。当他试图将你们赶出你们的教堂时，他使自己与基督徒的共融隔绝。当他拖曳并驱使许多人陷入与错误的认同中时，他以许多创伤刺伤了自己的灵魂，他是一个在所有方面、透过所有方面、并为所有方面而独自被定罪的人，因为他是一切人受到控告的原因。但是，弟兄，你虽以坚实的食粮被养育，很少需要这样的提醒，然而为履行我们职位的责任，按照宗徒的话，他说：\BibleQuote{除了这些外面的事，还有那每日压在我身上的，对众教会的挂虑。有谁软弱，我不软弱呢？有谁跌倒，我不焦心如焚呢？}\BibleRef{格后 11:28}我们相信，这个劝告尤其应当在此刻给予，即每当藉着神圣恩宠的服事，我们在教义的池中要么淹没那些外人，要么洁净他们时，我们绝不可从加采东大公会议上由圣神的天主性所提出的信德准则中偏离任何一点，并以万分谨慎衡量我们的言语，以避免新虚假教义的两种极端：不再（但愿不是如此）以辩论不确定之事的方式，而是以完全的权威提出既定的结论；因为我们知道，在从宗座发出并由整个神圣公会议一致赞同的信函中，汇集了如此多神圣授权的见证，以致没有人能再存有任何疑虑，除非他宁愿将自己包裹在错误的云雾中，而且公会议的程序——无论是我们读到制定信德定义的程序，还是弟兄你热切支持上述宗座信函的程序，特别是整个公会议致我们最虔诚的皇帝们的致辞——都得到过去如此多教父的见证的支持，以至于任何心虽愚昧顽固的人，只要他尚未因邪恶而与魔鬼一同被定罪，都会被说服。\par

% structure: p0012 source=h2 target=subsection reason=relative-heading-rank; base=h2

\subsection{五、\Person{狄奥多雷}的正统性已得到幸福而彻底的辩护。}

为此，我们也有责任防范教会的敌人，尽我们所能，不给他们留下诽谤我们的机会；同时，在反对聂斯多略派或欧迪奇派时，切勿显得我们已向另一方退让，而是应同等回避和谴责这两个基督的敌人，以便每当听众的利益需要时，我们都能以十分敏捷和清晰的态度，用应得的绝罚打击他们及其教义，免得我们若显得犹豫或迟缓，便被以为并非出于本意行事。尽管事实本身足以提醒你的智慧注意此事，但如今实际经验已使这一教训深入人心。然而，愿我们的天主受赞颂，祂那不可战胜的真理已使你在宗座（Apostolic See）的裁决中，被证明与一切异端无关。你将为此一切辛劳向祂报答应有的感谢，只要你保持自己作为普世教会的捍卫者，正如我们已证实并仍在证实你的那样。因为天主驱散了所有诽谤性的谬误，我们将此归于蒙福的\Person{伯多禄}对我们众人的奇妙关怀：在批准其宗座关于界定信仰的裁决之后，他不允许任何不祥的指控留在你们这些曾与我们一同为公教信仰而辛劳的人身上；因为圣神判定，那些其信仰已得胜的人，不可能不战胜其敌人。\par

% structure: p0014 source=h2 target=subsection reason=relative-heading-rank; base=h2

\subsection{六、他请求\Person{德奥多勒}继续合作，并提及他写给\Place{安提约基雅}主教的信。}

剩下的是我们劝勉你继续与宗座合作，因为我们得知，在你们中间仍有一些欧迪奇和聂斯多略错误的残余。因为我们的主基督所赐予祂教会的胜利，虽增加了我们的信心，却并未完全消除我们的焦虑，我们蒙恩并非为沉睡，而是为更平静地工作。因此，我们期望借助你的警醒关怀在这方面得到协助，即你尽快通过定期报告，将主的道理在那些地区进展如何告知宗座；以便我们根据经验所提示的任何方式，协助该地区的司铎。\par

关于在多次提及的会议中，非法反对可敬的\WorkTitle{尼西亚教规}所提出的那些问题，我们已致信我们的弟兄和同为主教者、\Place{安提约基雅}主教座的现任者，并附加了你通过代表口头告知我们的关于某些隐修士不法行为的内容，同时严格规定：任何人，无论是隐修士还是平信徒，凡自诩略懂某方面知识的，都不得擅自宣讲，除非是主的司铎。然而，我们希望这封信通过我们前述的弟兄和同为主教者\Person{玛克西穆}，使全教会都知道以造福普世教会；因此，我们认为不便在此信后附上其副本，因为我们确信我们对我们前述弟兄和同为主教者的指令会得到妥善执行。（另一手迹）愿天主保佑你平安，亲爱的弟兄。颁于显赫的\Person{奥庇利奥}执政之年（453年）6月11日。\par

\SourceRecord{Translated by Charles Lett Feltoe.；Nicene and Post-Nicene Fathers, Second Series；Vol. 12.；Edited by Philip Schaff and Henry Wace.；Buffalo, NY: Christian Literature Publishing Co.,；1895.；<http://www.newadvent.org/fathers/3604120.htm>.}

% structure: p0001 source=h1 target=section reason=page-title

\section{第一二三封书信}

\newline{}\Emph{致欧多基娅·奥古斯塔，关于巴勒斯坦的隐修士们。}\par

% structure: p0003 source=h2 target=subsection reason=relative-heading-rank; base=h2

\subsection{一、请求她运用自己的影响力使巴勒斯坦的隐修士们恢复秩序。}

\newline{}主教良致欧多基娅·奥古斯塔。\par

\newline{}我并不怀疑你的虔敬深知我对公教信仰有多么大的热忱，以及我负有怎样的责任——在天主助佑下——谨防真理的福音在任何时候受到无知或不忠之人的抗拒。因此，在向你致以我应当表示的、你的仁厚始终乐意从我这里接受的敬意之后，我恳求主以你平安的消息使我喜乐，并藉着你日益增进对信仰那一条文的维护——巴勒斯坦省内的某些隐修士对此条文深感不安——好使你这热诚的虔敬能彻底摧毁对这类异端乖谬的一切信赖。因为对于那些既不为天主奥秘的原则所动，也不为圣经的权威，甚至不为圣地的见证本身所动的人，除了彻底的毁灭，还能担心什么呢？那么，愿教会因天主的恩宠而获益——正如它确实在获益——也愿人类本身——圣言在降生成人时接纳了它——因你已怀有定居在那地方的意愿而获益，那里他奇妙作为的证据和他受难的标记向你宣告我们的主耶稣基督不仅是真天主，也是真人。\par

% structure: p0006 source=h2 target=subsection reason=relative-heading-rank; base=h2

\subsection{二、应该告诉他们，公教信仰同时摒弃欧提克派和奈斯托利派的极端。他希望得知她取得了多大成功。}

\newline{}如果上述之人尊敬并热爱\Quote{公教}之名，并愿意被列在主身体的肢体中，就让他们摈弃他们在鲁莽中犯下的扭曲错误，并为他们邪恶的亵渎和流血行为表示忏悔。为了他们灵魂的救恩，让他们服从在卡尔采东城所确认的教务会议法令。因为只有真正的信仰和谦卑的驯顺才能领悟人类救恩的奥秘，让他们相信他们在福音中所读到的、他们在信经中所宣认的，并且不让自己卷入不健全的教义。因为正如公教信仰谴责了奈斯托利——他曾胆敢在我们的主耶稣基督内主张两位——同样它也谴责了欧提克和狄约斯科若，他们否认独生圣言在童贞母的胎中取了真实的人性。\par

\newline{}如果你的劝诫在说服这些人方面获得任何成功——这将为你赢得永恒的荣耀——我恳请你通过书信告知我；使我得以知道你已经收获了善工之果，而他们藉着主的仁慈没有丧亡。日期为六月十五日，在显赫的奥皮利奥执政之年（四五三年）。\par

\SourceRecord{Translated by Charles Lett Feltoe.；Nicene and Post-Nicene Fathers, Second Series；Vol. 12.；Edited by Philip Schaff and Henry Wace.；Buffalo, NY: Christian Literature Publishing Co.,；1895.；<http://www.newadvent.org/fathers/3604123.htm>.}

% structure: p0001 source=h1 target=section reason=page-title

\section{\WorkTitle{第一二四封书信}}

\Emph{致巴勒斯坦的隐修士们。}\par

% structure: p0003 source=h2 target=subsection reason=relative-heading-rank; base=h2

\subsection{一、他们可能因关于降生成人致弗拉维安的信函的误译而受到误导。}

良主教，致定居于整个巴勒斯坦的全体隐修士。\par

我所欠于整个教会及其所有子女的焦虑关怀，已从多方得知，亲爱的诸位，由于我的翻译者——他们看来或出于无知，或出于恶意——使你们对某些陈述的理解与我的原意不同，因为他们未能将拉丁语准确译成希腊语，尽管在解释微妙和困难的问题时，即便用母语进行讨论的人也几乎难以使自己满意。然而这一点对我也有利，因为你们拒绝接受大公信仰所摒弃的，我们就知道你们更喜爱真理而非谬误；并且你们完全正确地拒绝相信那些连我自己也按照古老教义所憎恶的东西。因为虽然我写给神圣记忆的弗拉维安主教的书信本身已足够明确，无需任何修正或解释，但我其他的著作也与那封书信一致，在其中我的立场也会以类似的方式阐明。因为我必须在最仁慈的君王和神圣的主教会议以及君士坦丁堡教会面前与那些使许多基督子民陷入混乱的异端者辩论，从而根据福音和宗徒的教导，阐明我们应当如何思考和感受关于圣言降生成人的事，并且我丝毫没有偏离圣教父的信经：因为信仰是唯一的、真实的、独一的、大公的，不可加添，也不可减少：虽然聂斯多略首先，现今欧提基又从相反的角度，以同样不忠的方式试图攻击它，并企图将两种矛盾的异端强加于天主的教会，这导致他们两人都理所当然地被真理的门徒们谴责；因为他们以不同方式共同持有的错误观点是极其疯狂和亵渎的。\par

% structure: p0006 source=h2 target=subsection reason=relative-heading-rank; base=h2

\subsection{二、混淆位格的欧提基与分离位格的聂斯多略同样应被摒弃。}

因此，必须绝罚聂斯多略，因为他相信童贞玛利亚只是人性之母，从而将祂肉身的位格与神性的位格视为不同的实体，而不承认在圣言和肉身中的同一位基督，却将天主子与人之子说成是分离和区别的：尽管祂与圣父和圣神从永远且不受时间限制地共有的不变本质并未失去，但\BibleQuote{圣言成了血肉}是那样地发生在童贞女的胎中，以至于由于那一次的受孕和一次的生育，同时由于两个实体的结合，她既是主的婢女，又是主的母亲。这同样也是依撒伯尔所知道的，正如路加福音所记载的，当她说：\BibleQuote{这事怎能临到我呢？我主的母亲竟到我这里来？}\BibleRef{路 1:43}但是欧提基也必须受到同样的绝罚，因为他陷入了古代异端者的诡诈错误，选择了阿波利纳里的第三种教条，以致他否认祂真实的人性血肉和灵魂，并主张我们的主耶稣基督整个是同一本性，仿佛圣言的神性变成了血肉和灵魂；并且仿佛受孕、出生、哺乳、成长、被钉十字架、死亡、被埋葬、复活、升天、坐在圣父的右边，并要从那里来审判活人死人——仿佛所有这些事都属于那唯一不接纳其中任何一项（没有真实血肉）的本质：因为独生子的本性就是圣父的本性、圣神的本性，而永恒圣三不可分割的合一与同体的平等既是无苦的，也是不变的。但如果这位异端者从阿波利纳里的谬见中退出，以免他被证明认为神性是可受苦和可死的，却仍敢宣称降生成人的圣言——即成了血肉的圣言——的本性是单一的，那么他无疑就转入了摩尼和马西昂的疯狂观点，相信那个人——即那在人与天主之间的中保耶稣基督——所做的一切都是虚幻的，并没有一个真实的身体，而是一个幻影般的显现呈现在观看者眼前。\par

% structure: p0008 source=h2 target=subsection reason=relative-heading-rank; base=h2

\subsection{三、承认基督内的我们的本性是正统信仰的必要条件。}

因为这些不义的谎言曾被公教信仰所拒绝，并且这些人的亵渎为天下荣福教父的一致意见所谴责，无论这些人是谁，他们如此眼瞎，远离真理之光，以至于否认自降生之日起，圣言内存在人性——即我们的人性；他们必须表明，他们凭什么声称基督徒的名号，又以何种方式与真正的福音和谐一致，如果荣福童贞的生育要么产生了没有神性的肉身，要么产生了没有肉身的性体。因为正如不能否认\BibleQuote{圣言成了血肉，寄居在我们中间\BibleRef{若1:14},}同样不能否认\BibleQuote{天主在基督内，使世界与自己和好\BibleRef{格后5:19}.}但是，有什么和好，能藉此使天主为人性而和解，除非天主与人之间的中保担当了众人的案件？祂如何能合适地履行其中保职分，除非那以天主的形体与圣父同等者，也以奴隶的形体分享我们的本性：这样，那一位新人能更新旧人，并且因一个人的过犯而加于我们身上的死亡锁链，能因唯一那位对死亡一无所欠的人的死而松开。因为义人为不义者流血，其效果如此强大，赎价如此丰盈，以至如果我们全体囚徒只相信他们的救赎者，就无一人会被暴君的束缚所困；因为正如宗徒所说\BibleQuote{罪恶在哪里增多，恩宠在那里也更形增多\BibleRef{罗5:20}.}而且，我们生而负罪，却蒙受了新生于义的能力，自由的恩赐已变得比奴隶的债更强大。\par

% structure: p0010 source=h2 target=subsection reason=relative-heading-rank; base=h2

\subsection{四、只有真正分享基督死亡与复活的人，才因基督的血受益。}

那么，那些否认我们救主身体中人性位格的实在性的人，在这奥迹的效能中为自己留下什么希望呢？让他们说出，他们藉何种祭献而被和好，藉何种流血而被挽回。谁是\BibleQuote{祂为我们舍己，作为奉献和牺牲，成为天主馨香的祭品\BibleRef{弗5:2}:}或者，有什么祭献比真正的大司祭藉贞献自己的肉身放在十字架的祭台上的祭献更为神圣？因为虽然在上主眼中，许多圣人的死是宝贵的，但没有一个无辜者的死曾是世界的赎罪祭。义人接受了冠冕，而非赐予冠冕；从信友的坚忍中生出了忍耐的榜样，而非成义的恩赐。因为他们的死只影响他们自己，没有人能将自己的死偿还另一个人的债：在人类中只有一位，我们的主耶稣基督，卓然屹立，众人在祂内同被钉死，同被埋葬，同被复活。关于他们，祂亲自说\Quote{当我从地上被举起来时，我要吸引万有到我这里来。}真实信德，它使犯法者成义并成为义人，也被吸引到那位分享了他们人性者那里，并在祂内获得救恩；在祂内，人唯独发现自己不再有罪；因此可以自由地因祂的能力而夸耀，祂在我们血肉的卑微中与高傲的仇敌交战，并与那些在祂身体上得胜的人分享祂的胜利。\par

% structure: p0012 source=h2 target=subsection reason=relative-heading-rank; base=h2

\subsection{五、基督二性的行动必须保持区分。}

因此，虽然在我们唯一的主耶稣基督内，真天主亦真人，圣言与肉身的位格是同一的，且两种存在有其共同的行动：然而我们必须理解这些行动本身的特性，并通过真诚信德的默观，区分那些是由于祂弱点的谦卑而显出的行动，与那些是由于祂崇高权能而倾向的行动：即肉身没有圣言或圣言没有肉身所不能做的事。例如，没有圣言的权能，童贞女就不能怀孕和生子；没有肉身的实在，祂的婴儿期就不能被包裹在襁褓中。没有圣言的权能，贤士们就不会朝拜那颗新星所指出的婴孩；没有肉身的实在，那婴孩就不会被命令逃往埃及，以躲避黑落德的迫害。没有圣言的权能，父从天上发出的声音就不会说\Quote{这是我的爱子，我所喜悦的}；没有肉身的实在，若翰就不能指着祂说\Quote{请看天主的羔羊，请看除去世罪的那位。}没有圣言的权能，就不会有病愈康复，没有死人复活；没有肉身的实在，祂就不会饥饿需要食物，也不会疲倦需要休息。最后，没有圣言的权能，主就不会宣称自己与圣父同等；没有肉身的实在，祂也不会说父比祂大：因为公教信仰维护并捍卫这两种立场，相信天主的独生子既是人又是圣言，依据祂的神性与人性实体的独特属性。\par

% structure: p0014 source=h2 target=subsection reason=relative-heading-rank; base=h2

\subsection{六、在基督内，二性并无混淆。}

因此，虽然从起初，即在童贞玛利亚的胎中\Quote{圣言成了血肉}之时，天主性体与人性性体二者之间从未发生过任何分离，并且通过祂身体所有的成长和变化，始终是同一位的行动，然而我们绝不因将之混淆而搞混那些不可分离地完成的行为本身；而从这些行为的特征中，我们看出哪一性体所属。因为祂的天主性的行为并不影响祂的人性行为，祂的人性行为也不影响祂的天主性的行为，因为二者以这种方式且为了这个目的共同作用，以致在它们的运作中，祂的双重品质不被彼此吸收，祂的独一性也不被加倍。因此，让那些基督徒的幻象贩子告诉我们，是什么样的救主本性被钉在十字架的木头，躺在坟墓里，并在第三天当石头从墓口滚开时从肉身中复活；或者，耶稣向门徒们显现时，在门关闭的情况下进入他们的眼界，是什么样的一具身体：因为为了驱散观看者的不信，祂要求他们用眼睛细看，用手触摸仍然敞开的钉痕和祂被刺透的肋旁的肉伤。但如果尽管真理如此清晰，他们仍坚持异端，不肯离开黑暗中的立场，那么让他们表明自己从何处承诺永生的希望，而这希望，除了通过天主与人之间的中保——那人\Person{耶稣基督}，无人能够获得。因为\BibleQuote{除祂以外，在天下没有赐下别的名字，使我们必须赖以得救的}\BibleRef{宗 4:12}。也没有任何从囚禁中赎买人类的，除非藉着祂的血，\BibleQuote{祂为众人舍了自己作赎价}\BibleRef{弟前 2:6}：这正如蒙福的宗徒所宣告的：\BibleQuote{祂虽具有天主的形体，并没有以自己与天主同等为应当把持不舍的，却使自己空虚，取了奴仆的形体，与人相似，形状也一见如人；祂贬抑自己，听命至死，且死在十字架上。为此，天主极其举扬祂，赐给祂一个名字，超越其他所有的名字：致使上天、地上和地下的一切，一听到耶稣的名字，无不屈膝叩拜，一切唇舌无不明认主耶稣基督是在天主圣父的光荣中}\BibleRef{斐 2:6}。\par

% structure: p0016 source=h2 target=subsection reason=relative-heading-rank; base=h2

\subsection{VII. 祂之所以被举扬，因为是取了\Quote{奴仆的形体}，而非因为祂是圣子。}

因此，虽然主\Person{耶稣基督}是独一的，在祂内真实的天主性和真实的人性构成了绝对同一的位格，并且这结合的整体不能被任何分割所分离，然而其中所说的举扬，即\Quote{天主举扬祂}并\Quote{赐给祂一个名号超越一切名号}，我们理解属于那个需要因这光荣的增加而得到充实的形体。当然，在\OriginalTerm{天主的形体}{the form of God}中，圣子是与圣父平等的，在本质方面，圣父与独生子之间毫无区别，在尊威方面也毫无差异；而且，通过降生成人的奥迹，圣言也没有被剥夺任何需要圣父赐予恢复的东西。但是\OriginalTerm{奴仆的形体}{the form of a slave}（藉此，那不可受痛苦的天主性完成了一项伟大仁爱的保证）是人的软弱，这软弱被举扬到天主大能的光荣中，因为从童贞玛利亚受孕之初，天主性与人性就如此完全地结合，以至于没有人性，天主性的行为无法进行，没有天主性，人性的行为也无法进行。因此，正如光荣的主被说成被钉在十字架上，同样，从永恒中与天主平等的那一位被说成被举扬。基督由哪一性体被提及，并不重要，因为祂位格的合一不可分割地保留着，祂同时根据血肉是完全的人子，根据与圣父合一的天主性是完全的天主子。因此，基督在时间内所接受的一切，都是凭借祂的人性接受的，凡是祂以前所没有的，都加给了那人性。因为根据圣言的权能，\Quote{圣父所有的一切}，圣子也毫无差别地拥有，而且祂在\Quote{奴仆的形体}中从圣父所接受的，祂自己在圣父的形体中也给予了。祂在自身内同时既富有又贫穷；富有，因为\BibleQuote{在起初已有圣言，圣言与天主同在，圣言就是天主。这圣言在起初与天主同在。万物是藉着祂造的；凡受造的，没有一样不是藉着祂造的}；贫穷，因为\BibleQuote{圣言成了血肉，居住在我们中间}。但那空虚自己或贫穷，除了藉着那奴仆的形体（天主的尊威被隐藏于其中，人类救赎的计划得以实现）的接受，又是什么呢？因为我们被囚的原始枷锁无法解开，除非我们种族、我们本性的一个人出现，他不受旧债的先入之见所束缚，并以其无玷的宝血涂去死亡的桎梏，正如从起初就由天主预定的那样。因此，在所定的时期一满，那以多种方式宣告的应许得以达成其长久期待的应验，从而，那由一个又一个见证频繁宣布的，可能消除一切疑惑。\par

% structure: p0018 source=h2 target=subsection reason=relative-heading-rank; base=h2

\subsection{VIII. 反对他们在此事上的无信和不一致。}

因此，既然所有这些异端都已被摧毁，它们因着主持公会议的神圣教父们的虔诚而被从大公合一的肢体上剪除，并且它们理应被放逐于基督之外，因为它们将圣言的降生成人——这正是那些正确相信者的唯一救恩——变成了自己的绊脚石和跌倒的缘由，我惊讶于你们，亲爱的，在辨明真理之光上竟有任何困难。既然众多阐述已清楚表明，基督信仰在谴责聂斯多略和欧迪奇连同狄奥斯库若上是正确的，并且一个人若赞同其中任一方的亵渎意见，便不能被称作基督徒，我痛心的是，正如我所听闻的，你们竟藐视福音和宗徒的教导，煽动民众叛乱，扰乱教会，并对司铎和主教施加侮辱甚至死亡，以致你们因愤怒和凶残而迷失了你们的决心和宣认。你们温良和安静的规范在哪里？忍耐的宽宏在哪里？和平的宁静在哪里？爱与坚忍的坚实基础在哪里？什么邪恶的说服将你们夺去，什么迫害使你们与基督的福音分离？或者那欺骗者何等奇怪的诡计显现了出来，以至于你们忘记了先知和宗徒，忘记了你们在领受圣洗圣事时在许多见证人面前宣认的令人得救的信经和宣认，竟要投身于魔鬼的欺骗之中？如果异端者空洞的评论竟有如此分量，能如此猛烈地夺取你们信仰的纯洁，那么\Quote{爪刑}和其他残酷的折磨会对你们有什么影响呢？你们以为是在为信仰行事，却是在反对信仰。你们以教会的名义武装自己，却在攻击教会。这就是你们从先知、传福音者和宗徒所学到的吗？否认基督真正的肉体，使圣言的本体遭受痛苦和死亡，使我们的本性不同于那修复它者的本性，并且认为所有十字架高举的、长枪刺穿的、墓石接收又归还的，都只是神圣大能的作为，而非全然也是人类谦卑的作为？正是针对这种谦卑，宗徒说：\BibleQuote{因为我不以福音为耻}\BibleRef{罗 1:16}，因为他知道基督徒的敌人对他们投以何等的羞辱。因此，主也宣告说：\BibleQuote{凡在人前承认我的，我在我父前也必承认他}\BibleRef{玛 10:32}。因为那些在基督的肉体中唤不起尊敬的人，将不配获得圣子和圣父的承认；而那些羞于用嘴唇承认他们同意佩戴在额上之标记的人，将证明自己从十字架的标记中一无所获。\par

% structure: p0020 source=h2 target=subsection reason=relative-heading-rank; base=h2

\subsection{九、劝勉接受对降生成人的大公观点}

放弃吧，我的儿子们，放弃这些来自魔鬼的提示。天主的真理是任何事物都不能损害的，但真理若不藉着我们的血肉就不能拯救我们。因为，正如先知所说：\Quote{真理从地上生出}，童贞玛利亚以这样的方式怀了圣言，她以自己实体的血肉服侍了祂，使之与祂结合，既没有增加第二个人，也没有消失她的本性：因为那具有天主形体者，取了奴仆的形体，以这样的方式，基督在两种形体中是完全同一的：天主俯就于人的软弱，而人提升至神性的崇高，正如宗徒所说：\BibleQuote{圣祖也是他们的，并且基督按血肉说，也是从他们来的，祂是在万有之上，永远受赞美的天主。阿们。}\BibleRef{罗 9:5}\par

\SourceRecord{Translated by Charles Lett Feltoe.；Nicene and Post-Nicene Fathers, Second Series；Vol. 12.；Edited by Philip Schaff and Henry Wace.；Buffalo, NY: Christian Literature Publishing Co.,；1895.；<http://www.newadvent.org/fathers/3604124.htm>.}

% structure: p0001 source=h1 target=section reason=page-title

\section{第一二九函}

\Emph{致\Place{亚历山大里亚}主教\Person{普罗泰里乌斯}}\par

% structure: p0003 source=h2 target=subsection reason=relative-heading-rank; base=h2

\subsection{一、他赞扬他对着信德的持久忠诚。}

\Person{莱奥}致\Person{普罗泰里乌斯}，\Place{亚历山大里亚}主教。\par

你的来信，亲爱的弟兄，由我们的弟兄及同为主教的\Person{聶斯托利}及时带给我们，使我大为喜悦。因为由\Place{亚历山大里亚}教会首领致宗座的这种信函是合宜的，它显示出埃及人从一开始就从最蒙福的宗徒\Person{伯多禄}的教导，经由他蒙福的门徒\Person{马尔谷}，学到了罗马人所一致相信的：除了主耶稣基督以外，\BibleQuote{在天下人间，没有赐下别的名字，使我们赖以得救的}\BibleRef{宗 4:12}。但因\BibleQuote{并非所有的人都拥有信德}\BibleRef{得后 3:2}，且狡猾的诱惑者最喜伤害人心，莫过于用违背福音真理的错谬毒害他们不设防的心智，我们必须藉圣神大能的教导努力阻止基督信仰知识被魔鬼的谎言所歪曲。为防备这一危险，教会的主管者尤其应当警惕，并从纯朴者的思想中驱除那些被某种真理外表所粉饰的谎言。\BibleQuote{因为通往生命的门是窄的，路是狭的。}\BibleRef{玛 7:14}。他们并不主要是通过观察人的行为来诱捕人，而是通过语义上的细微区别，以极其微小的增删或改动来败坏句子的效力，因此有时一句原本有助于得救的陈述，因着微妙的变化而转为毁灭。但既然宗徒说：\BibleQuote{该当有异端，好叫那些经得起考验的人显明出来，}\BibleRef{格前 11:19}，这有助于整个教会的进步：每当邪恶在提出错误见解时显露自身，有害的事物就不再被隐藏，而最终必然导致毁灭的事物也就不会伤害他人的纯真。因此，那些鲁莽固执、宁愿以自己的耻辱为荣而不接受所提供的补救的人，必须将他们盲目的游荡和跌倒归咎于自己。弟兄，你对他们固执的不悦是正确的，我们赞扬你坚持那从蒙福的宗徒和圣教父传承给我们的教导。\par

% structure: p0006 source=h2 target=subsection reason=relative-heading-rank; base=h2

\subsection{二、让他通过公开朗读有关问题的教父引言和\WorkTitle{大卷}来坚固信友们。}

因为我写给可敬的\Person{弗拉维安}答复他关于我主耶稣基督降生成人的询问的书信中，并无新的宣讲；我丝毫未偏离你们的祖先和我们的祖先所公开持守的信德准则。如果\Person{狄奥斯库鲁}愿意追随和效法他们，他本会留在基督的身体内，拥有蒙福记忆的\Person{亚大纳削}的著作作为教导的材料，并拥有圣德回忆的\Person{德奥斐禄}和\Person{济利禄}的讲道作为可赞美地反对已被谴责的教义的手段，而不是选择与\Person{欧提克斯}同流合污于他的亵渎。因此，亲爱的弟兄，出于对我们共同信德的焦虑，我建议：因为基督十字架的仇敌窥伺我们所有的字句，我们不要给他们丝毫机会来错误地断言我们赞同聶斯托利派的教义。你必须如此勤勉地劝勉平信徒、圣职人员和所有弟兄会成员在信德上长进，以表明你并未教导任何新东西，而是将可敬的教父们以一致声明所教导、且我们的书信在一切方面都与之相符的那些道理，灌输到所有人的心中。这不仅要通过你的言语，也要通过实际地朗读先前的陈述来表明，好让天主的子民知道，教父们从前辈接受并传给后代的东西，在今日仍然灌输给他们。为此目的，在先读了上述司铎们的声明之后，最后也让我的著作被宣读，使信友们的耳朵可以见证，我们所宣讲的没有别的，正是我们从祖先那里接受的。因为他们的理解力在这些事上少有操练，至少让他们从教父的书信中学习，这邪恶是何等古老，如今我们在\Person{聶斯托利}和\Person{欧提克斯}身上都予以谴责，他们二人均以按照主的教导宣讲基督的福音为耻。\par

% structure: p0008 source=h2 target=subsection reason=relative-heading-rank; base=h2

\subsection{三、古代的成规应始终维持。}

因此，无论是在信德准则还是在纪律遵守上，都当始终维持古代的规范，亲爱的，要显示一位审慎统治者的坚定，使\Place{亚历山大里亚}教会从我对其古老特权的坚决维护、对抗某些人的无原则野心，以及我从所有都主教应保持其尊严不受减损的决心中获得益处。你可从我致圣教会议、至虔敬的皇帝或\Place{君士坦丁堡}主教的书信语气中看出；因为你会察知我特别留意不使主的教会的信德准则有任何偏离，也不使任何人的肆无忌惮导致其特权有任何减损。既然这样，弟兄，你要坚持你前任的惯例，对你们同省的主教保持应有的权威，按古代体制，他们是隶属于\Place{亚历山大里亚}宗座的；使他们不反抗教会惯例，也不拒绝在你主持下聚会，无论是在固定时间或任何合理原因需要时；若有任何在全体会议中讨论的事将有益于教会，弟兄们既已如此聚集，便可一致作出相应决议。因为没有什么应使他们脱离这种服从，我们既对你的信德和品行有如此好的认识，弟兄，我们将不允许你失去任何前任的权威，也不允许你受轻慢而不受惩罚。公元454年3月10日，在显赫的\Person{埃提乌斯}与\Person{斯图迪乌斯}执政时期。\par

\SourceRecord{Translated by Charles Lett Feltoe.；Nicene and Post-Nicene Fathers, Second Series；Vol. 12.；Edited by Philip Schaff and Henry Wace.；Buffalo, NY: Christian Literature Publishing Co.,；1895.；<http://www.newadvent.org/fathers/3604129.htm>.}

% structure: p0001 source=h1 target=section reason=page-title

\section{第一三九封书信}

\Emph{致耶路撒冷主教\Person{朱文纳尔}}\par

% structure: p0003 source=h2 target=subsection reason=relative-heading-rank; base=h2

\subsection{一、他因\Person{朱文纳尔}回归正统而喜乐，虽也责备其曾误入歧途。}

\Signature{城市罗马的主教\Person{良}，致耶路撒冷主教\Person{朱文纳尔}。}\par

当我收到你的来信，亲爱的弟兄，即我们的儿子——司铎\Person{安德肋}和执事\Person{伯多禄}——带给我的，我确实喜乐，因为你得以回归你的主教座；但当我忆起所有使你陷入如此巨大困苦的原因时，我悲伤地想到，你自己就是你逆境的根源，因为你未能坚持抵挡异端者：因为人们只能认为你没有足够的勇气去驳斥那些你曾与之一同犯错并自称满意的人。因为将可敬的\Person{弗拉维安}定罪，并接纳最邪恶的\Person{欧提克斯}，若不是否认我们的主耶稣基督在肉身内，又是什么呢？这乃是祂亲身以极大的仁慈使之倾覆，当祂藉着神圣的\Place{加采东}大公会议的权威，废止了\Place{厄弗所}会议的诅咒判决，而未阻止任何被定罪者通过改正而得到医治。因此，因为在长期忍耐中，你选择了回归智慧而非执迷不悟，我欢喜你如此寻求天上的良药，以至于最终成为了那受异端者攻击的信德的护卫者。因为，虽然任何司铎都不应不知晓他所宣讲的，但任何居住在\Place{耶路撒冷}的基督徒比所有无知者更不可原谅，因为他不仅通过文字，而且通过地方本身的见证，被教导理解福音的权能；而在别处或许不可不信的事，在那里不能不被看见。当眼睛成为人的导师时，理解为何陷入困境？当人类救恩的一切奥迹都呈现在视觉和触觉面前时，阅读或听闻的事物为何仍存疑？仿佛主仍然用祂的人性声音对每个疑惑的人说：\Quote{你们为什么惊慌？为什么心里起疑虑？看看我的手和我的脚，就知道是我本人了。摸摸我，看看，因为（或：因）鬼魂没有骨肉，如同你们看见我有一样。}\BibleRef{路 24:38}\par

% structure: p0006 source=h2 target=subsection reason=relative-heading-rank; base=h2

\subsection{二、愿他因所居之处的神圣联系而坚定信心。}

因此，亲爱的弟兄，你要运用这些公教信仰无可争议的证据，并用你所居住的圣地的见证来支持福音宣讲者的宣讲。在你的地区有\Place{白冷}，在那里救恩之光从\Person{达味}家族的童贞女的胎中升起，被襁褓包裹，被拥挤客栈的马槽接纳。在你的地区，救主的婴孩时期由天使宣告，由贤士朝拜，被\Person{黑落德}通过许多婴孩的死亡搜寻。在你的地区，祂的童年成长，青年成熟，祂真实的人性通过身体的增长达到完满的成人，不是没有饥饿的食物，不是没有休息的睡眠，不是没有怜悯的眼泪，不是没有畏惧和战栗：因为祂是同一位，在\OriginalTerm{天主性}{divinitas}中施行大能奇迹，在奴仆形象中遭受苦难的残酷。这十字架本身不停地向你述说；这墓穴的石头呼喊，主以人的状况躺在下面，并藉\OriginalTerm{天主性}{divinitas}力量从其中复活。当你接近\Place{橄榄山}，朝拜升天的地方时，难道天使的声音不在你耳边回响吗？他对那些在主升天时目瞪口呆的人说：\Quote{加里肋亚人啊，你们为什么站着望天呢？这位被接升天的耶稣，你们看见祂怎样升天，也要怎样降来。}\BibleRef{宗 1:11}\par

% structure: p0008 source=h2 target=subsection reason=relative-heading-rank; base=h2

\subsection{三、福音事实证实降生成人。}

因此，基督的真实诞生由真实的十字架所证实；因为祂是在我们的肉体内出生，祂也是在我们的肉体内被钉十字架。这肉体，既然没有罪进入其中，本不会死亡，除非它是我们族类的肉体。但为使祂恢复所有人的生命，祂承担了所有人的事业，并藉为所有人偿付，废除了旧债的效力，因为惟独祂在我们所有人中并不欠债：以便正如因一人的罪过所有人都成为罪人，同样因一人的义所有人都成为义人，公义由承担人性者赐予人。因为祂绝不处于我们真实肉体本性之外，福音作者在故事开头说：\Quote{耶稣基督的世代书：\Person{达味}之子，\Person{亚巴郎}之子}\BibleRef{玛 1:1}，蒙福的宗徒\Person{保禄}的教导与此一致，他说：\Quote{他们拥有祖先，基督按肉身也出自他们，祂是在万有之上永远受赞美的天主}\BibleRef{罗 9:5}，以及致\Person{弟茂德}书：\Quote{你要记住：耶稣基督从死者中复活了，出自\Person{达味}的后裔}\BibleRef{弟后 2:8}。\par

% structure: p0010 source=h2 target=subsection reason=relative-heading-rank; base=h2

\subsection{四、仍在错误中者当受到历史信德的彻底教导。}

但是，无论是新约还是旧约，有多少权威宣告了这一真理，正如适合您宗座古老那样，您清楚明白，因为教父们的信仰和写给可敬的\Person{弗拉维安}的信（您自己提及过），这封信已经由普世会议确认，对您已经足够。因此，亲爱的弟兄，你应当留心，不让任何人反对我们救赎与希望不可言喻的奥迹。但如果还有人仍在无知的黑暗或悖逆的纷争中，让他们受到那些在教会中其教导清晰并确定的人的权威教导，使他们认识到我们对于天主圣言降生成人的信仰与他们相同，并使他们不因固执而将自己置于基督身体之外，我们在其中与祂同死同生：因为无论是信德的忠诚还是奥迹的计划，都不容许\OriginalTerm{天主性}{divinitas}在其本质内可受苦难，也不容许在我们肉身的取纳中现实被篡改。日期：9月4日，在显赫的\Person{埃提乌斯}和\Person{斯图狄乌斯}执政之年（454年）。\par

\SourceRecord{Translated by Charles Lett Feltoe.；Nicene and Post-Nicene Fathers, Second Series；Vol. 12.；Edited by Philip Schaff and Henry Wace.；Buffalo, NY: Christian Literature Publishing Co.,；1895.；<http://www.newadvent.org/fathers/3604139.htm>.}

% structure: p0001 source=h1 target=section reason=page-title

\section{\WorkTitle{第一五六封书信}}

\Person{主教良}，致\Person{良·奥古斯都}。\par

% structure: p0003 source=h2 target=subsection reason=relative-heading-rank; base=h2

\subsection{一、现在没有必要重新开启教义的问题。}

我恭敬地接受了您充满刚毅信德和真理之光的谕旨，我希望能服从它，即使是在您认为所需的亲临之事上；因为那样我将从瞻仰您的辉煌中获得更大的益处。但我相信，当理由显示其更可取时，您会赞同我的看法。因为既然您以神圣和属灵的热忱始终维护教会的和平，且没有什么比坚持在圣神不断的引导下不可争辩地界定的那些事物更有利于维护信仰，那么，若我们似乎尽力要推翻这些法令，并应异端者的请求推翻普世教会所采纳的权威，从而消除教会冲突的一切界限，并完全放任叛乱，扩大而非平息争论，那将是极其不当的。因此，在厄弗所主教会议的可耻场景之后——在那里，因狄奥斯库若的邪恶，公教信仰被拒绝，而欧提奇斯的异端被接受——，为维护基督徒信仰，没有什么比神圣的加采东主教会议撤销他的邪恶行为，并如此关注天上教义，使任何人的思想中不再存有与先知或宗徒的言论相左之处，更为有益；当然，其中保持适度，只将顽固的叛乱者逐出教会合一，而不拒绝悔改者的宽恕。那么，还有什么比您陛下所颁布的、更符合人们期望或宗教精神的法令，即从今以后，任何人不得再攻击那些由神圣而非人的法令所决定之事，以免那些敢于怀疑天主真理的人，真正配得失去天主的恩赐呢？\par

% structure: p0005 source=h2 target=subsection reason=relative-heading-rank; base=h2

\subsection{二、重新考虑此问题的提议出自假基督或魔鬼本身。}

因此，既然普世教会通过那块原初磐石的建立而成为一块\OriginalTerm{磐石}{petra}，而首位宗徒、最蒙福的伯多禄听到了主的声音说：\Quote{你是伯多禄，在这\OriginalTerm{磐石}{petra}上，我要建立我的教会\BibleRef{玛 16:18}。} 那么，谁敢攻击如此不可动摇的力量呢？除非他是假基督或魔鬼，他执迷不悟，活在邪恶之中，焦急地借着适合他诡计的义怒之器散布谎言，同时以勤奋的假名假装寻求真理。他的无节制疯狂和盲目邪恶理所当然地给自己带来了轻蔑和恶名，以致当他以魔鬼般的意图攻击亚历山大的圣教会时，人们可以认识那些想要重新考虑加采东主教会议之人的品性。因为绝不可能在那里表达出与圣尼西亚主教会议相左的意见，正如异端者错误地坚持的那样，他们假装坚持尼西亚会议的信德。在那次会议上，我们的圣善可敬教父们聚集反对亚略，他们宣告的不是主的肉躯，而是圣子的神性与圣父\OriginalTerm{同质}{homoousion}；而在加采东主教会议上，为反对欧提奇斯的亵渎，它被界定为：主耶稣基督从童贞母亲的实体中接纳了我们身体的实际。\par

% structure: p0007 source=h2 target=subsection reason=relative-heading-rank; base=h2

\subsection{三、全体基督教世界的主教都在这点上与他一致。}

因此，在致我们至为虔诚的皇帝——他配被列在基督的护卫者之中——时，我运用公教信仰的自由，毫无畏惧地劝诫您与宗徒和先知为伍；坚决藐视和拒绝那些自绝于基督徒之名的人，并且不让那些亵渎的弑亲者（他们公认想取消信仰）在背信弃义的借口下讨论那信仰。因为主既然赐予您如此洞察祂奥秘的智慧，您应当毫不犹豫地认为，王权不仅是为治理世界而赐予您，更是为保护教会；即通过平息邪恶的企图，既能捍卫已被正确决定的事物，又能在有骚乱之处恢复真正的和平，也就是说，通过罢免篡夺他人权利者，并在亚历山大教座上恢复古代信仰，使天主的义怒藉您的改革得以平息，从而祂不再因他们的行为惩罚一个迄今虔诚的民族，而是赦免他们。可敬的皇帝，您要将此事置于您心眼前：普世所有主的司铎都在为那信仰恳求您，在那信仰中包含了全世界的救赎。其中那些维护宗徒信仰的人特别向您呼吁，他们曾治理过亚历山大教会，恳求您陛下不要让那些因其乖谬而合法被谴责的异端者继续篡位；因为无论您看他们错误的邪恶，还是考虑他们疯狂所行的事，他们不仅不能被接纳到司铎尊位，甚至配得被从基督徒之名中剪除。因为——我恳请您陛下原谅我这样说——当那些背信弃义的弑亲者敢于要求连无罪之人也不能合法获得的东西时，他们在某种程度上黯淡了您自己的光辉，最荣耀的皇帝。\par

% structure: p0009 source=h2 target=subsection reason=relative-heading-rank; base=h2

\subsection{四、呈递给皇帝的两份陈情书之间的区别。}

朕前曾呈递请愿书，其副本附于圣驾来函之中。然天主教徒所上之祈免书，内附署名名单；因其案情有理，故个人姓名乃至显赫品级皆坦然开列。反之，异端僭越者竟敢以乌合之众之暧昧许可，向我正统皇帝进呈之书，则隐匿一切特定姓名，其故非惟恐人数稀少，亦惧其品级暴露。盖彼等以为隐匿人数为便，虽其品级已显，且彼等自知当受谴责，故不敢宣扬其位份，亦非不当。故一文档含天主教徒之请愿，另一则陈异端之虚构。此处哀叹主之司铎、全体基督徒及隐修院之倾覆；彼处则展示巨恶之持续，以致本不应听闻之事竟得广为蔓延。\par

% structure: p0011 source=h2 target=subsection reason=relative-heading-rank; base=h2

\subsection{五、皇帝有机会彰显其信德。}

若亚历山大里亚教会，向为\Quote{祈祷之所}，而今竟不得为\BibleQuote{贼窝}\BibleRef{路 19:46}？盖因至残忍疯狂之暴行，天上奥迹之光尽灭，昭然若揭。祭献之奉献已绝，圣油之祝圣已废，一切奥迹皆从凶徒杀人之手隐退。凡此种种，彼等既行无可名状之亵渎，又流最尊贵司铎之血，焚其尸骨扬灰于风中以为戏，竟敢僭夺尊位之权，且于公议会前控告宗徒传承之不可侵犯信德，则当如何判决，毫无疑义。故尔，尔有绝佳良机，从主手加于尔冠冕之上亦信德之冠，向教会之敌奏凯：盖若能战胜敌对民族之军旅固可称颂，然使\Place{亚历山大里亚}教会脱离疯狂暴君，其荣耀何其大哉？此教会之苦难，乃全体基督徒之损伤。\par

% structure: p0013 source=h2 target=subsection reason=relative-heading-rank; base=h2

\subsection{六、许诺日后详陈信理，并请求纠正阿纳托略疏失之处。}

为使朕之函件于圣驾前如面谈然，朕以为凡关于我等共同信德之建议，须于后续信函中传达。且恐此信篇幅过长，故另函所述关于维护公教信德之要义，俾使宗座已发表之声明虽已足够，而此补充陈述亦能粉碎异端之网罗。盖圣驾具司祭与宗徒之心志，更当因此疫疠之恶而燃起义怒，此恶玷污了\Place{君士坦丁堡}教会之纯洁，其中发现若干神职人员，附和异端之解释，并在教会心脏之地以支持助长之。若朕之兄弟\Person{阿纳托略}因过分宽厚之仁慈而疏于清除彼等，伏请彰显圣驾信德，亦以此次疗救施于教会，使此类人等不仅被逐出神职之列，且不得居于城内。朕将圣驾之忠仆、主教\Person{尤利安}与司铎\Person{埃提乌斯}托付于尔，恳请俯允静听彼等为维护公教信德所陈之建议，盖彼等乃真诚之人，凡事皆有助于圣驾之信德。颁于十二月一日，显贵\Person{君士坦丁}与\Person{鲁弗斯}执政之年（四五七年）。\par

\SourceRecord{Translated by Charles Lett Feltoe.；Nicene and Post-Nicene Fathers, Second Series；Vol. 12.；Edited by Philip Schaff and Henry Wace.；Buffalo, NY: Christian Literature Publishing Co.,；1895.；<http://www.newadvent.org/fathers/3604156.htm>.}

% structure: p0001 source=h1 target=section reason=page-title

\section{第一五八号书信}

\Person{良}致书寄居在\Place{君士坦丁堡}的埃及公教主教们。\par

% structure: p0003 source=h2 target=subsection reason=relative-heading-rank; base=h2

\subsection{他鼓励他们为信仰受苦，并向皇帝恳求伸冤。}

我之前就因听闻在\Place{亚历山大}所犯的罪孽而悲伤，我的心灵因那暴行的残忍而创伤，不知该展示何等的泪水，发出何等的哀悼，且愿意用先知的话说：\BibleQuote{谁给予我的头水源，给我的眼睛泪泉？}然而，预见到你们的抱怨，亲爱的，我已向最仁慈的基督君王恳求补救这些大恶，并通过我们的儿子和助手\Person{格隆提乌斯}和\Person{奥林匹乌斯}在不同时间要求他尽快清除\Place{亚历山大}那城市教会中一个已被判罪的异端——在那里曾有许多公教教师兴盛——并且不允许那些无法无天的杀人者——他们毫无对场所或时间的敬畏，竟敢流统治者的血——从基督的仁慈中获得任何好处，尤其当他们想要重新审议\Place{加采东}大公会议以推翻信仰之时。因此，亲爱的，那使你们离开自己的\OriginalTerm{教座}{Sees}的理由，应当安慰你们的受苦；因为确知，为祂的名而遭受逆境的受苦灵魂，绝不会失去主的保护。因此勇敢地忍受，并记住那属于你们的家乡，为你们如今寄居异乡而喜乐。不要为流放而忧伤，不要为当前的疲惫而悲哀，你们知道宗徒曾以许多为信仰而临的险境夸耀。你们有一位知道你们的争战，并准备了补偿的赏报。不要让任何人退缩于这劳苦，其赏报是永远为王和永生。让所有战斗者的脚坚定地立在\Place{耶路撒冷}的厅堂中；因为在那报偿的希望中，他们将没有理由惧怕敌营或敌人的攻击。得胜从不困难，战胜卑劣敌人的残余——这敌人已被全世界一致击溃——的胜利也非艰难，尤其当你们看到其首领已经仆倒。因此，以不断的祈祷（正如我也没有停止过），恳求最基督的君王的恩宠，他在天主的仁慈中预备垂听：按照我已寄出的信件，愿他以我们确信他拥有的心思奋发，坚固共同信仰的事业，并以他的虔敬消除异端者的疯狂所编造的一切有害控告，安排你们的回归，亲爱的，并使每个省份和所有教会连同其司铎，因基督的平安而喜乐。日期：457年12月1日，\Person{君士坦丁}和\Person{鲁富斯}执政之年。\par

\SourceRecord{Translated by Charles Lett Feltoe.；Nicene and Post-Nicene Fathers, Second Series；Vol. 12.；Edited by Philip Schaff and Henry Wace.；Buffalo, NY: Christian Literature Publishing Co.,；1895.；<http://www.newadvent.org/fathers/3604158.htm>.}

% structure: p0001 source=h1 target=section reason=page-title

\section{第一五九书}

\Person{良}主教，致\Place{阿奎莱亚}主教\Person{尼塞塔}，请安。\par

% structure: p0003 source=h2 target=subsection reason=relative-heading-rank; base=h2

\subsection{一．序言。}

我儿\Person{阿德奥达图斯}，我教区的执事，在返回我们这里时，已转交了您的请求，可敬的，希望从我们这里获得宗座的权威，处理那些看似难以决定、但我们必须考虑到时代的需要而加以预备的事项，好使因敌人攻击所造成的创伤，主要通过宗教的力量得以医治。\par

% structure: p0005 source=h2 target=subsection reason=relative-heading-rank; base=h2

\subsection{二．关于丈夫被俘而再嫁的妇女。}

既然您说，由于战争的灾难和敌人严重的入侵，有些家庭已被拆散，以致丈夫被掳为囚，妻子被遗弃，而后者因认为自己丈夫已死或绝不可能从主人手中获释，遂在孤寂的压力下与他人缔结了婚姻；如今，借着主的帮助，情势已有改善，一些原本被认为已死去的人归来了。亲爱的弟兄，您似乎自然要怀疑，关于这些与他人结合的妇女，我们该如何定夺。但由于我们知道经上记载：\BibleQuote{贤惠的妻子是上主所赐}\BibleRef{箴 19:14}，并且我们也知晓那条诫命：\BibleQuote{天主所结合的，人不可拆散}\BibleRef{玛 19:6}，我们坚持认为，合法婚姻的契约必须重建，并且在敌人所加祸害消除后，各人合法拥有的一切必须归还；我们必须竭尽全力，使每个人收回自己的所有。\par

% structure: p0007 source=h2 target=subsection reason=relative-heading-rank; base=h2

\subsection{三．娶被俘者之妻者是否应受责难？}

但是，凡取代了被认为已不在世的丈夫而娶其妻者，不应被视为应受责难或侵犯他人权利者。因为许多属于被掳之人的东西已经转归他人所有，而他们回来后，财产理应恢复。如果这在奴隶、土地、房屋乃至动产上得到恰当遵守，那么对于妻子的恢复，岂不更应该如此，好使战争必需所扰乱之事，藉和平的补救得以恢复？\par

% structure: p0009 source=h2 target=subsection reason=relative-heading-rank; base=h2

\subsection{四．妻子必须归还其原夫。}

因此，如果长期被俘后归来的丈夫仍对妻子怀有如此深情，以致愿意她们回归婚姻关系，那么，因必需而生之事应予以忽略，并视为无可指责，而信德的要求则应予以满足。\par

% structure: p0011 source=h2 target=subsection reason=relative-heading-rank; base=h2

\subsection{五．拒绝返回的妇女应受绝罚。}

若有妇女深爱其后的丈夫，以致宁愿继续与他同居，而不愿回到合法的配偶身边，则她们理当受到谴责，甚至被剥夺教会的共融；因为在一件可宽恕的事上，她们选择了以罪恶玷污自己，显示出她们在失节中寻求自己的享乐，而正当的归回本能使她们获得宽恕。因此，让她们返回原来的状态，并自愿做出补赎，决不可让必需之境逼迫她们做的事，因其邪恶的欲望而转为耻辱；因为正如拒绝返回丈夫身边的妇女应被视为不圣洁，那些回到上主认可的情爱中的人，则理当受赞扬。\par

% structure: p0013 source=h2 target=subsection reason=relative-heading-rank; base=h2

\subsection{六．关于被迫吃祭肉的被俘者。}

关于那些据称在被掳期间被祭肉玷污的基督徒，我们认为对您的询问应作如下答复：亲爱的弟兄，他们应通过适当的补赎来洁净自己，补赎的衡量与其说是时间的长短，不如说是痛悔的强度。无论他们的顺从是由于恐怖还是饥饿，都无需犹豫予以开释，因为那食物是出于恐惧或需要而吃，并非出于迷信的敬意。\par

% structure: p0015 source=h2 target=subsection reason=relative-heading-rank; base=h2

\subsection{七．关于因恐惧或错误而重受洗者。}

至于您所认为，亲爱的，我们也应予以商议的那些人，他们或因恐惧所迫，或因错误而重复领受洗礼，如今却明白自己行径与公教信仰的规条相悖；对他们应持如此温和态度：准其与我们完全共融，但须以补赎的医治和主教的覆手为条件；补赎的期限（适度考虑）留给您判断，依您所见忏悔者心灵的倾向而定；在此事上，万勿忽略年老、疾病及其他风险。因为若有人处在如此危险的境地，以致生命无望，而仍在补赎之中，则司铎的慈祥关怀应赐予共融的恩助。\par

% structure: p0017 source=h2 target=subsection reason=relative-heading-rank; base=h2

\subsection{八．关于异端者施行的洗礼。}

凡此前未受洗者，若从异端者领受洗礼，只需藉呼求圣神而覆手以坚振，因为他们仅领受了洗礼的仪式，而未得圣化的德能。如您所知，我们要求在所有教会中遵守此规，使已进入的洗礼之泉不致因重复而受玷污，正如主所说：\BibleQuote{只有一个主，一个信德，一个洗礼}。并且那洗濯不致因重复而被亵渎，而只是呼求圣神的圣化，以使无人能从异端者领受的，可由公教司铎获得。我们这封回信，已答复了兄弟团体的垂询，您应将此信告知省内所有兄弟及同僚主教，使我们的权威一经给予，即可供普遍遵守。于三月二十一日，马约里安皇帝执政时（四五八年）。\par

\SourceRecord{Translated by Charles Lett Feltoe.；Nicene and Post-Nicene Fathers, Second Series；Vol. 12.；Edited by Philip Schaff and Henry Wace.；Buffalo, NY: Christian Literature Publishing Co.,；1895.；<http://www.newadvent.org/fathers/3604159.htm>.}

% structure: p0001 source=h1 target=section reason=page-title

\section{书信第一百六十二号}

由\Person{菲洛克塞努斯}\Emph{\OriginalTerm{事务代理人}{agens in rebus}}之手。\Person{利奥主教}致\Person{利奥·奥古斯都}。\par

% structure: p0003 source=h2 target=subsection reason=relative-heading-rank; base=h2

\subsection{一、加采东与尼西亚的法令是相同且终局的。}

我的心灵在主内欢跃，充满极大的喜乐，现在我因看出陛下至卓越的信仰在一切事上因天上的恩宠而增进，并以更大的勤勉体会到您内里司铎精神的虔敬，我大有感恩的理由。因为在陛下的文告中，无疑显明了圣神为整个教会的好处藉着您所成就的，以及全体信徒的祈祷如何热切期盼您的帝国在荣耀中处处扩展，因为您除了关心现世之事外，还如此恒久地运用宗教远见服务于神圣永恒之事：即那唯一赋予人类生命并使之圣化的天主教信仰得以存于同一宣认中，而源自地上意见多样的分歧得以从那筑有天主之城的坚固磐石上被驱散。最荣耀的皇帝啊，这些天主的恩赐终将来自祂，若我们不对所蒙恩赐忘恩负义，好像我们所得的全无价值，反而追求完全相反的事。因为寻求已发现的，重新考虑已完成的，推翻已界定的，若非是对已获恩赐不感恩，并放纵死欲之不洁渴望于禁树之果，还能是什么呢？因此，您屈尊显示对普世教会和平更周全的关怀，便清楚认识到异端者强大阴谋的企图：即期望在欧提克斯和狄奥斯库若的门徒与宗座代表之间进行更周详的讨论，仿佛尚未有任何界定，以使加采东神圣公会议在全世界天主教司铎欣然赞同下所决定之事失效，同时也损害最神圣的尼西亚大公会议。因为在我们时代于加采东关于我们的主耶稣基督降生成人的决定，也已由尼西亚的那神秘数目的教父如此界定，免得天主独生子在何事上与圣父不平等，或当祂成为人子时不曾拥有我们血肉与灵魂的真性体。\par

% structure: p0005 source=h2 target=subsection reason=relative-heading-rank; base=h2

\subsection{二、必须坚定不移地抵制异端者的邪恶图谋。}

因此，我们必须憎恶并恒常避免异端欺诈所力求获得的，也不应将已妥善充分界定之事再次置于讨论之下，免得我们自己在被定罪之人的意愿下，对显然与先知、福音作者和宗徒的权威完全一致的事产生怀疑。因此，若有任何人不同意这些受天启的决定，就让他们保留自己的意见，并带着他们所选的那邪恶派别离开教会的合一。因为那些敢于反对神圣奥迹的人绝不可能与我们有任何共融。让他们以空谈为荣，并夸耀他们反对信仰的论证之巧妙：我们却乐意听从宗徒的训诫，他说：\Quote{你们要小心，免得有人用哲学和虚妄的欺哄，\newline{}\BibleRef{哥 2:8}把你们掳去。}因为同一宗徒说：\Quote{若我重建我所拆毁的，就证明我是个过犯的人，\newline{}\BibleRef{迦 2:18}}并且使自己承受那不仅蒙福记忆的马尔西安皇帝的权威，而且我自己也以同意所接受的惩罚条件。因为正如您公正而真实地主张，完善不能增加，圆满不能添加。因此，既然我知道您，可敬的皇帝，被真理最纯净的光所浸透，在信仰的任何部分都不动摇，而以公正完美的判断分辨是非，并区分应接受和应拒绝的，我恳求您不要认为我的谦卑因缺乏信心而应受责备，因为我的谨慎不仅有利于普世教会，也有助于您自己的荣耀，即在您的统治下，异端者的肆无忌惮不致显得得势，而天主教徒的安全不致受扰。\par

% structure: p0007 source=h2 target=subsection reason=relative-heading-rank; base=h2

\subsection{三、他承诺派遣使节，不是与欧提克斯派讨论，而是向皇帝解释信仰。}

因此，虽然我对您心中的虔诚在所有事上非常信赖，并且看出因住在您内的天主之神，您已受到足够教导，没有任何错误能欺骗您的信仰，但我仍将尽力遵从您的命令，派遣几位兄弟在您面前代表我本人，陈述宗徒信仰准则是什么，尽管如我所说，这为您是众所周知的，并使一切事明确无疑：凡不接受可敬的尼西亚大公会议定义或加采东大公会议法令的人，绝不列于天主教徒之列，因为两道神圣法令显然出自福音与宗徒之源，凡非基督之水灌溉的，就如蛇毒之饮。陛下应预先知道，最可敬的皇帝，我所派遣之人将来自宗座，不是要与信仰的敌人争战，也不是要与任何人争论，因为已经由天主所喜悦的在尼西亚和加采东所确定的事，我们不敢再进入任何讨论，仿佛如此伟大的权威藉圣神所定的是可疑或软弱的。\par

% structure: p0009 source=h2 target=subsection reason=relative-heading-rank; base=h2

\subsection{四、必须迫使异端者放弃其僭取，并交托于天主的审判。}

但我们并不拒绝以我们的职务来辅助教诲我们的小孩子，他们在用奶哺养之后，渴望以更坚固的食物得到满足。正如我们不轻视朴实的人，同样我们也不与叛逆的异端者往来，牢记主的命令，祂说：\Quote{不要把圣物给狗，也不要把你们的珍珠丢在猪前}\BibleRef{玛 7:6}。确实，让圣神藉先知的口所描述的人——\Quote{外邦人的儿子向我撒谎}——参与自由讨论，是全然不当和不公的。因为即使他们不反对福音，他们却显示出是属于那些如经上所记载的人：\Quote{他们声称认识天主，但以行为否认祂}\BibleRef{铎 1:16}，同时义人亚伯尔的血仍在向恶人加音呼喊，加音受主斥责后，并未安静地从事悔改，反而爆发为杀人。我们希望他的惩罚抑留于主的审判，使这个无耻的掠夺者和嗜血的凶手被抛回自身，放弃属于我们的东西。我们也祈求你们，不再容忍亚历山大的圣教会的可悲被掳状态延续下去，它应藉你们的信德和正义的帮助重获自由，以使在埃及的各城中，教父的尊严和他们的司铎权利得以恢复。日期：三月二十一日，利奥与马约里安·奥古斯都执政时期（458年）。\par

\SourceRecord{Translated by Charles Lett Feltoe.；Nicene and Post-Nicene Fathers, Second Series；Vol. 12.；Edited by Philip Schaff and Henry Wace.；Buffalo, NY: Christian Literature Publishing Co.,；1895.；<http://www.newadvent.org/fathers/3604162.htm>.}

% structure: p0001 source=h1 target=section reason=page-title

\section{第一六四封信}

\Person{良}主教，致\Person{利奥}皇帝。\par

% structure: p0003 source=h2 target=subsection reason=relative-heading-rank; base=h2

\subsection{一、他派遣使节，但劝阻任何对信仰的新讨论。}

由于诸多明证使我确信，您以何等热忱关怀普世教会的利益，我未曾迟延，立即遵行陛下之命，派遣我的兄弟与同僚主教\Person{多米提安}和\Person{杰米尼安}。他们为促进我热切的祈祷，将恳求您平安接受福音训诲，并获得信仰的自由——您借圣神的教导，在此信仰中已显赫出众。如今基督的仇敌已被远远驱散，他们即使想隐藏疯狂也无法隐匿，因为主的羊群那神圣纯朴与披着羊皮隐藏的野兽之伪装截然不同，且因其至极的狂妄暴露无遗，他们再也无法借伪善渗入。因此，尊贵可敬的皇帝，请认清天主眷顾如何召唤您守护整个世界，并明白您欠您的母亲——以您为荣的教会——何等援助。已终结的争端绝不允许重新振作，来对抗全能者右手的胜利，尤其此事绝不可允许异端者，他们的企图早已被定罪，信友的劳苦正当地要求此结果：整个教会圆满地安全存于其合一的完整，且任何已经妥善奠定之事不得再行考虑，因为在神圣引导下合法制定的规章之后，仍愿争辩并非和平，而是悖逆之灵的表现，如宗徒所说：\Quote{因为以言辞争辩毫无益处，只会败坏听者。}\par

% structure: p0005 source=h2 target=subsection reason=relative-heading-rank; base=h2

\subsection{二、在信仰之事上，人的修辞学不合时宜。}

因为若人类幻想总能在争论中自恃，就绝不缺少敢于反对真理、倚靠此世智慧油滑言辞的人，而基督的信仰与智慧从主耶稣基督自身的教导中得知，应如何严格地避开这种最有害的虚妄。当基督即将召唤万民归信之时，祂不是从哲学家或演说家中拣选那些献身于宣讲福音的人，而是选取卑微的渔夫作为彰显自己的工具，免得那本身充满大能的天上教导似乎需要言辞的帮助。因此宗徒抗议说：\BibleQuote{因为基督派遣我，不是为施洗，而是为宣讲福音，且不用智慧的言辞，免得基督的十字架失去效力；因为十字架的道理，为那丧亡的是愚妄，为我们这些得救的，却是天主的德能。经上记载：我要摧毁智者的智慧，废弃贤者的聪明。智者在哪里？经师在哪里？这世代的探究者在哪里？难道天主不是使这世界的智慧成为愚拙了吗？\BibleRef{格前 1:17}}\par

% structure: p0007 source=h2 target=subsection reason=relative-heading-rank; base=h2

\subsection{三、\Person{欧提克斯}的教义已被圣经的见证所谴责，不能再予以考虑。}

然而，没有什么比那些被自己发明所欺骗的人更远离福音的光明，因为他们不认为主的降生成人真实地涉及人性——即我们的本性；仿佛天主的荣耀不配那不能受苦的圣言之尊威取了真实的人性，而人类的救恩本不能恢复，除非那具有天主形像者屈尊取了奴仆的形像。因此，既然加采东大公会议——罗马世界各省份均出席，其决定获得普世接受，且与最神圣的尼西亚会议完全一致——已将一切\Person{欧提克斯}教义的邪恶追随者从大公教会的共融中剪除，那么任何失足者若不通过完全的补赎净化自己，又如何能重获教会的平安呢？因为有何自由可允许他们讨论？他们已因公正神圣的判决应受定罪，从而最真实地落在蒙福宗徒的那句话下，他在初期教会之初就用这话推翻了基督十字架的仇敌，说：\BibleQuote{凡明认耶稣基督是在肉身内降世的神，便是出于天主；凡否认耶稣的神，不是出于天主，而是反基督的。}我们必须忠信坚定地运用这圣神预先赐予的教导，免得因接纳这些人的讨论而削弱天主启示的法令之权威，当在您王国的各地及地极，那在加采东所确认的信仰正建立在最稳固的平安基础之上，且任何自绝于与我们共融的人都不配称为基督徒。关于这些人，宗徒说：\BibleQuote{对异端的人，在第一次和第二次劝戒以后，就该弃绝，知道这样的人已背道，且自定己罪。\BibleRef{铎 3:10}}\par

% structure: p0009 source=h2 target=subsection reason=relative-heading-rank; base=h2

\subsection{四、若欲行使天主的仁慈，异端者必须完全停止其道路的错谬。}

因此，那不圣的弑亲者侵犯圣教会并残忍杀害其首领所犯下的罪行，不能因人的宽恕而赎补，除非那唯独能公正惩罚此类罪行、又唯独能以其不可言喻的慈悲赦免它们的主被平息。但我们虽不急于报复，却绝不能与魔鬼的仆役结盟。然而，若我们得知他们正离开异端的行列，悔改自己的错误，从纷争的武器转向悲哀的痛悔，我们也可以为他们代祷，以免他们永远丧亡，如此效法主的慈爱榜样——当祂被钉在十字架上时，曾为迫害者祈祷说：\BibleQuote{父啊，宽赦他们罢，因为他们不知道他们做的是什么。}\BibleRef{路 23:34}为使基督徒的爱能如此有益地为仇敌施行，邪恶的异端者必须停止骚扰天主始终虔诚敬拜的教会；他们不得再以谎言扰乱单纯者的心灵，以致于在所有以往时代纯正信仰得以兴盛之处，福音和宗徒的教导如今也能自由传扬；因为我们也尽己所能效法天主的慈悲，不愿任何人受正义的惩罚，而愿所有人都因慈悲获得释放。\par

% structure: p0011 source=h2 target=subsection reason=relative-heading-rank; base=h2

\subsection{五、他应召回避难的圣职人员和教友，并彻底弃绝坚持异端者。}

我恳求陛下的仁慈，请听我前述弟兄们的建议；正如我先前在一封前函中所言，我派遣他们并非要与受定罪者争辩，而只是为天主教信仰的稳固向陛下代祷。根据您对天主威严的信德和敬意，尤其应赐予这一点：彻底放下异端者的争论，仁慈地垂顾那些陷入如此厄运的人，在恢复亚历山大里亚教会至其原始状态后，在那里设立一位主教——他应持守加采东大公会议的法令，与福音的规条一致，并能为极度动荡的人民恢复和平。那些被不圣的弑亲者逐出教会的主教和圣职人员，也应奉陛下之命被召回；所有其他因类似恶意而被逐离住所的人，也应恢复其原有地位，使我们能因天主的恩宠和您的信德而充分、完全地喜乐，不再有任何纷争的噪音。因为若有人如此忘记基督徒的望德和自身的得救，竟敢以任何争论攻击圣加采东大公会议的福音和宗徒法令，从而也推翻最神圣的尼西亚大公会议，那么我们以同样的绝罚和相等的咒骂谴责他，并与所有对主耶稣基督的降生成人持有亵渎可憎观点的异端者一同定罪；如此，在未拒绝给予那些完全合法补偿者悔改之补救的前提下，基于真理的大公会议之判决仍可加于那些继续抵抗者。签署于458年8月17日，利奥和皇帝麦约里安执政之年。\par

\SourceRecord{Translated by Charles Lett Feltoe.；Nicene and Post-Nicene Fathers, Second Series；Vol. 12.；Edited by Philip Schaff and Henry Wace.；Buffalo, NY: Christian Literature Publishing Co.,；1895.；<http://www.newadvent.org/fathers/3604164.htm>.}

% structure: p0001 source=h1 target=section reason=page-title

\section{书信 166}

\Person{良}主教，致\Place{拉文纳}主教\Person{内奥}，问安。\par

% structure: p0003 source=h2 target=subsection reason=relative-heading-rank; base=h2

\subsection{一、那些在婴孩时期被掳走、无法回忆亦无法为自身圣洗提供见证的人，不应被拒绝此圣事。}

我们确实经常在圣神的教导下，当弟兄们的心在疑难问题的险滑处摇摆不定时，通过从圣经的教导或教父的规则中拟定答案来稳定他们；但最近在主教会议上，出现了一个新的、迄今闻所未闻的辩论议题。因为通过某些弟兄的提示，我们发现一些战俘在自由返回家园后——即那些在年幼时被掳走、无法对任何事情有确切认知的人——渴望获得圣洗的治愈之水，但由于婴儿时期无知，不能记得自己是否领受了圣洗的奥迹和礼仪；因此，在这种记忆缺失的不确定性中，他们的灵魂陷入危险，只要在看似谨慎的借口下，他们被拒绝了一项恩宠，而这恩宠之所以被扣留，是因为被认为已经赐予了。于是，由于某些弟兄基于并非不合理的恐惧而犹豫执行主之奥迹的礼仪，在一次主教会议（如我们所说）上，我们收到了关于此事的正式咨询请求，并在仔细讨论后，我们想权衡每位成员的意见，并以如此谨慎的方式处理，以便借助多人的知识确实抵达真理。因此，正是那些由神圣灵感进入我们心中的同样事物，得到了众多弟兄的同意和确认。所以，我们首先必须留意，免得当我们固守某种看似谨慎的态度时，却导致那些将要重生之人的灵魂丧失。因为谁会如此执着于猜疑，以至于认为仅凭猜测而无任何证据所怀疑的事情就是真的呢？因此，凡渴望重生的人自己无法回忆其圣洗，也无人能对其对天主的奉献作见证时，就没有引入罪过的可能性，因为就他们所知，奉献的施予者和领受者都无可指责。我们确实知道，每当有人根据异端者的制度（圣教父们已谴责了这些制度）被迫两次进入洗礼池——这洗礼池对那些将要重生的人来说只能用一次——就是犯下了不可饶恕的罪行，这与宗徒的教导相反，宗徒教导我们说：在天主圣三中只有一个天主性，在信德中只有一个宣认，在圣洗中只有一个圣事。但在当前情形中，无需担心类似的事情，因为完全未知是否曾做过的事，不构成重复的指控。因此，每当出现这种情况时，首先通过仔细调查来筛别，并花费相当的时间（除非他的临终时刻临近），去查问是否完全没有人能够以自己的证言帮助对方的无知。当确定那个需要圣洗圣事的人仅仅是由于毫无根据的怀疑而被阻止时，让他勇敢地前来获得这项恩宠吧，因为他自己在内心中没有发觉任何这恩宠的痕迹。我们无需害怕这样敞开救恩之门，因为此前并未显示有人已经进入过此门。\par

% structure: p0005 source=h2 target=subsection reason=relative-heading-rank; base=h2

\subsection{二、不可因再次施洗而使异端者所施行的圣洗圣事归于无效。}

但如果确定一个人曾由异端者施洗，绝不可在他身上重复重生的奥迹，而只应补行先前所欠缺的，使他能通过主教的覆手获得圣神的力量。亲爱的弟兄，我们希望此项决定能传达到你们全体，以免天主的仁慈因不当的畏惧而被拒绝给那些渴望得救的人。日期：十月廿四日，\Person{马约里安·奥古斯都}执政之年（458年）。\par

\SourceRecord{Translated by Charles Lett Feltoe.；Nicene and Post-Nicene Fathers, Second Series；Vol. 12.；Edited by Philip Schaff and Henry Wace.；Buffalo, NY: Christian Literature Publishing Co.,；1895.；<http://www.newadvent.org/fathers/3604166.htm>.}

% structure: p0001 source=h1 target=section reason=page-title

\section{Letter 167}

\Emph{致纳博讷高卢主教鲁斯蒂库斯，并就其提出的各项问题给予答复。}\par

% structure: p0003 source=h2 target=subsection reason=relative-heading-rank; base=h2

\subsection{一、劝勉他对两位冒犯他的主教持温和态度。}

主教良致纳博讷高卢主教鲁斯蒂库斯。\par

兄弟，你的信函由你的总执事赫尔墨斯带来，我欣然收悉；信中内容繁多，确实篇幅较长，但在我耐心细读之下，并未因四面八方的烦务而忽略任何一点。因此，在掌握了你的指控要点，并审查了主教及首领们调查过程中所发生的事之后，我们得知，司铎萨比尼安和良对你的行动缺乏信任，而且他们已不再有任何正当的抱怨理由，因为他们自愿退出了已开始的讨论。你应当以何种方式或何种程度的正义对待他们，我交由你自己斟酌；然而，我以爱德的劝勉提醒你，应对患病者施行属灵的药物，并记住圣经所言：\Quote{不要过于正义}，对这些人应温和行事，因为他们似乎出于对贞洁的热忱而超越了报复的限度，免得那欺骗了奸淫者的魔鬼，反而在奸淫的报复者身上得胜。\par

% structure: p0006 source=h2 target=subsection reason=relative-heading-rank; base=h2

\subsection{二、劝诫他不要因想放弃职守而表现出对天主应许的不信任。}

然而，亲爱的，我感到惊讶的是，你竟因各种原因引起的冒犯而如此困扰，以至于说宁可放下主教职务的辛劳，过安静闲适的生活，也不愿继续承担托付给你的职责。但是，主既说：\BibleQuote{唯独坚持到底的，才可得救}\BibleRef{玛 24:13}，这坚持到底的福分从何而来，若非来自忍耐的力量？因为正如宗徒所宣告的：\BibleQuote{凡是愿意在基督耶稣内虔敬生活的，都必遭受迫害}\BibleRef{弟后 3:12}。而且，并非只有刀剑、烈火或其他激烈手段针对基督信仰才算迫害；最残酷的迫害往往来自习俗的悖逆、持续的顽抗以及恶舌的刺伤。既然教会的所有肢体常受这些攻击，信友中无一不受诱惑，无论是安逸还是辛劳的生活都离不开危险，那么，若舵手擅离岗位，谁能引导船只在海浪中前行？若牧者自己不加警惕，谁能保护羊群免于豺狼的诡计？最后，若喜爱安逸的念头使望楼的守望者松懈了严格的警戒，谁能抵抗盗贼和强盗？因此，必须坚守所托付的职守和承担的任务。正义必须坚定不移地持守，仁慈必须满含爱德地施予。要憎恨的不是人，而是他们的罪。要斥责骄傲的人，容忍软弱的人；那些需要更严厉管教的罪，应以医治的意愿而非报复的精神来处理。如果我们遭遇更猛烈的苦难风暴，切莫惊恐，仿佛要靠自己的力量克服灾难，因为我们的谋略和能力都是基督，藉着祂我们能做一切，没有祂我们什么也不能做。祂为坚定福音的宣讲者和奥迹的施行者，说：\BibleQuote{看！我同你们天天在一起，直到今世的终结}\BibleRef{玛 28:20}。祂又说：\BibleQuote{我给你们讲了这些，是要你们在我内得到平安。在世界上你们要受苦难；然而你们放心，我已战胜了世界}\BibleRef{若 16:33}。这些应许如此清楚，我们不应让任何冒犯的原因削弱它们，免得我们对天主使我们成为祂特选的器皿忘恩负义，因为祂的援助大能，祂的应许真实。\par

% structure: p0008 source=h2 target=subsection reason=relative-heading-rank; base=h2

\subsection{三、所提的许多问题如果当面交谈会比书面更容易解决。}

亲爱的，关于你的总执事单独写给我的那些询问要点，如果你能让我们当面讨论，就每个问题更容易得出结论。因为有些问题似乎超出了常规考证的范围，我认为它们更适合谈话而非书写：既然有些事绝无可争议，也有许多事需要根据年龄或情况的需要加以调整；但务必记住，在可疑或不明之事上，应遵循既不违反福音命令又不违背圣教父法令的做法。\par

% structure: p0010 source=h3 target=subsubsection reason=relative-heading-rank; base=h2

\subsubsection{问题一}

\Emph{关于自称主教的司铎或执事，以及他们所祝圣的人。}\par

\Emph{答复}：决不允许将那些未经圣职界选举、未经平信徒要求、未经省区主教在总主教同意下祝圣的人算作主教。因此，既然常有人不当晋升的问题，谁还会怀疑，那显然未授予他们的东西绝不可能有效？如果这些假主教在自有主教的教会中祝圣了某些圣职人员，且祝圣得到了合法主教的同意和批准，则只要他们继续留在同一教会，可视为有效。否则，必须视为无效，因为它既无固定场所，又无任何权威依据。\par

% structure: p0013 source=h3 target=subsubsection reason=relative-heading-rank; base=h2

\subsubsection{问题二}

\Emph{关于司铎或执事，在其罪行暴露后请求公开补赎，是否应藉覆手礼准许之？}\par

答：凡已蒙受司铎尊位或执事品秩之人，若犯任何罪行，不得藉按手礼接受补赎的救治，此乃违反教会惯例之事；此项惯例无疑源自宗徒传统，正如经上所载：\BibleQuote{司铎若犯了罪，谁来为他祈祷？}\BibleRef{撒上 2:25}因此，这样的人失足后，为求得天主的仁慈，必须寻求私下的退隐之处，在那里他们的补赎才可能既有益又充足。\par

% structure: p0016 source=h3 target=subsubsection reason=relative-heading-rank; base=h2

\subsubsection{问题三}

关于在祭台供职且已有妻室者，是否可合法与她们同居？\par

答：祭台服务人员与主教及司铎在贞洁律上相同；他们身为平信徒或读经者时，可合法结婚生育。但一旦擢升至上述品级，先前合法之事便不再合法。因此，为使他们的婚姻成为属灵而非属肉体的，他们不应休妻，而应\Quote{有妻如无妻}，如此既可保持对妻子的情感，又终止婚姻的职能。\par

% structure: p0019 source=h3 target=subsubsection reason=relative-heading-rank; base=h2

\subsubsection{问题四}

关于司铎或执事将其未婚女儿嫁给已有女子结合且生有子女的男子。\par

答：并非每个与男子结合的女子都是他的妻子，正如并非每个儿子都是其父的继承人。婚姻盟约在生而自由者之间、在同等人之间才是合法的：这是主在罗马法律诞生之前早已制定的。因此妻子与妾不同，正如婢女与自由妇不同。为此宗徒为表明这二者的区别，引用了\WorkTitle{创世纪}中对亚巴郎说的话：\BibleQuote{把婢女和她的儿子赶走，因为婢女的儿子不能与我的儿子依撒格一同继承产业。}因此，既然婚姻盟约从一开始就如此设立，使其在性结合之外，还象征基督与祂的教会的奥秘结合，那么，在该女子身上若显示婚姻的奥秘未曾发生，则她毫无疑问与婚姻无分。因此，任何等级的圣职人员，若将女儿嫁给一个有妾的男子，不应被视为将其嫁给了已婚之人，除非那另一个女子似乎已获自由、有合法妆奁，并经公开婚礼而受到尊重。\par

% structure: p0022 source=h3 target=subsubsection reason=relative-heading-rank; base=h2

\subsubsection{问题五}

关于年轻女子嫁给有妾的男子。\par

答：凡由父命而嫁给丈夫的女子，若其丈夫先前所拥有的女子不属于婚姻关系，则她们无可指摘。\par

% structure: p0025 source=h3 target=subsubsection reason=relative-heading-rank; base=h2

\subsubsection{问题六}

关于离开与其生有子女的女子而另娶妻者。\par

答：既然妻子与妾不同，将婢女从卧榻上辞退，另娶一位确定生而自由的妻子，这不是重婚，而是可敬之举。\par

% structure: p0028 source=h3 target=subsubsection reason=relative-heading-rank; base=h2

\subsubsection{问题七}

关于那些在患病时接受补赎条件，痊愈后又拒绝遵守的人。\par

答：此等之人的疏忽应受责备，但不应最终放弃他们，好使他们因屡次劝勉而被激励，忠诚地履行他们在急难中所请求之事。因为只要人仍在此肉身内，就不可对任何人绝望；因为有时因年迈畏缩而推迟之事，藉更成熟的思虑得以成就。\par

% structure: p0031 source=h3 target=subsubsection reason=relative-heading-rank; base=h2

\subsubsection{问题八}

关于那些临终时承诺悔改，却在领受圣体前去世的人。\par

答：他们的情况应交由天主的审判，因为他们的死亡推迟到领圣体那一刻，是掌握在天主手中。但我们不能与那些在生前我们未与之共融的人，在死后共融。\par

% structure: p0034 source=h3 target=subsubsection reason=relative-heading-rank; base=h2

\subsubsection{问题九}

关于那些在剧痛之下请求给予补赎，而当司铎前来施行时，若痛苦稍减，便推辞拒绝接受所予之物的人。\par

答：这种推托并非出于对救治的轻视，而是出于害怕陷入更重的罪。因此，那被推迟的补赎，当更恳切地寻求时，不应被拒绝，好使受伤的灵魂能以任何方式获得赦免的医治。\par

% structure: p0037 source=h3 target=subsubsection reason=relative-heading-rank; base=h2

\subsubsection{问题十}

关于那些已宣示悔改的人，若开始在法庭上争讼。\par

答：追讨正当债务固然是一回事，因完美的爱德而置己有财产于不顾是另一回事。但祈求宽恕不法行为者，甚至应当戒除许多合法之事，如宗徒所说：\BibleQuote{凡事我都可行，但并非事事都有益处。}\BibleRef{格前 6:12}因此，若悔罪者有一件或许不该忽视的事，他最好诉诸教会的裁判，而非世俗法庭。\par

% structure: p0040 source=h3 target=subsubsection reason=relative-heading-rank; base=h2

\subsubsection{问题十一}

关于那些在补赎期间或之后从事商业活动的人。\par

答：其收益的性质或可为商贩开脱，或可定其罪，因为利润有可敬的与卑劣之别。然而，对悔罪者而言，蒙受损失比卷入贸易的风险更为有益，因为在买卖双方的交易中，很难避免不犯罪。\par

% structure: p0043 source=h3 target=subsubsection reason=relative-heading-rank; base=h2

\subsubsection{问题十二}

关于那些在补赎之后重返军役的人。\par

答：在补赎之后重返世俗军役，完全违反教会规条，如宗徒所说：\BibleQuote{天主的士兵不将俗务缠身。}\BibleRef{弟后 2:4}因此，那愿意缠身于世俗军役者，不能免于魔鬼的罗网。\par

% structure: p0046 source=h3 target=subsubsection reason=relative-heading-rank; base=h2

\subsubsection{问题十三}

关于那些在补赎之后娶妻或与妾结合的人。\par

\Emph{答}：若一青年因恐惧死亡或被俘的危险已行了补赎，此后又因害怕陷入少年时的淫欲而选择娶妻，以免犯奸淫，则其所行似乎可蒙宽恕，只要他除妻子外不曾认识任何女人。然而在此我们不定规条，仅表达一种意见，以为何者较少受指摘。因为按正确观点，对已行补赎者而言，最合适的莫过于持守心灵与身体的双重贞洁。\par

% structure: p0049 source=h3 target=subsubsection reason=relative-heading-rank; base=h2

\subsubsection{问题十四}

\Emph{论隐修士从军或结婚者}。\par

\Emph{答}：隐修士自愿或由己愿所发的誓愿，不可无罪放弃。因为人向天主所许的愿，也该偿还。因此，凡舍弃独身生活的修道并投身军旅或婚姻者，必须做补赎并公开表明自己，因为尽管军役可能无罪、婚姻生活可敬，但抛弃更高选择即是过犯。\par

% structure: p0052 source=h3 target=subsubsection reason=relative-heading-rank; base=h2

\subsubsection{问题十五}

\Emph{论少女曾穿着会衣一段时间但未受奉献，后出嫁者}。\par

\Emph{答}：少女若未经父母命令强迫，而是凭自己自由意志发了贞洁愿并穿会衣，后来若选择婚姻，则行事有错，即使未受奉献亦然；她们若坚守誓愿，无疑不会受欺骗而不得奉献。\par

% structure: p0055 source=h3 target=subsubsection reason=relative-heading-rank; base=h2

\subsubsection{问题十六}

\Emph{论曾有基督徒父母留下的婴儿，若找不到他们受洗的证明，是否应为他们施洗？}。\par

\Emph{答}：若在亲属、邻舍或神职人员中，并无证据可证明所问之人已受洗，则必须设法为他们施行再生之礼，免得他们显然丧亡；因为在此情况下，理由不允许将未证明已行之事视为重复。\par

% structure: p0058 source=h3 target=subsubsection reason=relative-heading-rank; base=h2

\subsubsection{问题十七}

\Emph{论曾被敌人掳掠、不确定是否受洗但记得父母屡次带他们去教堂者，当他们回到罗马领土时，是否可以或应该受洗？}。\par

\Emph{答}：凡能记得曾随父母去教堂者，也能记得是否领受了那常给父母的东西。但若连此事也已遗忘，则似乎应将那未知是否已施予之事施行于他们，因为在最忠信的谨慎下，不会有冒失的妄断。\par

% structure: p0061 source=h3 target=subsubsection reason=relative-heading-rank; base=h2

\subsubsection{问题十八}

\Emph{论来自非洲或毛里塔尼亚、不知自己曾在何派别受洗者，该如何处理？}。\par

\Emph{答}：这些人对自己的洗礼并不怀疑，只是声称不知施洗者的信仰；因此，既然他们已以某种方式领受了洗礼的形式，不应再受洗，而应藉按手礼、呼求圣神之能（他们无法从异端者领受）后，与天主教徒合一。\par

% structure: p0064 source=h3 target=subsubsection reason=relative-heading-rank; base=h2

\subsubsection{问题十九}

\Emph{论婴孩时受洗后被外邦人掳去，按外邦人方式与他们同住，当他们仍年轻回到罗马领土时，若求领圣体，该如何办理？}。\par

\Emph{答}：若他们仅与外邦人同住并吃了祭物，则可通过禁食和按手得以净化，日后禁绝祭偶像之物，便可分享基督的奥迹。但若他们曾拜偶像，或被凶杀或奸淫所玷污，则不得领受共融，除非经过公开补赎。\par

\SourceRecord{Translated by Charles Lett Feltoe.；Nicene and Post-Nicene Fathers, Second Series；Vol. 12.；Edited by Philip Schaff and Henry Wace.；Buffalo, NY: Christian Literature Publishing Co.,；1895.；<http://www.newadvent.org/fathers/3604167.htm>.}

% structure: p0001 source=h1 target=section reason=page-title

\section{\WorkTitle{Letter 169}}

\Person{Leo}主教，致\Person{Leo Augustus}皇帝。\par

% structure: p0003 source=h2 target=subsection reason=relative-heading-rank; base=h2

\subsection{Ⅰ. 他衷心感谢皇帝所作的一切，并请求他以任何可能的方式完成这一事业。}

若要竭尽全力称颂陛下为捍卫信德而作出的辉煌决断，以配得上这一伟大事件所要求的赞辞，我们必会发现，凭我这笨拙的口舌，实难表达感恩之情与普世教会的欢欣。然而，更相称的酬报正等待着您的行为和功绩——正是为了这事业，您展示了如此卓越的热忱，如今正荣耀地赢得所期望的结局。因此，陛下的仁爱理当知道：所有天主的教会都在赞美您、欢欣鼓舞，因为那不敬神的弑亲者已被从亚历山大里亚教会的颈项上甩脱；天主的人民，那可恶的强盗曾如此沉重地压迫他们，如今已恢复信德的古老自由，能藉着忠信司铎的宣讲重回救恩之路——当他们看到一切瘟疫的温床随着首恶本人被铲除而消失时。现在，既然您以坚定的决心和不变的意志完成了此事，请颁布合乎天主喜乐的法令，支持此城的公教统治者，他丝毫未染上那早已屡次被谴责的异端邪说，以此完成为信德而工作的事业：否则，那看似愈合、却仍潜伏于疤痕下的伤口可能会蔓延，而因您的公开行动已摆脱异端者乖谬的基督徒平信徒，将再次被致命的毒液吞噬。\par

% structure: p0005 source=h2 target=subsection reason=relative-heading-rank; base=h2

\subsection{Ⅱ. 司铎不仅需要信德的纯正，也需要善行。}

但是，可敬的皇帝，您看到并清楚了解：在考虑那将要被绝罚之人的案件时，不仅必须审视他信德的纯正；因为即使信德能借某种惩罚和告解得以涤除，并通过某些条件完全恢复，然而那些犯下的邪恶血腥行径，决不会因花言巧语的主张而消除：因为在天主的司祭中，尤其是在如此重大教会的司铎中，舌头的声响和嘴唇的发音是不够的；若天主以自己的声音宣告，而心灵被证实有亵渎之罪，则什么也无效。关于这样的人，圣神藉宗徒说：\BibleQuote{有虔敬的外貌，却否认了它的能力；}又在别处说：\BibleQuote{他们自称认识天主，但在行为上却否认祂。}因此，既然在教会的每一个成员身上，既要求真信德的纯正，又要求善行的丰盛，那么在首席司祭身上，这两者岂不更应突出，因为缺少其中一样，就不能与基督奥体结合。\par

% structure: p0007 source=h2 target=subsection reason=relative-heading-rank; base=h2

\subsection{Ⅲ. 不应允许蒂茂德以正统为由请求宽免。}

我们无需在此逐一列举使蒂茂德成为可诅咒者的一切事，因为通过他和因他所作的事，已充分且显著地为全世界所知；任何由暴民违反公义所犯的罪行，都归在他头上，因为狂乱的手是在满足他的意愿。因此，即使他在信德宣认中毫无遗漏、丝毫未欺骗我们，最相称陛下荣耀的做法是：完全将他排除在这项计划之外，因为在这座如此重大城市的主教身上，普世教会应以神圣的欢跃而喜乐，使得主的真平安不仅因信德的宣讲，也因人们行为的榜样而受到光荣。公元460年6月17日，在\Person{Magnus}与\Person{Apollonius}执政之年。（经由事务代理人\Person{Philoxenus}之手。）\par

\SourceRecord{Translated by Charles Lett Feltoe.；Nicene and Post-Nicene Fathers, Second Series；Vol. 12.；Edited by Philip Schaff and Henry Wace.；Buffalo, NY: Christian Literature Publishing Co.,；1895.；<http://www.newadvent.org/fathers/3604169.htm>.}

% structure: p0001 source=h1 target=section reason=page-title

\section{第一七一号书信}

良主教，致亚历山大里亚教会之天主教主教弟茂德。\par

% structure: p0003 source=h2 target=subsection reason=relative-heading-rank; base=h2

\subsection{一、祝贺当选，并劝勉赢回迷途者归栈}

宗徒所引之言光辉灿烂，显明\BibleQuote{凡爱天主的人，事事都有助于他们的获得益处}\BibleRef{罗 8:28}，且因天主怜悯的安排，在遭遇逆境之处，也赐予顺境。亚历山大里亚教会的经验正显示这一点：其中谦卑者的温和与长期忍耐，为他们的耐心积蓄了巨大的回报；因为\Quote{上主亲近心灵破碎的人，必拯救精神谦卑的人}，我们高贵君王的信德在万事中受光荣，藉此\Quote{上主的右手大显威能}，阻止了敌基督之可憎之物再占据诸圣教父的宝座；其亵渎伤害他人不如伤害自己，因为他虽引诱一些人成为其罪孽的同伙，却以自己的血使自己永不可赎。因此，关于在信德指引下，你的当选——由司铎、平信徒及全体信徒所促成——我向你保证，整个主的教会与我一同喜乐；我切愿神圣的怜悯以多种恩宠标志慈爱地坚定这喜乐，你的虔诚在一切事上协助，使你藉着教会的祈祷，殷勤地赢回那些迄今仍抗拒真理的人，使他们与天主和好，并作为热忱的牧者，引领他们与公教信德奥体合一，这奥体之整体不可分割，效法那位真实而温和的善牧，祂为自己的羊舍弃生命，当一只羊迷失时，祂不用鞭子驱回，而是亲自背在肩上带回羊栈。\par

% structure: p0005 source=h2 target=subsection reason=relative-heading-rank; base=h2

\subsection{二、警醒防备异端，并频繁向罗马呈报}

因此，亲爱的弟兄，你要谨慎，免得在天主子民中发现任何聂斯多略或欧迪克错误的痕迹；因为\BibleQuote{除已奠立了的根基，即耶稣基督外，任何人不能再奠立别的}\BibleRef{格前 3:11}；祂若未藉信德重生，在我们的肉身实质上收纳我们众人，就不会使整个世界与天主父和好。因此，每当有写信的机会，弟兄，你既已按惯例通过我们的儿子——司铎达尼尔和执事弟茂德——将你祝圣的消息呈报给我们，就请继续时常如此行，并尽可能频繁地向我们（我们对此殷切期盼）寄送有关和平进展的报告，以便通过定期往来，我们感受到\BibleQuote{天主的爱，藉着所赐与我们的圣神，已倾注在我们心中}\BibleRef{罗 5:5}。发于八月十八日，玛格努斯与亚波罗尼乌斯执政之年（460年）。\par

\SourceRecord{Translated by Charles Lett Feltoe.；Nicene and Post-Nicene Fathers, Second Series；Vol. 12.；Edited by Philip Schaff and Henry Wace.；Buffalo, NY: Christian Literature Publishing Co.,；1895.；<http://www.newadvent.org/fathers/3604171.htm>.}
